% GENERATED FILE -- DO NOT EDIT THE SEMANTIC BODY DIRECTLY.
% CANONICAL-MARKDOWN-PATH: ../research/formal-theorems-and-proofs.md
% CANONICAL-MARKDOWN-SHA256: 34c3e205ae13877db508d778a8fea1878fff597d011be819bf801f15881ebd32
% Regenerate with md/research/check-proof-sync.ps1 -UpdateProjections.
% Options for packages loaded elsewhere
\PassOptionsToPackage{unicode}{hyperref}
\PassOptionsToPackage{hyphens}{url}
\documentclass[
  english,
  11pt,
  a4paper,
]{article}
\usepackage{xcolor}
\usepackage[margin=2.5cm]{geometry}
\usepackage{amsmath,amssymb}
\setcounter{secnumdepth}{-\maxdimen} % remove section numbering
\usepackage{iftex}
\ifPDFTeX
  \usepackage[T1]{fontenc}
  \usepackage[utf8]{inputenc}
  \usepackage{textcomp} % provide euro and other symbols
\else % if luatex or xetex
  \usepackage{unicode-math} % this also loads fontspec
  \defaultfontfeatures{Scale=MatchLowercase}
  \defaultfontfeatures[\rmfamily]{Ligatures=TeX,Scale=1}
\fi
\usepackage{lmodern}
\ifPDFTeX\else
  % xetex/luatex font selection
\fi
% Use upquote if available, for straight quotes in verbatim environments
\IfFileExists{upquote.sty}{\usepackage{upquote}}{}
\IfFileExists{microtype.sty}{% use microtype if available
  \usepackage[]{microtype}
  \UseMicrotypeSet[protrusion]{basicmath} % disable protrusion for tt fonts
}{}
\makeatletter
\@ifundefined{KOMAClassName}{% if non-KOMA class
  \IfFileExists{parskip.sty}{%
    \usepackage{parskip}
  }{% else
    \setlength{\parindent}{0pt}
    \setlength{\parskip}{6pt plus 2pt minus 1pt}}
}{% if KOMA class
  \KOMAoptions{parskip=half}}
\makeatother
\usepackage{color}
\usepackage{fancyvrb}
\newcommand{\VerbBar}{|}
\newcommand{\VERB}{\Verb[commandchars=\\\{\}]}
\DefineVerbatimEnvironment{Highlighting}{Verbatim}{commandchars=\\\{\}}
% Add ',fontsize=\small' for more characters per line
\newenvironment{Shaded}{}{}
\newcommand{\AlertTok}[1]{\textcolor[rgb]{1.00,0.00,0.00}{\textbf{#1}}}
\newcommand{\AnnotationTok}[1]{\textcolor[rgb]{0.38,0.63,0.69}{\textbf{\textit{#1}}}}
\newcommand{\AttributeTok}[1]{\textcolor[rgb]{0.49,0.56,0.16}{#1}}
\newcommand{\BaseNTok}[1]{\textcolor[rgb]{0.25,0.63,0.44}{#1}}
\newcommand{\BuiltInTok}[1]{\textcolor[rgb]{0.00,0.50,0.00}{#1}}
\newcommand{\CharTok}[1]{\textcolor[rgb]{0.25,0.44,0.63}{#1}}
\newcommand{\CommentTok}[1]{\textcolor[rgb]{0.38,0.63,0.69}{\textit{#1}}}
\newcommand{\CommentVarTok}[1]{\textcolor[rgb]{0.38,0.63,0.69}{\textbf{\textit{#1}}}}
\newcommand{\ConstantTok}[1]{\textcolor[rgb]{0.53,0.00,0.00}{#1}}
\newcommand{\ControlFlowTok}[1]{\textcolor[rgb]{0.00,0.44,0.13}{\textbf{#1}}}
\newcommand{\DataTypeTok}[1]{\textcolor[rgb]{0.56,0.13,0.00}{#1}}
\newcommand{\DecValTok}[1]{\textcolor[rgb]{0.25,0.63,0.44}{#1}}
\newcommand{\DocumentationTok}[1]{\textcolor[rgb]{0.73,0.13,0.13}{\textit{#1}}}
\newcommand{\ErrorTok}[1]{\textcolor[rgb]{1.00,0.00,0.00}{\textbf{#1}}}
\newcommand{\ExtensionTok}[1]{#1}
\newcommand{\FloatTok}[1]{\textcolor[rgb]{0.25,0.63,0.44}{#1}}
\newcommand{\FunctionTok}[1]{\textcolor[rgb]{0.02,0.16,0.49}{#1}}
\newcommand{\ImportTok}[1]{\textcolor[rgb]{0.00,0.50,0.00}{\textbf{#1}}}
\newcommand{\InformationTok}[1]{\textcolor[rgb]{0.38,0.63,0.69}{\textbf{\textit{#1}}}}
\newcommand{\KeywordTok}[1]{\textcolor[rgb]{0.00,0.44,0.13}{\textbf{#1}}}
\newcommand{\NormalTok}[1]{#1}
\newcommand{\OperatorTok}[1]{\textcolor[rgb]{0.40,0.40,0.40}{#1}}
\newcommand{\OtherTok}[1]{\textcolor[rgb]{0.00,0.44,0.13}{#1}}
\newcommand{\PreprocessorTok}[1]{\textcolor[rgb]{0.74,0.48,0.00}{#1}}
\newcommand{\RegionMarkerTok}[1]{#1}
\newcommand{\SpecialCharTok}[1]{\textcolor[rgb]{0.25,0.44,0.63}{#1}}
\newcommand{\SpecialStringTok}[1]{\textcolor[rgb]{0.73,0.40,0.53}{#1}}
\newcommand{\StringTok}[1]{\textcolor[rgb]{0.25,0.44,0.63}{#1}}
\newcommand{\VariableTok}[1]{\textcolor[rgb]{0.10,0.09,0.49}{#1}}
\newcommand{\VerbatimStringTok}[1]{\textcolor[rgb]{0.25,0.44,0.63}{#1}}
\newcommand{\WarningTok}[1]{\textcolor[rgb]{0.38,0.63,0.69}{\textbf{\textit{#1}}}}
\usepackage{longtable,booktabs,array}
\newcounter{none} % for unnumbered tables
\usepackage{calc} % for calculating minipage widths
% Correct order of tables after \paragraph or \subparagraph
\usepackage{etoolbox}
\makeatletter
\patchcmd\longtable{\par}{\if@noskipsec\mbox{}\fi\par}{}{}
\makeatother
% Allow footnotes in longtable head/foot
\IfFileExists{footnotehyper.sty}{\usepackage{footnotehyper}}{\usepackage{footnote}}
\makesavenoteenv{longtable}
\ifLuaTeX
\usepackage[bidi=basic,shorthands=off]{babel}
\else
\usepackage[bidi=default,shorthands=off]{babel}
\fi
\ifLuaTeX
  \usepackage{selnolig} % disable illegal ligatures
\fi
\setlength{\emergencystretch}{3em} % prevent overfull lines
\providecommand{\tightlist}{%
  \setlength{\itemsep}{0pt}\setlength{\parskip}{0pt}}
\usepackage{bookmark}
\IfFileExists{xurl.sty}{\usepackage{xurl}}{} % add URL line breaks if available
\urlstyle{same}
% fallback for those not using the hyperref driver hyperxmp:
\makeatletter
\@ifundefined{xmpquote}{\newcommand{\xmpquote}[1]{#1}}{}
\makeatother
\hypersetup{
  pdftitle={Detailed Proof Report on Conditional Semantic Preservation for OCL-to-Cypher Validation},
  pdflang={en},
  hidelinks,
  pdfcreator={LaTeX via pandoc}}

\title{Detailed Proof Report on Conditional Semantic Preservation for
OCL-to-Cypher Validation}
\author{}
\date{}

\begin{document}
\maketitle

{
\setcounter{tocdepth}{3}
\tableofcontents
}
This document gives a paper-ready technical foundation for the
preservation claim of the framework:

\begin{Shaded}
\begin{Highlighting}[]
\NormalTok{The reference Cypher invariant query returns exactly the same violating}
\NormalTok{objects as the OCL validation semantics over the source UML object model.}
\end{Highlighting}
\end{Shaded}

The claim is intentionally conditional and scoped. The framework is not
claimed to be a complete OCL-to-Cypher compiler and this document does
not prove correctness for full OMG OCL.

The paper-safe wording is:

\begin{Shaded}
\begin{Highlighting}[]
\NormalTok{We prove conditional end{-}to{-}end preservation of violation sets for a finite{-}set}
\NormalTok{OCL validation fragment translated to Cypher under an explicit}
\NormalTok{UML{-}to{-}property{-}graph encoding contract.}
\end{Highlighting}
\end{Shaded}

\begin{center}\rule{0.5\linewidth}{0.5pt}\end{center}

\section{1. Scope and Assumptions}\label{1-scope-and-assumptions}

All theorem statements in this document are restricted to the following
scope.

\begin{Shaded}
\begin{Highlighting}[]
\NormalTok{A1. tInv is invariant text; ParseInvariant(tInv)=astInv succeeds;}
\NormalTok{    decodeInvariant\_val(astInv)=I=(cName,R,e), where cName is an unresolved}
\NormalTok{    source class name, and astInv belongs to the admitted}
\NormalTok{    representation ASTInv\_val. The expression e is pred(I); it is never a}
\NormalTok{    free top{-}level compilation root and is never identified with tInv or}
\NormalTok{    astInv.}
\NormalTok{A2. resolveClass\_MM(cName)=C and}
\NormalTok{    T\_BIND\_INV(astInv,MM)=BInvariant(C,R,b) succeed, and b is well typed as}
\NormalTok{    Boolean in BoundOCL\_val(MM) under Gamma0=\{self:C\}. At runtime MM may be}
\NormalTok{    supplied by a graph{-}backed metamodel view that is resolver{-}equivalent to}
\NormalTok{    the encoded M2 metamodel.}
\NormalTok{A3. Validation uses the explicitly defined certified finite{-}set semantics of}
\NormalTok{    OCL\_val, including its validation truth collapse and its Set{-}image collect;}
\NormalTok{    it does not use the full OMG null/invalid or collection{-}kind semantics.}
\NormalTok{A4. MM=M2 is a finite domain metamodel and M=M1 is a valid finite UML object}
\NormalTok{    model conforming to MM.}
\NormalTok{A5. GraphAdequate(MM,M,G) holds: G contains the required M2 schema, M1}
\NormalTok{    instance, and typing projections and satisfies R1{-}{-}R7, including}
\NormalTok{    resolver{-}equivalent M2 lookup. Writing G=Phi(MM,M) is permitted only when}
\NormalTok{    Phi has separately been shown to establish GraphAdequate.}
\NormalTok{A6. nva=T\_NORM(va) is the only normative normalized artifact,}
\NormalTok{    ValidNVA(nva) holds, plan=T\_CQM\^{}spec(C,R,nva) is defined, and}
\NormalTok{    T\_TEXT\^{}spec(plan)=(tq,pi) is defined,}
\NormalTok{    parse\_Cypher(tq)=q in CYPHER5\_val, the selected runtime satisfies}
\NormalTok{    CY1{-}{-}CY9, and SpecPlanAdequacy(MM,C,nva,plan,q,pi,G) holds.}
\NormalTok{    SpecPlanAdequacy is structural only. Returned{-}ID equality is the}
\NormalTok{    conclusion of SpecPlanSim\_sound, proved by constructor induction from the}
\NormalTok{    local NVA/CQM rules; it is not a field or premise of adequacy.}
\NormalTok{    Java T\_OPT, its planner, runtime tests, and performance measurements are}
\NormalTok{    non{-}normative implementation evidence and do not occur in this theorem.}
\NormalTok{    Every collection boundary implements finiteSet(bottom\_Set(tau))=empty}
\NormalTok{    rather than exposing backend null to a set operator.}
\NormalTok{A7. q has exactly one public result column named useId. returnedIds(q,G,pi) is}
\NormalTok{    defined only when every value in that column belongs to the target ObjectId}
\NormalTok{    carrier; there is no alternative node{-}returning interpretation. A7 does}
\NormalTok{    not assume that a returned ObjectId denotes a source object{-}{-}that no{-}ghost}
\NormalTok{    property is the conclusion of C6.}
\NormalTok{A8. ScalarClosed(C,e,M) holds: every scalar subevaluation for every admitted}
\NormalTok{    context object remains in the exact scalar profile of Section 3.3 (no}
\NormalTok{    overflow, non{-}finite result, inexact coercion, or zero divisor).}
\NormalTok{A9. For the generated pair \textasciigrave{}(q,pi)\textasciigrave{} fixed by A6,}
\NormalTok{    BottomSeparated(MM,M,G,e,pi) holds and Params(q) subseteq dom(pi): the reserved representation of}
\NormalTok{    semantic bottom is disjoint from every admitted literal, stored scalar,}
\NormalTok{    object id, class key, association key, attribute key, and qualifier payload}
\NormalTok{    reachable during this compilation and execution.}
\end{Highlighting}
\end{Shaded}

These assumptions deliberately separate four levels:

\begin{Shaded}
\begin{Highlighting}[]
\NormalTok{tInv : invariant text}
\NormalTok{astInv : parser representation}
\NormalTok{I=(cName,R,e) : decoded source invariant in OCL\_val}
\NormalTok{BInvariant(C,R,b) : resolved and typed bound invariant}
\end{Highlighting}
\end{Shaded}

Only \texttt{e} and \texttt{b} have denotations. Parsing is applied to
\texttt{tInv}; binding is applied to \texttt{astInv}. A certified
\texttt{let} or iterator declaration may omit its type or state it
explicitly. The parser AST retains the optional type name; the binder
resolves it against \texttt{MM}, checks initializer/element conformance,
and stores the resolved declaration type in BoundOCL, semantic IR, and
the Cypher plan. Likewise, \texttt{bottom}, Boolean collapse, and
Set-image \texttt{collect} below define OCL\_val; they must not be cited
as the denotation of full OMG OCL.

\subsection{1.0 Normative specification
pipeline}\label{10-normative-specification-pipeline}

The only theorem-bearing pipeline in this report is:

\begin{Shaded}
\begin{Highlighting}[]
\NormalTok{OCL\_val {-}\textgreater{} OVA {-}\textgreater{} T\_NORM {-}\textgreater{} NVA {-}\textgreater{} T\_CQM\^{}spec {-}\textgreater{} CQM {-}\textgreater{} T\_TEXT\^{}spec {-}\textgreater{} Cypher}
\end{Highlighting}
\end{Shaded}

\texttt{NVA} is the typed normal-form subtype described by
\texttt{specification/metamodel/Normalized-Validation-Algebra.emf}. Its
concrete grammar contains exactly the reachable rows marked
\emph{reachable} in the Theorem 4 constructor matrix. It has no
constructors for \texttt{ForAll}, \texttt{Implies}, \texttt{Xor},
\texttt{Reject}, \texttt{Includes}, \texttt{Excludes}, \texttt{Size},
\texttt{IsEmpty}, \texttt{NotEmpty}, or the two count-on-navigation
redexes. Define:

\begin{Shaded}
\begin{Highlighting}[]
\NormalTok{ValidNVA(nva)}
\NormalTok{iff EcoreConforms\_NVA(nva)}
\NormalTok{  and WF\_NVA(nva)}
\NormalTok{  and NF\_R(nva).}

\NormalTok{CertifiedNVA(va,nva)}
\NormalTok{iff ValidNVA(nva)}
\NormalTok{  and Ref\_OVA\_NVA(va,nva)}
\NormalTok{  and nva=T\_NORM(va)}
\NormalTok{  and [[nva]]\_I=[[va]]\_I for I in \{obj,graph\}.}
\end{Highlighting}
\end{Shaded}

The last two conjuncts are relational proof facts. They are never
represented by a mutable \texttt{correct:Boolean} model attribute.
\texttt{WF\_NVA} checks rooted finite containment, unique IDs,
constructor typing, nearest lexical binding,
\texttt{sourceCollectionType}, finite-Set shape, bottom safety, and
absence of every redex in \texttt{R}. Theorem 3 supplies determinism,
totality on certified OVA, termination, type/scope/alpha preservation,
bottom/finite-set preservation, and denotational preservation.

The scope excludes arbitrary OCL, the full OMG invalid/null and
collection-kind semantics, arbitrary Cypher, encodings outside
\texttt{GraphAdequate}, unbounded or unprobed Neo4j behavior, and
correctness of the Java optimizer. Java \texttt{T\_OPT}, the Java
planner/renderer, real-Neo4j runs, and performance measurements are
implementation evidence only. The optional proposition
\texttt{OptRefines(T\_OPT(va),T\_NORM(va))} is not used by Theorems
0-\/-6.

\subsection{1.1 Canonical Theorem
Contracts}\label{11-canonical-theorem-contracts}

The following dependency rows are normative. A theorem\textquotesingle s
prose statement may expand a row but may not silently remove an
assumption or dependency.

The structured JSON block immediately below is part of this formal
document and is the canonical machine-readable contract. It is hidden
from the rendered paper, but it is intentionally kept beside the
definitions and theorems that give it meaning. The external registry is
regenerated from this block; it is not an independent authority.

\{ "version": "PC-2026-07-22.3", "registrySchemaVersion": 4,
"claimStatusVocabulary": \{ "proved": "kernel-checked theorem under its
explicitly quantified premises", "conditional": "mathematical conclusion
requiring named graph/Cypher assumptions", "partial": "a proof boundary
with undisclosed local cases still prohibited", "empirical": "finite
implementation or runtime evidence, never universal proof" \},
"claimInventory": {[} \{ "id": "CLAIM-SPEC-NORM", "status": "proved",
"scope": "T\_NORM termination, relative normal form, typing and
denotation under Theorem 3 premises" \}, \{ "id": "CLAIM-SPEC-PLAN",
"status": "conditional", "scope": "NVA-to-CQM-to-reference-Cypher
realization under LR/C/BR/CY local agreements" \}, \{ "id":
"CLAIM-JAVA-OPT", "status": "partial", "scope": "optional
OptRefines(T\_OPT(va),T\_NORM(va)); excluded from Theorems 0-\/-6" \},
\{ "id": "CLAIM-NEO4J-RUNTIME", "status": "empirical", "scope": "finite
executions on the pinned selected runtime profile" \} {]}, "authority":
\{ "kind": "derived\_projection\_metadata", "canonicalPath":
"md/research/formal-theorems-and-proofs.md", "canonicalBlock":
"CANONICAL PROOF-CONTRACT REGISTRY JSON" \}, "assumptions": {[} \{ "id":
"A1", "key": "invariant\_parse\_decode\_and\_admission" \}, \{ "id":
"A2", "key": "successful\_context\_fixed\_invariant\_binding" \}, \{
"id": "A3", "key": "finite\_set\_validation\_semantics" \}, \{ "id":
"A4", "key": "valid\_finite\_conforming\_model" \}, \{ "id": "A5",
"key": "adequate\_graph\_encoding" \}, \{ "id": "A6", "key":
"reference\_nva\_cqm\_cypher\_and\_structural\_adequacy" \}, \{ "id":
"A7", "key": "single\_well\_typed\_use\_id\_result\_column" \}, \{ "id":
"A8", "key": "scalar\_closed" \}, \{ "id": "A9", "key":
"bottom\_separated" \} {]}, "scopeLemmas": {[} \{ "id": "M1", "title":
"Source-\/-AST Representation", "titleVi": "Biểu diễn Source-\/-AST" \},
\{ "id": "M2", "title": "Bound Subexpression Closure and Structural
Induction", "titleVi": "Tính đóng biểu thức Bound con và quy nạp cấu
trúc" \}, \{ "id": "M3", "title": "Canonical Type Uniqueness" \}, \{
"id": "M3a", "title": "Set-Join Functionality", "titleVi": "Tính hàm của
các phép join tập" \}, \{ "id": "M4", "title": "Admission and Binding
Soundness", "titleVi": "Tính đúng đắn của admission và binding" \}, \{
"id": "M5", "title": "Exclusion Stability" \} {]}, "theorems": {[} \{
"id": "PC-T0", "number": 0, "title": "Object-Graph Representation",
"titleVi": "Biểu diễn đối tượng-\/-đồ thị", "statement": "validation
observations of \texttt{(MM,M)} and every adequate \texttt{G} agree",
"statementVi": "Các quan sát kiểm tra trên đối tượng và đồ thị đủ điều
kiện trùng khớp", "assumptions": {[}"A4", "A5"{]}, "requires": {[}"R1",
"R2", "R3", "R4", "R5", "R6", "R7"{]} \}, \{ "id": "PC-T1", "number": 1,
"title": "Binder Soundness", "titleVi": "Tính đúng đắn của binder",
"statement": "decoded Source and uniquely bound denotations are related
and erase equally", "statementVi": "Ngữ nghĩa Source đã giải mã và Bound
duy nhất tương ứng sau khi xóa thẻ", "assumptions": {[}"A1", "A2", "A4",
"A8"{]}, "requires": {[}"M2", "M3", "M3a", "M4", "B0", "B1", "B2", "B3",
"B4", "SB0"{]} \}, \{ "id": "PC-T2", "number": 2, "title": "Validation
Algebra Abstraction", "titleVi": "Trừu tượng hóa Đại số Kiểm tra",
"statement": "erased Bound denotation equals typed object-VA
denotation", "statementVi": "Ngữ nghĩa Bound sau xóa thẻ bằng ngữ nghĩa
VA có kiểu phía đối tượng", "assumptions": {[}"A2", "A3", "A4", "A8"{]},
"requires": {[}"M2", "M3", "M3a", "VA1", "VA2", "VA3", "BV0", "B4"{]}
\}, \{ "id": "PC-T3", "number": 3, "title": "Normalization
Preservation", "titleVi": "Bảo toàn chuẩn hóa", "statement":
"deterministic bottom-up normalization terminates, preserves
type/denotation and reachable closure, and reaches \texttt{NF\_R} modulo
alpha-equivalence", "statementVi": "Chuẩn hóa bottom-up dừng, bảo toàn
kiểu/ngữ nghĩa/miền đóng và đạt NF\_R theo alpha", "assumptions":
{[}"A2", "A3", "A8"{]}, "requires": {[}"N0", "N1", "N2", "N3", "N4"{]}
\}, \{ "id": "PC-T4", "number": 4, "title": "Validation Preservation",
"titleVi": "Bảo toàn kiểm tra", "statement": "typed object and graph VA
denotations commute with \texttt{encodeValue} on reachable closed
evaluations", "statementVi": "Đánh giá VA có kiểu giao hoán với mã hóa
trên các đánh giá reachable đóng", "assumptions": {[}"A2", "A3", "A4",
"A5", "A8"{]}, "requires": {[}"T0", "G1", "G2", "G3", "G3a", "G4"{]} \},
\{ "id": "PC-T5", "number": 5, "title": "Reference Cypher Realization
under Cypher Assumptions", "titleVi": "Hiện thực hóa Cypher tham chiếu
dưới các giả thiết Cypher", "statement": "the reference
NVA-to-CQM-to-Cypher query returns exactly normalized graph-VA violation
IDs and no ghost IDs", "statementVi": "Truy vấn tham chiếu
NVA-\/-CQM-\/-Cypher trả về chính xác ID vi phạm VA đồ thị và không có
ID ma", "assumptions": {[}"A2", "A3", "A5", "A6", "A7", "A8", "A9"{]},
"requires": {[}"C1", "C2", "C3", "C3a", "C4", "C5", "C5a", "C5b", "C6",
"SpecPlanAdequacy", "SpecPlanSim\_sound", "TXT1", "TXT2", "TXT3",
"TXT4", "TXT5", "BR1-BR10", "CY1-CY9"{]} \}, \{ "id": "PC-T6", "number":
6, "title": "End-to-End Validation Equivalence", "titleVi": "Tương đương
kiểm tra đầu-cuối", "statement":
"\texttt{returnedIds(q,G,pi)=id{[}Viol\_OCL(e,C,M){]}}", "statementVi":
"Tập ID trả về bằng ảnh ID của tập vi phạm OCL", "assumptions": {[}"A1",
"A2", "A3", "A4", "A5", "A6", "A7", "A8", "A9"{]}, "requires": {[}"T0",
"T1", "T2", "T3", "T4", "T5"{]} \} {]}, "semanticFunctions": {[} \{
"id": "SF-01", "name": "source\_denotation", "signature":
"Expr\_OCLval(tau) x Model x Env\_S -\textgreater{} Val\_S(tau)",
"role": "independent Source OCL\_val denotation" \}, \{ "id": "SF-02",
"name": "bound\_denotation", "signature": "Expr\_Bound(tau) x Model x
Env\_B -\textgreater{} Val\_B(tau)", "role": "independent Bound OCL
denotation" \}, \{ "id": "SF-03", "name": "erase\_source", "signature":
"Val\_S(tau) -\textgreater{} ValBottom\_obj(tau)", "role": "forgets
Source semantic tags" \}, \{ "id": "SF-04", "name": "erase\_bound",
"signature": "Val\_B(tau) -\textgreater{} ValBottom\_obj(tau)", "role":
"forgets Bound semantic tags" \}, \{ "id": "SF-05", "name":
"oclval\_denotation", "signature": "Expr\_OCLval(tau) x Model x Env\_obj
-\textgreater{} ValBottom\_obj(tau)", "role": "erase\_source after
tagging an object environment; text and AST have no denotation" \}, \{
"id": "SF-06", "name": "erased\_bound\_denotation", "signature":
"Expr\_Bound(tau) x Model x Env\_obj -\textgreater{}
ValBottom\_obj(tau)", "role": "erase\_bound after tagging an object
environment" \}, \{ "id": "SF-07", "name": "object\_va\_denotation",
"signature": "Expr\_VA(tau) x Model x Env\_obj -\textgreater{}
ValBottom\_obj(tau)", "role": "object-side VA denotation" \}, \{ "id":
"SF-08", "name": "graph\_va\_denotation", "signature": "Expr\_VA(tau) x
Graph x Env\_graph -\textgreater{} ValBottom\_G(tau)", "role":
"graph-side VA denotation" \}, \{ "id": "SF-09", "name":
"graph\_denotation\_alias", "signature": "PrimitiveAccess\_VA(tau) x
Graph x Env\_graph -\textgreater{} ValBottom\_G(tau)", "role":
"restriction of SF-08, not a new semantics" \}, \{ "id": "SF-10",
"name": "encode\_value", "signature": "ValBottom\_obj(tau)
-\textgreater{} ValBottom\_G(tau)", "role": "typed object-to-graph value
encoding" \}, \{ "id": "SF-11", "name": "finite\_set", "signature":
"ValBottom\_I(Set(sigma)) -\textgreater{} P\_fin(ValBottom\_I(sigma))",
"role": "whole-collection-bottom-to-empty completion, retaining element
bottom" \}, \{ "id": "SF-12", "name": "navigation\_source\_lifting",
"signature": "ValBottom\_I(C) -\textgreater{} P\_fin(Val\_I(C))",
"role": "scalar entity/bottom navigation lifting" \}, \{ "id": "SF-13",
"name": "validation\_truth", "signature": "ValBottom\_I(Boolean)
-\textgreater{} Boolean", "role": "true iff represented value is Boolean
true" \}, \{ "id": "SF-14", "name": "typed\_equality", "signature":
"ValBottom\_I(tau) x ValBottom\_I(tau) -\textgreater{} Boolean", "role":
"total equality on one tagged type" \}, \{ "id": "SF-15", "name":
"attribute\_interpretation", "signature": "ValBottom\_I(C) x
Attribute(C,tau) -\textgreater{} ValBottom\_I(tau)", "role": "typed
attribute observation" \}, \{ "id": "SF-16", "name":
"navigation\_interpretation", "signature": "Val\_I(C) x
NavigationMetadata -\textgreater{} P\_fin(Val\_I(D))", "role":
"directional qualified target-set primitive used by distinct ONE/MANY
observations" \}, \{ "id": "SF-17", "name": "class\_interpretation",
"signature": "Class(C) -\textgreater{} P\_fin(Val\_I(C))", "role":
"finite allInstances interpretation" \}, \{ "id": "SF-18", "name":
"cypher\_representation", "signature": "ValBottom\_G(tau)
-\textgreater{} CyVal\_C(tau)", "role": "type-indexed injective target
representation" \}, \{ "id": "SF-19", "name": "cypher\_abstraction",
"signature": "CyVal\_C(tau) -\textgreater{} ValBottom\_G(tau)", "role":
"type-indexed left inverse of SF-18" \}, \{ "id": "SF-20", "name":
"cypher\_expression\_evaluation", "signature": "CExpr(tau) x Row
-\textgreater{} CyVal\_C(tau)", "role": "selected-dialect typed
expression evaluation" \}, \{ "id": "SF-21", "name":
"cypher\_row\_evaluation", "signature": "CQuery x Bag(Row)
-\textgreater{} Bag(Row)", "role": "selected-dialect correlated row
evaluation" \}, \{ "id": "SF-22", "name": "cypher\_query\_evaluation",
"signature": "CQuery x Graph x ParamEnv -\textgreater{} Bag(Row)",
"role": "closed-query target evaluation" \}, \{ "id": "SF-23", "name":
"prelude\_execution", "signature": "List(CClause) x Graph x ParamEnv x
Row -\textgreater{} Bag(Row)", "role": "scalar-plan prelude execution
under an explicit graph and parameter environment" \}, \{ "id": "SF-24",
"name": "semantic\_set\_projection", "signature": "Alias(sigma) x
Bag(Row) -\textgreater{} P\_fin(ValBottom\_G(sigma))", "role":
"abstracts one typed column and removes multiplicity" \}, \{ "id":
"SF-25", "name": "generated\_parameter\_map", "signature":
"CompiledQuery -\textgreater{} ParamEnv", "role": "second projection of
a generated query/parameter artifact" \}, \{ "id": "SF-26", "name":
"parameter\_name\_set", "signature": "CQuery -\textgreater{}
P\_fin(ParamName)", "role": "parameter names occurring syntactically in
the raw query" \}, \{ "id": "SF-27", "name":
"returned\_identifier\_set", "signature": "CQuery x Graph x ParamEnv
-\textgreater{} P\_fin(ObjectId)", "role": "set abstraction of returned
identifiers" \}, \{ "id": "SF-28", "name": "object\_violation\_set",
"signature": "Expr\_OCLval(Boolean) x Class x Model -\textgreater{}
P\_fin(Obj(M))", "role": "source validation violation set" \}, \{ "id":
"SF-29", "name": "to\_one\_navigation\_collection\_view", "signature":
"ValBottom\_I(C) -\textgreater{} P\_fin(Val\_I(C))", "role": "explicit
singleton/empty finite-set view used only when a collection operator
directly consumes a multiplicity-one navigation" \} {]}, "evidence": {[}
\{ "id": "EV-PO01-VOCABULARY", "kind": "test", "path":
"neo4j-tgg/src/test/java/org/uet/dse/neo4jtgg/experiment/RendererAccessorAgreementTest.java",
"symbol": "class RendererAccessorAgreementTest" \}, \{ "id":
"EV-PO01-PGMM", "kind": "source", "path": "md/research/model/PGMM.emf",
"symbol": "package pgmm;" \}, \{ "id": "EV-PO02-CERTIFIED", "kind":
"test", "path":
"neo4j-tgg/src/test/java/org/uet/dse/neo4jtgg/service/impl/CertifiedValidationCompilationTest.java",
"symbol": "class CertifiedValidationCompilationTest" \}, \{ "id":
"EV-PO02-OVA-METAMODEL", "kind": "source", "path":
"md/research/model/OCL-Validation-Algebra.emf", "symbol": "package ova;"
\}, \{ "id": "EV-PO02-CQM-METAMODEL", "kind": "source", "path":
"md/research/model/Cypher-Query-Model.emf", "symbol": "package cqm;" \},
\{ "id": "EV-PO02-CERTIFIED-VALIDATOR", "kind": "source", "path":
"neo4j-tgg/src/main/java/org/uet/dse/neo4jtgg/ocl/ir/OclCertifiedModelValidator.java",
"symbol": "class OclCertifiedModelValidator" \}, \{ "id":
"EV-PO02-CERTIFIED-VALIDATOR-GUARD", "kind": "test", "path":
"neo4j-tgg/src/test/java/org/uet/dse/neo4jtgg/ocl/ir/OclCertifiedModelValidatorTest.java",
"symbol": "class OclCertifiedModelValidatorTest" \}, \{ "id":
"EV-PO03-KEYS", "kind": "test", "path":
"neo4j-tgg/src/test/java/org/uet/dse/neo4jtgg/experiment/CanonicalEncodingIndependentContractTest.java",
"symbol": "class CanonicalEncodingIndependentContractTest" \}, \{ "id":
"EV-PO04-08-REPRESENTATION", "kind": "test", "path":
"neo4j-tgg/src/test/java/org/uet/dse/neo4jtgg/experiment/RepresentationAdequacyEvaluatorTest.java",
"symbol": "class RepresentationAdequacyEvaluatorTest" \}, \{ "id":
"EV-PO09-RENDERER", "kind": "test", "path":
"neo4j-tgg/src/test/java/org/uet/dse/neo4jtgg/experiment/RendererAccessorAgreementTest.java",
"symbol": "class RendererAccessorAgreementTest" \}, \{ "id":
"EV-PO10-BOTTOM", "kind": "test", "path":
"neo4j-tgg/src/test/java/org/uet/dse/neo4jtgg/ocl/OclBottomTokenContractTest.java",
"symbol": "class OclBottomTokenContractTest" \}, \{ "id":
"EV-PO11-SCALAR", "kind": "test", "path":
"neo4j-tgg/src/test/java/org/uet/dse/neo4jtgg/ocl/OclScalarClosureCheckerTest.java",
"symbol": "class OclScalarClosureCheckerTest" \}, \{ "id":
"EV-PO12-ALIAS", "kind": "test", "path":
"neo4j-tgg/src/test/java/org/uet/dse/neo4jtgg/experiment/BottomSafeReceiverAgreementTest.java",
"symbol": "class BottomSafeReceiverAgreementTest" \}, \{ "id":
"EV-PO13-RAW-AST", "kind": "test", "path":
"neo4j-tgg/src/test/java/org/uet/dse/neo4jtgg/ocl/ir/RawCypherAstTotalityTest.java",
"symbol": "class RawCypherAstTotalityTest" \}, \{ "id":
"EV-PO14-DIRECT-TEXT-BRIDGE", "kind": "test", "path":
"neo4j-tgg/src/test/java/org/uet/dse/neo4jtgg/experiment/Neo4jCypherAstBridgeTest.java",
"symbol": "class Neo4jCypherAstBridgeTest" \}, \{ "id":
"EV-PO14-BOTTOM-ROUNDTRIP", "kind": "test", "path":
"neo4j-tgg/src/test/java/org/uet/dse/neo4jtgg/experiment/BottomSafeReceiverAgreementTest.java",
"symbol": "bottomSafeReceiverBranchesRoundTripThroughBothParsers" \}, \{
"id": "EV-PO14-CONSTRUCTOR-MATRIX", "kind": "file", "path":
"verification/coverage/cypher\_plan\_constructor\_matrix.csv" \}, \{
"id": "EV-PO14-CONSTRUCTOR-MATRIX-GUARD", "kind": "test", "path":
"neo4j-tgg/src/test/java/org/uet/dse/neo4jtgg/experiment/OclCypherPlanConstructorEvidenceMatrixTest.java",
"symbol":
"everySealedPlanConstructorHasFactoryPlannerRendererFormalRuleAndRoundTripEvidence"
\}, \{ "id": "EV-PO15-RUNTIME-MANIFEST", "kind": "file", "path":
"verification/evidence/cypher5-val-runtime-2026-08-01.tsv" \}, \{ "id":
"EV-PO15-RUNTIME-GUARD", "kind": "test", "path":
"neo4j-tgg/src/test/java/org/uet/dse/neo4jtgg/experiment/Cypher5ValRuntimeEvidenceManifestTest.java",
"symbol": "class Cypher5ValRuntimeEvidenceManifestTest" \}, \{ "id":
"EV-PO16-OCL47-MANIFEST", "kind": "file", "path":
"verification/evidence/oclval47-nonvacuity-runtime-2026-08-01.tsv" \},
\{ "id": "EV-PO16-OCL47-GUARD", "kind": "test", "path":
"neo4j-tgg/src/test/java/org/uet/dse/neo4jtgg/experiment/OclVal47NonVacuityEvidenceTest.java",
"symbol": "class OclVal47NonVacuityEvidenceTest" \}, \{ "id":
"EV-CANONICAL-RUNTIME-SUPPLEMENT", "kind": "file", "path":
"verification/evidence/canonical-profile-runtime-2026-08-09.tsv" \}, \{
"id": "EV-CANONICAL-RUNTIME-SUPPLEMENT-GUARD", "kind": "test", "path":
"neo4j-tgg/src/test/java/org/uet/dse/neo4jtgg/experiment/CanonicalProfileRuntimeEvidenceManifestTest.java",
"symbol": "class CanonicalProfileRuntimeEvidenceManifestTest" \}, \{
"id": "EV-PO17-CHECKER", "kind": "source", "path":
"verification/scripts/check-verification-contract.ps1", "symbol":
"Machine verification contract PASS" \}, \{ "id": "EV-PO17-MUTATIONS",
"kind": "source", "path":
"verification/scripts/check-verification-contract-mutations.ps1",
"symbol": "Machine verification mutation tests PASS" \}, \{ "id":
"EV-PO18-LEAN", "kind": "source", "path":
"verification/lean/Ocl2CypherProof.lean", "symbol": "theorem
java\_ir\_eval\_refinement" \}, \{ "id": "EV-PO18-PRODPLAN-INDUCTION",
"kind": "source", "path": "verification/lean/Ocl2CypherProof.lean",
"symbol": "theorem prod\_plan\_sim\_sound" \}, \{ "id":
"EV-SPEC-NVA-GRAMMAR", "kind": "source", "path":
"verification/lean/Ocl2CypherProof.lean", "symbol": "theorem
certified\_nva\_grammar\_complete" \}, \{ "id": "EV-SPEC-PLANSIM",
"kind": "source", "path": "verification/lean/Ocl2CypherProof.lean",
"symbol": "theorem spec\_plan\_sim\_sound" \}, \{ "id":
"EV-SPEC-NVA-METAMODEL", "kind": "source", "path":
"md/research/specification/metamodel/Normalized-Validation-Algebra.emf",
"symbol": "package nva;" \}, \{ "id": "EV-SPEC-NVA-ECORE", "kind":
"file", "path":
"md/research/specification/metamodel/Normalized-Validation-Algebra.ecore"
\}, \{ "id": "EV-SPEC-NVA-INSTANCE", "kind": "file", "path":
"verification/instances/nva-certified-v1.xmi" \}, \{ "id":
"EV-SPEC-NVA-VALIDATOR", "kind": "source", "path":
"neo4j-tgg/src/main/java/org/uet/dse/neo4jtgg/ocl/ir/NvaCertifiedModelValidator.java",
"symbol": "class NvaCertifiedModelValidator" \}, \{ "id":
"EV-SPEC-NVA-VALIDATOR-GUARD", "kind": "test", "path":
"neo4j-tgg/src/test/java/org/uet/dse/neo4jtgg/ocl/ir/DynamicEmfModelValidatorTest.java",
"symbol": "nvaStructuralValidityDoesNotForgeSemanticCertification" \},
\{ "id": "EV-SPEC-PLANSIM-CATALOG", "kind": "source", "path":
"neo4j-tgg/src/main/java/org/uet/dse/neo4jtgg/ocl/ir/SpecPlanSimContract.java",
"symbol": "class SpecPlanSimContract" \}, \{ "id":
"EV-SPEC-TOTAL-FEATURES", "kind": "file", "path":
"verification/coverage/ova\_cqm\_total\_feature\_refinement.csv" \}, \{
"id": "EV-SPEC-TOTAL-FEATURES-GUARD", "kind": "test", "path":
"neo4j-tgg/src/test/java/org/uet/dse/neo4jtgg/ocl/ir/MetamodelFeatureRefinementCoverageTest.java",
"symbol": "everyDeclaredOvaAndCqmFeatureHasOneExplicitPolicy" \}, \{
"id": "EV-SPEC-PGMM-EMF", "kind": "test", "path":
"neo4j-tgg/src/test/java/org/uet/dse/neo4jtgg/experiment/CanonicalPgmmEmfConformanceTest.java",
"symbol": "canonicalInstanceConformsToFullPgmmAndJavaEncoding" \}, \{
"id": "EV-PO18-CHECKER", "kind": "source", "path":
"verification/scripts/check-mechanized-proof.ps1", "symbol": "Mechanized
proof check PASS" \}, \{ "id": "EV-PO18-MUTATIONS", "kind": "source",
"path": "verification/scripts/check-mechanized-proof-mutations.ps1",
"symbol": "Mechanized proof mutation tests PASS" \}, \{ "id":
"EV-PO18-NORM-BOUNDARY", "kind": "test", "path":
"neo4j-tgg/src/test/java/org/uet/dse/neo4jtgg/ocl/ir/OclRewritePreservationTest.java",
"symbol":
"formalNormalizerAndGraphOptimizerAreSemanticNotSyntacticCounterparts"
\}, \{ "id": "EV-PO18-CAPTURE-GUARD", "kind": "test", "path":
"neo4j-tgg/src/test/java/org/uet/dse/neo4jtgg/ocl/ir/OclRewritePreservationTest.java",
"symbol": "fusionFallsBackWhenIteratorRenameWouldCaptureNestedBinder"
\}, \{ "id": "EV-PO18-JAVA-MATRIX", "kind": "file", "path":
"verification/coverage/java\_ir\_refinement\_matrix.csv" \}, \{ "id":
"EV-PO18-JAVA-GUARD", "kind": "test", "path":
"neo4j-tgg/src/test/java/org/uet/dse/neo4jtgg/experiment/OclIrJavaRefinementCoverageTest.java",
"symbol":
"everyProductionOptimizedConstructorHasLeanAndJavaRefinementEvidence"
\}, \{ "id": "EV-PO18-JAVA-REFINEMENT", "kind": "source", "path":
"neo4j-tgg/src/test/java/org/uet/dse/neo4jtgg/experiment/PipelineRefinementVerifier.java",
"symbol": "verifyOptimizedToPlan" \}, \{ "id": "EV-PO18-JAVA-WITNESSES",
"kind": "test", "path":
"neo4j-tgg/src/test/java/org/uet/dse/neo4jtgg/experiment/OclIrPlanPayloadRefinementTest.java",
"symbol":
"allSixteenOptimizedConstructorsHaveExecutableFieldPreservingWitnesses"
\}, \{ "id": "EV-PO18-METAMODEL-REFINEMENT", "kind": "source", "path":
"neo4j-tgg/src/main/java/org/uet/dse/neo4jtgg/ocl/ir/OclMetamodelRefinement.java",
"symbol": "class OclMetamodelRefinement" \}, \{ "id":
"EV-PO18-OVA-REFINEMENT-MATRIX", "kind": "file", "path":
"verification/coverage/ova\_metamodel\_java\_refinement.csv" \}, \{
"id": "EV-PO18-CQM-REFINEMENT-MATRIX", "kind": "file", "path":
"verification/coverage/cqm\_metamodel\_java\_refinement.csv" \}, \{
"id": "EV-PO18-METAMODEL-REFINEMENT-GUARD", "kind": "test", "path":
"neo4j-tgg/src/test/java/org/uet/dse/neo4jtgg/ocl/ir/OclMetamodelRefinementContractTest.java",
"symbol": "class OclMetamodelRefinementContractTest" \}, \{ "id":
"EV-PO18-FIELD-REFINEMENT", "kind": "file", "path":
"verification/coverage/ova\_cqm\_field\_refinement.csv" \}, \{ "id":
"EV-PO18-FIELD-REFINEMENT-GUARD", "kind": "test", "path":
"neo4j-tgg/src/test/java/org/uet/dse/neo4jtgg/ocl/ir/OclFieldRefinementContractTest.java",
"symbol": "class OclFieldRefinementContractTest" \}, \{ "id":
"EV-PO18-AST-LOWERING-CATALOG", "kind": "file", "path":
"verification/coverage/cqm\_ast\_lowering\_rules.csv" \}, \{ "id":
"EV-PO18-AST-LOWERING-GUARD", "kind": "test", "path":
"neo4j-tgg/src/test/java/org/uet/dse/neo4jtgg/experiment/CqmAstLoweringContractTest.java",
"symbol": "class CqmAstLoweringContractTest" \}, \{ "id":
"EV-PO18-CQM-AST-SHAPE", "kind": "test", "path":
"neo4j-tgg/src/test/java/org/uet/dse/neo4jtgg/experiment/CqmCypherAstStructuralAgreementTest.java",
"symbol": "class CqmCypherAstStructuralAgreementTest" \}, \{ "id":
"EV-PO18-METAMODEL-ECORE", "kind": "source", "path":
"verification/scripts/Generate-EcoreFromEmfatic.ps1", "symbol":
"Generate-EcoreFromEmfatic.ps1" \}, \{ "id": "EV-PO18-OVA-ECORE",
"kind": "file", "path": "md/research/model/OCL-Validation-Algebra.ecore"
\}, \{ "id": "EV-PO18-CQM-ECORE", "kind": "file", "path":
"md/research/model/Cypher-Query-Model.ecore" \}, \{ "id":
"EV-PO18-METAMODEL-INSTANCES", "kind": "file", "path":
"verification/instances/ova-certified-v1.xmi" \}, \{ "id":
"EV-PO18-CQM-INSTANCE", "kind": "file", "path":
"verification/instances/cqm-certified-v1.xmi" \}, \{ "id":
"EV-PO18-METAMODEL-INSTANCES-GUARD", "kind": "test", "path":
"neo4j-tgg/src/test/java/org/uet/dse/neo4jtgg/experiment/MetamodelInstanceXmlContractTest.java",
"symbol": "class MetamodelInstanceXmlContractTest" \}, \{ "id":
"EV-PO18-PGMM-CONFORMANCE", "kind": "source", "path":
"neo4j-tgg/src/main/java/org/uet/dse/neo4jtgg/ocl/ir/CanonicalPgmmConformance.java",
"symbol": "class CanonicalPgmmConformance" \}, \{ "id":
"EV-PO18-PGMM-CONFORMANCE-GUARD", "kind": "test", "path":
"neo4j-tgg/src/test/java/org/uet/dse/neo4jtgg/experiment/CanonicalPgmmConformanceTest.java",
"symbol": "class CanonicalPgmmConformanceTest" \}, \{ "id":
"EV-PO18-PGMM-INSTANCE", "kind": "file", "path":
"verification/pgmm/PGMM-canonical-v1.tsv" \}, \{ "id":
"EV-PO18-RUNTIME-CONSTRUCTORS", "kind": "test", "path":
"neo4j-tgg/src/test/java/org/uet/dse/neo4jtgg/experiment/OclValRealNeo4jCoverageTest.java",
"symbol": "allAdmittedConstructsAgreeWithUseOnRealNeo4j" \}, \{ "id":
"EV-RUNTIME-RECAPTURE-BLOCKED-2026-08-21", "kind": "file", "path":
"verification/evidence/runtime-recapture-blocked-2026-08-21.md" \}, \{
"id": "EV-RUNTIME-RECAPTURE-2026-08-22", "kind": "file", "path":
"verification/evidence/runtime-recapture-2026-08-22.tsv" \}, \{ "id":
"EV-PO21-LEAN", "kind": "source", "path":
"verification/lean/Ocl2CypherProof.lean", "symbol": "theorem pa\_comp"
\}, \{ "id": "EV-PO21-CERTIFICATE", "kind": "source", "path":
"neo4j-tgg/src/main/java/org/uet/dse/neo4jtgg/experiment/AdapterAdequacyCertificate.java",
"symbol": "record AdapterAdequacyCertificate" \}, \{ "id":
"EV-PO21-SNAPSHOT-READER", "kind": "source", "path":
"neo4j-tgg/src/main/java/org/uet/dse/neo4jtgg/experiment/AdapterAdequacySnapshotReader.java",
"symbol": "class AdapterAdequacySnapshotReader" \}, \{ "id":
"EV-PO21-SERVICE-GATE", "kind": "source", "path":
"neo4j-tgg/src/main/java/org/uet/dse/neo4jtgg/service/impl/DefaultOclValidationService.java",
"symbol": "AdapterAdequacyCertificate.issue" \}, \{ "id":
"EV-PO21-RESEARCH-GATE", "kind": "source", "path":
"neo4j-tgg/src/main/java/org/uet/dse/neo4jtgg/ui/ResearchToolDialog.java",
"symbol": "AdapterAdequacyCertificate.issue" \}, \{ "id":
"EV-PO21-MATRIX", "kind": "file", "path":
"verification/coverage/adapter\_adequacy\_matrix.csv" \}, \{ "id":
"EV-PO21-MATRIX-GUARD", "kind": "test", "path":
"neo4j-tgg/src/test/java/org/uet/dse/neo4jtgg/experiment/AdapterAdequacyEvidenceMatrixTest.java",
"symbol": "everyPaCompPremiseHasLeanJavaObservationAndEvidence" \}, \{
"id": "EV-PO21-MUTATIONS", "kind": "test", "path":
"neo4j-tgg/src/test/java/org/uet/dse/neo4jtgg/experiment/AdapterAdequacyCertificateTest.java",
"symbol": "rejectsEveryObservationMutationAndSnapshotMixing" \}, \{
"id": "EV-PO21-LABEL", "kind": "test", "path":
"neo4j-tgg/src/test/java/org/uet/dse/neo4jtgg/experiment/AdapterAdequacyCertificateTest.java",
"symbol":
"graphSnapshotUsesRendererLabelsAndScansCanonicalKeysAcrossWrongModelTags"
\}, \{ "id": "EV-PO21-RUNTIME", "kind": "test", "path":
"neo4j-tgg/src/test/java/org/uet/dse/neo4jtgg/experiment/AdapterAdequacyCertificateRealNeo4jTest.java",
"symbol": "certificateAndQueriesObserveOneNeo4jSnapshot" \}, \{ "id":
"EV-PO19-ARTIFACT-GATE", "kind": "source", "path":
".github/workflows/maven.yml", "symbol":
"proof-contract-PC-2026-07-22.3" \}, \{ "id": "EV-PO19-REPORT", "kind":
"source", "path": "verification/scripts/write-proof-report.ps1",
"symbol": "proof-report.schema.v1" \}, \{ "id":
"EV-PO22-TYPE-REGRESSION", "kind": "test", "path":
"neo4j-tgg/src/test/java/org/uet/dse/neo4jtgg/experiment/OclValFragmentCoverageTest.java",
"symbol": "nativeUseAndBoundTreeAgreeOnToOneTypeBeforeCollectionView"
\}, \{ "id": "EV-PO22-ADMISSION", "kind": "test", "path":
"neo4j-tgg/src/test/java/org/uet/dse/neo4jtgg/experiment/OclValNegativeAdmissionCoverageTest.java",
"symbol":
"toOneNavigationKeepsUseScalarTypeAndIsNotSilentlyCertifiedAsASet" \},
\{ "id": "EV-PO22-LOWERING", "kind": "test", "path":
"neo4j-tgg/src/test/java/org/uet/dse/neo4jtgg/experiment/OclValFragmentCoverageTest.java",
"symbol": "toOneCollectionViewSurvivesBoundIrPlanAndDirectTextRendering"
\}, \{ "id": "EV-PO22-LIFT1", "kind": "source", "path":
"verification/lean/Ocl2CypherProof.lean", "symbol":
"ToOneLift.lift1\_consumer\_agreement" \}, \{ "id": "EV-PO22-RUNTIME",
"kind": "test", "path":
"neo4j-tgg/src/test/java/org/uet/dse/neo4jtgg/experiment/OclValSemanticDiscriminatorRealNeo4jTest.java",
"symbol": "nullAndDuplicateCollectFollowTheCanonicalProfile" \}, \{
"id": "EV-PO23-VOID", "kind": "test", "path":
"neo4j-tgg/src/test/java/org/uet/dse/neo4jtgg/experiment/OclValFragmentCoverageTest.java",
"symbol":
"nullLiteralHasInternalVoidTypeButInvariantRemainsBooleanInUseAndBoundTrees"
\}, \{ "id": "EV-PO23-MATRIX", "kind": "file", "path":
"verification/coverage/void\_context\_matrix.csv" \}, \{ "id":
"EV-PO23-MATRIX-GUARD", "kind": "test", "path":
"neo4j-tgg/src/test/java/org/uet/dse/neo4jtgg/experiment/OclVoidContextMatrixTest.java",
"symbol": "everyVoidContextAgreesWithTheDeclaredUseAndCertifiedBoundary"
\}, \{ "id": "EV-PO23-RUNTIME", "kind": "test", "path":
"neo4j-tgg/src/test/java/org/uet/dse/neo4jtgg/experiment/OclValSemanticDiscriminatorRealNeo4jTest.java",
"symbol": "nullAndDuplicateCollectFollowTheCanonicalProfile" \}, \{
"id": "EV-PO24-REFINEMENT", "kind": "source", "path":
"neo4j-tgg/src/test/java/org/uet/dse/neo4jtgg/experiment/PipelineRefinementVerifier.java",
"symbol": "verify" \}, \{ "id": "EV-PO24-PAYLOAD", "kind": "test",
"path":
"neo4j-tgg/src/test/java/org/uet/dse/neo4jtgg/experiment/OclIrPlanPayloadRefinementTest.java",
"symbol": "payloadAndLoweringMutationsAreRejected" \}, \{ "id":
"EV-PO24-MATRIX", "kind": "file", "path":
"verification/coverage/bound\_va\_refinement\_matrix.csv" \}, \{ "id":
"EV-PO24-MATRIX-GUARD", "kind": "test", "path":
"neo4j-tgg/src/test/java/org/uet/dse/neo4jtgg/experiment/BoundVaProductionRefinementCoverageTest.java",
"symbol": "everyProductionBoundConstructorHasExactVaAbstractionEvidence"
\}, \{ "id": "EV-PO24-LEAN", "kind": "source", "path":
"verification/lean/Ocl2CypherProof.lean", "symbol":
"BoundVaAbstraction.bound\_va\_abstraction" \} {]}, "proofObligations":
{[} \{ "id": "PO-01", "title": "Vocabulary/schema agreement",
"classification": "required", "status": "discharged", "scope": "the
canonical PGMM and implementation use UmlClass plus joint
modelKey/classKey scope across the graph contract, production
invariant/allInstances/type/cast accessors, formal rules, and
generated-query contract", "evidenceIds": {[}"EV-PO01-VOCABULARY",
"EV-PO01-PGMM"{]} \}, \{ "id": "PO-02", "title": "Certified negative
admission", "classification": "required", "status": "discharged",
"scope": "the normative OVA/CQM metamodels expose only the finite-set
certified fragment; executable WF\_OVA, CertifiedOVA, WF\_CQM, and
CertifiedCQM predicates reject non-Boolean roots, invalid
scope/type/stage metadata, general-profile constructors, and any
violation policy other than NOT\_VALIDATION\_TRUE", "evidenceIds":
{[}"EV-PO02-CERTIFIED", "EV-PO02-OVA-METAMODEL",
"EV-PO02-CQM-METAMODEL", "EV-PO02-CERTIFIED-VALIDATOR",
"EV-PO02-CERTIFIED-VALIDATOR-GUARD"{]} \}, \{ "id": "PO-03", "title":
"Key agreement/injectivity", "classification": "required", "status":
"discharged", "evidenceIds": {[}"EV-PO03-KEYS"{]} \}, \{ "id": "PO-04",
"title": "Object/id exactness", "classification": "required", "status":
"discharged", "evidenceIds": {[}"EV-PO04-08-REPRESENTATION"{]} \}, \{
"id": "PO-05", "title": "Type exactness/inheritance", "classification":
"required", "status": "discharged", "evidenceIds":
{[}"EV-PO04-08-REPRESENTATION"{]} \}, \{ "id": "PO-06", "title": "Scalar
attribute exactness", "classification": "required", "status":
"discharged", "evidenceIds": {[}"EV-PO04-08-REPRESENTATION"{]} \}, \{
"id": "PO-07", "title": "Link/qualifier exactness", "classification":
"required", "status": "discharged", "evidenceIds":
{[}"EV-PO04-08-REPRESENTATION"{]} \}, \{ "id": "PO-08", "title":
"allInstances agreement", "classification": "required", "status":
"discharged", "evidenceIds": {[}"EV-PO04-08-REPRESENTATION"{]} \}, \{
"id": "PO-09", "title": "Renderer/accessor agreement", "classification":
"required", "status": "discharged", "evidenceIds":
{[}"EV-PO09-RENDERER"{]} \}, \{ "id": "PO-10", "title": "BottomSeparated
enforcement", "classification": "required", "status": "discharged",
"evidenceIds": {[}"EV-PO10-BOTTOM"{]} \}, \{ "id": "PO-11", "title":
"ScalarClosed enforcement", "classification": "required", "status":
"discharged", "evidenceIds": {[}"EV-PO11-SCALAR"{]} \}, \{ "id":
"PO-12", "title": "Bottom-safe alias preservation", "classification":
"required", "status": "discharged", "evidenceIds": {[}"EV-PO12-ALIAS"{]}
\}, \{ "id": "PO-13", "title": "Raw AST closure/totality",
"classification": "required", "status": "discharged", "evidenceIds":
{[}"EV-PO13-RAW-AST"{]} \}, \{ "id": "PO-14", "title": "Admitted
direct-text parse/normalization bridge", "classification": "required",
"status": "discharged", "scope": "all 17 sealed plan constructors and 52
checked queries agree with the reviewed canonical full-tree manifest
after UmlClass normalization; both the independent parser and pinned
Neo4j 2026.06.0 Cypher 5 parser accept/project all 52, with bottom-safe
and wrong-label mutations checked", "evidenceIds":
{[}"EV-PO14-DIRECT-TEXT-BRIDGE", "EV-PO14-BOTTOM-ROUNDTRIP",
"EV-PO14-CONSTRUCTOR-MATRIX", "EV-PO14-CONSTRUCTOR-MATRIX-GUARD"{]} \},
\{ "id": "PO-15", "title": "CY1-\/-CY9 runtime profile",
"classification": "required", "status": "discharged", "scope": "fresh
real-Neo4j capture at revision ebc1a4da verified CY1-\/-CY9=9/9 on Neo4j
2026.06.0, Cypher 5, enterprise/demo, with current renderer and source
hashes", "evidenceIds": {[}"EV-PO15-RUNTIME-MANIFEST",
"EV-PO15-RUNTIME-GUARD", "EV-RUNTIME-RECAPTURE-2026-08-22"{]} \}, \{
"id": "PO-16", "title": "47-case differential equivalence",
"classification": "required", "status": "discharged", "scope": "fresh
real-Neo4j differential capture at revision ebc1a4da agrees with USE for
all 47 admitted OCL\_val cases, including 19 profile-tautology and 28
non-vacuous mixed cases", "evidenceIds": {[}"EV-PO16-OCL47-MANIFEST",
"EV-PO16-OCL47-GUARD", "EV-RUNTIME-RECAPTURE-2026-08-22"{]} \}, \{ "id":
"PO-17", "title": "Semantic proof synchronization", "classification":
"required", "status": "discharged", "evidenceIds": {[}"EV-PO17-CHECKER",
"EV-PO17-MUTATIONS"{]} \}, \{ "id": "PO-18", "title": "Mechanized core
lemmas and optional Java optimization refinement", "classification":
"recommended", "status": "partial", "scope": "the normative NVA grammar
is enumerated and spec\_plan\_sim\_sound composes every recursive
NVA/CQM case under explicit local AlgebraAgreement; executable NVA
Ecore/XMI/WF\_NVA, total OVA/CQM feature coverage, SpecPlanSim catalog,
and full PGMM EMF conformance are checked. The partial remainder is only
the optional non-normative Java OptRefines(T\_OPT(va),T\_NORM(va)) and
universal Java/Neo4j instantiation; it is not a premise of Theorems
0-\/-6", "evidenceIds": {[}"EV-PO18-LEAN", "EV-PO18-PRODPLAN-INDUCTION",
"EV-SPEC-NVA-GRAMMAR", "EV-SPEC-PLANSIM", "EV-SPEC-NVA-METAMODEL",
"EV-SPEC-NVA-ECORE", "EV-SPEC-NVA-INSTANCE", "EV-SPEC-NVA-VALIDATOR",
"EV-SPEC-NVA-VALIDATOR-GUARD", "EV-SPEC-PLANSIM-CATALOG",
"EV-SPEC-TOTAL-FEATURES", "EV-SPEC-TOTAL-FEATURES-GUARD",
"EV-SPEC-PGMM-EMF", "EV-PO18-CHECKER", "EV-PO18-MUTATIONS",
"EV-PO18-NORM-BOUNDARY", "EV-PO18-CAPTURE-GUARD", "EV-PO18-JAVA-MATRIX",
"EV-PO18-JAVA-GUARD", "EV-PO18-JAVA-REFINEMENT",
"EV-PO18-JAVA-WITNESSES", "EV-PO18-METAMODEL-REFINEMENT",
"EV-PO18-OVA-REFINEMENT-MATRIX", "EV-PO18-CQM-REFINEMENT-MATRIX",
"EV-PO18-METAMODEL-REFINEMENT-GUARD", "EV-PO18-FIELD-REFINEMENT",
"EV-PO18-FIELD-REFINEMENT-GUARD", "EV-PO18-AST-LOWERING-CATALOG",
"EV-PO18-AST-LOWERING-GUARD", "EV-PO18-CQM-AST-SHAPE",
"EV-PO18-METAMODEL-ECORE", "EV-PO18-OVA-ECORE", "EV-PO18-CQM-ECORE",
"EV-PO18-METAMODEL-INSTANCES", "EV-PO18-CQM-INSTANCE",
"EV-PO18-METAMODEL-INSTANCES-GUARD", "EV-PO18-PGMM-CONFORMANCE",
"EV-PO18-PGMM-CONFORMANCE-GUARD", "EV-PO18-PGMM-INSTANCE",
"EV-PO18-RUNTIME-CONSTRUCTORS"{]} \}, \{ "id": "PO-19", "title": "Clean
machine-verification artifact", "classification": "required", "status":
"discharged", "scope": "discharged 2026-08-09 for CI/artifact
provenance: GitHub Actions completed successfully for clean commit
7b7bdd2538434d64e2e77c9321b93b12feee1a42 and uploaded
proof-contract-PC-2026-07-22.3 plus the build package; the recorded
clean aggregate baseline is rooted at
f09385197460d3a9b91fc878354583062bacdfb1 with contract/mutations, Lean
32/32, Lean mutations 6/6, 311/311 selected conformance tests, and zero
failures/errors/skips; this discharge does not expand the conditional
semantic scope", "evidenceIds": {[}"EV-PO19-ARTIFACT-GATE",
"EV-PO19-REPORT"{]} \}, \{ "id": "PO-20", "title": "Paper publication
build (local only)", "classification": "out\_of\_scope", "status":
"open", "scope": "paper sources remain local and ignored under md/ by
explicit repository policy; GitHub Actions neither compiles nor uploads
the paper, and publication evidence is excluded from the
machine-correctness contract", "evidenceIds": {[}{]} \}, \{ "id":
"PO-21", "title": "AdapterAdequate composition from PA1-\/-PA9",
"classification": "required", "status": "discharged", "scope":
"PA1-\/-PA9 static composition and fresh real-Neo4j certificate agree on
one shared snapshot: observations=6, plans=2, and wrong-label mutation
is killed under canonical joint model scope", "evidenceIds":
{[}"EV-PO21-LEAN", "EV-PO21-CERTIFICATE", "EV-PO21-SNAPSHOT-READER",
"EV-PO21-SERVICE-GATE", "EV-PO21-RESEARCH-GATE", "EV-PO21-MATRIX",
"EV-PO21-MATRIX-GUARD", "EV-PO21-MUTATIONS", "EV-PO21-LABEL",
"EV-PO21-RUNTIME", "EV-CANONICAL-RUNTIME-SUPPLEMENT",
"EV-CANONICAL-RUNTIME-SUPPLEMENT-GUARD",
"EV-RUNTIME-RECAPTURE-2026-08-22"{]} \}, \{ "id": "PO-22", "title":
"Multiplicity-preserving to-one navigation lift", "classification":
"required", "status": "discharged", "scope": "LIFT1,
sourceCollectionType propagation, admission, static regressions, and
fresh real-Neo4j discriminator capture agree for the three to-one
collection-view observations", "evidenceIds":
{[}"EV-PO22-TYPE-REGRESSION", "EV-PO22-ADMISSION", "EV-PO22-LOWERING",
"EV-PO22-LIFT1", "EV-PO22-RUNTIME", "EV-CANONICAL-RUNTIME-SUPPLEMENT",
"EV-CANONICAL-RUNTIME-SUPPLEMENT-GUARD",
"EV-RUNTIME-RECAPTURE-2026-08-22"{]} \}, \{ "id": "PO-23", "title":
"Internal Void/null typing and semantic boundary", "classification":
"required", "status": "discharged", "scope": "formal Void rules, 22-row
static matrix, typed v1\textbar V bottom representation, and fresh
real-Neo4j capture agree for void contexts=13 and duplicate collect=1",
"evidenceIds": {[}"EV-PO23-VOID", "EV-PO23-MATRIX",
"EV-PO23-MATRIX-GUARD", "EV-PO23-RUNTIME",
"EV-CANONICAL-RUNTIME-SUPPLEMENT",
"EV-CANONICAL-RUNTIME-SUPPLEMENT-GUARD",
"EV-RUNTIME-RECAPTURE-2026-08-22"{]} \}, \{ "id": "PO-24", "title":
"Production Bound/VA abstraction agreement", "classification":
"required", "status": "discharged", "scope": "discharged 2026-08-09 for
the production abstraction boundary: an 11-row reflection matrix covers
every BoundExpression record and all 12 SemanticExpression targets,
including the disjoint BoundProperty attribute/navigation cases and
every record component; all targets are witnessed across the 47-case
corpus; the axiom-free bound\_va\_abstraction theorem proves recursive
evaluation equality under every shared payload-parametric primitive
algebra and explicitly preserves sourceCollectionType; this does not
instantiate the primitives with JVM or Neo4j semantics, which remains
PO-18", "evidenceIds": {[}"EV-PO24-REFINEMENT", "EV-PO24-PAYLOAD",
"EV-PO24-MATRIX", "EV-PO24-MATRIX-GUARD", "EV-PO24-LEAN"{]} \} {]},
"runtimeProfiles": {[} \{ "id": "RP-NEO4J-2026.06-CYPHER5-CANONICAL-V1",
"serverVersion": "2026.06.0", "edition": "enterprise", "cypherVersion":
"5", "database": "demo", "encodingProfile": "canonical-v1",
"parserArtifact": "org.neo4j:cypher-parser-factory:2026.06.0",
"evidenceIds": {[}"EV-PO15-RUNTIME-MANIFEST", "EV-PO16-OCL47-MANIFEST",
"EV-CANONICAL-RUNTIME-SUPPLEMENT"{]} \} {]}, "constructorCoverage": \{
"matrixPath": "verification/coverage/oclval\_coverage\_matrix.csv",
"admittedCaseSourcePath":
"neo4j-tgg/src/test/java/org/uet/dse/neo4jtgg/experiment/OclValFragmentCoverageTest.java",
"features": {[} "self", "lexical variable", "Boolean literal", "Integer
literal", "Real literal", "String literal", "attribute", "navigation
forward", "navigation reverse", "qualified navigation", "allInstances",
"not", "and", "or", "implies", "equals", "not equals", "less than",
"less or equal", "greater than", "greater or equal", "addition",
"subtraction", "multiplication", "division", "if", "let", "exists",
"forAll", "select", "reject", "collect scalar", "includes", "excludes",
"includesAll", "excludesAll", "size", "isEmpty", "notEmpty",
"oclIsKindOf", "oclAsType", "xor", "Set literal", "union",
"intersection", "asSet", "isUnique" {]}, "requiredColumns": \{ "parse":
"OclDocumentParser", "bind": "OclSemanticBinder", "typing":
"BoundTypeRules", "va": "OclIrBuilder", "normalize": "OclIrOptimizer",
"cypher": "OclCypherPlannerRenderer", "proof\_case": "T1-T5",
"positive\_test": "OclValFragmentCoverageTest", "negative\_test":
"OclValNegativeAdmissionCoverageTest", "real\_neo4j":
"OclValRealNeo4jCoverageTest" \}, "artifactPaths": {[}
"neo4j-tgg/src/main/java/org/uet/dse/neo4jtgg/service/impl/OclDocumentParser.java",
"neo4j-tgg/src/main/java/org/uet/dse/neo4jtgg/ocl/OclSemanticBinder.java",
"neo4j-tgg/src/main/java/org/uet/dse/neo4jtgg/ocl/ir/OclIrBuilder.java",
"neo4j-tgg/src/main/java/org/uet/dse/neo4jtgg/ocl/ir/OclIrOptimizer.java",
"neo4j-tgg/src/main/java/org/uet/dse/neo4jtgg/ocl/ir/OclCypherPlanner.java",
"neo4j-tgg/src/main/java/org/uet/dse/neo4jtgg/ocl/ir/OclCypherRenderer.java",
"neo4j-tgg/src/test/java/org/uet/dse/neo4jtgg/experiment/OclValNegativeAdmissionCoverageTest.java",
"neo4j-tgg/src/test/java/org/uet/dse/neo4jtgg/experiment/OclValRealNeo4jCoverageTest.java"
{]} \}, "planConstructorCoverage": \{ "matrixPath":
"verification/coverage/cypher\_plan\_constructor\_matrix.csv",
"constructors": {[} "VariablePlan", "LiteralPlan", "SetLiteralPlan",
"NotPlan", "BinaryPlan", "AttributeAccessPlan", "NavigationAccessPlan",
"MethodCallPlan", "CollectionOperationPlan", "IfPlan", "LetPlan",
"IteratorOperationPlan", "ExistsSubqueryPlan", "NotExistsSubqueryPlan",
"CountSubqueryComparisonPlan", "NavigationAggregationPlan",
"NavigationUniquenessPlan" {]}, "requiredColumns": {[} "constructor",
"ir\_constructor", "factory\_method", "planner\_method",
"renderer\_method", "formal\_rule", "evidence\_test", "evidence\_method"
{]} \}, "javaIrRefinementCoverage": \{ "matrixPath":
"verification/coverage/java\_ir\_refinement\_matrix.csv",
"constructors": {[} "Variable", "Literal", "SetLiteral", "Not",
"Binary", "AttributeAccess", "NavigationAccess", "MethodCall",
"CollectionOperation", "IteratorOperation", "If", "Let",
"NavigationPredicateCheck", "NavigationCountComparison",
"NavigationAggregation", "NavigationUniquenessCheck" {]},
"requiredColumns": {[} "java\_constructor", "java\_payload\_fields",
"plan\_constructors", "lean\_constructor", "lean\_agreement\_case",
"planner\_method", "java\_refinement\_method", "evidence\_test",
"evidence\_method", "witness\_test", "witness\_method" {]} \},
"boundVaRefinementCoverage": \{ "matrixPath":
"verification/coverage/bound\_va\_refinement\_matrix.csv",
"boundConstructors": {[} "BoundVariable", "BoundLiteral",
"BoundSetLiteral", "BoundNot", "BoundIf", "BoundLet", "BoundBinary",
"BoundProperty", "BoundMethodCall", "BoundCollectionOperation",
"BoundIterator" {]}, "vaConstructors": {[} "Variable", "Literal",
"SetLiteral", "Not", "If", "Let", "Binary", "AttributeAccess",
"NavigationAccess", "MethodCall", "CollectionOperation",
"IteratorOperation" {]}, "requiredColumns": {[} "bound\_constructor",
"bound\_payload\_fields", "va\_constructors", "lean\_constructors",
"builder\_method", "java\_refinement\_method", "evidence\_test",
"evidence\_method" {]} \}, "implementationVocabulary": {[} \{ "term":
"ObjectInstanceOf", "role": "M1 object to runtime UmlClass typing link",
"path":
"neo4j/src/main/java/org/uet/dse/neo4j/encoding/CanonicalGraphVocabulary.java",
"literal": "public static final String OBJECT\_INSTANCE\_OF =
"ObjectInstanceOf";" \}, \{ "term": "InstanceOf", "role": "M2 schema
element to MetaNode typing link", "path":
"neo4j/src/main/java/org/uet/dse/neo4j/encoding/CanonicalGraphVocabulary.java",
"literal": "public static final String SCHEMA\_INSTANCE\_OF =
"InstanceOf";" \}, \{ "term": "modelKey", "role": "repository model
scope conjoined with every validation-visible canonical lookup", "path":
"neo4j-tgg/src/main/java/org/uet/dse/neo4jtgg/ocl/ir/OclCypherRenderer.java",
"literal": "\{modelKey: \$" \}, \{ "term": "classKey", "role": "exact
UmlClass identity used jointly with modelKey by invariant
context/allInstances", "path":
"neo4j-tgg/src/main/java/org/uet/dse/neo4jtgg/ocl/ir/OclCypherRenderer.java",
"literal": ""UmlClass", "classKey", classParam, state)" \}, \{ "term":
"attributeKey", "role": "exact scalar attribute-slot identity", "path":
"neo4j-tgg/src/main/java/org/uet/dse/neo4jtgg/ocl/ir/OclCypherRenderer.java",
"literal": ".attributeKey = \$" \}, \{ "term": "associationKey", "role":
"exact association identity for navigation", "path":
"neo4j-tgg/src/main/java/org/uet/dse/neo4jtgg/ocl/ir/OclCypherRenderer.java",
"literal": ".associationKey = \$" \}, \{ "term": "sourceQualifiers",
"role": "qualifier payload at the source association end", "path":
"neo4j-tgg/src/main/java/org/uet/dse/neo4jtgg/ocl/ir/OclCypherRenderer.java",
"literal": "? "sourceQualifiers" : binding.sourceQualifierProperty();"
\}, \{ "term": "targetQualifiers", "role": "qualifier payload at the
target association end", "path":
"neo4j-tgg/src/main/java/org/uet/dse/neo4jtgg/ocl/ir/OclCypherRenderer.java",
"literal": "case INCOMING -\textgreater{} binding == null ?
"targetQualifiers" : binding.targetQualifierProperty();" \}, \{ "term":
"scalarCodecV1", "role": "typed scalar wire format separating bottom
from ordinary strings and numeric domains", "path":
"neo4j/src/main/java/org/uet/dse/neo4j/sync/helper/CanonicalScalarValueCodec.java",
"literal": "public static final String BOTTOM = "v1\textbar V";" \}, \{
"term": "\_\_oclBottom", "role": "reserved semantic bottom marker",
"path":
"neo4j-tgg/src/main/java/org/uet/dse/neo4jtgg/ocl/OclBottomToken.java",
"literal": "public static final String MARKER = "\_\_oclBottom";" \}
{]}, "mechanization": \{ "framework": "Lean", "version": "4.32.2",
"toolchain": "leanprover/lean4:v4.32.2", "releaseAsset":
"lean-4.32.2-windows.zip", "releaseAssetBytes": 832110051,
"releaseAssetSha256":
"369c2b480a2a6f8bfb727af42c333c894c4872a73b3503099abad7bef67549fa",
"sourcePath": "verification/lean/Ocl2CypherProof.lean", "checkerPath":
"verification/scripts/check-mechanized-proof.ps1",
"mutationCheckerPath":
"verification/scripts/check-mechanized-proof-mutations.ps1",
"projectPaths": {[} "verification/lean/lakefile.lean",
"verification/lean/lean-toolchain",
"verification/lean/lake-manifest.json" {]}, "externalChecker": \{
"kind": "nanoda-ndjson", "sourcePath":
"verification/scripts/check-nanoda.sh", "configPath":
"verification/contract/nanoda-config.json", "workflowPath":
".github/workflows/maven.yml", "lean4exportCommit":
"9fb131bb100eb32ccf6836f14e4f8328d13b6792", "nanodaCommit":
"418320295890faed83a96fd97907b12a3b6728c2", "upstreamIssue":
"\url{https://github.com/leanprover/lean-action/issues/169}" \},
"requiredTheorems": {[} "image\_reflects\_membership",
"image\_preserves\_subset", "exists\_over\_image",
"forall\_over\_image", "lift1\_source\_bound",
"lift1\_bound\_validation", "lift1\_present", "lift1\_absent",
"lift1\_consumer\_agreement", "encodeValue\_injective",
"implies\_rewrite", "forall\_rewrite", "notEmpty\_rewrite",
"normalize\_preserves\_eval", "normalize\_reaches\_redex\_free",
"implies\_root\_strictly\_decreases", "all\_rewrite\_semantics",
"normalize\_reaches\_normal\_form", "normalize\_idempotent",
"root\_rewrite\_strictly\_decreases", "typed\_rewrite\_preserves\_type",
"scoped\_rename\_preserves\_binder\_boundary",
"named\_to\_scoped\_semantic\_correspondence",
"java\_capture\_guard\_sound",
"java\_guarded\_rename\_preserves\_scoped\_semantics",
"structural\_preservation", "java\_ir\_eval\_refinement",
"prod\_plan\_sim\_sound", "certified\_nva\_grammar\_complete",
"spec\_plan\_sim\_sound", "bound\_va\_abstraction", "pa\_comp",
"theorem6\_forward", "theorem6\_backward", "theorem6\_at\_object" {]},
"axiomAudit": \{ "allowedCoreAxioms": {[}"propext", "Quot.sound"{]},
"theorems": {[} \{ "name": "Ocl2CypherProof.encodeValue\_injective",
"expected": {[}"propext"{]} \}, \{ "name":
"Ocl2CypherProof.ToOneLift.lift1\_consumer\_agreement", "expected":
{[}{]} \}, \{ "name":
"Ocl2CypherProof.BoolExpr.normalize\_preserves\_eval", "expected":
{[}"propext"{]} \}, \{ "name":
"Ocl2CypherProof.BoolExpr.normalize\_reaches\_redex\_free", "expected":
{[}"propext"{]} \}, \{ "name":
"Ocl2CypherProof.Formula.structural\_preservation", "expected":
{[}"propext"{]} \}, \{ "name":
"Ocl2CypherProof.JavaIrRefinement.java\_ir\_eval\_refinement",
"expected": {[}"propext"{]} \}, \{ "name":
"Ocl2CypherProof.JavaIrRefinement.prod\_plan\_sim\_sound", "expected":
{[}"propext"{]} \}, \{ "name":
"Ocl2CypherProof.SpecificationPlanRefinement.certified\_nva\_grammar\_complete",
"expected": {[}"propext"{]} \}, \{ "name":
"Ocl2CypherProof.SpecificationPlanRefinement.spec\_plan\_sim\_sound",
"expected": {[}{]} \}, \{ "name":
"Ocl2CypherProof.BoundVaAbstraction.bound\_va\_abstraction", "expected":
{[}{]} \}, \{ "name": "Ocl2CypherProof.AdapterComposition.pa\_comp",
"expected": {[}{]} \}, \{ "name":
"Ocl2CypherProof.theorem6\_at\_object", "expected": {[}{]} \}, \{
"name": "Ocl2CypherProof.Normalization.all\_rewrite\_semantics",
"expected": {[}{]} \}, \{ "name":
"Ocl2CypherProof.Normalization.typed\_rewrite\_preserves\_type",
"expected": {[}"propext"{]} \}, \{ "name":
"Ocl2CypherProof.Normalization.Scoped.scoped\_rename\_preserves\_binder\_boundary",
"expected": {[}"propext"{]} \}, \{ "name":
"Ocl2CypherProof.Normalization.NamedBridge.named\_to\_scoped\_semantic\_correspondence",
"expected": {[}"propext", "Quot.sound"{]} \}, \{ "name":
"Ocl2CypherProof.Normalization.NamedBridge.java\_capture\_guard\_sound",
"expected": {[}"propext"{]} \}, \{ "name":
"Ocl2CypherProof.Normalization.NamedBridge.java\_guarded\_rename\_preserves\_scoped\_semantics",
"expected": {[}"propext", "Quot.sound"{]} \}, \{ "name":
"Ocl2CypherProof.Normalization.RewriteExpr.normalize\_reaches\_normal\_form",
"expected": {[}"propext"{]} \}, \{ "name":
"Ocl2CypherProof.Normalization.RewriteExpr.normalize\_idempotent",
"expected": {[}"propext"{]} \}, \{ "name":
"Ocl2CypherProof.Normalization.RewriteExpr.RootRule.root\_rewrite\_strictly\_decreases",
"expected": {[}"propext"{]} \} {]} \}, "coverage": {[} \{ "id":
"MK-FINITE-SET", "status": "mechanized", "scope": "predicate-set
image/subset/exists/forall and injective membership reflection" \}, \{
"id": "MK-LIFT1", "status": "mechanized", "scope": "axiom-free
absent-to-empty and present-to-singleton agreement among source, Bound,
and validation-algebra collection views, including the asSet, size,
isEmpty, and notEmpty observations admitted for scalar to-one receivers"
\}, \{ "id": "MK-ENCODE", "status": "mechanized", "scope": "flat typed
bottom/scalar/entity/finite-set value encoding" \}, \{ "id":
"MK-NORMALIZE", "status": "partial", "scope": "all eleven listed
rewrite-family semantic equations over validation Bool/finite List,
intrinsic Bool/Nat/Elem/Set N2 indexes, lexical named-to-de-Bruijn
semantic correspondence, soundness and scoped-semantic preservation of
the exact three-disjunct Java capture guard over the Boolean binder
kernel, relative normal form, idempotence, and lexicographic root
decrease; concrete semantic premises for every production optimization
rule remain open" \}, \{ "id": "MK-T4", "status": "mechanized", "scope":
"payload-parametric relational structural induction over all 16
production OclIr.OptimizedExpression constructors, including
sourceCollectionType on collection and iterator operations, every other
non-recursive record payload, recursive expression lists, and optional
predicates, under one explicit primitive-agreement premise per
constructor" \}, \{ "id": "MK-PRODPLAN", "status": "mechanized",
"scope": "named prod\_plan\_sim\_sound theorem composes the
constructor-wise evaluation relation under the same explicit local
AlgebraAgreement premise used by the formal ProdPlanSim contract" \}, \{
"id": "MK-SPECPLAN", "status": "mechanized", "scope": "axiom-free
structural induction over independent finite NVA trees whose tag ranges
over exactly the 26 certified constructors; composes SpecPlanSim under
an explicit LR/C/BR/CY local equation for every constructor tag" \}, \{
"id": "MK-BOUND-VA", "status": "mechanized", "scope": "axiom-free
recursive evaluation equality for the alpha abstraction from all eleven
production BoundExpression record families, with BoundProperty split
into the twelve SemanticExpression targets and every
sourceCollectionType payload retained, under an arbitrary shared
primitive algebra" \}, \{ "id": "MK-PA-COMP", "status": "mechanized",
"scope": "six AdapterAdequate observation equalities derived from
exact-M2 and PA1-\/-PA9 over one explicit shared source/graph snapshot
premise" \}, \{ "id": "MK-T6", "status": "mechanized", "scope":
"forward/backward pointwise violation-ID inclusions under agreement and
ID injectivity" \} {]}, "openScope": {[} "nested collection encodeValue
injectivity", "metamodel-specific OCL subtyping/coercion and concrete
semantic agreement for every Java T\_OPT normalization rule",
"instantiation of the 16-constructor relational algebra with the actual
object and graph evaluators, including Java payload semantics",
"instantiation of the 26-constructor NVA algebra with the concrete graph
and reference-Cypher evaluators; CLAIM-SPEC-PLAN remains conditional on
the registered LR/C/BR/CY local agreements", "Neo4j/Cypher semantics
beyond the selected runtime-profile assumptions and
non-invariant/fallback entry points", "universal Theorem 6 composition
beyond the abstract agreement premises" {]} \}, "artifactPolicy": \{
"kind": "machine-verification", "reportSchema":
"proof-report.schema.v1", "reportGeneratorPath":
"verification/scripts/write-proof-report.ps1", "baselineEvidencePath":
"verification/evidence/proof-baseline-2026-08-04.json",
"paperPublication": "external-local-only" \}, "claimPolicy": \{
"prototypeCorrectness": "conditional", "activeValue": "active",
"blockingClassifications": {[}"required", "recommended"{]},
"blockingStatuses": {[}"open", "partial"{]} \} \}

{\def\LTcaptype{none} % do not increment counter
\begin{longtable}[]{@{}llll@{}}
\toprule\noalign{}
Contract ID & Statement kernel & Assumptions & Required results \\
\midrule\noalign{}
\endhead
\bottomrule\noalign{}
\endlastfoot
PC-T0 & validation observations of \texttt{(MM,M)} and every adequate
\texttt{G} agree & A4, A5 & R1, R2, R3, R4, R5, R6, R7 \\
PC-T1 & decoded Source and uniquely bound denotations are related and
erase equally & A1, A2, A4, A8 & M2, M3, M3a, M4, B0, B1, B2, B3, B4,
SB0 \\
PC-T2 & erased Bound denotation equals typed object-VA denotation & A2,
A3, A4, A8 & M2, M3, M3a, VA1, VA2, VA3, BV0, B4 \\
PC-T3 & deterministic bottom-up normalization terminates, preserves
type/denotation and reachable closure, and reaches \texttt{NF\_R} modulo
alpha-equivalence & A2, A3, A8 & N0, N1, N2, N3, N4 \\
PC-T4 & typed object and graph VA denotations commute with
\texttt{encodeValue} on reachable closed evaluations & A2, A3, A4, A5,
A8 & T0, G1, G2, G3, G3a, G4 \\
PC-T5 & the reference NVA-to-CQM-to-Cypher query returns exactly
normalized graph-VA violation IDs and no ghost IDs & A2, A3, A5, A6, A7,
A8, A9 & C1, C2, C3, C3a, C4, C5, C5a, C5b, C6, SpecPlanAdequacy,
SpecPlanSim\_sound, TXT1, TXT2, TXT3, TXT4, TXT5, BR1-BR10, CY1-CY9 \\
PC-T6 & \texttt{returnedIds(q,G,pi)=id{[}Viol\_OCL(e,C,M){]}} & A1, A2,
A3, A4, A5, A6, A7, A8, A9 & T0, T1, T2, T3, T4, T5 \\
\end{longtable}
}

\subsection{1.2 Current Theory-\/-Implementation Alignment Audit
(2026-08-22)}\label{12-current-theory--implementation-alignment-audit-2026-08-22}

This subsection is the current-state index for the rest of the report.
It was checked against the production Java sources, the proof registry,
the coverage matrices, the Lean source, and
\texttt{verification/report/proof-report.json}. Older dated passages
later in the document are historical evidence and must not be read as
overriding this index.

\subsubsection{1.2.1 Current Counts and Claim
Status}\label{121-current-counts-and-claim-status}

{\def\LTcaptype{none} % do not increment counter
\begin{longtable}[]{@{}lrl@{}}
\toprule\noalign{}
Item & Current value & Interpretation \\
\midrule\noalign{}
\endhead
\bottomrule\noalign{}
\endlastfoot
Contract assumptions / scope lemmas / paper theorem contracts & 9 / 6 /
7 & A1-\/-A9, M1-\/-M5 including M3a, and PC-T0-\/-PC-T6 \\
Registered semantic functions / evidence records & 29 / 57 & Registry
inventory, not 29 fully mechanized functions \\
Proof obligations & 24 & 22 discharged (91.7\%), 1 partial (4.2\%), 1
open (4.2\%) \\
Obligation classifications & 22 required, 1 recommended, 1 out-of-scope
& Runtime-dependent rows are discharged by the 2026-08-22 real-Neo4j
capture; only the recommended universal Java/Neo4j refinement row
remains partial \\
Remaining correctness blockers & PO-18 & Universal primitive commutation
and concrete Java optimizer/planner refinement remain outside the
mechanized proof; all required finite-profile obligations are
discharged \\
Publication-only row & PO-20, \texttt{open}, out-of-scope & Local paper
compilation is intentionally outside the machine-correctness contract \\
Frozen OCL / Cypher-plan / Java-IR coverage rows & 47 / 17 / 16 &
Admitted feature rows, sealed plan constructors, and production
optimized-expression constructors \\
Bound-\/-VA / adapter / Void-context rows & 11 / 11 / 22 & Production
abstraction, PA-COMP observations, and 13 admitted plus 9 rejected Void
contexts \\
Current local Java conformance / compiler mutations & 606
\texttt{neo4j-tgg} + 29 \texttt{neo4j} discovered / 20/20 & All
non-provenance Java contracts pass; opt-in real-server tests and the
three checked-in runtime-manifest guards pass with clean-evidence
mode \\
Recorded parser trees / runtime CY rows / differential cases & 52/52 /
9/9 / 47/47 & Finite evidence for the pinned Cypher 5 runtime profile,
not universal Neo4j semantics \\
Recorded Lean theorems / Lean mutations & 33/33 / 6/6 & Selected
mechanized kernels under their explicit premises, including named
\texttt{prod\_plan\_sim\_sound} constructor composition \\
\end{longtable}
}

The source and finite-tree contracts pass locally, including 20/20
compiler mutations. The 2026-08-22 selected-runtime execution against
Neo4j 2026.06.0 passed CY1-\/-CY9, OCL\_val-47, adapter/shared-snapshot,
Void/bottom/scalar/ receiver, and the Families/Persons case study. The
three publication manifests were regenerated against revision
\texttt{ebc1a4da} and pass their clean-evidence guards. This discharges
the finite runtime profile, but does not broaden the paper claim to full
OMG OCL, arbitrary Cypher, or a universal Java optimizer proof. The
strongest publication wording therefore remains \textbf{conditional
correctness under the certified profile}.

\subsubsection{1.2.2 Audit Result and Priority
Scale}\label{122-audit-result-and-priority-scale}

The source-level audit currently tracks \textbf{34 issue families,
AI-01-\/-AI-34}. The word \emph{family} is deliberate: one row can
contain several manifestations of the same broken boundary. This is the
complete list found by the 2026-08-14 audit through 2026-08-15 of the
named production and proof artifacts; it is not a mathematical claim
that no defect can exist in an unexecuted program state.

\begin{Shaded}
\begin{Highlighting}[]
\NormalTok{P0  can admit an out{-}of{-}language term, change a certified result, make a}
\NormalTok{    runtime premise check pass vacuously, or make a certified query fail;}
\NormalTok{P1  formal/implementation scope mismatch or certificate blind spot that must}
\NormalTok{    be closed or explicitly excluded before the corresponding claim is used;}
\NormalTok{P2  structural, profile, or finite{-}evidence boundary that must be stated}
\NormalTok{    accurately but is not by itself a demonstrated wrong result.}
\end{Highlighting}
\end{Shaded}

\subsubsection{1.2.3 Original Rows Rechecked
(AI-01-\/-AI-10)}\label{123-original-rows-rechecked-ai-01--ai-10}

{\def\LTcaptype{none} % do not increment counter
\begin{longtable}[]{@{}lllll@{}}
\toprule\noalign{}
ID & Priority & Boundary & Rechecked production evidence & Current
assessment and required action \\
\midrule\noalign{}
\endhead
\bottomrule\noalign{}
\endlastfoot
AI-01 & P2 & Invariant root & A1, T1, T6, and the binder rules now
separate invariant text, parser AST, decoded source, and bound invariant
& \textbf{Closed in the proof specification.} Expression-local lemmas
use \texttt{e} only as \texttt{pred(I)} \\
AI-02 & P1 & \texttt{let} syntax and binding & \texttt{ASTLet} retains
the optional declaration type; B-Let resolves it, checks
\texttt{type(init)\ \textless{}=\ declaredType}, and stores the resolved
type through BoundOCL, IR, and plan & \textbf{Closed by implementation
and proof.} Inferred and typed \texttt{let} share one rule through
\texttt{declType} \\
AI-03 & P1 & \texttt{if} branch join & \texttt{ifJoin} is exactly the
implemented partial function: equal branch types, or one \texttt{Void}
branch & \textbf{Closed by specification narrowing.} Numeric/class LUB
branches remain outside the certified fragment \\
AI-04 & P1 & Equality admission & \texttt{eqCompatible} now means equal
types, two numeric scalar types, or either side \texttt{Void} &
\textbf{Closed by specification narrowing.} General class-upcast
equality is not certified \\
AI-05 & P0 & Operation arity and argument typing & Binder, source
admission, and bound admission enforce exact arity, receiver kind,
class-reference position, and element compatibility; 45
certified-negative cases pass & \textbf{Closed in code and static
evidence.} No admitted argument is silently erased \\
AI-06 & P0 & Common-type Set coercion & The admitted mixed numeric
domain is guarded by exact Integer-to-Real64 convertibility; optimizer
and \texttt{ScalarClosed} use exact arithmetic and tests cover values
beyond \texttt{2\^{}53} & \textbf{Closed conditionally under A8.}
Inexact mixed values fail the premise instead of entering the theorem
domain \\
AI-07 & P2 & Raw target AST & Formal \texttt{let} now materializes the
initializer once under an environment alias and formal set-valued
branches use correlated raw-AST schemes, while production may use
different \texttt{CASE}/list-comprehension shapes & \textbf{Bounded
structural bridge only.} \texttt{PlanAdequacy} is now an explicit T5
premise; finite parser/tree evidence is not its universal proof \\
AI-08 & P0 & \texttt{ScalarClosed} executable witness & Stored values
are decoded by declared UML type; malformed, legacy, and cross-tag
payloads fail; empty required numeric observations cannot pass vacuously
& \textbf{Closed for the certified profile; real-Neo4j manifest
CY1-\/-CY9 and OCL47 guards pass} \\
AI-09 & P2 & Scalar \texttt{collect} oracle & The profile uses
non-flattening finite-Set image semantics, whereas native USE
\texttt{collect} is bag/sequence-oriented and the duplicate
discriminator compares with an explicit \texttt{asSet} reference &
\textbf{Profile/oracle boundary remains.} Call it a ``profile-adjusted
source reference,'' not evaluation of an unchanged general-OMG-OCL
invariant \\
AI-10 & P1 & Bottom separation & Stored bottom is `v1 &
V\texttt{;\ legal\ strings\ use\ the\ disjoint\ }v1 \\
\end{longtable}
}

\subsubsection{1.2.4 Additional Confirmed Rows
(AI-11-\/-AI-34)}\label{124-additional-confirmed-rows-ai-11--ai-34}

{\def\LTcaptype{none} % do not increment counter
\begin{longtable}[]{@{}lllll@{}}
\toprule\noalign{}
ID & Priority & Boundary & Concrete mismatch or counterexample family &
Required action \\
\midrule\noalign{}
\endhead
\bottomrule\noalign{}
\endlastfoot
AI-11 & P1 & Typed iterator syntax & The certified binder resolves the
optional iterator type, checks
\texttt{elementType\ \textless{}=\ declaredType}, scopes the body with
the declared type, and preserves that type through BoundOCL, IR, and
plan & \textbf{Closed by implementation and proof for the
single-iterator finite-Set fragment} \\
AI-12 & P0 & Class-reference values & Class references are admitted only
as the receiver of \texttt{allInstances} or the sole type argument of a
type operation; equality, Set, \texttt{let}, and iterator uses are
certified-negative & \textbf{Closed} \\
AI-13 & P1 & Implicit collection property projection & The certified
path rejects \texttt{collection.property}; explicit \texttt{collect} is
required & \textbf{Closed by certified exclusion} \\
AI-14 & P0 & Integer constant folding & Integer folding uses exact
integral arithmetic and checked Int64 bounds; mixed Real conversion is
explicit and exact under A8 & \textbf{Closed; boundary tests pass} \\
AI-15 & P0 & Optimizer lexical scope & Shadowed bindings are removed
from substitution environments; free-variable/capture guards block
unsafe fusion; retained nested-\texttt{let} witnesses pass &
\textbf{Closed for implemented rewrites; PO-18 universal mechanization
remains} \\
AI-16 & P0 & Cypher alias freshness & Renderer maintains a
source-to-target alpha environment plus a globally fresh generated-alias
set & \textbf{Closed; collision and nested-shadowing tests pass} \\
AI-17 & P0 & Bottom-valued navigation receiver & Every navigation owner
is normalized through bottom-safe entity-source expansion before
\texttt{MATCH} & \textbf{Closed statically and by the canonical
real-Neo4j supplement} \\
AI-18 & P0 & Extensional Set navigation & Navigation targets are
deduplicated by object identity; cardinality uses
\texttt{size(COLLECT\ \{\ ...\ RETURN\ DISTINCT\ target\ \})} &
\textbf{Closed; duplicate-target and count formal-tree/mutation
contracts pass} \\
AI-19 & P0 & Scalar attribute codec & \texttt{CanonicalScalarValueCodec}
provides typed tags, exact escaping, strict decoding, Int64/finite-Real
checks, and an independently implemented expected-payload calculation in
the adapter & \textbf{Closed statically and by the canonical real-Neo4j
supplement} \\
AI-20 & P0 & Qualifier codec and bottom & Writer and renderer use the
typed wire contract; the adapter derives expected payloads
independently; graph pull decodes payloads by association-end type; a
bottom qualifier has an explicit non-null guard and matches no link row
& \textbf{Closed statically and by the canonical real-Neo4j
supplement} \\
AI-21 & P1 & Model isolation and certificate visibility & Context,
allInstances, attribute, navigation, type, cast, re-identification, M2
synchronization, graph snapshots/pull, and fallback lookups constrain
\texttt{modelKey}; current-key scans see wrong-model clones and
cross-model edges; class/attribute keys are unique & \textbf{Closed
statically and by the fresh adapter certificate} \\
AI-22 & P1 & Association-class boundary & Metamodel resolution rejects
association-class navigation with
\texttt{ASSOCIATION\_CLASS\_UNSUPPORTED}; dedicated admission/compiler
tests pass & \textbf{Closed by certified exclusion} \\
AI-23 & P2 & Semantic/optimized IR separation & \texttt{InvariantQuery}
carries \texttt{SEMANTIC}/\texttt{OPTIMIZED} stage and producer version;
only the optimizer creates the opaque current artifact accepted by
expression planning; invariant planning validates the optimizer version
& \textbf{Closed at the Java API boundary} \\
AI-24 & P0 & Nested navigation-iterator planning &
\texttt{Department::HasMinimumWitness} produced a correlated
\texttt{EXISTS\ \{\ UNWIND\ COLLECT\ \{\ ...\ NOT\ EXISTS\ ...\ \}\ \}}
whose Neo4j \texttt{EXPLAIN} alone ran for more than three minutes and
had to be terminated & \textbf{Closed for the demonstrated family.} The
optimizer now retains/materializes the outer iterator when its predicate
already contains an iterator/navigation query; the focused real case
fell to 6.732 s and the full 30-case suite completed in 21.85 s with
exact IDs. PO-18 still requires universal rewrite preservation \\
AI-25 & P0 & Persisted scalar type aliases & M1 repeat comparison failed
after a correct \texttt{v1\textbar{}I\textbar{}18} write because graph
metadata said \texttt{Int} while the strict decoder accepted only
\texttt{Integer}; \texttt{Double}/\texttt{Real} had the same latent
boundary & \textbf{Closed.} The codec normalizes only the declared
database aliases \texttt{Int\ -\textgreater{}\ Integer} and
\texttt{Double\ -\textgreater{}\ Real}; focused codec/snapshot tests and
the real research repeat-diff pass \\
AI-26 & P0 & Canonicality of real-test fixtures & Bottom-receiver and
research/scale fixtures still wrote legacy untagged scalars; the scale
bulk loader also omitted \texttt{slotKey}/\texttt{linkKey}, so a nominal
canonical-profile experiment could observe bottom or a weakened graph
instead of the stated domain & \textbf{Closed in the exercised
harnesses.} Fixtures use the production typed codec or exact wire tags;
scale assertions reject malformed scalar slots, missing binary-link
metadata, and non-injective link keys before timing \\
AI-27 & P1 & Mutation/profile evidence validity & The KEY mutation used
an unlabeled node outside the snapshot domain, changing it to
\texttt{:UmlClass} was blocked by the correct uniqueness constraint, and
the research manifest selected the first DBMS component (\texttt{5})
rather than \texttt{Neo4j\ Kernel} (\texttt{2026.06.0}) & \textbf{Closed
in the harness and clean checked-in evidence.} KEY is mutated through
association metadata inside the observation domain, all R/PA/M2/KEY
mutants are killed, and manifest extraction selects/asserts the pinned
kernel \\
AI-28 & P0 & Whole-collection bottom at a primary set receiver &
\texttt{finiteSet(bottom\_Set(tau))=empty}; the discriminator
\texttt{(if\ self.flag\ then\ null\ else\ Set\{1\}\ endif)-\textgreater{}isEmpty()}
requires either runtime bottom representation to behave as
\texttt{{[}{]}} before the set observation & \textbf{Closed for the
certified profile.} \texttt{renderCollectionView} maps both Cypher
\texttt{null} and the reserved bottom token to \texttt{{[}{]}}; focused
and clean selected-runtime cases pass \\
AI-29 & P0 & Whole-collection bottom at secondary collection boundaries
& Normalizing only the primary receiver left RHS operands unsafe. The
admitted counterexample
\texttt{Set\{1\}-\textgreater{}includesAll(if\ true\ then\ null\ else\ Set\{1\}\ endif)}
reached \texttt{all(...\ IN\ null\ ...)}. Direct bound/plan construction
can additionally place a token-shaped collection cell at set equality,
\texttt{flatten}, collection definition checks, entity-source
\texttt{UNWIND}, or collection attribute projection & \textbf{Repaired
with an explicit scope distinction.} Every collection cell at these
renderer boundaries now passes through the same null-or-token-to-empty
\texttt{CASE}; the admitted null-RHS discriminator passes on the
selected runtime. Token-shaped plan regressions are defensive
renderer-totality evidence only: their current producers use
\texttt{any} or a collection-valued \texttt{collect} body and are
rejected by frozen OCL\_val admission, so they are not cited as
certified theorem cases. Universal \texttt{PlanAdequacy} remains
pending \\
AI-30 & P0 & Optimized navigation alias capture & The optimized
navigation path bound its target before rendering the owner/qualifiers.
In
\texttt{self.children-\textgreater{}exists(x\ \textbar{}\ x.children-\textgreater{}exists(x\ \textbar{}\ true))},
the inner owner named \texttt{x} could resolve to the new, not-yet-bound
target alias instead of the outer iterator & \textbf{Closed for the
demonstrated family.} Owner and qualifier expressions are rendered in
the outer lexical scope before the target binding is entered; the
nested-shadowing optimized-navigation regression passes. PO-18 still
requires a constructor-universal scope proof \\
AI-31 & P0 & Real64 signed-zero refinement & Formal numeric equality
identifies \texttt{-0.0} and \texttt{+0.0}, while
\texttt{Double.compare} distinguished them during constant folding; the
writer/oracle and qualifier text could also emit two wire encodings.
\texttt{(0.0\ /\ (0\ -\ 1))\ =\ 0.0} was a concrete optimizer
counterexample & \textbf{Closed for the demonstrated family.} Optimizer
folding canonicalizes zero, the scalar codec and independent adapter
oracle emit only \texttt{v1\textbar{}R\textbar{}0.0}, non-canonical
\texttt{v1\textbar{}R\textbar{}-0.0} is rejected, and qualifier
rendering canonicalizes zero. Focused optimizer/codec/qualifier/renderer
tests pass \\
AI-32 & P0 & Validation truth of a non-null bottom token &
\texttt{bool\_val(v)} is true iff \texttt{v=true}, but production used
\texttt{coalesce(v,false)}. When a bottom element was represented by the
non-null map token, the expression remained a map and could reach
Boolean operators or \texttt{WHERE} & \textbf{Closed for the
demonstrated family.} \texttt{validationTruth(v)} is now exactly
\texttt{coalesce((v)=true,false)}, so false, null, and every non-Boolean
bottom representation collapse to false. The 52 generated trees, the
335-case implementation-conformance gate, and the selected-runtime
\texttt{Set\{null,true\}-\textgreater{}exists(x\textbar{}x)}
discriminator pass \\
AI-33 & P0 & Bottom-token scalar operations & A bottom element of a
finite Set is materialized as the reserved non-null token. Scalar
equality recognized only Cypher \texttt{null}, while ordering,
arithmetic, and qualifier serialization could receive the token
directly. Thus the admitted
\texttt{Set\{null,1\}-\textgreater{}exists(x\textbar{}x=null)} disagreed
with \texttt{eq\_val(bottom,bottom)}, and arithmetic could raise a
backend type error & \textbf{Closed for the demonstrated certified
family.} Scalar equality now uses the same total \texttt{isBottom}
predicate as the formal \texttt{semEq}; ordering and arithmetic return
Cypher \texttt{null} whenever either operand is either bottom
representation; qualifier guards and serialization also recognize both
forms. Focused compiler tests and real-Neo4j discriminators for
equality, ordering, and arithmetic pass \\
AI-34 & P0 & \texttt{Void}-\/-\texttt{Set} equality dispatch &
\texttt{eqJoin(Void,Set(sigma))=Set(sigma)} and the certified set
observation maps collection bottom to the empty set, but production
selected extensional set equality only when both operand records were
already collection-typed. The admitted
\texttt{null\ =\ Set\{1\}-\textgreater{}select(x\textbar{}false)}
therefore used scalar bottom equality and returned false instead of true
& \textbf{Closed for the demonstrated family.} Equality selects the
extensional set path when at least one operand is collection-typed and
both operands are collection-or-\texttt{Void}; each operand then crosses
C5a before mutual inclusion. The focused certified compiler check and
selected real-Neo4j discriminator pass \\
\end{longtable}
}

\subsubsection{1.2.5 What the Passing Gates Do and Do Not
Establish}\label{125-what-the-passing-gates-do-and-do-not-establish}

The discriminator families from the previous audit are now represented
by negative-admission, exact-numeric, optimizer-scope, alias,
bottom-receiver, duplicate-navigation, codec, model-isolation,
association-class, and IR-stage tests. The static suite and the 52 exact
generated-tree cases pass, and the compiler mutation score is 20/20.
Together with the focused collection-bottom regressions and the
2026-08-22 real-Neo4j run, this closes the demonstrated AI-01-\/-AI-34
families under their stated boundaries and discharges the finite runtime
obligations PO-15, PO-16, and PO-21-\/-PO-23. It does not close PO-18:
MK-T4 remains payload-parametric and MK-NORMALIZE does not universally
prove every concrete production rewrite. Runtime evidence therefore
supports only the certified finite profile, not a universal Java/Neo4j
implementation proof.

\subsubsection{1.2.6 Repair Order}\label{126-repair-order}

\begin{enumerate}
\def\labelenumi{\arabic{enumi}.}
\tightlist
\item
  \textbf{Completed --- capture the repaired semantic boundaries
  cleanly:} AI-28 and AI-29-\/-AI-34 are included in the common
  selected-runtime evidence.
\item
  \textbf{Completed --- publish clean runtime evidence:} the stabilized
  implementation is committed at \texttt{ebc1a4da}; the Cypher 5/Neo4j
  profile was rerun for adapter adequacy, OCL\_val-47,
  Void/bottom/scalar, and Families/Persons, and all three manifests use
  the same source revision and pass clean-evidence guards.
\item
  \textbf{Remaining research task --- close proof universality:}
  instantiate the primitive-agreement obligations and prove the concrete
  Java optimizer/renderer refinement needed to close PO-18.
\item
  \textbf{Completed for the certified fragment --- finish structural
  cleanup:} retain the invariant-root separation and either prove or
  replace the bounded raw-AST bridge (AI-07). Keep the profile-adjusted
  \texttt{collect} boundary (AI-09) explicit.
\item
  \textbf{Completed --- synchronize publication artifacts:} registry,
  Lean hash, runtime manifests, and plan status are synchronized; only
  PO-18 remains explicitly partial.
\end{enumerate}

The following are outside the main theorem unless explicitly added later
with separate semantics, implementation evidence, and proof obligations.

\begin{Shaded}
\begin{Highlighting}[]
\NormalTok{full OMG OCL}
\NormalTok{full null/invalid four{-}valued OCL semantics}
\NormalTok{Bag, Sequence, OrderedSet preservation}
\NormalTok{ordered collection semantics}
\NormalTok{non{-}binary association navigation}
\NormalTok{fallback{-}only constructs}
\NormalTok{arbitrary Cypher}
\NormalTok{arbitrary property{-}graph encodings}
\end{Highlighting}
\end{Shaded}

Validation truth is deliberately a validation-level policy:

\begin{Shaded}
\begin{Highlighting}[]
\NormalTok{bool\_val(v) = true  iff v = true}
\NormalTok{bool\_val(v) = false otherwise}
\end{Highlighting}
\end{Shaded}

Thus the certified \texttt{bottom}, which deliberately collapses the
admitted null/undefined/invalid-like outcomes of this profile, is not
validation-true. An unsupported backend value is outside the theorem
domain; it is not silently reclassified as \texttt{bottom}. This policy
is suitable for invariant violation queries but is not a proof of full
OMG OCL truth tables.

\begin{center}\rule{0.5\linewidth}{0.5pt}\end{center}

\section{2. Supported OCL Validation Fragment
OCL\_val}\label{2-supported-ocl-validation-fragment-ocl_val}

\texttt{OCL\_val} is the theorem-scope source-language fragment and is
defined before its parser representation. It fixes the supported source
constructors and validation-level semantics, but it does not yet contain
resolved UML references. Optional source type annotations remain
unresolved names at this level. The parser represents an
\texttt{OCL\_val} source expression as an OCL AST. The partial
structural decoder
\texttt{decode\_val\ :\ OCL\_AST\ -\textgreater{}partial\ OCL\_val}
recognizes exactly the AST forms of this source fragment. Successful
binding then produces a different model, \texttt{BoundOCL\_val(MM)},
whose names, navigation metadata, scopes, and types have been resolved
against \texttt{MM}.

Thus the specification order is:

\begin{Shaded}
\begin{Highlighting}[]
\NormalTok{OCL\_val source metamodel}
\NormalTok{  {-}\textgreater{} AST\_val representation}
\NormalTok{  {-}\textgreater{} BoundOCL\_val(MM)}
\end{Highlighting}
\end{Shaded}

while the runtime pipeline remains:

\begin{Shaded}
\begin{Highlighting}[]
\NormalTok{OCL Text {-}\textgreater{} Parse {-}\textgreater{} OCL AST {-}\textgreater{} T\_BIND {-}\textgreater{} Bound OCL model.}
\end{Highlighting}
\end{Shaded}

Parser support for a construct does not imply theorem coverage.

\subsection{2.1 Grammar-Level Fragment}\label{21-grammar-level-fragment}

The proved source fragment \texttt{OCL\_val} contains invariant
expressions built from the following constructs. Expressions must be
typable for binding to succeed, but the source metamodel itself stores
unresolved names rather than resolved UML elements. The metavariable
\texttt{e} in this grammar denotes a decoded predicate expression; it is
neither the invariant text \texttt{tInv} nor the parser tree
\texttt{astInv}:

\begin{Shaded}
\begin{Highlighting}[]
\NormalTok{e ::= self}
\NormalTok{    | x}
\NormalTok{    | literal}
\NormalTok{    | Set\{e1,...,en\}                 where n \textgreater{}= 1 and elements have one canonical non{-}collection type}
\NormalTok{    | e.a                               where a has a primitive scalar type}
\NormalTok{    | e.role}
\NormalTok{    | e.role[q1,...,qk]}
\NormalTok{    | C.allInstances()}
\NormalTok{    | not e}
\NormalTok{    | e and e}
\NormalTok{    | e or e}
\NormalTok{    | e xor e}
\NormalTok{    | e implies e}
\NormalTok{    | e = e}
\NormalTok{    | e \textless{}\textgreater{} e}
\NormalTok{    | e \textless{} e | e \textless{}= e | e \textgreater{} e | e \textgreater{}= e}
\NormalTok{    | e + e | e {-} e | e * e | e / e}
\NormalTok{    | if e then e else e endif}
\NormalTok{    | let x [: tauD] = e1 in e2}
\NormalTok{    | e{-}\textgreater{}exists(x [: sigmaD] | e)}
\NormalTok{    | e{-}\textgreater{}forAll(x [: sigmaD] | e)}
\NormalTok{    | e{-}\textgreater{}select(x [: sigmaD] | e)}
\NormalTok{    | e{-}\textgreater{}reject(x [: sigmaD] | e)}
\NormalTok{    | e{-}\textgreater{}collect(x [: sigmaD] | e)}
\NormalTok{    | e{-}\textgreater{}isUnique(x [: sigmaD] | e)  where the body is not collection{-}valued}
\NormalTok{    | e{-}\textgreater{}includes(e)}
\NormalTok{    | e{-}\textgreater{}excludes(e)}
\NormalTok{    | e{-}\textgreater{}includesAll(e)}
\NormalTok{    | e{-}\textgreater{}excludesAll(e)}
\NormalTok{    | e{-}\textgreater{}union(e)}
\NormalTok{    | e{-}\textgreater{}intersection(e)}
\NormalTok{    | e{-}\textgreater{}asSet()}
\NormalTok{    | e{-}\textgreater{}size()}
\NormalTok{    | e{-}\textgreater{}isEmpty()}
\NormalTok{    | e{-}\textgreater{}notEmpty()}
\NormalTok{    | e.oclIsKindOf(C)}
\NormalTok{    | e.oclAsType(C)}
\end{Highlighting}
\end{Shaded}

Navigation is limited to ordinary binary associations represented by
direct \texttt{Link*} relationships. Association classes are excluded
because the repository encodes their instances as link-object nodes and
spokes, not as this direct relationship shape. Qualified navigation is
included only for scalar primitive qualifier payloads whose typed
serialization is fixed by Section 4.

For omitted declarations, the binder infers \texttt{tauD} from the
initializer and \texttt{sigmaD} from the source element type. For
explicit declarations, the named type must resolve uniquely in
\texttt{MM} and the actual type must conform to it. Class names remain
position-restricted as expression values: only \texttt{C.allInstances()}
and the single class argument of an admitted type operation may contain
a class reference. A class name used as a declaration type is type
syntax, not a first-class class-reference value. Implicit
collection-property projection (\texttt{collection.property}) is not
desugared in the certified path; the source must spell the corresponding
\texttt{collect} explicitly.

\texttt{OCL\_val} preserves the source OCL/USE result kind of binary
navigation. An ordinary unqualified end with upper multiplicity one has
static type \texttt{D}; a native collection result is normalized to
\texttt{Set(D)} only in the certified finite-set profile. The native
result binding is authoritative because a qualified navigation can
remain collection-typed even when its target end is upper-one. When
\texttt{-\textgreater{}} applies one of \texttt{asSet}, \texttt{size},
\texttt{isEmpty}, or \texttt{notEmpty} directly to a native-scalar
to-one navigation, binding records a separate \texttt{Set(D)} collection
view at that operator boundary: a present target becomes a singleton and
an absent target becomes the empty set. An iterator never performs this
implicit lift; it requires a native \texttt{Set(D)} source or an
explicit preceding \texttt{asSet()}. The property node itself remains
typed \texttt{D}. This separation is necessary for the
binding-preservation statement below.

Type operations are included under these conditions:

\begin{Shaded}
\begin{Highlighting}[]
\NormalTok{oclIsKindOf(C) uses the conformance relation.}
\NormalTok{oclAsType(C) preserves values when conformance holds and yields bottom}
\NormalTok{otherwise.}
\end{Highlighting}
\end{Shaded}

\texttt{size()} denotes finite-set cardinality. The internal VA operator
\texttt{Count(S)} also denotes finite-set cardinality and is used by
normalization and Cypher realization. OCL
\texttt{collection-\textgreater{}count(element)} is not part of
\texttt{OCL\_val}.

\texttt{collect} is interpreted only as finite-set image construction
with a non-collection body type. The theorem does not preserve Bag
multiplicity, ordering, implicit flattening, or nested collection
results of full OMG OCL \texttt{collect}.

This restriction is executable, not merely documentary. The certified
compiler first applies the closed syntactic admission policy, then binds
with the finite-set profile, and finally applies a bound/type admission
policy. The certified binder preserves a native-scalar to-one
navigation\textquotesingle s \texttt{Entity} type and stores
\texttt{sourceCollectionType=Set(Entity)} only on a directly consuming
collection operation/iterator; every native collection-valued navigation
and \texttt{collect} use certified \texttt{Set} results. The bound
policy rejects a non-Boolean invariant root, an arbitrary scalar
collection source, every residual \texttt{Bag}, \texttt{Sequence}, or
\texttt{OrderedSet}, a collection-valued \texttt{collect} body, and
every collection- or reference-valued attribute access. The only
admitted scalar collection-source view is a resolved native-scalar
\texttt{{[}0..1{]}}/\texttt{{[}1{]}} navigation with matching target
entity type. Qualified collection-valued navigation is not misclassified
merely from its target-end multiplicity. The general \texttt{compile()}
API may still bind those broader forms, but it produces no
\texttt{OCL\_val} certificate and is outside Theorems 0-\/-6.
Enumeration literals and enumeration-valued attributes are treated the
same way: the general compiler retains them, while bound admission
excludes them until the static type grammar, typed encoding, equality,
and PA5 cases are extended together.

\subsection{2.1.1 Source, AST, and Bound Metamodel
Certificate}\label{211-source-ast-and-bound-metamodel-certificate}

The grammar above defines the source metamodel \texttt{OCL\_val}. Its
parser representation and bound target are distinct:

\begin{Shaded}
\begin{Highlighting}[]
\NormalTok{encode\_ast : OCL\_val {-}\textgreater{} AST\_val}
\NormalTok{decode\_val : OCL\_AST {-}\textgreater{}partial OCL\_val}
\NormalTok{decode\_val(encode\_ast(s)) = s}

\NormalTok{encode\_inv : Inv\_val {-}\textgreater{} ASTInv\_val}
\NormalTok{decodeInvariant\_val(encode\_inv(I)) = I}

\NormalTok{T\_BIND\_EXPR : AST\_val x MM x Gamma {-}\textgreater{}partial BoundOCL\_val(MM)}
\NormalTok{T\_BIND\_INV  : ASTInv\_val x MM {-}\textgreater{}partial BoundValInvariant(MM)}

\NormalTok{Adm\_MM(astInv) iff}
\NormalTok{  RawAdm(astInv)}
\NormalTok{  and decodeInvariant\_val(astInv)=I is defined}
\NormalTok{  and T\_BIND\_INV(astInv,MM)=bInv is defined}
\NormalTok{  and BoundAdm\_MM(bInv) holds.}
\end{Highlighting}
\end{Shaded}

The metaclasses below describe \texttt{BoundOCL\_val(MM)}, not source
\texttt{OCL\_val} and not the parser OCL AST metamodel.

{\def\LTcaptype{none} % do not increment counter
\begin{longtable}[]{@{}lll@{}}
\toprule\noalign{}
Metaclass & Structural features & Scope condition \\
\midrule\noalign{}
\endhead
\bottomrule\noalign{}
\endlastfoot
\texttt{BoundValModule} & containment
\texttt{invariants\ :\ BoundValInvariant{[}*{]}} & invariant roots
only \\
\texttt{BoundValInvariant} & resolved \texttt{contextClass},
\texttt{name}, containment \texttt{predicate} & predicate has type
Boolean \\
\texttt{BoundValExpr\textless{}tau\textgreater{}} & abstract
\texttt{staticType\ :\ tau} & unique canonical type \\
\texttt{VSelf}, \texttt{VVariable}, \texttt{VLiteral} &
context/declaration/literal data & typed base expressions \\
\texttt{VSetLiteral} & nonempty contained element list & one canonical
non-collection element type \\
\texttt{VAttribute} & \texttt{source}, resolved \texttt{attribute} &
attribute belongs to \texttt{MM} and has primitive scalar type \\
\texttt{VNavigation} & \texttt{source}, resolved association, roles,
direction, qualifiers & binary, fixed qualifier signature \\
\texttt{VAllInstances} & resolved class & finite \texttt{Set(C)}
result \\
\texttt{VNot}, \texttt{VBinary}, \texttt{VIf} & typed child containments
& admitted scalar/Boolean operators only \\
\texttt{VLet} & declaration, resolved \texttt{variableType},
initializer, body & initializer type conforms to \texttt{variableType};
declaration scopes only the body \\
\texttt{VIterator} & operation, declaration, resolved
\texttt{iteratorVariableType}, source, body & source element type
conforms to \texttt{iteratorVariableType}; operation in
exists/forAll/select/reject/collect/isUnique \\
\texttt{VCollectionOp} & operation, source, optional argument &
finite-set operation only \\
\texttt{VTypeOp} & operation, source, resolved class & kindOf or cast
only \\
\end{longtable}
}

An instance is in \texttt{BoundOCL\_val(MM)} only when: containment is
finite, rooted, acyclic, and unshared; all UML references are resolved;
every expression has its canonical type; navigation is binary and
direction-explicit; collection types are finite \texttt{Set};
iterator/let scope uses nearest binding; collect has a non-collection
body; qualifiers are fixed-arity ordered primitive scalar lists; and
only the closed operator sets in the grammar occur. In particular,
\texttt{oclIsTypeOf} has no certified metaclass/operation literal.
Collection- and reference-valued UML attributes likewise have no
certified \texttt{VAttribute} instance; support in the general compiler
does not extend this bound language.

\textbf{M1 Source-\/-AST Representation.} \texttt{encode\_ast} and
\texttt{decode\_val} form a retraction between source \texttt{OCL\_val}
expressions and \texttt{AST\_val}:

\begin{Shaded}
\begin{Highlighting}[]
\NormalTok{decode\_val(encode\_ast(s)) = s.}
\end{Highlighting}
\end{Shaded}

Every \texttt{ast\ in\ AST\_val} has one decoded source expression, up
to capture-avoiding alpha-renaming of bound variables.

Proof. Structural recursion maps every source constructor to its unique
parser AST constructor and recursively maps its children. Conversely,
\texttt{decode\_val} has exactly one equation for each supported AST
form and is undefined for all other forms. Induction on the finite
source tree proves the retraction equation; induction on an admitted AST
proves uniqueness. Binder names may differ only by capture-avoiding
alpha-renaming.

\textbf{M2 Bound Subexpression Closure and Structural Induction.} Every
contained expression of \texttt{BoundOCL\_val(MM)} is a well-formed
typed bound expression under the lexical environment induced by its
ancestors, and strict containment is well founded.

Proof. Finite acyclic containment is well founded. Child typing premises
are required before parent admission. Iterator and let bodies receive
exactly one displayed environment extension; other children retain the
parent environment. Induction on the unique root-to-child path proves
closure.

\textbf{M3 Canonical Type Uniqueness.} If
\texttt{MM;Gamma\ \textbar{}-\ e:tau1} and
\texttt{MM;Gamma\ \textbar{}-\ e:tau2} are canonical typings of an
admitted expression, then \texttt{tau1=tau2}.

Proof. Induct on the metaclass. Resolved references and fixed-result
operators determine their result types. \texttt{arithResult},
\texttt{ifJoin}, \texttt{eqJoin}, \texttt{memberJoin},
\texttt{setJoinN}, and \texttt{setJoin2} are partial functions. Iterator
and let cases use the induction hypotheses under the unique lexical
extension selected by \texttt{declType}. There is no general subsumption
rule: conformance is used only at the declared binder boundary, whose
resolved type is therefore unique.

\textbf{M4 Admission and Binding Soundness.} If
\texttt{decode\_val(ast)=s} and the certified binder derives
\texttt{MM;Gamma\ \textbar{}-\ ast\ =\textgreater{}\ b:tau}, then
\texttt{s\ in\ OCL\_val}, \texttt{ast\ in\ AST\_val},
\texttt{b\ in\ BoundOCL\_val(MM)}, and
\texttt{MM;Gamma\ \textbar{}-\ b:tau}.

Proof. Induct on the binding derivation. Each rule creates one listed
metaclass; child hypotheses establish finite typed containments,
resolver functionality establishes references/navigation metadata, and
binder rules establish scope. No certified rule exists for an excluded
construct.

\textbf{M5 Exclusion Stability.} Neither an admitted source/AST pair nor
its \texttt{BoundOCL\_val(MM)} result contains an excluded operation,
ordered/multiplicity-sensitive collection sort, non-binary navigation,
or non-invariant root.

Proof. M1 gives no \texttt{decode\_val} equation for an excluded AST
constructor. By M2, every bound subexpression has a listed bound
metaclass. Closed operator and type domains, binary navigation metadata,
and the single invariant root leave no source, AST, or bound
representation case for an excluded construct.

M1-\/-M5, including M3a, are scope lemmas. M2 justifies structural
induction, M3 supplies type functionality, M4 connects parser admission
to the theorem domain, and F1 below establishes the converse
erase-\/-rebind result under stable resolvers.

\subsection{2.2 Excluded or Partial
Constructs}\label{22-excluded-or-partial-constructs}

The following constructs are not included in Theorem 6:

\begin{Shaded}
\begin{Highlighting}[]
\NormalTok{any          unless a deterministic choice policy is fixed and proved}
\NormalTok{oclIsTypeOf unless a direct runtime{-}class accessor is added and proved}
\NormalTok{one          unless rewritten to finite{-}set cardinality and proved separately}
\NormalTok{count(e)     unless represented by a separate CountElem(S,E) operator}
\NormalTok{sortedBy     unless ordered collection semantics is added}
\NormalTok{iterate      outside current fragment}
\NormalTok{closure      outside current fragment}
\NormalTok{Tuple        outside current fragment}
\NormalTok{Message      outside current fragment}
\NormalTok{State        outside current fragment}
\NormalTok{pre/post/body/init/derive constraints}
\NormalTok{Bag, Sequence, OrderedSet multiplicity/order preservation}
\NormalTok{association{-}class navigation through link{-}object/spoke encoding}
\NormalTok{class references used as first{-}class values}
\NormalTok{implicit collection{-}property projection}
\NormalTok{fallback{-}only constructs}
\end{Highlighting}
\end{Shaded}

Implementations may support some of these constructs experimentally.
Such support is outside the theorem unless the construct is added to
\texttt{OCL\_val} with typing, denotation, realization, and proof cases.

The prototype enforces this distinction with a closed-world
\texttt{OclValAdmissionPolicy} on the instrumented proof/conformance
path. The policy admits exactly the source constructor families listed
in Section 2.1 and rejects any other iterator, collection operation,
method call, or unknown AST node with stable diagnostic
\texttt{OCL\_VAL\_EXCLUDED\_CONSTRUCT} before binding. The general
\texttt{compile()} API may continue to expose explicitly experimental
translations; success there is not an \texttt{OCL\_val} admission
certificate and is not consumed by Theorem 6 evidence. All instrumented
and real-Neo4j differential theorem suites pass through the closed-world
policy.

In particular, \texttt{one} is outside Theorem 6 unless it is rewritten
to \texttt{Compare(EQ,\ Count(Select(S,x,P)),\ 1)} and the rewrite is
proved. \texttt{count(element)} is outside Theorem 6 unless it is
translated to a separate \texttt{CountElem(S,E)} operator under the
chosen finite-set membership semantics.

\subsection{2.3 Static Types, Resolver Contracts, and Typing
Judgments}\label{23-static-types-resolver-contracts-and-typing-judgments}

Let \texttt{sigma} range over non-collection value types and
\texttt{tau} over every type in the proved fragment:

\begin{Shaded}
\begin{Highlighting}[]
\NormalTok{sigma ::= Boolean | Integer | Real | String | C}
\NormalTok{delta ::= Void | sigma}
\NormalTok{tau   ::= Void | sigma | Set(sigma)}
\end{Highlighting}
\end{Shaded}

where \texttt{C} ranges over classes of \texttt{MM}. \texttt{Void} is
the binder-internal type of the literal \texttt{null}; its only
denotation is runtime \texttt{bottom}. It is not a claim that the full
OMG \texttt{OclVoid}/\texttt{OclInvalid} lattice is supported. Other
sources of runtime \texttt{bottom} remain value-policy results rather
than acquiring type \texttt{Void}. Let \texttt{tau1\ \textless{}=\ tau2}
be the least conformance relation that contains
\texttt{Void\ \textless{}=\ tau} for every admitted \texttt{tau}, UML
class conformance, \texttt{Integer\ \textless{}=\ Real}, reflexivity,
and the covariant finite-set rule
\texttt{Set(sigma1)\ \textless{}=\ Set(sigma2)} whenever
\texttt{sigma1\ \textless{}=\ sigma2}. \texttt{Set(Void)} and nested
sets are not types of \texttt{OCL\_val}. A \texttt{Void} expression can
nevertheless be lifted to a non-Void branch or set-element type by the
explicit partial functions below. No other implicit collection
conversion is used.

Let \texttt{T?} be either an omitted annotation or a surface type name.
Declaration resolution is the deterministic partial function:

\begin{Shaded}
\begin{Highlighting}[]
\NormalTok{declType\_MM(omitted,tauActual) = tauActual}

\NormalTok{declType\_MM(T,tauActual) = tauDeclared}
\NormalTok{  iff resolveType\_MM(T)=tauDeclared}
\NormalTok{  and tauActual \textless{}= tauDeclared.}
\end{Highlighting}
\end{Shaded}

\texttt{resolveType\_MM} covers the admitted primitive types, UML
classes, and finite \texttt{Set(sigma)} types. An unknown name or a
failed conformance premise makes binding undefined. The canonical
embedding associated with a successful premise is written
\texttt{iota{[}tauActual-\textgreater{}tauDeclared{]}}; it is identity
for reflexivity and UML class upcast, maps \texttt{Void} to typed
bottom, maps Integer to Real under A8, and lifts componentwise to finite
sets.

Define the deterministic partial type functions:

\begin{Shaded}
\begin{Highlighting}[]
\NormalTok{arithResult(PLUS, Integer, Integer)  = Integer}
\NormalTok{arithResult(MINUS,Integer, Integer)  = Integer}
\NormalTok{arithResult(TIMES,Integer, Integer)  = Integer}
\NormalTok{arithResult(DIV,  Integer, Integer)  = Real}

\NormalTok{arithResult(op,tau1,tau2) = Real}
\NormalTok{  when op in \{PLUS,MINUS,TIMES,DIV\},}
\NormalTok{       tau1,tau2 in \{Integer,Real\},}
\NormalTok{       and at least one operand type is Real.}

\NormalTok{arithJoin(PLUS, Integer, Integer)  = Integer}
\NormalTok{arithJoin(MINUS,Integer, Integer)  = Integer}
\NormalTok{arithJoin(TIMES,Integer, Integer)  = Integer}
\NormalTok{arithJoin(DIV,  Integer, Integer)  = Real}
\NormalTok{arithJoin(op,tau1,tau2)            = Real}
\NormalTok{  when op in \{PLUS,MINUS,TIMES,DIV\},}
\NormalTok{       tau1,tau2 in \{Integer,Real\},}
\NormalTok{       and at least one operand type is Real.}

\NormalTok{orderJoin(Integer,Integer) = Integer}
\NormalTok{orderJoin(Real,Real)       = Real}
\NormalTok{orderJoin(Integer,Real)    = Real}
\NormalTok{orderJoin(Real,Integer)    = Real}

\NormalTok{ifJoin(tau,tau)  = tau}
\NormalTok{ifJoin(Void,tau) = tau}
\NormalTok{ifJoin(tau,Void) = tau}
\NormalTok{ifJoin is undefined in every other case.}

\NormalTok{eqJoin(tau,tau)         = tau}
\NormalTok{eqJoin(Void,tau)        = tau}
\NormalTok{eqJoin(tau,Void)        = tau}
\NormalTok{eqJoin(Integer,Real)    = Real}
\NormalTok{eqJoin(Real,Integer)    = Real}
\NormalTok{eqJoin is undefined in every other case.}

\NormalTok{eqCompatible(tau1,tau2) iff eqJoin(tau1,tau2) is defined.}

\NormalTok{baseJoin(sigma,sigma)        = sigma}
\NormalTok{baseJoin(Integer,Real)       = Real}
\NormalTok{baseJoin(Real,Integer)       = Real}
\NormalTok{baseJoin(C,D)                = lub\_MM(C,D), when that least common superclass is unique}

\NormalTok{setJoin2(sigma1,sigma2) = baseJoin(sigma1,sigma2).}

\NormalTok{memberJoin(Void,sigma) = sigma}
\NormalTok{memberJoin(sigma,Void) = sigma}
\NormalTok{memberJoin(sigma1,sigma2) = baseJoin(sigma1,sigma2)}
\NormalTok{memberJoin(Void,Void) is undefined.}

\NormalTok{setJoinN([tau1,...,taun]) = sigma, where n\textgreater{}=1, iff:}
\NormalTok{  (i) every tau\_i is Void or a non{-}collection type sigma\_i;}
\NormalTok{  (ii) at least one tau\_i is non{-}Void; and}
\NormalTok{  (iii) folding baseJoin over the non{-}Void entries, in source order, is}
\NormalTok{        defined and returns sigma.}

\NormalTok{\textasciigrave{}setJoinN\textasciigrave{} is undefined for an empty list, a Void{-}only list, any collection}
\NormalTok{element, incompatible primitive types, primitive/class mixtures, or class}
\NormalTok{types without one unique least common superclass. \textasciigrave{}setJoin2\textasciigrave{} is used by union}
\NormalTok{and intersection, whose operands already have non{-}Void element types;}
\NormalTok{\textasciigrave{}setJoinN\textasciigrave{} is used only by a nonempty set literal.}

\NormalTok{compatibleOrdered(tau1,tau2)}
\NormalTok{  iff orderJoin(tau1,tau2) is defined.}
\end{Highlighting}
\end{Shaded}

If a static type function is undefined, the expression is not admitted.
The object evaluator may complete division by zero to \texttt{bottom},
but such an evaluation fails \texttt{ScalarClosed}; T3-\/-T6 do not
compare it with a backend that may raise an error instead.

\textbf{M3a Functionality of set joins.} If
\texttt{setJoinN({[}tau1,...,taun{]})=sigma1} and
\texttt{setJoinN({[}tau1,...,taun{]})=sigma2}, then
\texttt{sigma1=sigma2}. The analogous property holds for
\texttt{setJoin2}. Whenever defined, the result is the unique canonical
least common supertype of all non-Void element types.

For every lifting selected by \texttt{ifJoin}, \texttt{eqJoin},
\texttt{memberJoin}, \texttt{setJoinN}, \texttt{setJoin2},
\texttt{arithJoin}, or \texttt{orderJoin}, define the canonical typed
embedding

\begin{Shaded}
\begin{Highlighting}[]
\NormalTok{iota\_I[tau {-}\textgreater{} upsilon] : ValBottom\_I(tau) {-}\textgreater{}partial ValBottom\_I(upsilon)}

\NormalTok{iota\_I[tau {-}\textgreater{} tau](v) = v}
\NormalTok{iota\_I[Void {-}\textgreater{} tau](bottom\_Void) = bottom\_tau, when tau != Void}
\NormalTok{iota\_I[C {-}\textgreater{} D](o) = o, when C conformsTo D}
\NormalTok{iota\_I[C {-}\textgreater{} D](bottom\_C) = bottom\_D, when C conformsTo D}
\NormalTok{iota\_I[Integer {-}\textgreater{} Real](n) = exactBinary64(n), when abs(n) \textless{}= 2\^{}53}
\NormalTok{iota\_I[Integer {-}\textgreater{} Real](bottom\_Integer) = bottom\_Real}
\end{Highlighting}
\end{Shaded}

For a finite set, \texttt{upSet\_I{[}tau\ -\textgreater{}\ upsilon{]}}
applies this embedding elementwise. Set literals, \texttt{union}, and
\texttt{intersection} first lift both operands to their unique common
element type and only then apply typed equality, duplicate elimination,
or membership. Thus no untyped equality is used between distinct element
domains.

If \texttt{memberJoin(sigma,delta)=upsilon}, scalar membership uses:

\begin{Shaded}
\begin{Highlighting}[]
\NormalTok{memberEq\_I[sigma,delta](u,v)}
\NormalTok{  = eq\_val\^{}upsilon(}
\NormalTok{      iota\_I[sigma {-}\textgreater{} upsilon](u),}
\NormalTok{      iota\_I[delta {-}\textgreater{} upsilon](v))}

\NormalTok{memberSet\_I[sigma {-}\textgreater{} upsilon](S)}
\NormalTok{  = upSet\_I[sigma {-}\textgreater{} upsilon](S).}
\end{Highlighting}
\end{Shaded}

Likewise, a scalar equality with \texttt{eqJoin(tau,sigma)=upsilon}
compares the two canonically embedded operands in
\texttt{ValBottom\_I(upsilon)}. This is intentionally narrower than
arbitrary UML conformance: two different class types are not
equality-compatible merely because one conforms to the other, matching
the certified binder.

Proof. The domains of the \texttt{baseJoin} equations are disjoint and
each returns one result. The class equation is applicable only when
\texttt{lub\_MM} is unique; incomparable minimal candidates make the
function undefined. Removing \texttt{Void} entries is deterministic and
leaves at least one entry by premise (ii). Induction over the remaining
fold, using functionality of the preceding fold and of
\texttt{baseJoin}, proves equality of results. The upper-bound and
leastness properties are preserved at every step. Case analysis on the
selected embedding proves that \texttt{iota\_I} is total on values
reachable under \texttt{ScalarClosed}: identity and class upcast
preserve the represented value, Integer-to-Real conversion is exact only
for \texttt{abs(n)\textless{}=2\^{}53}, and typed bottom maps to typed
bottom. Elementwise image construction therefore makes \texttt{upSet\_I}
functional as well. These are typed helpers of SF-01-\/-SF-08, not a
second semantic evaluator.

The binder relies on deterministic partial resolver functions:

\begin{Shaded}
\begin{Highlighting}[]
\NormalTok{resolveClass\_MM(name)                 {-}\textgreater{} C}
\NormalTok{resolveAttribute\_MM(C,name)           {-}\textgreater{} (a,tau)}
\NormalTok{resolveNavigation\_MM(C,role,qTypes)   {-}\textgreater{}}
\NormalTok{  (A,fromRole,toRole,direction,targetClass,qualifierTypes,resultKind)}
\NormalTok{resolveProperty\_MM(C,name,qTypes)      {-}\textgreater{}}
\NormalTok{  Attr(a,tau) or}
\NormalTok{  Nav(A,fromRole,toRole,direction,targetClass,qualifierTypes,resultKind)}
\NormalTok{resolveTypeOperation\_MM(name)         {-}\textgreater{} supported type operation}
\end{Highlighting}
\end{Shaded}

Here \texttt{resultKind} is \texttt{ONE} or \texttt{MANY}. It is part of
the resolved metadata and is never reconstructed from the target class
after binding.

Each resolver is adequate and functional on its domain. Adequacy means
that a successful result denotes exactly the visible UML element
selected by the source OCL name under inherited-feature lookup.
Functionality means that two successful results for the same arguments
are equal. Ambiguous, inaccessible, or unsupported lookups are undefined
and therefore not admitted.

The typing environment \texttt{Gamma} is a finite stack of bindings.
Lookup selects the nearest binding, thereby defining lexical shadowing.
The principal typing judgment is:

\begin{Shaded}
\begin{Highlighting}[]
\NormalTok{MM ; Gamma |{-} b : tau}
\end{Highlighting}
\end{Shaded}

The partial source-view function used by both source typing and binding
is:

\begin{Shaded}
\begin{Highlighting}[]
\NormalTok{collectionView(op,S)=Set(tau)                      if type(S)=Set(tau)}
\NormalTok{collectionView(op,S)=Set(D)                        if type(S)=D, S is a}
\NormalTok{  directly consumed resolved native{-}scalar to{-}one navigation to D, and}
\NormalTok{  op is one of \{asSet,size,isEmpty,notEmpty\}}
\NormalTok{collectionView(op,S)=undefined                     otherwise.}
\end{Highlighting}
\end{Shaded}

Thus an iterator never obtains a collection merely because its source
has upper multiplicity one. The following rules define all
theorem-relevant source typing cases. Premises named
\texttt{resolved(...)} refer to the resolver functions above.

\begin{Shaded}
\begin{Highlighting}[]
\NormalTok{(T{-}Self)      Gamma(self)=C}
\NormalTok{              {-}{-}{-}{-}{-}{-}{-}{-}{-}{-}{-}{-}{-}{-}{-}{-}{-}{-}{-}{-}{-}{-}{-}{-}{-}}
\NormalTok{              MM;Gamma |{-} self : C}

\NormalTok{(T{-}Var)       nearest(Gamma,x)=tau}
\NormalTok{              {-}{-}{-}{-}{-}{-}{-}{-}{-}{-}{-}{-}{-}{-}{-}{-}{-}{-}{-}{-}{-}{-}{-}{-}{-}}
\NormalTok{              MM;Gamma |{-} x : tau}

\NormalTok{(T{-}Lit)       typeOfLiteral(c)=tau}
\NormalTok{              {-}{-}{-}{-}{-}{-}{-}{-}{-}{-}{-}{-}{-}{-}{-}{-}{-}{-}{-}{-}{-}{-}{-}{-}{-}}
\NormalTok{              MM;Gamma |{-} c : tau}

\NormalTok{              typeOfLiteral(null)=Void}

\NormalTok{(T{-}SetLit)    n\textgreater{}=1}
\NormalTok{              MM;Gamma |{-} ei : tau\_i for every i}
\NormalTok{              setJoinN([tau\_1,...,tau\_n])=sigma}
\NormalTok{              {-}{-}{-}{-}{-}{-}{-}{-}{-}{-}{-}{-}{-}{-}{-}{-}{-}{-}{-}{-}{-}{-}{-}{-}{-}{-}{-}{-}{-}{-}{-}{-}{-}{-}{-}{-}{-}{-}}
\NormalTok{              MM;Gamma |{-} Set\{e1,...,en\} : Set(sigma)}

\NormalTok{(T{-}Attr)      MM;Gamma |{-} e : D}
\NormalTok{              resolveAttribute\_MM(D,a)=(a,tau)}
\NormalTok{              {-}{-}{-}{-}{-}{-}{-}{-}{-}{-}{-}{-}{-}{-}{-}{-}{-}{-}{-}{-}{-}{-}{-}{-}{-}{-}{-}{-}{-}{-}{-}{-}{-}}
\NormalTok{              MM;Gamma |{-} e.a : tau}

\NormalTok{(T{-}Nav{-}One)   MM;Gamma |{-} e : D}
\NormalTok{              MM;Gamma |{-} qi : sigma\_i for every qualifier qi}
\NormalTok{              resolveNavigation\_MM(D,r,[sigma\_i])=}
\NormalTok{                (A,fr,tr,dir,E,[sigma\_i],ONE)}
\NormalTok{              {-}{-}{-}{-}{-}{-}{-}{-}{-}{-}{-}{-}{-}{-}{-}{-}{-}{-}{-}{-}{-}{-}{-}{-}{-}{-}{-}{-}{-}{-}{-}{-}{-}{-}{-}{-}{-}{-}{-}{-}{-}{-}{-}{-}{-}{-}{-}{-}{-}}
\NormalTok{              MM;Gamma |{-} e.r[q1,...,qk] : E}

\NormalTok{(T{-}Nav{-}Many)  MM;Gamma |{-} e : D}
\NormalTok{              MM;Gamma |{-} qi : sigma\_i for every qualifier qi}
\NormalTok{              resolveNavigation\_MM(D,r,[sigma\_i])=}
\NormalTok{                (A,fr,tr,dir,E,[sigma\_i],MANY)}
\NormalTok{              {-}{-}{-}{-}{-}{-}{-}{-}{-}{-}{-}{-}{-}{-}{-}{-}{-}{-}{-}{-}{-}{-}{-}{-}{-}{-}{-}{-}{-}{-}{-}{-}{-}{-}{-}{-}{-}{-}{-}{-}{-}{-}{-}{-}{-}{-}{-}{-}{-}}
\NormalTok{              MM;Gamma |{-} e.r[q1,...,qk] : Set(E)}

\NormalTok{There is deliberately no rule whose navigation receiver has type \textasciigrave{}Set(D)\textasciigrave{};}
\NormalTok{implicit \textasciigrave{}collection.property\textasciigrave{} projection is outside the certified binder.}
\NormalTok{\textasciigrave{}resultKind=ONE\textasciigrave{} is used exactly when the native resolved navigation is scalar}
\NormalTok{(normally an unqualified upper{-}one end); otherwise the kind is \textasciigrave{}MANY\textasciigrave{}.}
\NormalTok{Multiplicity alone never overrides the native resolved result kind.}

\NormalTok{(T{-}AllInst)   resolveClass\_MM(name)=C}
\NormalTok{              {-}{-}{-}{-}{-}{-}{-}{-}{-}{-}{-}{-}{-}{-}{-}{-}{-}{-}{-}{-}{-}{-}{-}{-}{-}{-}{-}{-}{-}{-}{-}{-}{-}}
\NormalTok{              MM;Gamma |{-} C.allInstances() : Set(C)}

\NormalTok{(T{-}Not)       MM;Gamma |{-} p : Boolean}
\NormalTok{              {-}{-}{-}{-}{-}{-}{-}{-}{-}{-}{-}{-}{-}{-}{-}{-}{-}{-}{-}{-}{-}{-}{-}{-}{-}}
\NormalTok{              MM;Gamma |{-} not p : Boolean}

\NormalTok{(T{-}Bool2)     MM;Gamma |{-} p1 : Boolean}
\NormalTok{              MM;Gamma |{-} p2 : Boolean}
\NormalTok{              op in \{and,or,xor,implies\}}
\NormalTok{              {-}{-}{-}{-}{-}{-}{-}{-}{-}{-}{-}{-}{-}{-}{-}{-}{-}{-}{-}{-}{-}{-}{-}{-}{-}}
\NormalTok{              MM;Gamma |{-} p1 op p2 : Boolean}

\NormalTok{(T{-}Eq)        MM;Gamma |{-} e1 : tau1}
\NormalTok{              MM;Gamma |{-} e2 : tau2}
\NormalTok{              eqJoin(tau1,tau2)=upsilon}
\NormalTok{              {-}{-}{-}{-}{-}{-}{-}{-}{-}{-}{-}{-}{-}{-}{-}{-}{-}{-}{-}{-}{-}{-}{-}{-}{-}{-}{-}{-}{-}{-}{-}{-}}
\NormalTok{              MM;Gamma |{-} e1 (= or \textless{}\textgreater{}) e2 : Boolean}

\NormalTok{(T{-}Ord)       MM;Gamma |{-} e1 : tau1}
\NormalTok{              MM;Gamma |{-} e2 : tau2}
\NormalTok{              compatibleOrdered(tau1,tau2)}
\NormalTok{              {-}{-}{-}{-}{-}{-}{-}{-}{-}{-}{-}{-}{-}{-}{-}{-}{-}{-}{-}{-}{-}{-}{-}{-}{-}{-}{-}{-}{-}{-}{-}{-}}
\NormalTok{              MM;Gamma |{-} e1 op e2 : Boolean}
\NormalTok{              where op in \{\textless{},\textless{}=,\textgreater{},\textgreater{}=\}}

\NormalTok{(T{-}Arith)     MM;Gamma |{-} e1 : tau1}
\NormalTok{              MM;Gamma |{-} e2 : tau2}
\NormalTok{              arithResult(op,tau1,tau2)=tau}
\NormalTok{              {-}{-}{-}{-}{-}{-}{-}{-}{-}{-}{-}{-}{-}{-}{-}{-}{-}{-}{-}{-}{-}{-}{-}{-}{-}{-}{-}{-}{-}{-}{-}{-}}
\NormalTok{              MM;Gamma |{-} e1 op e2 : tau}

\NormalTok{(T{-}If)        MM;Gamma |{-} c : Boolean}
\NormalTok{              MM;Gamma |{-} t : tau1}
\NormalTok{              MM;Gamma |{-} f : tau2}
\NormalTok{              ifJoin(tau1,tau2)=tau}
\NormalTok{              {-}{-}{-}{-}{-}{-}{-}{-}{-}{-}{-}{-}{-}{-}{-}{-}{-}{-}{-}{-}{-}{-}{-}{-}{-}{-}{-}{-}{-}{-}{-}{-}}
\NormalTok{              MM;Gamma |{-} if c then t else f : tau}

\NormalTok{(T{-}Let)       MM;Gamma |{-} init : tau1}
\NormalTok{              declType\_MM(T?,tau1)=tauD}
\NormalTok{              MM;Gamma[x:tauD] |{-} body : tau2}
\NormalTok{              {-}{-}{-}{-}{-}{-}{-}{-}{-}{-}{-}{-}{-}{-}{-}{-}{-}{-}{-}{-}{-}{-}{-}{-}{-}{-}{-}{-}{-}{-}{-}{-}}
\NormalTok{              MM;Gamma |{-} let x[:T?]=init in body : tau2}

\NormalTok{(T{-}Exists)    MM;Gamma |{-} S : Set(tauE)}
\NormalTok{              declType\_MM(T?,tauE)=tauD}
\NormalTok{              MM;Gamma[x:tauD] |{-} P : Boolean}
\NormalTok{              {-}{-}{-}{-}{-}{-}{-}{-}{-}{-}{-}{-}{-}{-}{-}{-}{-}{-}{-}{-}{-}{-}{-}{-}{-}{-}{-}{-}{-}{-}{-}{-}}
\NormalTok{              MM;Gamma |{-} S{-}\textgreater{}exists(x[:T?]|P) : Boolean}

\NormalTok{(T{-}ForAll)    same premises as T{-}Exists}
\NormalTok{              {-}{-}{-}{-}{-}{-}{-}{-}{-}{-}{-}{-}{-}{-}{-}{-}{-}{-}{-}{-}{-}{-}{-}{-}{-}{-}{-}{-}{-}{-}{-}{-}}
\NormalTok{              MM;Gamma |{-} S{-}\textgreater{}forAll(x|P) : Boolean}

\NormalTok{(T{-}Select)    MM;Gamma |{-} S : Set(tauE)}
\NormalTok{              declType\_MM(T?,tauE)=tauD}
\NormalTok{              MM;Gamma[x:tauD] |{-} P : Boolean}
\NormalTok{              {-}{-}{-}{-}{-}{-}{-}{-}{-}{-}{-}{-}{-}{-}{-}{-}{-}{-}{-}{-}{-}{-}{-}{-}{-}{-}{-}{-}{-}{-}{-}{-}}
\NormalTok{              MM;Gamma |{-} S{-}\textgreater{}select(x[:T?]|P) : Set(tauE)}

\NormalTok{(T{-}Reject)    same premises and result type as T{-}Select}

\NormalTok{(T{-}Collect)   MM;Gamma |{-} S : Set(tauE)}
\NormalTok{              declType\_MM(T?,tauE)=tauD}
\NormalTok{              MM;Gamma[x:tauD] |{-} E : sigma}
\NormalTok{              sigma is not of the form Set(\_)}
\NormalTok{              {-}{-}{-}{-}{-}{-}{-}{-}{-}{-}{-}{-}{-}{-}{-}{-}{-}{-}{-}{-}{-}{-}{-}{-}{-}{-}{-}{-}{-}{-}{-}{-}}
\NormalTok{              MM;Gamma |{-} S{-}\textgreater{}collect(x|E) : Set(sigma)}

\NormalTok{(T{-}IsUnique)  MM;Gamma |{-} S : Set(tauE)}
\NormalTok{              declType\_MM(T?,tauE)=tauD}
\NormalTok{              MM;Gamma[x:tauD] |{-} E : sigma}
\NormalTok{              sigma is not of the form Set(\_)}
\NormalTok{              {-}{-}{-}{-}{-}{-}{-}{-}{-}{-}{-}{-}{-}{-}{-}{-}{-}{-}{-}{-}{-}{-}{-}{-}{-}{-}{-}{-}{-}{-}{-}{-}}
\NormalTok{              MM;Gamma |{-} S{-}\textgreater{}isUnique(x|E) : Boolean}

\NormalTok{(T{-}Member)    collectionView(op,S)=Set(sigma)}
\NormalTok{              MM;Gamma |{-} e : delta}
\NormalTok{              memberJoin(sigma,delta)=upsilon}
\NormalTok{              op in \{includes,excludes\}}
\NormalTok{              {-}{-}{-}{-}{-}{-}{-}{-}{-}{-}{-}{-}{-}{-}{-}{-}{-}{-}{-}{-}{-}{-}{-}{-}{-}{-}{-}{-}{-}{-}{-}{-}}
\NormalTok{              MM;Gamma |{-} S{-}\textgreater{}op(e) : Boolean}

\NormalTok{(T{-}SetRel)    collectionView(op,S)=Set(tau)}
\NormalTok{              MM;Gamma |{-} T : Set(sigma)}
\NormalTok{              memberJoin(tau,sigma)=upsilon}
\NormalTok{              op in \{includesAll,excludesAll\}}
\NormalTok{              {-}{-}{-}{-}{-}{-}{-}{-}{-}{-}{-}{-}{-}{-}{-}{-}{-}{-}{-}{-}{-}{-}{-}{-}{-}{-}{-}{-}{-}{-}{-}{-}}
\NormalTok{              MM;Gamma |{-} S{-}\textgreater{}op(T) : Boolean}

\NormalTok{(T{-}SetComb)   collectionView(op,S)=Set(tau)}
\NormalTok{              MM;Gamma |{-} T : Set(sigma)}
\NormalTok{              setJoin2(tau,sigma)=upsilon}
\NormalTok{              op in \{union,intersection\}}
\NormalTok{              {-}{-}{-}{-}{-}{-}{-}{-}{-}{-}{-}{-}{-}{-}{-}{-}{-}{-}{-}{-}{-}{-}{-}{-}{-}{-}{-}{-}{-}{-}{-}{-}}
\NormalTok{              MM;Gamma |{-} S{-}\textgreater{}op(T) : Set(upsilon)}

\NormalTok{(T{-}AsSet)     collectionView(asSet,S)=Set(tau)}
\NormalTok{              {-}{-}{-}{-}{-}{-}{-}{-}{-}{-}{-}{-}{-}{-}{-}{-}{-}{-}{-}{-}{-}{-}{-}{-}{-}{-}{-}{-}{-}{-}{-}{-}}
\NormalTok{              MM;Gamma |{-} S{-}\textgreater{}asSet() : Set(tau)}

\NormalTok{(T{-}Card)      collectionView(size,S)=Set(tau)}
\NormalTok{              {-}{-}{-}{-}{-}{-}{-}{-}{-}{-}{-}{-}{-}{-}{-}{-}{-}{-}{-}{-}{-}{-}{-}{-}{-}{-}{-}{-}{-}{-}{-}{-}}
\NormalTok{              MM;Gamma |{-} S{-}\textgreater{}size() : Integer}

\NormalTok{(T{-}Empty)     collectionView(op,S)=Set(tau)}
\NormalTok{              op in \{isEmpty,notEmpty\}}
\NormalTok{              {-}{-}{-}{-}{-}{-}{-}{-}{-}{-}{-}{-}{-}{-}{-}{-}{-}{-}{-}{-}{-}{-}{-}{-}{-}{-}{-}{-}{-}{-}{-}{-}}
\NormalTok{              MM;Gamma |{-} S{-}\textgreater{}op() : Boolean}

\NormalTok{(T{-}TypePred)  MM;Gamma |{-} e : tau   tau=Void or tau is a class type D}
\NormalTok{              resolveClass\_MM(name)=C}
\NormalTok{              {-}{-}{-}{-}{-}{-}{-}{-}{-}{-}{-}{-}{-}{-}{-}{-}{-}{-}{-}{-}{-}{-}{-}{-}{-}{-}{-}{-}{-}{-}{-}{-}}
\NormalTok{              MM;Gamma |{-} e.oclIsKindOf(C) : Boolean}

\NormalTok{(T{-}Cast)      MM;Gamma |{-} e : tau   tau=Void or tau is a class type D}
\NormalTok{              resolveClass\_MM(name)=C}
\NormalTok{              {-}{-}{-}{-}{-}{-}{-}{-}{-}{-}{-}{-}{-}{-}{-}{-}{-}{-}{-}{-}{-}{-}{-}{-}{-}{-}{-}{-}{-}{-}{-}{-}}
\NormalTok{              MM;Gamma |{-} e.oclAsType(C) : C}
\end{Highlighting}
\end{Shaded}

The cast typing rule gives the successful result type; failed runtime
conformance yields \texttt{bottom} according to the validation policy.
\texttt{oclIsTypeOf} has no typing rule in the certified fragment.

The binder judgment records construction and typing simultaneously:

\begin{Shaded}
\begin{Highlighting}[]
\NormalTok{MM ; Gamma |{-} ast =\textgreater{} b : tau}
\end{Highlighting}
\end{Shaded}

It is syntax-directed. Every rule first recursively binds its children,
invokes the unique resolver required by the constructor, stores the
resolved UML reference and static type in \texttt{b}, and then applies
the corresponding typing rule above. Iterator and \texttt{let} bodies
are bound under \texttt{Gamma{[}x:tauD{]}}, where
\texttt{declType\_MM(T?,tauActual)=tauD}; after the body is bound, the
previous stack is restored. The environment is an argument of expression
binding; at invariant level the context fixes it:

\begin{Shaded}
\begin{Highlighting}[]
\NormalTok{T\_BIND\_EXPR(ast,MM,Gamma)=b}
\NormalTok{  iff exists unique tau: MM;Gamma |{-} ast =\textgreater{} b:tau.}

\NormalTok{T\_BIND\_INV(astInv,MM)=BInvariant(C,R,b)}
\NormalTok{  iff decodeInvariantShape(astInv)=(contextName,R,predAst)}
\NormalTok{  and resolveClass\_MM(contextName)=C}
\NormalTok{  and T\_BIND\_EXPR(predAst,MM,\{self:C\})=b}
\NormalTok{  and MM;\{self:C\} |{-} b:Boolean.}
\end{Highlighting}
\end{Shaded}

\textbf{B0 Binding Functionality.} If
\texttt{MM;Gamma\ \textbar{}-\ ast\ =\textgreater{}\ b1:tau1} and
\texttt{MM;Gamma\ \textbar{}-\ ast\ =\textgreater{}\ b2:tau2}, then
\texttt{b1=b2} and \texttt{tau1=tau2}.

Proof. Induct on the AST constructor. Child results are unique by the
induction hypotheses. Resolver results, \texttt{arithResult},
\texttt{ifJoin}, \texttt{eqJoin}, \texttt{setJoinN}, and
\texttt{setJoin2} are partial functions, so the parent constructor,
stored metadata, and result type are unique. \texttt{resolveType\_MM}
and \texttt{declType\_MM} are partial functions, so the declared lexical
extension is also unique for inferred and annotated binders. These cases
exhaust the successful rules.

\subsection{2.4 Complete Binder Construction
Rules}\label{24-complete-binder-construction-rules}

The binder output syntax used by the proof is explicit:

\begin{Shaded}
\begin{Highlighting}[]
\NormalTok{BInv ::= BInvariant(C,R,b)}

\NormalTok{b ::= BSelf(C) | BVar(x,tau) | BLiteral(c,tau)}
\NormalTok{    | BSetLiteral([b1,...,bn],Set(sigma))}
\NormalTok{    | BAttribute(b,a,tau)}
\NormalTok{    | BNavigation(b,A,fr,tr,dir,[bq1,...,bqk],resultKind,tauNav)}
\NormalTok{    | BAllInstances(C)}
\NormalTok{    | BNot(b) | BBinary(op,b,b,tau)}
\NormalTok{    | BIf(b,b,b,tau) | BLet(x,tauD,b,b,tau)}
\NormalTok{    | BIterator(op,b,kappa,x,tauD,b,tau)}
\NormalTok{    | BCollectionOp(op,b,kappa,[b1,...,bk],tau)}
\NormalTok{    | BTypeOp(op,b,C,tau)}
\end{Highlighting}
\end{Shaded}

\texttt{resolveProperty\_MM} is the disjoint classifier used for parser
\texttt{PropertyCallExp} nodes. It delegates to attribute or navigation
resolution and is functional: a successful admitted call has exactly one
classification. The following table is the complete definition of the
successful binding judgment. In every row, all displayed child binding
judgments are premises, and failure of any typing or resolver premise
makes the rule inapplicable.

{\def\LTcaptype{none} % do not increment counter
\begin{longtable}[]{@{}lll@{}}
\toprule\noalign{}
Rule & Source AST form and premises & Bound result and type \\
\midrule\noalign{}
\endhead
\bottomrule\noalign{}
\endlastfoot
B-InvRoot & \texttt{resolveClass\_MM(contextName)=C}; bind
\texttt{predicate=\textgreater{}b:Boolean} under
\texttt{Gamma0=\{self:C\}} &
\texttt{ContextInvariant(contextName,R,predicate)\ =\textgreater{}\ BInvariant(C,R,b)} \\
B-Self & \texttt{Gamma(self)=C} &
\texttt{MM;Gamma\ proves\ SelfExp\ =\textgreater{}\ BSelf(C)\ :\ C} \\
B-Var & \texttt{nearest(Gamma,x)=tau} &
\texttt{VariableExp(x)\ =\textgreater{}\ BVar(x,tau)\ :\ tau} \\
B-Lit & \texttt{typeOfLiteral(c)=tau} &
\texttt{LiteralExp(c)\ =\textgreater{}\ BLiteral(c,tau)\ :\ tau} \\
B-SetLit & \texttt{n\textgreater{}=1};
\texttt{ei=\textgreater{}bi:tau\_i};
\texttt{setJoinN({[}tau\_1,...,tau\_n{]})=tau}; no element type is a
collection &
\texttt{Set\{e1,...,en\}\ =\textgreater{}\ BSetLiteral({[}b1,...,bn{]},Set(tau))\ :\ Set(tau)} \\
B-Attr & \texttt{ast\ =\textgreater{}\ b:D};
\texttt{resolveProperty\_MM(D,n,{[}{]})\ =\ Attr(a,tau)} &
\texttt{PropertyCall(ast,n,{[}{]})\ =\textgreater{}\ BAttribute(b,a,tau)\ :\ tau} \\
B-Nav-One & \texttt{ast\ =\textgreater{}\ b:D}; each
\texttt{qi\ =\textgreater{}\ bqi:sigma\_i};
\texttt{resolveProperty\_MM(D,r,{[}sigma\_i{]})\ =\ Nav(A,fr,tr,dir,E,{[}sigma\_i{]},ONE)}
&
\texttt{PropertyCall(ast,r,{[}qi{]})\ =\textgreater{}\ BNavigation(b,A,fr,tr,dir,{[}bqi{]},ONE,E)\ :\ E} \\
B-Nav-Many & \texttt{ast\ =\textgreater{}\ b:D}; each
\texttt{qi\ =\textgreater{}\ bqi:sigma\_i};
\texttt{resolveProperty\_MM(D,r,{[}sigma\_i{]})\ =\ Nav(A,fr,tr,dir,E,{[}sigma\_i{]},MANY)}
&
\texttt{PropertyCall(ast,r,{[}qi{]})\ =\textgreater{}\ BNavigation(b,A,fr,tr,dir,{[}bqi{]},MANY,Set(E))\ :\ Set(E)} \\
B-AllInst & \texttt{resolveClass\_MM(name)=C} &
\texttt{MethodCall(ClassRef(name),allInstances,{[}{]})\ =\textgreater{}\ BAllInstances(C)\ :\ Set(C)} \\
B-Not & \texttt{p\ =\textgreater{}\ bp:Boolean} &
\texttt{NotExp(p)\ =\textgreater{}\ BNot(bp)\ :\ Boolean} \\
B-Bool & \texttt{ei\ =\textgreater{}\ bi:Boolean};
\texttt{op\ in\ \{and,or,xor,implies\}} &
\texttt{BinaryExp(op,e1,e2)\ =\textgreater{}\ BBinary(op,b1,b2,Boolean)\ :\ Boolean} \\
B-Eq & \texttt{ei\ =\textgreater{}\ bi:tau\_i};
\texttt{eqCompatible(tau1,tau2)};
\texttt{op\ in\ \{=,\textless{}\textgreater{}\}} &
\texttt{BinaryExp(op,e1,e2)\ =\textgreater{}\ BBinary(op,b1,b2,Boolean)\ :\ Boolean} \\
B-Ord & \texttt{ei\ =\textgreater{}\ bi:tau\_i};
\texttt{compatibleOrdered(tau1,tau2)};
\texttt{op\ in\ \{\textless{},\textless{}=,\textgreater{},\textgreater{}=\}}
&
\texttt{BinaryExp(op,e1,e2)\ =\textgreater{}\ BBinary(op,b1,b2,Boolean)\ :\ Boolean} \\
B-Arith & \texttt{ei\ =\textgreater{}\ bi:tau\_i};
\texttt{arithResult(op,tau1,tau2)=tau} &
\texttt{BinaryExp(op,e1,e2)\ =\textgreater{}\ BBinary(op,b1,b2,tau)\ :\ tau} \\
B-If & \texttt{c=\textgreater{}bc:Boolean};
\texttt{t=\textgreater{}bt:tau1}; \texttt{f=\textgreater{}bf:tau2};
\texttt{ifJoin(tau1,tau2)=tau} &
\texttt{IfExp(c,t,f)\ =\textgreater{}\ BIf(bc,bt,bf,tau)\ :\ tau} \\
B-Let & \texttt{init=\textgreater{}bi:tau1};
\texttt{declType\_MM(T?,tau1)=tauD}; \texttt{body=\textgreater{}bb:tau2}
under \texttt{Gamma{[}x:tauD{]}} &
\texttt{LetExp(x,T?,init,body)\ =\textgreater{}\ BLet(x,tauD,bi,bb,tau2)\ :\ tau2} \\
B-Exists & \texttt{S=\textgreater{}bS:Set(tauE)};
\texttt{kappa=Set(tauE)}; \texttt{declType\_MM(T?,tauE)=tauD};
\texttt{P=\textgreater{}bP:Boolean} under \texttt{Gamma{[}x:tauD{]}} &
\texttt{IteratorExp(exists,S,x,T?,P)\ =\textgreater{}\ BIterator(exists,bS,kappa,x,tauD,bP,Boolean)\ :\ Boolean} \\
B-ForAll & same iterator premises as B-Exists &
\texttt{BIterator(forAll,bS,kappa,x,tauD,bP,Boolean)\ :\ Boolean} \\
B-Select & same source/body environment; Boolean body &
\texttt{BIterator(select,bS,kappa,x,tauD,bP,Set(tauE))\ :\ Set(tauE)} \\
B-Reject & same source/body environment; Boolean body &
\texttt{BIterator(reject,bS,kappa,x,tauD,bP,Set(tauE))\ :\ Set(tauE)} \\
B-Collect & \texttt{S=\textgreater{}bS:Set(tauE)};
\texttt{kappa=Set(tauE)}; \texttt{declType\_MM(T?,tauE)=tauD};
\texttt{E=\textgreater{}bE:sigma} under \texttt{Gamma{[}x:tauD{]}};
\texttt{sigma\ !=\ Set(\_)} &
\texttt{BIterator(collect,bS,kappa,x,tauD,bE,Set(sigma))\ :\ Set(sigma)} \\
B-IsUnique & same declaration/source premises as B-Collect;
non-collection body &
\texttt{BIterator(isUnique,bS,kappa,x,tauD,bE,Boolean)\ :\ Boolean} \\
B-Member & \texttt{collectionView(op,S)=kappa=Set(sigma)};
\texttt{e=\textgreater{}be:delta};
\texttt{memberJoin(sigma,delta)=upsilon};
\texttt{op\ in\ \{includes,excludes\}} &
\texttt{CollectionOperationExp(op,S,{[}e{]})\ =\textgreater{}\ BCollectionOp(op,bS,kappa,{[}be{]},Boolean)\ :\ Boolean} \\
B-SetRel & \texttt{collectionView(op,S)=kappa=Set(tau)};
\texttt{T=\textgreater{}bT:Set(sigma)};
\texttt{memberJoin(tau,sigma)=upsilon};
\texttt{op\ in\ \{includesAll,excludesAll\}} &
\texttt{BCollectionOp(op,bS,kappa,{[}bT{]},Boolean)\ :\ Boolean} \\
B-SetComb & \texttt{collectionView(op,S)=kappa=Set(tau)};
\texttt{T=\textgreater{}bT:Set(sigma)};
\texttt{setJoin2(tau,sigma)=upsilon};
\texttt{op\ in\ \{union,intersection\}} &
\texttt{BCollectionOp(op,bS,kappa,{[}bT{]},Set(upsilon))\ :\ Set(upsilon)} \\
B-AsSet & \texttt{collectionView(asSet,S)=kappa=Set(tau)} &
\texttt{BCollectionOp(asSet,bS,kappa,{[}{]},Set(tau))} \\
B-Size & \texttt{collectionView(size,S)=kappa=Set(tau)} &
\texttt{CollectionOperationExp(size,S,{[}{]})\ =\textgreater{}\ BCollectionOp(size,bS,kappa,{[}{]},Integer)\ :\ Integer} \\
B-Empty & \texttt{collectionView(op,S)=kappa=Set(tau)};
\texttt{op\ in\ \{isEmpty,notEmpty\}} &
\texttt{BCollectionOp(op,bS,kappa,{[}{]},Boolean)\ :\ Boolean} \\
B-KindOf & \texttt{e=\textgreater{}be:tau}; \texttt{tau=Void} or class
type; \texttt{resolveClass\_MM(name)=C} &
\texttt{MethodCall(e,oclIsKindOf,{[}C{]})\ =\textgreater{}\ BTypeOp(kindOf,be,C,Boolean)\ :\ Boolean} \\
B-Cast & \texttt{e=\textgreater{}be:tau}; \texttt{tau=Void} or class
type; \texttt{resolveClass\_MM(name)=C} &
\texttt{MethodCall(e,oclAsType,{[}C{]})\ =\textgreater{}\ BTypeOp(cast,be,C,C)\ :\ C} \\
\end{longtable}
}

Here \texttt{comparable\ element\ types} abbreviates the premises of
T-Member or T-SetRel. The iterator body premise is derived only under
the displayed stack extension; it is not visible while binding the
source or after the rule returns. The table is exhaustive for
\texttt{OCL\_val}. Parser AST constructors or operation names not
matching a row have no successful binding derivation.
\texttt{TopLevelQuery} may be handled by the prototype, but it is not an
invariant root and is outside Theorem 6. Prototype parsing of
\texttt{oclIsTypeOf} is an experimental extension and has no successful
certified binding derivation.

The binder uses exactly the \texttt{collectionView} function defined in
Section 2.3; it does not infer a second, implementation-specific view.

Certified-closure side condition: a \texttt{BNavigation(...,D)} with
native scalar result may occur only as the immediate source child of a
\texttt{BCollectionOp} for one of the four operations above, whose
stored view is \texttt{Set(D)}. A direct iterator such as
\texttt{self.manager-\textgreater{}select(...)} is rejected because
native USE also rejects a scalar iterator source; the admitted spelling
is
\texttt{self.manager-\textgreater{}asSet()-\textgreater{}select(...)}.
Scalar chaining such as \texttt{self.manager.name}, scalar-navigation
equality, or applying a type operation directly to that navigation
remains outside the certified profile until it has separate binding,
lowering, and receiver-safety cases. The experimental general compiler
may still support those forms.

The proof constructors are an abstraction of the production records, not
a claim that Java declares classes named \texttt{BSelf},
\texttt{BAllInstances}, or \texttt{BTypeOp}. The total abstraction on
certified production Bound records is fixed by:

\begin{Shaded}
\begin{Highlighting}[]
\NormalTok{alpha\_B(BoundVariable("self",C))                   = BSelf(C)}
\NormalTok{alpha\_B(BoundVariable(x,tau)), x != "self"         = BVar(x,tau)}
\NormalTok{alpha\_B(BoundMethodCall(ClassRef(C),allInstances)) = BAllInstances(C)}
\NormalTok{alpha\_B(BoundMethodCall(e,oclIsKindOf,[C]))        = BTypeOp(kindOf,alpha\_B(e),C,Boolean)}
\NormalTok{alpha\_B(BoundMethodCall(e,oclAsType,[C]))          = BTypeOp(cast,alpha\_B(e),C,C)}
\end{Highlighting}
\end{Shaded}

The remaining production records map homomorphically to the
correspondingly named proof cases. In particular
\texttt{BoundLet.variableType} maps to \texttt{tauD}, while
\texttt{BoundIterator.sourceCollectionType} and
\texttt{BoundIterator.iteratorVariableType} map to \texttt{kappa} and
\texttt{tauD}. The semantic IR and Cypher plan retain the same
declaration-type payload. Java refinement evidence must check these
payloads rather than comparing constructor names alone.

\subsection{2.5 Complete Source/Bound Denotation
Correspondence}\label{25-complete-sourcebound-denotation-correspondence}

Source OCL and Bound OCL are given independent semantic algebras. Their
value domains are disjointly tagged copies of the theorem-supported
carriers:

\begin{Shaded}
\begin{Highlighting}[]
\NormalTok{Val\_S ::= SBool(b) | SInt(i) | SReal(r) | SString(s)}
\NormalTok{        | SEntity(o) | SSet(A) | SBottom}

\NormalTok{Val\_B ::= BBool(b) | BInt(i) | BReal(r) | BString(s)}
\NormalTok{        | BEntity(o) | BSet(A) | BBottom.}
\end{Highlighting}
\end{Shaded}

\texttt{Expr\_OCLval(tau)} and \texttt{Expr\_Bound(tau)} denote the
fibres selected by their unique typing derivations. Likewise
\texttt{Val\_S(tau)} and \texttt{Val\_B(tau)} are the corresponding
tagged value fibres: they contain the constructor for \texttt{tau} plus
the layer-local bottom tag; at \texttt{Set(sigma)} they contain finite
sets of \texttt{Val\_S(sigma)} or \texttt{Val\_B(sigma)}. These names
are typed restrictions of the displayed disjoint sums, not additional
value constructors.

The source semantics is the unique structural function:

\begin{Shaded}
\begin{Highlighting}[]
\NormalTok{[[.]]\_S\^{}tau : Expr\_OCLval(tau) x Model x Env\_S {-}\textgreater{} Val\_S(tau),}
\end{Highlighting}
\end{Shaded}

using source-local operations \texttt{truth\_S}, \texttt{not\_S},
\texttt{bool\_S}, \texttt{eq\_S}, \texttt{ord\_S}, \texttt{arith\_S},
\texttt{finite\_S}, \texttt{sources\_S}, \texttt{attr\_S},
\texttt{nav\_S}, and \texttt{type\_S}. The Bound semantics is separately
the unique structural function:

\begin{Shaded}
\begin{Highlighting}[]
\NormalTok{[[.]]\_B\^{}tau : Expr\_Bound(tau) x Model x Env\_B {-}\textgreater{} Val\_B(tau),}
\end{Highlighting}
\end{Shaded}

using distinct operations \texttt{truth\_B}, \texttt{not\_B},
\texttt{bool\_B}, \texttt{eq\_B}, \texttt{ord\_B}, \texttt{arith\_B},
\texttt{finite\_B}, \texttt{sources\_B}, \texttt{attr\_B},
\texttt{nav\_B}, and \texttt{type\_B}. Each family is defined directly
by recursion/case analysis on its own tagged domain. For example:

The \texttt{Model} argument carries its unique conformance metamodel
\texttt{metamodelOf(M)}. Under A4, \texttt{metamodelOf(M)=MM}. A
name-bearing source case first invokes the deterministic source resolver
over \texttt{metamodelOf(M)}; it does not read a resolved UML reference
from the source node. The corresponding Bound case reads the reference
stored by \texttt{T\_BIND}. B1/B3 prove that these two routes select the
same class, attribute, or navigation metadata. This implicit
model-to-metamodel projection is why the public source-evaluator
signature need not duplicate \texttt{MM} as a separate argument.

\begin{Shaded}
\begin{Highlighting}[]
\NormalTok{truth\_S(SBool(true)) = true;       truth\_S(v) = false otherwise}
\NormalTok{truth\_B(BBool(true)) = true;       truth\_B(v) = false otherwise}

\NormalTok{finite\_S(SSet(A)) = A;             finite\_S(SBottom) = empty}
\NormalTok{finite\_B(BSet(A)) = A;             finite\_B(BBottom) = empty}

\NormalTok{not\_S(v) = SBool(not truth\_S(v));  not\_B(w) = BBool(not truth\_B(w)).}

\NormalTok{view\_S(SSet(A))=A; view\_B(BSet(A))=A}
\NormalTok{view\_S(SBottom)=empty; view\_B(BBottom)=empty}
\NormalTok{view\_S(SEntity(o))=\{SEntity(o)\}; view\_B(BEntity(o))=\{BEntity(o)\}}

\NormalTok{collView\_S(v,Set(sigma))=finite\_S(v)}
\NormalTok{collView\_B(v,Set(sigma))=finite\_B(v)}
\NormalTok{collView\_S(v,D)=view\_S(v); collView\_B(v,D)=view\_B(v)}
\NormalTok{  only when D is in Class(MM) and the operator carries the certified}
\NormalTok{  direct{-}ONE collection{-}view marker}

\NormalTok{oneOrBottom\_S(empty)=SBottom}
\NormalTok{oneOrBottom\_S(\{SEntity(o)\})=SEntity(o)}
\NormalTok{oneOrBottom\_B(empty)=BBottom}
\NormalTok{oneOrBottom\_B(\{BEntity(o)\})=BEntity(o)}
\end{Highlighting}
\end{Shaded}

The uniform \texttt{collView\_X} boundary uses \texttt{finite\_X} for a
native \texttt{Set(sigma)} and uses the entity equations of
\texttt{view\_X} only at a certified direct collection boundary over a
native \texttt{ONE} navigation. It is not an implicit
collection-property navigation rule. \texttt{oneOrBottom\_X} is applied
only when the well-formed-model multiplicity premise gives a target set
of cardinality at most one; it is undefined otherwise, which places an
invalid model outside A4.

The remaining Boolean, scalar, collection, navigation, and type
operations are defined by the equations in the two corresponding columns
below, not by calling VA semantics. Define the cross-layer logical
relation recursively:

\begin{Shaded}
\begin{Highlighting}[]
\NormalTok{SBottom \textasciitilde{}=SB BBottom}
\NormalTok{SBool(b) \textasciitilde{}=SB BBool(b); SInt(i) \textasciitilde{}=SB BInt(i)}
\NormalTok{SReal(r) \textasciitilde{}=SB BReal(r); SString(s) \textasciitilde{}=SB BString(s)}
\NormalTok{SEntity(o) \textasciitilde{}=SB BEntity(o)}
\NormalTok{SSet(A) \textasciitilde{}=SB BSet(B)}
\NormalTok{  iff every a in A has a unique \textasciitilde{}=SB{-}related element in B and conversely.}
\end{Highlighting}
\end{Shaded}

Environments satisfy \texttt{rho\_S\ \textasciitilde{}=SB\ rho\_B}
pointwise. Define two total erasures with disjoint domains:

\begin{Shaded}
\begin{Highlighting}[]
\NormalTok{erase\_S\^{}tau : Val\_S(tau) {-}\textgreater{} ValBottom\_obj(tau)}
\NormalTok{erase\_B\^{}tau : Val\_B(tau) {-}\textgreater{} ValBottom\_obj(tau)}

\NormalTok{erase\_S(SBottom)=bottom             erase\_B(BBottom)=bottom}
\NormalTok{erase\_S(SBool(b))=b                 erase\_B(BBool(b))=b}
\NormalTok{erase\_S(SInt(i))=i                  erase\_B(BInt(i))=i}
\NormalTok{erase\_S(SReal(r))=r                 erase\_B(BReal(r))=r}
\NormalTok{erase\_S(SString(s))=s               erase\_B(BString(s))=s}
\NormalTok{erase\_S(SEntity(o))=o               erase\_B(BEntity(o))=o}
\NormalTok{erase\_S(SSet(A))=\{erase\_S(a)|a in A\}}
\NormalTok{erase\_B(BSet(B))=\{erase\_B(b)|b in B\}.}
\end{Highlighting}
\end{Shaded}

By structural induction on \texttt{\textasciitilde{}=SB}:

\begin{Shaded}
\begin{Highlighting}[]
\NormalTok{v\_S \textasciitilde{}=SB v\_B implies erase\_S(v\_S)=erase\_B(v\_B).        (SB{-}Erase)}
\end{Highlighting}
\end{Shaded}

For an untagged object environment \texttt{rho}, \texttt{tag\_S(rho)}
and \texttt{tag\_B(rho)} apply the corresponding tag recursively to each
bound value. Define the public object-side denotations as derived
notation:

\begin{Shaded}
\begin{Highlighting}[]
\NormalTok{[[s]]\_OCLval(M,rho)}
\NormalTok{  := erase\_S([[s]]\_S(M,tag\_S(rho)))}

\NormalTok{[[b]]\_Bound(M,rho)}
\NormalTok{  := erase\_B([[b]]\_B(M,tag\_B(rho))).}
\end{Highlighting}
\end{Shaded}

Thus equality in the theorem chain is well-typed equality in
\texttt{ValBottom\_obj(tau)}, not literal equality of independently
tagged values.

Write \texttt{S(s,rho\_S)={[}{[}s{]}{]}\_S(M,rho\_S)} and
\texttt{B(b,rho\_B)={[}{[}b{]}{]}\_B(M,rho\_B)}. Iterator and
\texttt{let} equations quantify only over the well-typed environments
established by the corresponding binder rule. In the source column
below, a resolved metavariable such as \texttt{a}, \texttt{A}, or
\texttt{C} means the unique result of resolving the displayed syntactic
name against \texttt{metamodelOf(M)} at that evaluation case; it is not
a field stored in the source expression.

{\def\LTcaptype{none} % do not increment counter
\begin{longtable}[]{@{}lll@{}}
\toprule\noalign{}
Source OCL\_val / Bound form & Source denotation & Bound denotation and
equality reason \\
\midrule\noalign{}
\endhead
\bottomrule\noalign{}
\endlastfoot
invariant root / \texttt{BInvariant(C,R,b)} &
\texttt{truth\_S(S(p,rho\_S{[}self-\textgreater{}SEntity(o){]}))} &
\texttt{truth\_B(B(b,rho\_B{[}self-\textgreater{}BEntity(o){]}))} \\
\texttt{SelfExp} / \texttt{BSelf(C)} & \texttt{rho\_S(self)} &
\texttt{rho\_B(self)} \\
\texttt{VariableExp(x)} / \texttt{BVar(x,tau)} & nearest
\texttt{rho\_S(x)} & resolved \texttt{rho\_B(x)} \\
\texttt{LiteralExp(c)} / \texttt{BLiteral(c,tau)} & source-tagged
literal \texttt{tag\_S(c)} & bound-tagged literal \texttt{tag\_B(c)} \\
\texttt{Set\{e1,...,en\}} / \texttt{BSetLiteral(...)} & \texttt{SSet} of
the canonically lifted child values using \texttt{setJoinN} &
\texttt{BSet} of the same independently lifted bound values; duplicate
elimination uses typed equality \\
attribute call / \texttt{BAttribute(b,a,tau)} &
\texttt{attr\_S(S(src,rho\_S),a)} by direct object-slot lookup &
\texttt{attr\_B(B(b,rho\_B),a)} by the resolved UML attribute
reference \\
directional qualified navigation with \texttt{resultKind=MANY} /
\texttt{BNavigation(...,MANY,Set(E))} & finite union of \texttt{nav\_S}
targets after evaluating qualifiers in declared order & corresponding
finite union of \texttt{nav\_B} targets using the stored association,
roles, direction, kind, and ordered qualifiers \\
directional qualified navigation with \texttt{resultKind=ONE} /
\texttt{BNavigation(...,ONE,E)} & \texttt{oneOrBottom\_S} of the same
target set & \texttt{oneOrBottom\_B} of the corresponding target set;
valid multiplicity ensures cardinality at most one \\
\texttt{C.allInstances()} / \texttt{BAllInstances(C)} &
\texttt{SSet(\{SEntity(o)\ such\ that\ classOf(o)\ conformsTo\ C\})} &
\texttt{BSet(\{BEntity(o)\ such\ that\ classOf(o)\ conformsTo\ resolved\ C\})} \\
\texttt{not\ p} / \texttt{BNot(bp)} & \texttt{not\_S(S(p,rho\_S))} &
\texttt{not\_B(B(bp,rho\_B))} \\
Boolean binary operation / \texttt{BBinary(op,...)} &
\texttt{bool\_S(op,S(e1,rho\_S),S(e2,rho\_S))} &
\texttt{bool\_B(op,B(b1,rho\_B),B(b2,rho\_B))} \\
equality/ordering / \texttt{BBinary(op,...)} & \texttt{eq\_S} or
\texttt{ord\_S} by source operand tags and source scalar table &
\texttt{eq\_B} or \texttt{ord\_B} by bound static types and bound scalar
table \\
arithmetic / \texttt{BBinary(op,...)} &
\texttt{arith\_S(op,S(e1,rho\_S),S(e2,rho\_S))} &
\texttt{arith\_B(op,B(b1,rho\_B),B(b2,rho\_B))} \\
\texttt{if\ c\ then\ t\ else\ f} / \texttt{BIf(...,tau)} & select the
source branch by \texttt{truth\_S}, then embed its value from the branch
type to \texttt{ifJoin} result \texttt{tau} & select the Bound branch by
\texttt{truth\_B} and apply the corresponding typed embedding; in
particular \texttt{Void} bottom becomes \texttt{bottom\_tau} \\
\texttt{let\ x{[}:T?{]}=init\ in\ body} / \texttt{BLet(x,tauD,...)} &
\texttt{S(body,rho\_S{[}x-\textgreater{}iota\_S{[}tau1-\textgreater{}tauD{]}(S(init,rho\_S)){]})}
&
\texttt{B(bb,rho\_B{[}x-\textgreater{}iota\_B{[}tau1-\textgreater{}tauD{]}(B(bi,rho\_B)){]})},
where \texttt{declType\_MM(T?,tau1)=tauD} \\
\texttt{exists} / \texttt{BIterator(exists,...,x,tauD,...)} &
\texttt{SBool(true)} iff some \texttt{s} in
\texttt{finite\_S(S(Src,rho\_S))} makes the body true under
\texttt{x-\textgreater{}iota\_S{[}tauE-\textgreater{}tauD{]}(s)} & the
corresponding bound quantification uses
\texttt{x-\textgreater{}iota\_B{[}tauE-\textgreater{}tauD{]}(s)}, where
\texttt{declType\_MM(T?,tauE)=tauD} \\
\texttt{forAll} / \texttt{BIterator(forAll,...)} & source universal
quantification over \texttt{finite\_S} & bound universal quantification
over \texttt{finite\_B} \\
\texttt{select}/\texttt{reject} / corresponding \texttt{BIterator} &
source set comprehension using \texttt{truth\_S} & bound set
comprehension using \texttt{truth\_B} \\
finite-set \texttt{collect} / corresponding \texttt{BIterator} &
\texttt{SSet} image under the source body & \texttt{BSet} image under
the bound body; no flattening or multiplicity \\
\texttt{isUnique} / corresponding \texttt{BIterator} & pairwise equality
of projected source values implies equality of source elements & the
identical injectivity test over the bound finite set and typed
equality \\
membership and set relations / \texttt{BCollectionOp} & source typed
membership, subset, or disjointness over \texttt{finite\_S} & bound
typed membership, subset, or disjointness over \texttt{finite\_B} \\
\texttt{union}/\texttt{intersection} / corresponding
\texttt{BCollectionOp} & lift both finite sets through
\texttt{setJoin2}, then take extensional union/intersection & perform
the same typed lift and finite-set operation on the independently tagged
bound domain \\
\texttt{asSet} / corresponding \texttt{BCollectionOp} &
\texttt{SSet(collView\_S(S(src,rho\_S),kappa))}, where \texttt{kappa} is
the binder-certified source type/view &
\texttt{BSet(collView\_B(B(bsrc,rho\_B),kappa))}, using the same
recorded \texttt{sourceCollectionType} \\
size and emptiness / \texttt{BCollectionOp} & source
cardinality/emptiness of \texttt{collView\_S(S(src,rho\_S),kappa)} &
bound cardinality/emptiness of
\texttt{collView\_B(B(bsrc,rho\_B),kappa)} \\
\texttt{oclIsKindOf(C)} / \texttt{BTypeOp(kindOf,...)} & source object
conformance to the syntactically denoted class & bound object
conformance to the resolved class reference \\
\texttt{oclAsType(C)} / \texttt{BTypeOp(cast,...)} & source entity if
conforming, otherwise \texttt{SBottom} & bound entity if conforming,
otherwise \texttt{BBottom} \\
\end{longtable}
}

\textbf{SB0 Independent-Operator Correspondence.} For every admitted
primitive operator of result type \texttt{tau}, if corresponding source
and Bound operands are pairwise \texttt{\textasciitilde{}=SB}-related,
the binder premises are derivable, and the reached operation is
\texttt{EvalClosed\_obj} (so every selected partial embedding and
cardinality representation is defined), then applying the source-local
and Bound-local operator produces \texttt{\textasciitilde{}=SB}-related
results. This statement includes Boolean operators (including
\texttt{xor}), typed scalar coercions and comparisons,
\texttt{SetLiteral}, the certified collection view, \texttt{union},
\texttt{intersection}, \texttt{asSet}, \texttt{isUnique}, attribute
access, \texttt{Navigation(ONE)}, \texttt{Navigation(MANY)},
\texttt{allInstances}, kind-of, and cast.

Proof. Use exhaustive case analysis for validation truth and each
admitted scalar row. Typed-lifting correspondence plus finite-set
extensionality proves the set cases. Resolver adequacy proves that both
navigation evaluators use the same metadata and target links; the
\texttt{ONE} case additionally commutes \texttt{oneOrBottom} by its
empty/singleton equations. Type cases use equality of the syntactic and
resolved UML denotations. No VA semantic function occurs in this proof.

For every table row, related child denotations and SB0 imply related
parent denotations. This is the local induction step used in Theorem 1;
B3 and B4 discharge navigation determinism and lexical scope.
\texttt{oclIsTypeOf} is not in the certified fragment because the
prototype has no independently validated direct runtime-class accessor;
it may be reintroduced only with \texttt{directType\_G} and a new
source/bound/graph realization case.

\begin{center}\rule{0.5\linewidth}{0.5pt}\end{center}

\section{3. Core Notation}\label{3-core-notation}

\subsection{3.1 Models and Encodings}\label{31-models-and-encodings}

\begin{Shaded}
\begin{Highlighting}[]
\NormalTok{MM = M2        UML domain metamodel}
\NormalTok{M = M1         UML object model conforming to MM}
\NormalTok{metamodelOf(M) unique UML metamodel to which M conforms; under A4 it is MM}
\NormalTok{G              repository property graph}
\NormalTok{G\_M2           schema projection containing encoded M2 declarations}
\NormalTok{G\_M1           instance projection containing encoded M1 objects and links}
\NormalTok{G\_typing       projection connecting M1 objects to M2 classes}
\NormalTok{Phi(MM,M)      repository encoding of M2 and M1}
\NormalTok{Obj(M)         finite set of UML objects in M}
\NormalTok{Class(MM)      finite set of classes in MM}
\NormalTok{Node(G)        finite set of graph nodes}
\NormalTok{Rel(G)         finite set of graph relationships}
\NormalTok{id(o)          stable object identifier}
\NormalTok{node(o)        graph node representing object o}
\NormalTok{classNode(C)   graph node representing class C}
\NormalTok{Rep\_G(o,n)     graph node n directly represents object o}
\NormalTok{ObjNode\_val(G) validation{-}visible graph nodes representing UML objects}
\end{Highlighting}
\end{Shaded}

The successful transformation chain is rooted at one complete invariant:

\begin{Shaded}
\begin{Highlighting}[]
\NormalTok{tInv                         : invariant text}
\NormalTok{astInv = ParseInvariant(tInv): invariant parser AST}
\NormalTok{I=(cName,R,e) = decodeInvariant\_val(astInv)}
\NormalTok{resolveClass\_MM(cName)=C}
\NormalTok{BInvariant(C,R,b) = T\_BIND\_INV(astInv,MM)}
\NormalTok{va  = T\_VA(b)}
\NormalTok{nva = T\_NORM(va)}
\NormalTok{plan = T\_CQM\^{}spec(C,R,nva)}
\NormalTok{(tq,pi) = T\_TEXT\^{}spec(plan)}
\NormalTok{q = parse\_Cypher(tq)}
\NormalTok{GraphAdequate(MM,M,G)}
\NormalTok{SpecPlanAdequacy(MM,C,nva,plan,q,pi,G)}
\end{Highlighting}
\end{Shaded}

Writing \texttt{G=Phi(MM,M)} is an optional specialization only after
\texttt{Phi} has been proved to establish the displayed
\texttt{GraphAdequate} premise.

Source denotation is defined on \texttt{OCL\_val}, not on parser AST
nodes:

\begin{Shaded}
\begin{Highlighting}[]
\NormalTok{[[e]]\_OCLval(M,rho)}
\NormalTok{  := erase\_S([[e]]\_S(M,tag\_S(rho))).}
\end{Highlighting}
\end{Shaded}

Parsing adequacy supplies
\texttt{decodeInvariant\_val(ParseInvariant(text(I)))=I}; this equation
connects text and AST representation to the source invariant without
assigning a denotation to either representation.

The source fragment, parser representation, and bound result are three
different levels:

\begin{Shaded}
\begin{Highlighting}[]
\NormalTok{encode\_ast : OCL\_val {-}\textgreater{} AST\_val}
\NormalTok{decode\_val : OCL\_AST {-}\textgreater{}partial OCL\_val}
\NormalTok{decode\_val(encode\_ast(s)) = s}

\NormalTok{AST\_val = \{ ast in OCL\_AST | decode\_val(ast) is defined \}}
\NormalTok{T\_BIND\_EXPR : AST\_val x MM x Gamma {-}\textgreater{}partial BoundOCL\_val(MM)}
\NormalTok{T\_BIND\_INV  : ASTInv\_val x MM {-}\textgreater{}partial BoundValInvariant(MM)}

\NormalTok{Adm\_MM(astInv) iff}
\NormalTok{  RawAdm(astInv)}
\NormalTok{  and decodeInvariant\_val(astInv)=I is defined}
\NormalTok{  and T\_BIND\_INV(astInv,MM)=bInv is defined}
\NormalTok{  and BoundAdm\_MM(bInv).}
\end{Highlighting}
\end{Shaded}

Logically, \texttt{G\ =\ G\_M2\ union\ G\_M1\ union\ G\_typing}. These
are projections of one repository graph, not necessarily separate
physical databases. Define \texttt{GraphMMView(G\_M2)} as the resolver
interface obtained from explicitly encoded M2 metadata. A5 requires it
to agree with \texttt{MM} for every reachable \texttt{resolveClass},
\texttt{resolveAttribute}, \texttt{resolveNavigation},
\texttt{resolveInheritance}, and \texttt{resolveQualifier} call.
Consequently \texttt{T\_BIND\_INV(astInv,GraphMMView(G\_M2))} and
\texttt{T\_BIND\_INV(astInv,MM)} produce isomorphic bound results. This
is schema recovery, not schema inference from M1 links.

The source fragment \texttt{OCL\_val} is defined before parsing.
\texttt{AST\_val} represents it in the parser metamodel. The optional
projection
\texttt{erase\_bound\ :\ BoundOCL\_val(MM)\ -\textgreater{}\ AST\_val}
removes resolved references and typing metadata. The operational
pipeline still runs from parsed AST to the bound model; these
definitions add no transformation stage.

\subsection{3.2 Canonical Semantic
Registry}\label{32-canonical-semantic-registry}

\textbf{Proof contract version: \texttt{PC-2026-07-22.3}.} This table is
the sole normative declaration of semantic-function symbols, domains,
and codomains. Later sections may give defining equations or
abbreviations, but may not overload a symbol with a different type. Its
structured source is the canonical JSON block in Section 1.1 of this
file. The external
\texttt{verification/contract/proof-contract-registry.json} and the
LaTeX table are synchronized projections.

{\def\LTcaptype{none} % do not increment counter
\begin{longtable}[]{@{}llll@{}}
\toprule\noalign{}
ID & Symbol & Total signature on the theorem domain & Canonical role \\
\midrule\noalign{}
\endhead
\bottomrule\noalign{}
\endlastfoot
SF-01 & \texttt{{[}{[}.{]}{]}\_S} &
\texttt{Expr\_OCLval(tau)\ x\ Model\ x\ Env\_S\ -\textgreater{}\ Val\_S(tau)}
& independently tagged admitted-source evaluator \\
SF-02 & \texttt{{[}{[}.{]}{]}\_B} &
\texttt{Expr\_Bound(tau)\ x\ Model\ x\ Env\_B\ -\textgreater{}\ Val\_B(tau)}
& independently tagged Bound OCL evaluator \\
SF-03 & \texttt{erase\_S} &
\texttt{Val\_S(tau)\ -\textgreater{}\ ValBottom\_obj(tau)} & erase
source tags into the common typed object-value universe \\
SF-04 & \texttt{erase\_B} &
\texttt{Val\_B(tau)\ -\textgreater{}\ ValBottom\_obj(tau)} & erase Bound
tags into the common typed object-value universe \\
SF-05 & \texttt{{[}{[}.{]}{]}\_OCLval} &
\texttt{Expr\_OCLval(tau)\ x\ Model\ x\ Env\_obj\ -\textgreater{}\ ValBottom\_obj(tau)}
& \texttt{erase\_S} after \texttt{tag\_S}; text and AST have no
denotation \\
SF-06 & \texttt{{[}{[}.{]}{]}\_Bound} &
\texttt{Expr\_Bound(tau)\ x\ Model\ x\ Env\_obj\ -\textgreater{}\ ValBottom\_obj(tau)}
& \texttt{erase\_B} after \texttt{tag\_B} \\
SF-07 & \texttt{{[}{[}.{]}{]}\_obj\^{}tau} &
\texttt{Expr\_VA(tau)\ x\ Model\ x\ Env\_obj\ -\textgreater{}\ ValBottom\_obj(tau)}
& typed object interpretation of VA \\
SF-08 & \texttt{{[}{[}.{]}{]}\_graph\^{}tau} &
\texttt{Expr\_VA(tau)\ x\ Graph\ x\ Env\_graph\ -\textgreater{}\ ValBottom\_G(tau)}
& typed graph interpretation of VA \\
SF-09 & \texttt{Denote\_graph\^{}tau} &
\texttt{PrimitiveAccess\_VA(tau)\ x\ Graph\ x\ Env\_graph\ -\textgreater{}\ ValBottom\_G(tau)}
& restriction of SF-08 to primitive graph access used by C3 \\
SF-10 & \texttt{encodeValue\_tau} &
\texttt{ValBottom\_obj(tau)\ -\textgreater{}\ ValBottom\_G(tau)} &
recursive object/scalar/set/bottom encoding \\
SF-11 & \texttt{finiteSet\_sigma} &
\texttt{ValBottom\_I(Set(sigma))\ -\textgreater{}\ P\_fin(ValBottom\_I(sigma))}
& collection-position bottom-to-empty policy \\
SF-12 & \texttt{asSources\_C} &
\texttt{ValBottom\_I(C)\ -\textgreater{}\ P\_fin(Val\_I(C))} & scalar
entity/bottom navigation-source lifting \\
SF-13 & \texttt{bool\_val} &
\texttt{ValBottom\_I(Boolean)\ -\textgreater{}\ Boolean} & true iff the
input is exactly validation true \\
SF-14 & \texttt{eq\_val\^{}tau} &
\texttt{ValBottom\_I(tau)\ x\ ValBottom\_I(tau)\ -\textgreater{}\ Boolean}
& typed equality including the declared bottom policy \\
SF-15 & \texttt{I\_attr\^{}tau} &
\texttt{ValBottom\_I(C)\ x\ Attribute(C,tau)\ -\textgreater{}\ ValBottom\_I(tau)}
& primitive attribute observation; bottom receiver maps to bottom \\
SF-16 & \texttt{I\_nav} &
\texttt{Val\_I(C)\ x\ NavigationMetadata\ -\textgreater{}\ P\_fin(Val\_I(D))}
& primitive binary directional navigation \\
SF-17 & \texttt{I\_class} &
\texttt{Class(C)\ -\textgreater{}\ P\_fin(Val\_I(C))} & finite
all-instances/conformance observation \\
SF-18 & \texttt{repr\_C\^{}tau} &
\texttt{ValBottom\_G(tau)\ -\textgreater{}\ CyVal\_C(tau)} &
type-indexed injective target representation \\
SF-19 & \texttt{abs\_C\^{}tau} &
\texttt{CyVal\_C(tau)\ -\textgreater{}\ ValBottom\_G(tau)} &
type-indexed target abstraction satisfying
\texttt{abs\_C\^{}tau(repr\_C\^{}tau(v))=v} \\
SF-20 & \texttt{EvalExpr\_G,pi} &
\texttt{CExpr(tau)\ x\ Row\ -\textgreater{}\ CyVal\_C(tau)} &
selected-dialect typed expression evaluation \\
SF-21 & \texttt{EvalRows\_G,pi} &
\texttt{CQuery\ x\ Bag(Row)\ -\textgreater{}\ Bag(Row)} &
selected-dialect row semantics \\
SF-22 & \texttt{Eval\_Cypher} &
\texttt{CQuery\ x\ Graph\ x\ ParamEnv\ -\textgreater{}\ Bag(Row)} &
\texttt{EvalRows\_G,pi(q,\{emptyRow\})} \\
SF-23 & \texttt{ExecPrelude} &
\texttt{List(CClause)\ x\ Graph\ x\ ParamEnv\ x\ Row\ -\textgreater{}\ Bag(Row)}
& \texttt{EvalRows\_G,pi(Seq(L),\{row\})} \\
SF-24 & \texttt{projectSemanticSet\_sigma} &
\texttt{Alias(sigma)\ x\ Bag(Row)\ -\textgreater{}\ P\_fin(ValBottom\_G(sigma))}
& abstracts one typed projected column and forgets row multiplicity \\
SF-25 & \texttt{paramsOf} &
\texttt{CompiledQuery\ -\textgreater{}\ ParamEnv} & second projection of
a generated \texttt{(CQuery,ParamEnv)} artifact \\
SF-26 & \texttt{Params} &
\texttt{CQuery\ -\textgreater{}\ P\_fin(ParamName)} & parameter names
occurring syntactically in the raw query \\
SF-27 & \texttt{returnedIds} &
\texttt{CQuery\ x\ Graph\ x\ ParamEnv\ -\textgreater{}\ P\_fin(ObjectId)}
& set of returned \texttt{use\_id} values, never a row table \\
SF-28 & \texttt{Viol\_OCL} &
\texttt{Expr\_OCLval(Boolean)\ x\ Class\ x\ Model\ -\textgreater{}\ P\_fin(Obj(M))}
& object-side invariant violation set \\
SF-29 & \texttt{navCollectionView} &
\texttt{ValBottom\_I(C)\ -\textgreater{}\ P\_fin(Val\_I(C))} &
singleton/empty finite-set view used only when a collection operator
directly consumes a to-one navigation \\
\end{longtable}
}

Here \texttt{CompiledQuery\ =\ CQuery\ x\ ParamEnv}, \texttt{I} ranges
over \texttt{\{obj,graph\}}, \texttt{P\_fin} denotes finite powerset,
and \texttt{Val\_G\^{}C}/\texttt{CyVal\_C} mean the supported
graph/Cypher value domains from Section 6.4, not arbitrary backend
values. Superscripts may be omitted only when the static typing
derivation uniquely determines them. In particular, \texttt{finiteSet},
\texttt{navCollectionView}, \texttt{asSources}, \texttt{eq\_val},
\texttt{I\_attr}, and both VA denotations are typed abbreviations, not
untyped overloaded functions.

The renderer produces a pair \texttt{(q,pi)}, not a parameter map
determined by query syntax. Two compilations may use alpha-equivalent
raw queries with different literal values, so no function
\texttt{CQuery\ -\textgreater{}\ ParamEnv} is assumed. When a compact
artifact notation is useful, it means only the pair projection:

\begin{Shaded}
\begin{Highlighting}[]
\NormalTok{paramsOf((q,pi)) = pi}
\NormalTok{returnedIds((q,pi),G) := returnedIds(q,G,pi).}
\end{Highlighting}
\end{Shaded}

Compound Cypher realization is not built into SF-09 or SF-22. It is
derived constructor-by-constructor in TXT4 and Theorem 5.

\subsection{3.3 Values and
Environments}\label{33-values-and-environments}

Validation values are:

\begin{Shaded}
\begin{Highlighting}[]
\NormalTok{Val ::= Boolean | Scalar | Entity | FiniteSet(ValueOrElementBottom)}
\NormalTok{      | WholeValueBottom}
\end{Highlighting}
\end{Shaded}

where \texttt{Entity} is interpreted as a UML object in the object
interpretation and as a graph node in the graph interpretation.

The raw grammar above is restricted by the static type. For
interpretation \texttt{I\ in\ \{obj,graph\}}, define typed values
recursively:

\begin{Shaded}
\begin{Highlighting}[]
\NormalTok{Val\_I(Void)    = empty}
\NormalTok{Val\_I(Boolean) = \{true,false\}}
\NormalTok{Val\_I(Integer) = Int64}
\NormalTok{Val\_I(Real)    = Real64}
\NormalTok{Val\_I(String)  = String\_C}
\NormalTok{Val\_obj(C)     = \{ o in Obj(M) | classOf(o) conformsTo C \}}
\NormalTok{Val\_graph(C)   = \{ node(o) | o in Val\_obj(C) \}}

\NormalTok{ValBottom\_I(Void)  = \{bottom\_Void\}}
\NormalTok{ValBottom\_I(sigma) = Val\_I(sigma) disjoint{-}union \{bottom\_sigma\}}

\NormalTok{Val\_I(Set(sigma))}
\NormalTok{  = P\_fin(ValBottom\_I(sigma))}

\NormalTok{ValBottom\_I(Set(sigma))}
\NormalTok{  = Val\_I(Set(sigma)) disjoint{-}union \{bottom\_Set(sigma)\}}
\end{Highlighting}
\end{Shaded}

Thus three values that must not be conflated are:

\begin{Shaded}
\begin{Highlighting}[]
\NormalTok{empty                         the empty finite set}
\NormalTok{\{bottom\_sigma\}                a set containing one bottom element}
\NormalTok{bottom\_Set(sigma)             a bottom value of collection type}
\end{Highlighting}
\end{Shaded}

The collection observation is totalized by:

\begin{Shaded}
\begin{Highlighting}[]
\NormalTok{finiteSet\_sigma(A)                   = A, for A in Val\_I(Set(sigma))}
\NormalTok{finiteSet\_sigma(bottom\_Set(sigma))   = empty.}
\end{Highlighting}
\end{Shaded}

Theorem environments and denotations are homogeneous by static type.
Retaining \texttt{bottom\_sigma} inside a finite-set image is the stated
\texttt{collect} policy; mapping the distinct whole-collection value
\texttt{bottom\_Set(sigma)} to the empty set is the collection-boundary
policy. A subscript on bottom or \texttt{finiteSet} may be omitted only
when the static typing derivation fixes it uniquely.

The selected scalar profile is the following mathematical domain:

\begin{Shaded}
\begin{Highlighting}[]
\NormalTok{Int64 = \{ n in Z | {-}2\^{}63 \textless{}= n \textless{}= 2\^{}63{-}1 \}}

\NormalTok{Real64 = \{ finite IEEE{-}754 binary64 values \}}
\NormalTok{         / identify {-}0.0 and +0.0 for numeric equality}

\NormalTok{String\_C = finite Unicode scalar{-}value sequences compared by exact code{-}point}
\NormalTok{           equality; no String ordering is admitted.}
\end{Highlighting}
\end{Shaded}

\texttt{NaN}, infinities, temporal/spatial values, byte arrays, locale
collation, and backend-specific scalar values are excluded. Define
\texttt{round64} as IEEE-754 round-to-nearest, ties-to-even. The
admitted operators are total only under the displayed side conditions:

{\def\LTcaptype{none} % do not increment counter
\begin{longtable}[]{@{}ll@{}}
\toprule\noalign{}
Operation & Exact theorem-covered definition and side condition \\
\midrule\noalign{}
\endhead
\bottomrule\noalign{}
\endlastfoot
Integer \texttt{+,-,*} & mathematical result in \texttt{Int64};
otherwise the evaluation is not \texttt{ScalarClosed} \\
Real \texttt{+,-,*} & \texttt{round64} of the exact real operation,
provided the result is finite \\
\texttt{/} & result type \texttt{Real64}; divisor is numerically
non-zero; operands are converted as below; rounded result must be
finite \\
Integer-to-Real & admitted only for \texttt{n} with
\texttt{abs(n)\ \textless{}=\ 2\^{}53}, so conversion is exact \\
mixed numeric operation & binder selects \texttt{Real64}; the Integer
operand satisfies exact conversion; then use the Real rule \\
numeric equality/order & compare after the same admitted common-type
conversion; \texttt{-0.0} equals \texttt{+0.0} \\
Boolean equality & ordinary two-valued equality \\
String equality & exact Unicode code-point-sequence equality \\
Entity equality & source object identity and corresponding graph-node
identity \\
\end{longtable}
}

For an invariant predicate \texttt{e} in context \texttt{C}, define:

\begin{Shaded}
\begin{Highlighting}[]
\NormalTok{ScalarClosed(C,e,M)}
\NormalTok{iff}
\NormalTok{for every o in Obj(M) with classOf(o) conformsTo C and every scalar}
\NormalTok{subexpression s of e,}
\NormalTok{  evaluation of s under every iterator/let environment reachable from}
\NormalTok{  self{-}\textgreater{}o satisfies the applicable row above and returns a profile value or}
\NormalTok{  the explicitly supported bottom value.}
\end{Highlighting}
\end{Shaded}

For a local typed evaluation at Source, Bound, or VA stage,
\texttt{EvalClosed\_I(v,rho)} abbreviates the same condition for every
arithmetic operation, canonical coercion, and finite-set cardinality
actually reached while evaluating \texttt{v} under \texttt{rho}; a
cardinality result must lie in \texttt{Int64}.
\texttt{ScalarClosed(C,e,M)} entails the object-side local predicate for
every VA term and environment reachable from the admitted invariant
chain; graph-side closure is transferred by G2 for finite-set
cardinality, G3 for scalar operations, and G3a for coercions.

The object evaluator and selected Cypher profile are assumed to
implement these same equations. This is now an explicit dynamic domain
premise, not the circular condition ``both sides happen to return the
same result.'' Scalar serialization is injective on this domain;
\texttt{semEq}, \texttt{semOrd}, and \texttt{semArith} may use native
Cypher operators only after typing and \texttt{ScalarClosed} discharge
their premises.

Encoding of values is defined recursively at one fixed static type:

\begin{Shaded}
\begin{Highlighting}[]
\NormalTok{encodeValue\_C(o) = node(o), if o is a UML object of type C}
\NormalTok{encodeValue\_sigma(s) = s, if s is a non{-}bottom scalar of type sigma}
\NormalTok{encodeValue\_tau(bottom\_tau) = bottom\_tau\^{}G}
\NormalTok{encodeValue\_Set(sigma)(A) = \{ encodeValue\_sigma(v) | v in A \}}
\NormalTok{encodeValue\_Set(sigma)(bottom\_Set(sigma)) = bottom\_Set(sigma)\^{}G}
\NormalTok{encodeSet\_sigma(S) = \{ encodeValue\_sigma(v) | v in S \}}
\end{Highlighting}
\end{Shaded}

There is no unindexed cross-type \texttt{encodeValue} equality. A bare
occurrence is an abbreviation whose unique type index comes from the
typing derivation.

An object environment satisfies a typing environment when:

\begin{Shaded}
\begin{Highlighting}[]
\NormalTok{rho |=\_M Gamma}
\NormalTok{iff}
\NormalTok{dom(rho) = dom(Gamma)}
\NormalTok{and for every nearest binding x:tau in Gamma,}
\NormalTok{    rho(x) in ValBottom\_obj(tau).}
\end{Highlighting}
\end{Shaded}

Define \texttt{eta\ \textbar{}=\_G\ Gamma} analogously using
\texttt{ValBottom\_G}. For \texttt{I\ in\ \{obj,graph\}}, the notation
\texttt{zeta\ \textbar{}=\_I\ Gamma} means respectively
\texttt{zeta\ \textbar{}=\_M\ Gamma} or
\texttt{zeta\ \textbar{}=\_G\ Gamma}; it introduces no third environment
relation. Object and graph environments are related by
\texttt{rho\ \textasciitilde{}Phi\_Gamma\ eta} iff their domains equal
\texttt{dom(Gamma)} and, for every nearest visible binding
\texttt{x:tau} in \texttt{Gamma}:

\begin{Shaded}
\begin{Highlighting}[]
\NormalTok{eta(x) = encodeValue\_tau(rho(x)).}
\end{Highlighting}
\end{Shaded}

Environment extension preserves the relation:

\begin{Shaded}
\begin{Highlighting}[]
\NormalTok{rho \textasciitilde{}Phi\_Gamma eta}
\NormalTok{and v in ValBottom\_obj(tau)}
\NormalTok{implies}
\NormalTok{rho[x {-}\textgreater{} v] \textasciitilde{}Phi\_Gamma[x:tau]}
\NormalTok{  eta[x {-}\textgreater{} encodeValue\_tau(v)].}
\end{Highlighting}
\end{Shaded}

Navigation uses the source-lifting helper \texttt{asSources}:

\begin{Shaded}
\begin{Highlighting}[]
\NormalTok{asSources(v) = \{ v \} if v is an Entity}
\NormalTok{asSources(bottom) = empty}
\end{Highlighting}
\end{Shaded}

Its domain is the scalar receiver fibre \texttt{ValBottom\_I(C)},
because certified navigation has no \texttt{Set(C)} receiver rule.
Consequently every non-bottom value returned by \texttt{asSources\_C} is
an entity. This helper is distinct from \texttt{finiteSet}: singleton
lifting is used for navigation receivers, whereas collection operators
require a collection value or \texttt{bottom}.

For collection operators, the proof uses the validation-level helpers
\texttt{finiteSet} and \texttt{navCollectionView}:

\begin{Shaded}
\begin{Highlighting}[]
\NormalTok{finiteSet(Set(S)) = S}
\NormalTok{finiteSet(bottom) = empty}
\NormalTok{navCollectionView(o) = \{o\}, when o is the present result of a to{-}one navigation}
\NormalTok{navCollectionView(bottom) = empty, for an absent optional to{-}one navigation}
\end{Highlighting}
\end{Shaded}

Here \texttt{kappa=Set(tau)} equals the source type at every iterator
boundary and at every collection-operation boundary whose source is a
native Set. Only the four operations \texttt{asSet}, \texttt{size},
\texttt{isEmpty}, and \texttt{notEmpty} may instead record the
singleton/empty \texttt{Set(D)} view of a directly consumed
native-scalar to-one navigation. This payload must be preserved by
Bound-to-VA, optimization, planning, and rendering; it is never
reconstructed from multiplicity.

For ordinary well-typed collection expressions in \texttt{OCL\_val},
collection-valued denotations are either finite sets or \texttt{bottom},
so \texttt{finiteSet} is used by the collection operators. If one of
those four operations records a direct to-one navigation view, it uses
\texttt{navCollectionView} instead. Treating an absent optional target
as the empty set in that position is part of the validation-level policy
of this theorem, not a claim about unrestricted OMG OCL null/invalid
semantics.

The object and graph value universes are disjoint tagged unions. In
particular, object entities, graph entities, scalars, sets, and
\texttt{bottom} cannot collide across value sorts. Equality on
theorem-supported values is typed equality:

\begin{Shaded}
\begin{Highlighting}[]
\NormalTok{eq\_val(v1,v2) = true iff v1 and v2 have the same value sort and}
\NormalTok{                         denote the same value.}

\NormalTok{For object entities: eq\_val(o1,o2) iff o1=o2.}
\NormalTok{For graph entities:  eq\_val(n1,n2) iff n1=n2.}
\end{Highlighting}
\end{Shaded}

Ordered comparisons remain restricted to compatible ordered scalar
types.

\subsection{3.4 Violation Sets}\label{34-violation-sets}

Object-side violation set:

\begin{Shaded}
\begin{Highlighting}[]
\NormalTok{Viol\_OCL(e,C,M)}
\NormalTok{= \{ o in Obj(M)}
\NormalTok{    | classOf(o) conformsTo C}
\NormalTok{      and bool\_val([[e]]\_OCLval(M, self {-}\textgreater{} o)) = false \}.}
\end{Highlighting}
\end{Shaded}

Cypher-side returned-id set:

\begin{Shaded}
\begin{Highlighting}[]
\NormalTok{returnedIds(q,G,pi)}
\NormalTok{= set of use\_id values returned by Eval\_Cypher(q,G,pi).}
\end{Highlighting}
\end{Shaded}

Cypher rows may contain duplicates. The comparison is always set-based:

\begin{Shaded}
\begin{Highlighting}[]
\NormalTok{id(o) in returnedIds(q,G,pi)}
\end{Highlighting}
\end{Shaded}

means membership in the set of returned identifiers. Implementations
should use \texttt{RETURN\ DISTINCT\ self.use\_id} for invariant
queries.

\begin{center}\rule{0.5\linewidth}{0.5pt}\end{center}

\section{4. Graph Encoding Contract}\label{4-graph-encoding-contract}

The representation theorem assumes an adequate graph encoding:

\begin{Shaded}
\begin{Highlighting}[]
\NormalTok{Phi : (MM,M) {-}\textgreater{} G.}
\end{Highlighting}
\end{Shaded}

The implementation can be tested against this contract, but the
mathematical theorems take it as an explicit assumption.

\subsection{4.0 Repository Levels and Metamodel
View}\label{40-repository-levels-and-metamodel-view}

The canonical repository graph stores both modeling levels used by
validation:

\begin{Shaded}
\begin{Highlighting}[]
\NormalTok{G = G\_M2 union G\_M1 union G\_typing.}
\end{Highlighting}
\end{Shaded}

Let \texttt{modelKey\_MM} be the nonblank repository scope of
\texttt{MM} (the production profile uses \texttt{MM.name}). Every
validation-visible schema node, object node, attribute-value node, and
\texttt{Link*} relationship created for this model carries
\texttt{modelKey=modelKey\_MM}. Every generated accessor conjuncts that
property with its canonical-key lookup. A row declared under another
model but carrying or referencing a current canonical key is an encoding
violation, not an invisible foreign row; the adapter snapshot therefore
scans both the declared model scope and the current canonical-key
namespace.

\texttt{G\_M2} contains the M2 declarations needed by semantic binding:
injective class, attribute, and association keys; attribute types;
binary association ends and roles; qualifier declarations; and
inheritance. \texttt{G\_M1} contains object nodes, attribute values, and
association-link instances. \texttt{G\_typing} contains the materialized
object-to-class membership used by context enumeration and
\texttt{allInstances}.

A graph-backed binder may inspect schema relationships in \texttt{G\_M2}
to resolve an association end, role, target type, qualifier, or
inheritance fact. It must not inspect M1 link instances to guess a
property meaning or specialize generated Cypher to the current dataset.
M1 relationships are traversed only when the generated query executes.
The original UML file need not remain available once
\texttt{GraphMMView(G\_M2)} has been shown resolver-equivalent to
\texttt{MM}.

\subsection{4.1 Class Nodes}\label{41-class-nodes}

For every class \texttt{C\ in\ Class(MM)}, \texttt{Phi} creates exactly
one class node:

\begin{Shaded}
\begin{Highlighting}[]
\NormalTok{classNode(C)}
\NormalTok{label(classNode(C), "UmlClass")}
\NormalTok{prop(classNode(C), "classKey") = key\_MM(C)}
\NormalTok{prop(classNode(C), "name") = C.name}
\NormalTok{prop(classNode(C), "modelKey") = modelKey\_MM}
\end{Highlighting}
\end{Shaded}

No two distinct classes share a class node.

\texttt{key\_MM} is an injective, namespace-aware class key:

\begin{Shaded}
\begin{Highlighting}[]
\NormalTok{key\_MM(C1) = key\_MM(C2) iff C1 = C2.}
\end{Highlighting}
\end{Shaded}

It may be a qualified name or stable metamodel identifier. An
implementation that uses the unqualified \texttt{name} property as the
lookup key is covered only when class names are globally unique in
\texttt{MM}; otherwise generated context, \texttt{allInstances}, and
type-operation patterns must use \texttt{classKey}.

\subsection{4.2 Object Nodes and
Identity}\label{42-object-nodes-and-identity}

Define the representation relation:

\begin{Shaded}
\begin{Highlighting}[]
\NormalTok{Rep\_G(o,n)}
\NormalTok{iff}
\NormalTok{  n is an object node generated by Phi from o}
\NormalTok{  and prop(n,"use\_id") = id(o)}
\NormalTok{  and prop(n,"modelKey") = modelKey\_MM.}
\end{Highlighting}
\end{Shaded}

For every object \texttt{o\ in\ Obj(M)}, \texttt{Phi} creates exactly
one representing node:

\begin{Shaded}
\begin{Highlighting}[]
\NormalTok{exists unique n in Node(G): Rep\_G(o,n).}
\end{Highlighting}
\end{Shaded}

Conversely, every validation-visible object node represents exactly one
source object. Define \texttt{ObjNode\_val(G)} as the nodes that can be
enumerated as UML objects by invariant contexts, navigation, or
\texttt{allInstances}. The exactness condition is:

\begin{Shaded}
\begin{Highlighting}[]
\NormalTok{for every n in ObjNode\_val(G),}
\NormalTok{  exists unique o in Obj(M): Rep\_G(o,n).}

\NormalTok{Equivalently:}
\NormalTok{ObjNode\_val(G) = \{ node(o) | o in Obj(M) \}.}
\end{Highlighting}
\end{Shaded}

This no-spurious-object-node direction is required by the reverse
implication of Theorem 6.

\texttt{node(o)} denotes this unique node. Therefore:

\begin{Shaded}
\begin{Highlighting}[]
\NormalTok{Rep\_G(o,node(o))}
\NormalTok{prop(node(o),"use\_id") = id(o).}
\end{Highlighting}
\end{Shaded}

The identifier function is injective:

\begin{Shaded}
\begin{Highlighting}[]
\NormalTok{id(o1) = id(o2) iff o1 = o2.}
\end{Highlighting}
\end{Shaded}

Therefore \texttt{node} is injective over object nodes.

\subsection{4.3 Type Membership and
Inheritance}\label{43-type-membership-and-inheritance}

The prototype has two intentionally distinct relationship names. They
are not aliases and must not be interchanged:

{\def\LTcaptype{none} % do not increment counter
\begin{longtable}[]{@{}lll@{}}
\toprule\noalign{}
Relationship type & Model-level meaning & Visible to OCL type
accessors \\
\midrule\noalign{}
\endhead
\bottomrule\noalign{}
\endlastfoot
\texttt{ObjectInstanceOf} & M0 \texttt{Object} conforms to an M1
\texttt{UmlClass}, including materialized supertypes & yes \\
\texttt{InstanceOf} & generic schema instantiation, principally M1
elements to M2 \texttt{MetaNode}; also \texttt{AttributeValue} to its
\texttt{Attribute} definition & no \\
\end{longtable}
}

\texttt{CanonicalGraphVocabulary.OBJECT\_INSTANCE\_OF} and
\texttt{CanonicalGraphVocabulary.SCHEMA\_INSTANCE\_OF} are the
executable vocabulary constants. In particular, context matching,
\texttt{oclIsKindOf}, \texttt{oclAsType}, and \texttt{allInstances} must
use \texttt{ObjectInstanceOf}. A pattern
\texttt{(o:Object)-{[}:InstanceOf{]}-\textgreater{}(c:UmlClass)} is
outside the canonical profile and is rejected by generated-query
conformance checks.

Type membership is materialized. For every object \texttt{o} whose
runtime class is \texttt{D}, and for every class \texttt{C} such that
\texttt{D\ conformsTo\ C}, the graph contains:

\begin{Shaded}
\begin{Highlighting}[]
\NormalTok{(node(o)){-}[:ObjectInstanceOf]{-}\textgreater{}(classNode(C)).}
\end{Highlighting}
\end{Shaded}

No-spurious-type-edge condition:

\begin{Shaded}
\begin{Highlighting}[]
\NormalTok{Every ObjectInstanceOf edge used by validation is generated from a true}
\NormalTok{conformance fact classOf(o) conformsTo C.}
\end{Highlighting}
\end{Shaded}

Thus:

\begin{Shaded}
\begin{Highlighting}[]
\NormalTok{o conformsTo C}
\NormalTok{iff}
\NormalTok{(node(o)){-}[:ObjectInstanceOf]{-}\textgreater{}(classNode(C)).}
\end{Highlighting}
\end{Shaded}

The graph interpretation of \texttt{allInstances} is:

\begin{Shaded}
\begin{Highlighting}[]
\NormalTok{allInstances\_graph(C)}
\NormalTok{= \{ n | (n){-}[:ObjectInstanceOf]{-}\textgreater{}(classNode(C)) \}.}
\end{Highlighting}
\end{Shaded}

No transitive class-hierarchy traversal is needed in the proof. If an
implementation stores only direct runtime type edges, it must
materialize supertype memberships before validation or it is outside
this proof for inherited \texttt{allInstances}.

For exact type checks, the graph encoding must additionally provide a
direct runtime-class accessor:

\begin{Shaded}
\begin{Highlighting}[]
\NormalTok{directType\_G(node(o)) = classNode(classOf(o)).}
\end{Highlighting}
\end{Shaded}

If this accessor is not implemented, \texttt{oclIsTypeOf} is outside the
theorem and only \texttt{oclIsKindOf} is covered by materialized
conformance edges.

\subsection{4.4 Attribute Encoding}\label{44-attribute-encoding}

Attributes are accessed through one canonical accessor:

\begin{Shaded}
\begin{Highlighting}[]
\NormalTok{value\_G(node(o), a) = attrVal\_M(o,a).}
\end{Highlighting}
\end{Shaded}

Each admitted UML attribute has an injective, namespace-aware key
\texttt{attributeKey\_MM(a)}:

\begin{Shaded}
\begin{Highlighting}[]
\NormalTok{attributeKey\_MM(a1)=attributeKey\_MM(a2) iff a1=a2.}
\end{Highlighting}
\end{Shaded}

The concrete graph may implement \texttt{value\_G} as direct node
properties, attribute-value nodes, or another fixed representation. If
multiple storage forms coexist, they must be observationally consistent
through \texttt{value\_G}.

The current canonical profile stores a scalar slot as tagged text and
defines \texttt{value\_G} by strict decoding with the declared UML
attribute type:

\begin{Shaded}
\begin{Highlighting}[]
\NormalTok{encScalar\_tau(bottom) = "v1|V"}
\NormalTok{encScalar\_String(s)   = "v1|S|" ++ esc(s)}
\NormalTok{encScalar\_Integer(i)  = "v1|I|" ++ canonicalInt64(i)}
\NormalTok{encScalar\_Real(r)     = "v1|R|" ++ canonicalFiniteReal64(r)}
\NormalTok{encScalar\_Boolean(b)  = "v1|B|" ++ (b ? "true" : "false")}
\NormalTok{encScalar\_Enum(x)     = "v1|E|" ++ esc(x)   {-}{-} reserved extension domain}

\NormalTok{esc(s) = replace(replace(s,"\%","\%25"),"|","\%7C")}
\NormalTok{value\_G(n,a) = decScalar\_type(a)(storedPayload(n,a)).}
\end{Highlighting}
\end{Shaded}

\texttt{decScalar} accepts \texttt{v1\textbar{}V} as bottom for every
supported scalar type and otherwise requires the tag selected by the
static type. It rejects null text, legacy sentinel strings, malformed
escapes, cross-type tags, noncanonical or nonfinite Real64 text, and
integers outside Int64. \texttt{esc} preserves quotes, whitespace, and
Unicode exactly. Hence ordinary Strings such as \texttt{Undefined},
\texttt{COLLECTION\_EMPTY}, and \texttt{O\textquotesingle{}Brien},
including values containing \texttt{\%} or \texttt{\textbar{}}, are
disjoint from bottom and round-trip without normalization. Enum storage
is defined by the codec but remains outside \texttt{OCL\_val} until its
typing/equality cases are admitted.

\subsection{4.5 Association Link
Encoding}\label{45-association-link-encoding}

For each binary association link, qualifiers are retained separately for
the two declared association ends. The predicate and its
association-indexed extension are:

\begin{Shaded}
\begin{Highlighting}[]
\NormalTok{link\_M(A, o1, sourceRole, sourceQs, o2, targetRole, targetQs)}

\NormalTok{links\_M(A)}
\NormalTok{  = \{ (o1,sourceRole,sourceQs,o2,targetRole,targetQs)}
\NormalTok{      | link\_M(A,o1,sourceRole,sourceQs,o2,targetRole,targetQs) \}.}
\end{Highlighting}
\end{Shaded}

\texttt{Phi} creates at least one validation-visible graph relationship
\texttt{r} such that:

\begin{Shaded}
\begin{Highlighting}[]
\NormalTok{src(r) = node(o1)}
\NormalTok{trg(r) = node(o2)}
\NormalTok{type(r) starts with "Link"}
\NormalTok{prop(r,"modelKey") = modelKey\_MM}
\NormalTok{linkAssociation\_G(r) = associationKey\_MM(A)}
\NormalTok{linkSourceRole\_G(r) = sourceRole}
\NormalTok{linkTargetRole\_G(r) = targetRole}
\NormalTok{linkQualifiers\_G(r,forward) = encodeQualifierList(sourceQs)}
\NormalTok{linkQualifiers\_G(r,reverse) = encodeQualifierList(targetQs)}
\end{Highlighting}
\end{Shaded}

\texttt{associationKey\_MM} is injective over admitted associations. It
may be a qualified association name. The proof uses the logical accessor
\texttt{linkAssociation\_G}; a concrete property called \texttt{name} or
\texttt{associationName} is not part of this identity contract.

The logical link accessors are total on validation-visible
\texttt{Link*} relationships. In the current prototype adapter they are
realized by:

\begin{Shaded}
\begin{Highlighting}[]
\NormalTok{linkAssociation\_G(r)          = prop(r,"associationKey")}
\NormalTok{linkSourceRole\_G(r)           = prop(r,"sourceRole")}
\NormalTok{linkTargetRole\_G(r)           = prop(r,"targetRole")}
\NormalTok{linkQualifiers\_G(r,forward)   = prop(r,"sourceQualifiers")}
\NormalTok{linkQualifiers\_G(r,reverse)   = prop(r,"targetQualifiers")}
\end{Highlighting}
\end{Shaded}

The domain quantified by representation and navigation lemmas is not all
graph relationships:

\begin{Shaded}
\begin{Highlighting}[]
\NormalTok{RelLink\_val(G,modelKey\_MM)}
\NormalTok{  = \{ r in Rel(G)}
\NormalTok{      | type(r) starts with "Link"}
\NormalTok{        and prop(r,"modelKey")=modelKey\_MM}
\NormalTok{        and src(r),trg(r) in ObjNode\_val(G) \}.}
\end{Highlighting}
\end{Shaded}

For semantic link observations, \texttt{decodeQualifierList} is the
strict typed inverse of \texttt{encodeQualifierList} on the admitted
payload domain. A malformed, wrong-sort, or wrong-arity payload makes
\texttt{GraphAdequate} false; raw storage text is never compared
directly with a source scalar in \texttt{Obs\_graph}.

No-spurious-link condition:

\begin{Shaded}
\begin{Highlighting}[]
\NormalTok{Every Link* relationship used by OCL\_val navigation is generated from exactly}
\NormalTok{one UML association link in M.}
\end{Highlighting}
\end{Shaded}

The theorem does not require relationship-level uniqueness. If the
concrete encoding contains duplicate relationships for the same semantic
UML link, Cypher realization must remove duplicate semantic results
according to the distinctness discipline. No-spurious-link is still
required in the reverse direction.

Association classes are not instances of this direct-link rule. Their
current repository representation is a link-object plus spoke
relationships, so navigation through an association class has no
admitted \texttt{nav\_graph} case and is rejected before planning.

Each result of \texttt{linkQualifiers\_G(r,d)} is an ordered list whose
length equals the qualifier arity of the association end selected by
\texttt{d}. For every supported primitive scalar qualifier sort,
serialization is \texttt{encScalar\_tau} above and is therefore typed,
sort-preserving, and injective; list serialization is componentwise and
length-preserving:

\begin{Shaded}
\begin{Highlighting}[]
\NormalTok{encodeQualifierList(qs1) = encodeQualifierList(qs2) iff qs1 = qs2.}
\end{Highlighting}
\end{Shaded}

If any evaluated qualifier is bottom, \texttt{encodeQualifierList} is
undefined for navigation matching and the generated predicate includes a
non-null guard; therefore no relationship row matches. The
adapter\textquotesingle s expected source payload is computed by an
independent implementation of these equations rather than by calling the
writer codec.

Thus matching a direction-selected encoded qualifier list both preserves
and reflects object-model qualifier equality. In particular, equality
cannot hold between lists of different arity, and component \texttt{i}
always denotes declared qualifier \texttt{i} of that end.

Under the OCL/USE association-end convention, forward navigation
compares the arguments with \texttt{linkQualifiers\_G(r,forward)} and
reverse navigation compares them with
\texttt{linkQualifiers\_G(r,reverse)}. This choice is fixed by resolved
direction metadata, not inferred from surface role names during
rendering.

\subsection{4.6 Navigation
Interpretation}\label{46-navigation-interpretation}

The binder normalizes every resolved navigation to metadata:

\begin{Shaded}
\begin{Highlighting}[]
\NormalTok{(association A, fromRole, toRole, direction, [q1,...,qk])}
\end{Highlighting}
\end{Shaded}

Graph navigation is defined by the link encoding and this normalized
metadata. For forward navigation:

\begin{Shaded}
\begin{Highlighting}[]
\NormalTok{nav\_graph(n,A,fromRole,toRole,forward,encodedQ?)}
\NormalTok{= \{ trg(r)}
\NormalTok{    | src(r)=n}
\NormalTok{      and prop(n,"modelKey") = modelKey\_MM}
\NormalTok{      and prop(trg(r),"modelKey") = modelKey\_MM}
\NormalTok{      and prop(r,"modelKey") = modelKey\_MM}
\NormalTok{      and type(r) starts with "Link"}
\NormalTok{      and linkAssociation\_G(r) = associationKey\_MM(A)}
\NormalTok{      and linkSourceRole\_G(r) = fromRole}
\NormalTok{      and linkTargetRole\_G(r) = toRole}
\NormalTok{      and qualifierMatches(encodedQ?,linkQualifiers\_G(r,forward)) \}.}
\end{Highlighting}
\end{Shaded}

For reverse navigation:

\begin{Shaded}
\begin{Highlighting}[]
\NormalTok{nav\_graph(n,A,fromRole,toRole,reverse,encodedQ?)}
\NormalTok{= \{ src(r)}
\NormalTok{    | trg(r)=n}
\NormalTok{      and prop(n,"modelKey") = modelKey\_MM}
\NormalTok{      and prop(src(r),"modelKey") = modelKey\_MM}
\NormalTok{      and prop(r,"modelKey") = modelKey\_MM}
\NormalTok{      and type(r) starts with "Link"}
\NormalTok{      and linkAssociation\_G(r) = associationKey\_MM(A)}
\NormalTok{      and linkSourceRole\_G(r) = toRole}
\NormalTok{      and linkTargetRole\_G(r) = fromRole}
\NormalTok{      and qualifierMatches(encodedQ?,linkQualifiers\_G(r,reverse)) \}.}
\end{Highlighting}
\end{Shaded}

Navigation Preservation applies only to navigations whose resolved
direction is represented by this fixed \texttt{nav\_graph} contract.
Here \texttt{encodedQ} is the empty list when the association end is
unqualified; otherwise it is computed as:

\begin{Shaded}
\begin{Highlighting}[]
\NormalTok{rawQs = [ [[q1]]\_I(rho), ..., [[qk]]\_I(rho) ]}
\NormalTok{encodedQ = encodeQualifierList(rawQs).}
\end{Highlighting}
\end{Shaded}

The list order is the declared qualifier order of the resolved
association end. Admission requires the number and static types of
qualifier expressions to equal that declaration.
\texttt{qualifierMatches({[}{]},stored)} is true exactly when the
selected end is unqualified and \texttt{stored={[}{]}}; it does not
silently ignore a non-empty stored qualifier list.

\subsection{4.7 Validation Observation and
Equivalence}\label{47-validation-observation-and-equivalence}

Theorem 0 does not reconstruct arbitrary UML serialization. It compares
the following typed observation structures:

\begin{Shaded}
\begin{Highlighting}[]
\NormalTok{Obs\_obj(MM,M) =}
\NormalTok{  ( Obj(M), Id\_obj, Type\_obj, Attr\_obj, Link\_obj,}
\NormalTok{    NavOne\_obj, NavMany\_obj, AllInst\_obj )}

\NormalTok{Obs\_graph(MM,G) =}
\NormalTok{  ( ObjNode\_val(G), Id\_graph, Type\_graph, Attr\_graph,}
\NormalTok{    Link\_graph, NavOne\_graph, NavMany\_graph, AllInst\_graph )}
\end{Highlighting}
\end{Shaded}

The components are restricted to classes, attributes, binary
associations, roles, directions, and primitive qualifiers admitted by
\texttt{OCL\_val}. \texttt{Link\_graph} is the semantic set obtained
from validation-visible \texttt{Link*} relationships, so duplicate
relationships representing the same UML link do not create distinct
semantic links. Its qualifier components are produced by strict typed
\texttt{decodeQualifierList}; they are not raw encoded property values.

Let \texttt{lift\_node} map object-valued components pointwise through
\texttt{node} and leave supported scalar values unchanged. Define
validation equivalence by:

\begin{Shaded}
\begin{Highlighting}[]
\NormalTok{(MM,M) ==\_val G}
\NormalTok{iff}
\NormalTok{lift\_node(Obs\_obj(MM,M)) = Obs\_graph(MM,G).}
\end{Highlighting}
\end{Shaded}

Equality is typed componentwise and extensional for relation-, set-, and
function-valued components. The partial reader \texttt{Psi\_val} is
defined exactly on graphs satisfying the encoding contract and returns
this observation structure; it is not claimed to reconstruct
non-validation UML metadata.

\subsection{\texorpdfstring{4.8 Constructive Status of
\texttt{Phi}}{4.8 Constructive Status of Phi}}\label{48-constructive-status-of-phi}

Theorem 0 is an adequacy theorem for the explicit contract above. It
does not quantify over arbitrary encoders. For readers who prefer a
constructive presentation, define the canonical encoding
\texttt{Phi\_star(MM,M)} as the least graph closed under exactly these
generation rules:

\begin{Shaded}
\begin{Highlighting}[]
\NormalTok{P{-}Class:}
\NormalTok{  for each C in Class(MM), add exactly one node c with}
\NormalTok{    label(c,"UmlClass"), prop(c,"classKey")=key\_MM(C),}
\NormalTok{    prop(c,"name")=C.name, prop(c,"modelKey")=modelKey\_MM;}
\NormalTok{    define classNode(C)=c.}

\NormalTok{P{-}Metamodel:}
\NormalTok{  for every admitted M2 attribute, binary association, association end,}
\NormalTok{  qualifier declaration, and inheritance fact, add its canonical schema node}
\NormalTok{  or relationship with an injective qualified key and enough metadata for the}
\NormalTok{  corresponding resolver; add no ambiguous resolver{-}visible declaration.}

\NormalTok{P{-}Object:}
\NormalTok{  for each o in Obj(M), add exactly one node n with}
\NormalTok{    label(n,"Object"), prop(n,"use\_id")=id(o),}
\NormalTok{    prop(n,"modelKey")=modelKey\_MM; define node(o)=n.}

\NormalTok{P{-}Type:}
\NormalTok{  for each pair (o,C), add exactly one relationship}
\NormalTok{    node(o){-}[:ObjectInstanceOf]{-}\textgreater{}classNode(C)}
\NormalTok{  iff classOf(o) conformsTo C.}

\NormalTok{P{-}Attr:}
\NormalTok{  for each admitted slot (o,a), add exactly one fresh node av and exactly one}
\NormalTok{    node(o){-}[:ObjectHasAttribute]{-}\textgreater{}av with}
\NormalTok{      label(av,"AttributeValue"),}
\NormalTok{      prop(av,"modelKey")=modelKey\_MM,}
\NormalTok{      prop(av,"attributeKey")=attributeKey\_MM(a),}
\NormalTok{      prop(av,"value")=storeScalar(attrVal\_M(o,a));}
\NormalTok{  add no other validation{-}visible AttributeValue node for (o,a).}

\NormalTok{P{-}Link:}
\NormalTok{  for each admitted binary link}
\NormalTok{    l=(A,o1,sourceRole,sourceQs,o2,targetRole,targetQs),}
\NormalTok{  add exactly one relationship r of concrete type Link from node(o1) to node(o2)}
\NormalTok{  with prop(r,"modelKey")=modelKey\_MM,}
\NormalTok{       linkAssociation\_G(r)=associationKey\_MM(A),}
\NormalTok{       linkSourceRole\_G(r)=sourceRole,}
\NormalTok{       linkTargetRole\_G(r)=targetRole,}
\NormalTok{       linkQualifiers\_G(r,forward)=encodeQualifierList(sourceQs),}
\NormalTok{       linkQualifiers\_G(r,reverse)=encodeQualifierList(targetQs).}

\NormalTok{P{-}Exact:}
\NormalTok{  the validation{-}visible subgraph is the least subgraph generated by these}
\NormalTok{  rules and contains no other validation{-}visible object, class, type,}
\NormalTok{  attribute, or link fact.}
\end{Highlighting}
\end{Shaded}

\texttt{storeScalar} is injective on the scalar profile, maps supported
bottom to its reserved representation, and satisfies
\texttt{readScalar(storeScalar(v))=v}. The canonical accessor is
therefore:

\begin{Shaded}
\begin{Highlighting}[]
\NormalTok{value\_G(n,a)}
\NormalTok{  = readScalar(the unique av.value such that}
\NormalTok{      (n){-}[:ObjectHasAttribute]{-}\textgreater{}(av:AttributeValue)}
\NormalTok{      and av.modelKey=modelKey\_MM}
\NormalTok{      and av.attributeKey=attributeKey\_MM(a)).}
\end{Highlighting}
\end{Shaded}

``Least'' and P-Exact make the backward directions of R2-\/-R5
immediate; the forward directions follow from the corresponding
generation rule. R6 and R7 then follow extensionally from R5 and R3. A
concrete implementation may use a different storage layout, but Theorem
0 applies to it only after an observation mapping shows that it
satisfies the same contract. Consequently the strongest paper-safe
statement remains ``representation adequacy under the encoding
contract,'' not ``every implementation of \texttt{Phi} is faithful.''

\subsection{4.9 Concrete-Encoding Adapter
Obligation}\label{49-concrete-encoding-adapter-obligation}

The canonical profile above is intentionally not identified
syntactically with the current prototype schema. Let \texttt{G\_impl} be
a graph produced by a concrete encoder and let \texttt{P} define the
concrete observations:

\begin{Shaded}
\begin{Highlighting}[]
\NormalTok{class\_P(C)              concrete class lookup}
\NormalTok{metamodel\_P             concrete M2 resolver view}
\NormalTok{objects\_P               concrete object{-}node/id lookup}
\NormalTok{type\_P(n,C)             concrete conformance lookup}
\NormalTok{value\_P(n,a)            concrete attribute accessor}
\NormalTok{nav\_P(n,A,fr,tr,d,qbar) concrete directional navigation accessor}
\NormalTok{allInstances\_P(C)       concrete instance lookup}
\end{Highlighting}
\end{Shaded}

The adapter is defined through observations rather than by requiring the
physical property names of \texttt{G\_impl} to be identical to the
logical names used by the canonical profile. For the current prototype,
the intended physical accessor projection is:

\begin{Shaded}
\begin{Highlighting}[]
\NormalTok{classKey\_P(C)                  = prop(class\_P(C),"classKey")}
\NormalTok{attributeKey\_P(a)              = prop(attribute\_P(a),"attributeKey")}
\NormalTok{associationKey\_P(r)            = prop(r,"associationKey")}
\NormalTok{qualifiers\_P(r,forward)        = prop(r,"sourceQualifiers")}
\NormalTok{qualifiers\_P(r,reverse)        = prop(r,"targetQualifiers")}
\NormalTok{objectId\_P(n)                  = prop(n,"use\_id")}
\end{Highlighting}
\end{Shaded}

These are concrete realizations of the logical accessors
\texttt{linkAssociation\_G} and \texttt{linkQualifiers\_G}; neither
\texttt{associationName} nor a single direction-free
\texttt{qualifierPayload} belongs to the normative observation
interface. The adapter obligation is:

\begin{Shaded}
\begin{Highlighting}[]
\NormalTok{associationKey\_P(r) = associationKey\_MM(A),}
\NormalTok{qualifiers\_P(r,d) = encodeQualifierList(qualifiers\_M(A,d)),}
\end{Highlighting}
\end{Shaded}

with direction, arity, component order, and scalar serialization
preserved and reflected. Display-only \texttt{name} properties are
outside these identity observations.

\subsubsection{PA1 --- canonical-v1 key agreement and
injectivity}\label{pa1--canonical-v1-key-agreement-and-injectivity}

For the prototype adapter, let \texttt{norm(x)=trim(x)} and admit a key
component only when \texttt{norm(x)} is non-empty and does not contain
the reserved delimiter \texttt{::}. The executable \texttt{canonical-v1}
constructors are:

\begin{Shaded}
\begin{Highlighting}[]
\NormalTok{classKey(m,c)       = norm(m) ++ "::class::" ++ norm(c)}
\NormalTok{attributeKey(m,c,a) = norm(m) ++ "::attribute::" ++ norm(c) ++ "::" ++ norm(a)}
\NormalTok{associationKey(m,A) = norm(m) ++ "::association::" ++ norm(A)}
\end{Highlighting}
\end{Shaded}

\texttt{CanonicalGraphEncoding} is the single production implementation
called by the metamodel/object graph writers and by
\texttt{OclCypherRenderer}. The binder does not construct text keys: it
resolves a class, the defining owner of an attribute, or an association
declaration; the renderer applies the same constructor to those resolved
semantic names. Thus encoder and renderer agree on every admitted
declaration, including inherited attributes, whose key uses the defining
owner rather than the invariant context class.

\textbf{Lemma PA1 (KeyAgreement).} For admitted normalized components,
each of the three constructors is injective on its typed tuple, and
their output namespaces are pairwise disjoint.

\textbf{Proof.} None of the components contains \texttt{::}, so every
output has a unique decomposition at the displayed delimiters. Equality
of two keys in one namespace therefore implies componentwise equality.
The fixed tags \texttt{class}, \texttt{attribute}, and
\texttt{association} differ, so outputs from distinct namespaces cannot
be equal. The implementation rejects rather than escapes a component
containing the delimiter, preventing an invalid component from aliasing
a different admitted tuple. QED.

This lemma proves constructor agreement and key-level injectivity only.
It does not prove that an encoded graph contains exactly one
node/relationship for each key, nor that it contains no spurious
observations; those remain PA2-\/-PA9.

\subsubsection{PA2/PA3 --- prototype object and identifier
exactness}\label{pa2pa3--prototype-object-and-identifier-exactness}

For one model \texttt{m}, define the concrete object observation:

\begin{Shaded}
\begin{Highlighting}[]
\NormalTok{objectObs\_P(n) = (elementId(n), prop(n,"use\_id"), prop(n,"objectKey"))}
\NormalTok{expectedIdentity\_m(o) = (id(o), objectKey(m,id(o)))}
\end{Highlighting}
\end{Shaded}

The checked synchronization profile has the following explicit
preconditions:

\begin{enumerate}
\def\labelenumi{\arabic{enumi}.}
\tightlist
\item
  source object identifiers are non-blank and injective within
  \texttt{m};
\item
  synchronization completes atomically with an exact diff: every
  source-only or mismatched object is upserted and every graph-only
  object is deleted;
\item
  validation observes only \texttt{:Object\ \{modelKey:m\}} nodes after
  that transaction.
\end{enumerate}

The production writer derives both fields from the same source object:
it \texttt{MERGE}s on \texttt{CanonicalGraphEncoding.objectKey(m,id(o))}
and sets \texttt{use\_id=id(o)}. Schema profile
\texttt{canonical-v5-object-identity} replaces legacy indexes on
\texttt{:Object(objectKey)} with a uniqueness constraint. Repeating an
upsert for the same source identity therefore selects the same graph
node and refreshes the same \texttt{use\_id}; a second node with that
canonical key is rejected. The synchronization transaction also
deep-deletes every object reported as graph-only.

\textbf{Lemma PA2 (ObjectExactness, checked synchronization profile).}
Under the three preconditions above, the post-state has exactly one
validation-visible object node for every source object and no
validation-visible object node not represented in the source snapshot.

\textbf{Proof.} Exact diff coverage sends every missing source object to
the upsert. PA1-style injectivity of \texttt{objectKey(m,id)} and the
unique constraint give at most one node per source identity, while
\texttt{MERGE} gives at least one. Exact diff reflection sends every
graph-only identity to deep deletion. Atomic commit establishes all
three facts in the same post-state. QED, conditional on the stated
synchronization preconditions.

\textbf{Lemma PA3 (IdentifierInjectivity, checked synchronization
profile).} In the same post-state, \texttt{use\_id} exists on every
validation-visible object, is stable under repeated upsert, and is
injective within \texttt{m}.

\textbf{Proof.} The upsert unconditionally sets \texttt{use\_id=id(o)}.
Stability follows because the same canonical \texttt{objectKey} selects
the same node. If two visible nodes had the same \texttt{use\_id},
source-ID injectivity maps both to the same source object and hence the
same \texttt{objectKey}; schema uniqueness and PA2 then force the nodes
to be equal. QED, conditional on the stated synchronization
preconditions.

\texttt{RepresentationAdequacyEvaluator} checks these conclusions
independently from query rendering. PA2 compares
\texttt{(use\_id,objectKey)} facts in both directions and rejects
duplicate or malformed object observations. PA3 additionally rejects
missing fields, duplicate \texttt{use\_id}, duplicate
\texttt{objectKey}, and unstable source-to-graph mappings. These
executable checks are implementation conformance evidence; they do not
remove the premises of the two lemmas.

\subsubsection{PA4 --- materialized type
exactness}\label{pa4--materialized-type-exactness}

For the same checked synchronization post-state, define:

\begin{Shaded}
\begin{Highlighting}[]
\NormalTok{typeFact\_M(o,C) iff C = classOf(o) or C is in allParents(classOf(o))}
\NormalTok{typeFact\_P(o,C) iff node(o){-}[:ObjectInstanceOf]{-}\textgreater{}classNode(C)}
\NormalTok{                     and classNode(C).classKey = classKey(m,C)}
\end{Highlighting}
\end{Shaded}

The batch writer constructs \texttt{row.classKeys} from the runtime
class followed by all transitive parents. Before materializing this set,
it deletes every existing \texttt{ObjectInstanceOf} membership whose
target \texttt{classKey} is not in \texttt{row.classKeys}; it then
\texttt{MERGE}s a membership to every \texttt{UmlClass} whose exact
canonical key is in the set.

\textbf{Lemma PA4 (TypeExactness, checked synchronization profile).}
Assuming PA1, PA2, an exact metamodel snapshot, complete
\texttt{allParents}, successful execution of the production batch
synchronization path, and no concurrent external graph mutation:

\begin{Shaded}
\begin{Highlighting}[]
\NormalTok{typeFact\_P(o,C) iff typeFact\_M(o,C).}
\end{Highlighting}
\end{Shaded}

\textbf{Proof.} For the forward implication, a membership retained or
created by the batch path has a target key in \texttt{row.classKeys};
construction of that list means the target is exactly the runtime class
or one of its transitive parents. For the reverse implication, every
runtime/parent class contributes its canonical key to
\texttt{row.classKeys}, exact M2 lookup finds its \texttt{UmlClass}, and
\texttt{MERGE} materializes the edge. Stale memberships are removed in
the same transaction. Relationship multiplicity is immaterial to this
extensional predicate and to the renderer\textquotesingle s
existence/distinct observations. QED under the stated premises.

The renderer\textquotesingle s invariant context, \texttt{oclIsKindOf},
and \texttt{allInstances} patterns all observe this predicate through
exact \texttt{classKey} matches. PA4 establishes the storage predicate;
the complete end-to-end \texttt{allInstances} accessor remains the
separate PA8 obligation.

\subsubsection{PA5 --- attribute slot and value
exactness}\label{pa5--attribute-slot-and-value-exactness}

\texttt{attributeKey(m,C,a)} identifies the declared attribute and is
intentionally shared by every object slot for that attribute. The
prototype therefore uses a separate physical slot identity:

\begin{Shaded}
\begin{Highlighting}[]
\NormalTok{slotKey(m,o,a) = objectKey(m,id(o)) ++ "::slot::" ++ attributeKey(m,owner(a),a)}
\end{Highlighting}
\end{Shaded}

\texttt{canonical-v6-attribute-slot-identity} migrates existing
reachable \texttt{AttributeValue} nodes to this key and installs a
uniqueness constraint on \texttt{AttributeValue.slotKey}. Both scalar
and complex-attribute production writers derive the key with
\texttt{CanonicalGraphEncoding.attributeSlotKey}, merge the value node
by \texttt{slotKey}, set the exact \texttt{attributeKey}, and connect it
to \texttt{node(o)} by \texttt{ObjectHasAttribute}. Undefined USE values
are represented by the explicit encoded payload \texttt{Undefined}, not
by omitting the slot.

\textbf{Lemma PA5 (Scalar AttributeExactness, checked synchronization
profile).} Assume PA1-\/-PA4, a fixed exact M2 attribute set during the
object-sync transaction, complete change detection, successful schema
migration, and no concurrent external graph mutation. For every source
object \texttt{o} and admitted attribute \texttt{a} of primitive scalar
type in its conforming class, there is exactly one validation-visible
\texttt{AttributeValue} node \texttt{v} such that:

\begin{Shaded}
\begin{Highlighting}[]
\NormalTok{node(o){-}[:ObjectHasAttribute]{-}\textgreater{}v}
\NormalTok{v.slotKey      = slotKey(m,o,a)}
\NormalTok{v.attributeKey = attributeKey(m,owner(a),a)}
\NormalTok{v.value        = encodeAttributeValue(attrVal\_M(o,a)).}
\end{Highlighting}
\end{Shaded}

No other validation-visible slot for \texttt{(o,a)} exists.

\textbf{Proof.} PA1/PA3 make both components of \texttt{slotKey} stable
and injective, so equality of slot keys implies equality of the object
and declared attribute. The writer covers every current attribute when
an object is introduced or reported changed, and \texttt{MERGE} supplies
existence while the schema constraint supplies at-most-one physical
keyed slot. The writer sets the payload from the same source value and
exact defining-owner key. Under fixed M2 and complete change detection,
unchanged slots remain equal and changed slots are refreshed; the
no-external-mutation premise excludes unkeyed duplicate observations.
QED under the stated premises.

The renderer follows \texttt{ObjectHasAttribute} from the already
selected receiver and filters by exact \texttt{attributeKey}. PA5 makes
the resulting list singleton, so its scalar \texttt{head(...)} accessor
is deterministic and equals the proof-level \texttt{value\_G}.
\texttt{RepresentationAdequacyEvaluator} additionally compares
\texttt{(objectId,attributeKey,value,slotKey)} in both directions and
rejects missing, spurious, duplicate, malformed, or wrong-value slots.

The writer and schema also use \texttt{slotKey} for complex attributes,
but PA5 and the current \texttt{VAttribute} proof case do not claim
semantic decoding correctness for collection- or reference-valued
payloads. Those accesses are rejected by certified bound admission and
remain experimental until separate value-domain, typing, lowering, and
realization lemmas are supplied.

\subsubsection{PA6/PA7 --- binary link and directional qualifier
exactness}\label{pa6pa7--binary-link-and-directional-qualifier-exactness}

For a validation-visible binary link, the prototype now gives the
relationship an identity derived from the complete semantic tuple rather
than merging only by association and endpoints:

\begin{Shaded}
\begin{Highlighting}[]
\NormalTok{linkKey(m,A,o1,o2,sourceQs,targetQs)}
\NormalTok{  = modelKey(m) ++ "::link::"}
\NormalTok{    ++ frame(name(A)) ++ frame(id(o1)) ++ frame(id(o2))}
\NormalTok{    ++ frameList(encodeQualifierList(sourceQs))}
\NormalTok{    ++ frameList(encodeQualifierList(targetQs))}

\NormalTok{frame(x)        = length(x) ++ ":" ++ x}
\NormalTok{frameList(xs)   = length(xs) ++ ":" ++ concat(map(frame,xs)).}
\end{Highlighting}
\end{Shaded}

Length-prefixing makes tuple decomposition unique even when ids or
serialized payloads contain punctuation. The diff
engine\textquotesingle s \texttt{LinkState.buildIdentity} uses the same
model-independent suffix, so detection, deletion, and writing cannot
disagree about binary-link identity. The production writer
\texttt{MERGE}s each of the three supported binary relationship labels
by this \texttt{linkKey}, then sets the exact \texttt{associationKey},
endpoints, roles, and both ordered qualifier lists.
\texttt{canonical-v7.1-binary-link-exactness} migrates existing binary
relationships, normalizes absent qualifier properties to empty lists,
indexes \texttt{linkKey} for all three labels, and refuses installation
if a model-scoped relationship is malformed or if the migrated graph
contains a duplicate semantic key.

\textbf{Lemma PA6 (LinkExactness, checked binary-link synchronization
profile).} Assume PA1-\/-PA5, a fixed exact binary-association M2 during
synchronization, complete source/graph link change detection, successful
v7 migration including its duplicate-key preflight, successful
production synchronization, and no concurrent external graph mutation.
For every admitted source binary link \texttt{l} there is exactly one
validation-visible relationship \texttt{r} with the canonical
\texttt{linkKey(l)}, exact \texttt{associationKey}, source/target
objects, and declared roles; and every validation-visible
\texttt{LinkAssociateWith}, \texttt{LinkAggregates}, or
\texttt{LinkComposeOf} relationship corresponds to exactly one such
source link.

\textbf{Proof.} Injectivity of the length-prefixed tuple encoding makes
equal \texttt{linkKey}s imply equality of association, ordered
endpoints, and both ordered qualifier payload lists. Complete change
detection selects every missing or mismatched source link and every
graph-only link. Deletion removes graph-only keys; keyed \texttt{MERGE}
supplies a witness for each source key and refreshes all observed
metadata. The migration preflight establishes the initial at-most-one
property; transactional keyed merges preserve it under the
no-concurrent- mutation premise. Thus source and graph link observations
are equal in both directions. QED under the stated premises.

\textbf{Lemma PA7 (QualifierAgreement, checked binary-link
synchronization profile).} Under PA6 and the admitted primitive
qualifier codec\textquotesingle s sort-preserving injectivity, for each
link \texttt{r}:

\begin{Shaded}
\begin{Highlighting}[]
\NormalTok{linkQualifiers\_G(r,forward) = encodeQualifierList(sourceQs)}
\NormalTok{linkQualifiers\_G(r,reverse) = encodeQualifierList(targetQs),}
\end{Highlighting}
\end{Shaded}

with equal declared arity, component order, and component payloads.

\textbf{Proof.} The writer enumerates \texttt{MLink.getQualifier()} by
association-end index and serializes each inner list in iteration order.
It stores index zero as \texttt{sourceQualifiers} and index one as
\texttt{targetQualifiers}. The renderer selects the former for outgoing
navigation and the latter for incoming navigation. Length and component
boundaries are retained by lists, while codec injectivity reflects
component equality. Therefore equality of the direction-selected stored
list is equivalent to equality of the corresponding ordered source
qualifier tuple. QED under the stated premises.

The independent evaluator records physical relationship tokens and
compares link witnesses and directional qualifier facts in both
directions. PA6 rejects missing, spurious, malformed, and duplicate-key
relationships; PA7 separately rejects missing properties and wrong
direction, arity, order, or payload.

\subsubsection{\texorpdfstring{PA8 --- inherited \texttt{allInstances}
accessor
agreement}{PA8 --- inherited allInstances accessor agreement}}\label{pa8--inherited-allinstances-accessor-agreement}

Define the prototype accessor observation independently of the stored
type-fact snapshot:

\begin{Shaded}
\begin{Highlighting}[]
\NormalTok{allInstances\_P(C)}
\NormalTok{  = \{ o.use\_id |}
\NormalTok{      (o){-}[:ObjectInstanceOf]{-}\textgreater{}(c)}
\NormalTok{      and c.classKey = classKey(m,C) \}.}
\end{Highlighting}
\end{Shaded}

\textbf{Lemma PA8 (AllInstancesAgreement, checked synchronization and
renderer profile).} Assume PA1-\/-PA4, successful synchronization, exact
\texttt{classKey} construction, complete transitive-parent
materialization, and no concurrent external mutation. For every admitted
class \texttt{C}:

\begin{Shaded}
\begin{Highlighting}[]
\NormalTok{allInstances\_P(C) = \{ id(o) | classOf(o) conformsTo C \}.}
\end{Highlighting}
\end{Shaded}

\textbf{Proof.} By PA1, the rendered parameter selects exactly the class
node for \texttt{C}. By PA4, an \texttt{ObjectInstanceOf} witness to
that node exists iff the source object conforms to \texttt{C}, including
indirect supertypes. PA2/PA3 identify every matched validation-visible
object with exactly one source identifier. These facts prove both
inclusions. The renderer realizes the finite-set quotient with
\texttt{COLLECT\ \{\ MATCH\ ...\ RETURN\ DISTINCT\ o\ \}}, so duplicate
physical membership edges cannot change membership or cardinality. QED
under the stated premises.

The executable PA8 observation is deliberately collected by a separate
query, not reconstructed from the evaluator\textquotesingle s
\texttt{typeFacts}; therefore a renderer or accessor mutation can fail
PA8 while PA4 still passes.

\subsubsection{PA9 --- renderer accessor
agreement}\label{pa9--renderer-accessor-agreement}

Let \texttt{Accessors(plan)} be the set of context, identity, type,
scalar-attribute, binary-navigation, qualifier, and
\texttt{allInstances} accessors reachable in an admitted
\texttt{OclCypherPlan}. Define
\texttt{RendererAccessorAgreement(plan,text,params)} to require the
following exact realization matrix:

{\def\LTcaptype{none} % do not increment counter
\begin{longtable}[]{@{}lll@{}}
\toprule\noalign{}
Plan constructor & Required text/parameter observation & Supporting
adapter lemma \\
\midrule\noalign{}
\endhead
\bottomrule\noalign{}
\endlastfoot
invariant context/result & \texttt{ObjectInstanceOf}, exact canonical
context \texttt{classKey}, \texttt{RETURN\ DISTINCT\ self.use\_id} &
PA1-\/-PA4 \\
scalar \texttt{AttributeAccessPlan} & bottom-safe canonical
\texttt{objectKey} re-identification, then \texttt{ObjectHasAttribute}
plus exact canonical \texttt{attributeKey}; never suffix or display-name
matching & PA1/PA2/PA5 \\
outgoing/incoming \texttt{NavigationAccessPlan} & direction-correct
\texttt{Link*} pattern, exact \texttt{associationKey}, source role and
target role & PA1/PA6 \\
qualified navigation & \texttt{sourceQualifiers{[}i{]}} outgoing and
\texttt{targetQualifiers{[}i{]}} incoming, in declared index order &
PA7 \\
\texttt{oclIsKindOf}/\texttt{oclAsType} & bottom-safe canonical
\texttt{objectKey} re-identification, then exact
\texttt{ObjectInstanceOf} and canonical target \texttt{classKey} &
PA1/PA2/PA4 \\
\texttt{allInstances} & exact \texttt{ObjectInstanceOf}/canonical
\texttt{classKey} inside \texttt{COLLECT}, with local
\texttt{RETURN\ DISTINCT} & PA8 \\
\end{longtable}
}

\textbf{Lemma PA9 (RendererAccessorAgreement, checked
admitted-constructor profile).} Assume successful certified
binding/planning, PA1-\/-PA8, and the selected Cypher dialect contract.
For every accessor in an admitted plan, the current renderer text and
parameter environment implement the corresponding prototype observation
used by PA1-\/-PA8.

\textbf{Proof.} By case analysis over the finite table. The invariant
wrapper emits the conformance edge, exact class parameter, and injective
result projection. The attribute case emits the unique PA5 slot path and
parameter obtained from the binder-resolved defining owner. The three
navigation directions emit the PA6 relationship filter and swap endpoint
roles only in the incoming case; qualifier indices select the PA7
directional property. Type operations reuse the PA4 conformance
predicate. \texttt{allInstances} is exactly the PA8 DISTINCT set
accessor. Parameter closure checks establish that every textual
\texttt{\$p} has one bound value and no extra value is silently
substituted. These cases exhaust the accessor-bearing plan constructors.
QED under the stated premises.

PA9 is not a parser-level proof that arbitrary rendered text is
alpha-equivalent to the formal raw Cypher AST. PA14a-\/-PA14d verify the
checked 17-constructor direct-text profile, including the bottom-safe
receiver cases, against the reviewed finite oracle. The canonical graph
contract, raw AST, production text, and oracle now use
\texttt{UmlClass}, discharging PO-01/PO-14 for this finite profile.
Universal full Neo4j-AST equality remains outside the checked result.

\subsubsection{PA10 --- checked bottom
separation}\label{pa10--checked-bottom-separation}

\textbf{Lemma PA10 (BottomSeparated, checked tagged-map execution
profile).} Assume certified binding, the canonical scalar/key/qualifier
domains, Neo4j\textquotesingle s property value boundary, and execution
through the checked validation entry points. The prototype parameter
bound as bottom is disjoint from every theorem-visible literal, stored
scalar, object identifier, metamodel key, and qualifier.

\textbf{Proof.} \texttt{OclBottomToken} is the sole production
constructor and returns the immutable tagged map
\texttt{\{\_\_oclBottom:true\}}. Certified literals and qualifiers are
admitted primitive scalars, identifiers and canonical keys are strings,
and persisted Neo4j properties cannot be maps. Hence none equals the
token. The runtime parameter guard recursively rejects the token in
externally supplied operation environments. Before document or
single-rule validation, the model-scoped graph checker rejects use of
the reserved marker on a node or relationship. Finally, generated
finite-set plans convert Cypher null to the non-null token before
uniqueness/cardinality processing. Therefore BR3 and the admission
portion of BR7 hold for each execution reporting
\texttt{BottomSeparated=PASS}. The selected-server probe establishes
BR5/BR6 for the recorded runtime profile. QED under the stated premises.

This is a conditional implementation-conformance lemma: bypassing the
checked entry points, changing the token/property profile, or using an
unprobed Cypher runtime reopens the obligation. Bottom-safe
entity-receiver branching and live alias preservation are handled
separately by PA12.

\subsubsection{PA11 --- checked scalar
closure}\label{pa11--checked-scalar-closure}

\textbf{Lemma PA11 (ScalarClosed, checked observed-domain profile).} Let
a certified run supply complete values for every primitive attribute
reachable from the invariant, and let \texttt{OclScalarClosureChecker}
return \texttt{PASS}. Then every checked scalar literal, stored value,
parameter, coercion, and arithmetic result belongs to the scalar domain
of Section 3.3.

\textbf{Proof.} Admission limits scalar types to Boolean, Integer, Real,
String and bottom. The checker accepts only Int64 integral values,
finite binary floating values, Boolean values, and strings containing no
unpaired UTF-16 surrogate. For each arithmetic IR node it recursively
obtains the finite observed operand domains. Attribute domains are
collected from the same USE fixture; literal, set and iterator domains
are propagated structurally. It evaluates their Cartesian product:
Integer \texttt{+,-,*} must remain in Int64, every Real result must be
finite, every divisor must be non-zero, and every Integer coerced to
Real must have magnitude at most \texttt{2\^{}53}. Thus the product is a
conservative superset of reachable operand pairs, so passing it implies
all reachable results satisfy the Section 3.3 side conditions. Missing
domains or a product larger than the checked bound yields
\texttt{OUT\_OF\_SCOPE}, never \texttt{PASS}. The graph and
runtime-parameter scans establish the same scalar boundary at production
validation entry points. QED under completeness of the stated
observations.

PA11 does not claim arbitrary external Cypher execution is
scalar-closed. A general validation run whose dynamic arithmetic has not
supplied complete operand observations cannot use this lemma; it remains
\texttt{OUT\_OF\_SCOPE} for the theorem even though graph/parameter
boundary admission is still enforced.

\subsubsection{PA12 --- bottom-safe receiver and alias
agreement}\label{pa12--bottom-safe-receiver-and-alias-agreement}

Define the prototype receiver guard for entity-or-bottom expression
\texttt{E} as:

\begin{Shaded}
\begin{Highlighting}[]
\NormalTok{WITH E.objectKey AS receiverKey}
\NormalTok{WHERE receiverKey IS NOT NULL}
\NormalTok{MATCH (receiver:Object \{objectKey: receiverKey\}) ...}
\end{Highlighting}
\end{Shaded}

The \texttt{WITH} occurs inside the expression subquery and may
reference every live outer alias occurring in \texttt{E}.

\textbf{Lemma PA12 (BottomSafeReceiverAgreement, checked canonical-key
profile).} Assume PA2, PA4/PA5 as applicable, \texttt{BottomSeparated},
and the selected Cypher expression-subquery scoping contract. The
prototype lowering of attribute, \texttt{oclIsKindOf}, and
\texttt{oclAsType} agrees extensionally with the formal
\texttt{withEntityReceiver}, never supplies bottom to an entity pattern,
evaluates the receiver identity once, and preserves all live aliases
used by the receiver.

\textbf{Proof.} If \texttt{E} is bottom, it is either Cypher null or the
tagged map. Neither has a non-null \texttt{objectKey}, so the
\texttt{WHERE} removes the row before \texttt{MATCH}; the enclosing
\texttt{head(COLLECT\ \{...\})} yields null for attribute/cast and
\texttt{EXISTS} yields false for kind-of. If \texttt{E} is an entity,
PA2 gives exactly one canonical \texttt{objectKey} and exactly one
matching validation-visible object node. PA5 then gives the attribute
result, while PA4 gives kind/cast conformance; hence the entity result
equals the formal arm. The receiver expression occurs once in the first
\texttt{WITH}. Correlated expression-subquery scope resolves
\texttt{self}, the current iterator, and every outer iterator referenced
by \texttt{E} without renaming; the real nested fixture distinguishes
loss of either \texttt{p} or \texttt{self}. These cases exhaust the
three entity-consuming constructors. QED under the stated premises.

The key guard is an implementation refinement of the formal
receiver-alias guard, not a change to its semantics. Replacing canonical
\texttt{objectKey} by display name or unscoped \texttt{use\_id} would
invalidate the PA2 step.

\subsubsection{PA13 --- prototype raw-AST datatype and printer
totality}\label{pa13--prototype-raw-ast-datatype-and-printer-totality}

Let \texttt{RawJava} be the sealed Java datatype \texttt{RawCypherAst},
whose four sealed families are \texttt{Expr}, \texttt{Pattern},
\texttt{Clause}, and \texttt{Query}. Its records correspond to the raw
grammar of Section 6.6, with Java names \texttt{ListExpr},
\texttt{CaseExpr}, \texttt{NodePattern}, \texttt{RelPattern}, and
\texttt{Seq} where needed to avoid host-language name collisions.

\textbf{Lemma PA13a (RawDatatypeClosure).} Every constructible
\texttt{RawJava} value is a finite tree over typed identifier atoms,
finite syntax enums, parameter atoms, and child \texttt{RawJava} values;
it contains no constructor for arbitrary model text.

\textbf{Proof.} The four roots are sealed. Alias and parameter atoms
enforce the certified identifier grammar. Property keys, labels,
relationship types, directions, unary/binary operators, and functions
are enums. Lists are copied immutably, required children are non-null,
duplicate property keys/imports and empty clauses that have no Cypher
meaning are rejected. Model strings have no literal constructor and must
be represented by \texttt{Param}. Structural induction over the sealed
permitted-subclass lists gives the result. QED.

\textbf{Lemma PA13b (RawRendererTotality).}
\texttt{RawCypherRenderer.render} terminates and returns concrete text
for every constructible \texttt{RawJava.Query}.

\textbf{Proof.} Each renderer branch consumes one sealed constructor and
recursively renders only proper children. Enum switches are exhaustive.
\texttt{Seq} maps over a finite non-empty clause list; \texttt{UnionAll}
recurses into its two proper children. Expression, pattern, and clause
helpers cover exactly their sealed permitted subclasses. \texttt{Call}
deterministically injects its unique typed imports as the leading
\texttt{WITH} of every \texttt{UnionAll} arm (or of the sole
\texttt{Seq}), matching Cypher\textquotesingle s branch-local scope.
Therefore the recursion terminates and has no missing constructor case.
The executable totality test compares instantiated coverage against
Java\textquotesingle s \texttt{getPermittedSubclasses()} for all four
roots. QED.

PA13a/b do \textbf{not} prove prototype \texttt{BuildCQ/Expand} closure.
The selected production architecture deliberately keeps
\texttt{OclCypherRenderer} as a direct \texttt{OclCypherPlan}-to-text
transformation; \texttt{RawJava} is a reference datatype/test oracle,
not an execution stage. The implementation bridge therefore becomes
PO-14: independently parse and normalize generated text, relate its
normalized tree to the expected formal constructor tree, and check it on
the selected Neo4j parser. The reviewed finite bridge is exact for the
admitted profile after the \texttt{UmlClass} alignment; TXT5 remains
conditional outside that finite profile.

\subsubsection{PA14 --- checked direct-Cypher syntax
closure}\label{pa14--checked-direct-cypher-syntax-closure}

Let \texttt{parseGen} be the independent test parser
\texttt{GeneratedCypherSyntaxTree.parse} for the closed token/clause
fragment emitted by \texttt{OclCypherRenderer}, and let
\texttt{printGen} be its whitespace-normalizing structural printer.

\textbf{Lemma PA14a (CheckedGeneratedSyntaxClosure).} For each of the 47
admitted coverage plans \texttt{P}, if \texttt{text(P)} is the
production output, then \texttt{parseGen(text(P))} is defined and
\texttt{parseGen(printGen(parseGen(text(P))))\ =\ parseGen(text(P))}.
Moreover, the parameter atoms occurring in \texttt{text(P)} are exactly
the keys of the returned parameter environment.

\textbf{Evidence.} \texttt{OclValFragmentCoverageTest} checks the
equation and deterministic recompilation for all 47 cases;
\texttt{GeneratedCypherContractVerifier} checks exact parameter-domain
equality. \texttt{GeneratedCypherSyntaxTreeTest} additionally rejects
unbalanced groups, malformed clauses, placeholders, comments, invalid
parameters, characters outside the fragment, bare \texttt{UNION}, and
empty union arms. The focused run passed 99/99 tests (95 coverage tests
plus four parser tests); the complete implementation-conformance run
passed 235/235 with mutation score 20/20 on 2026-07-31.

PA14a is normalization stability for an independent closed-fragment
parser. It does \textbf{not} by itself prove constructor agreement,
alias alpha-equivalence, or acceptance/equivalence by the selected Neo4j
parser.

\textbf{Lemma PA14b (CheckedPlanFormalTreeAgreement).} Let
\texttt{Kinds5} be the checked AST-kind projection taken from the local
Cypher 5 grammar and AST constructor sources under
\texttt{md/neo4j/community/cypher/front-end}. For every sealed
\texttt{OclCypherPlan.ExpressionPlan} constructor \texttt{K}, the test
oracle defines an expected witness \texttt{W\_K} consisting of the
required Cypher AST kinds, tokens, directions, canonical parameter
suffixes, and forbidden alternatives. If \texttt{P} is a checked
instance of \texttt{K}, then parsing the independently rendered text of
\texttt{P} yields all observations in \texttt{W\_K}. Recursing over the
typed plan children yields an \texttt{AgreementNode} tree with one
successful node for every plan node.

\textbf{Evidence.} \texttt{ExpectedFormalCypherTreeVerifier} contains an
explicit rule registry equal to
\texttt{ExpressionPlan.getPermittedSubclasses()}, currently 17/17.
\texttt{OclCypherPlanFormalTreeAgreementTest} checks all 47 admitted
plans, adds three general-pipeline fixtures for renderer constructors
not fully reached by the frozen admission corpus, and adds one typed
manual \texttt{LetPlan} fixture because the optimizer eliminates or
inlines \texttt{let} before planning. The general navigation-aggregation
fixtures are test-only renderer evidence and remain outside frozen
\texttt{OCL\_val}. Four lexically valid mutations are rejected:
replacing \texttt{attributeKey} by \texttt{name}, reversing a navigation
arrow, replacing \texttt{COUNT} by \texttt{COLLECT}, and deleting
\texttt{NOT} from \texttt{NOT\ EXISTS}.

\texttt{OpenCypherFrontendReferenceContractTest} pins the projection
vocabulary to the checked-in Cypher 5 \texttt{Cypher5Parser.g4} and
Scala AST constructors. This check exposed and corrected one
classification error: renderer navigation syntax is a
\texttt{PatternComprehension}, not a \texttt{ListComprehension}. The
combined parser, coverage, source-reference, and plan-agreement run
passed 102/102 tests on 2026-08-01. The complete
implementation-conformance gate then passed 238/238 tests with compiler
mutation score 20/20.

PA14b proves expected-witness agreement for the checked constructor
profile; it is deliberately weaker than full-tree equality or equality
with Neo4j\textquotesingle s complete internal AST.

\textbf{Lemma PA14c (CheckedCanonicalFullTreeAgreement).} Let
\texttt{canonGen} retain every atom and delimiter group in
\texttt{parseGen(text)} while replacing only variable and alias
identifiers by deterministic lexical-scope names. It preserves property
and map keys, labels, relationship types, function names, literals,
operators, and parameters exactly. Let \texttt{Expected52(i)} be the
collision-free serialized canonical token/group tree stored for checked
case \texttt{i}. For each of the 52 checked cases,

\begin{Shaded}
\begin{Highlighting}[]
\NormalTok{serialize(canonGen(parseGen(text(P\_i)))) = Expected52(i).}
\end{Highlighting}
\end{Shaded}

The 52 cases comprise the 47 admitted queries, four general-pipeline
renderer fixtures, and one manual \texttt{LetPlan} renderer fixture. One
general fixture is a model-typed OCL invariant with nested iterators
that deliberately shadow the name \texttt{x}; it traverses the OCL
parser, binder, IR builder, optimizer, planner, and direct renderer.

\textbf{Evidence.} \texttt{GeneratedCypherCanonicalTreeTest} proves on
targeted pairs that valid nested renaming preserves the canonical tree
while variable capture does not. It also shows that changing direction,
relationship type, label, property/map key, or parameter changes the
tree. The checked-in \texttt{formal-cypher-trees.tsv} stores gzip/base64
only as an encoding; the test inflates the complete canonical strings
and compares them directly, so hash collision is not a premise.
\texttt{GeneratedCypherFormalTreeManifestTest} checks all 52 names and
complete strings. The focused source/constructor/full-tree run passed
7/7 and the complete implementation-conformance gate passed 242/242 with
compiler mutation score 20/20 on 2026-08-01.

PA14c is an exact finite-corpus statement about the independent
normalized token/group tree, not a universal theorem for every possible
plan and not an assertion about Neo4j\textquotesingle s internal AST.
The manifest must not be regenerated automatically after renderer
changes: constructor witnesses and semantic review must justify any
expected-tree update. PA14d below adds the selected actual Neo4j-parser
projection; universal closure and full internal-AST equality remain
separate obligations.

\textbf{Lemma PA14d (CheckedNeo4jParserAstProjectionAgreement).} Let
\texttt{parseN5} be \texttt{AstParserFactory(CypherVersion.Cypher5)}
from test-scope artifact
\texttt{org.neo4j:cypher-parser-factory:2026.06.0}, and let
\texttt{projN5} recursively project the resulting Scala products while
erasing source positions and object identities only. It preserves
constructor/field order, schema and property names, function names,
parameters, literals, directions, and aliases. For each of the 52
checked cases \texttt{P\_i}, \texttt{parseN5(text(P\_i))} is defined.
Moreover, for every AST-kind requirement in the constructor witness
\texttt{W\_K} of PA14b, both the independent \texttt{parseGen}
observation and a corresponding node/type in
\texttt{projN5(parseN5(text(P\_i)))} exist.

\textbf{Evidence.} \texttt{Neo4jCypherAstBridge} calls the official
factory and \texttt{singleStatement()}, checks that the loaded artifact
is exactly \texttt{cypher-parser-factory-2026.06.0.jar}, and produces a
deterministic canonical projection. \texttt{Neo4jCypherAstBridgeTest}
parses all 52 formal-tree cases, of which 47 are the certified corpus;
repeated projection is identical. It additionally checks two nested
iterator scope nodes and retained shadowed alias names, rejects
malformed query/CASE syntax, and distinguishes outgoing from incoming
relationship direction and \texttt{CountExpression} from
\texttt{CollectExpression}. \texttt{ExpectedFormalCypherTreeVerifier}
now accepts an \texttt{AstKind} requirement only if both independent
inference and the selected Neo4j AST witness it. The valid direction
mutation was corrected from the formerly malformed
\texttt{-{[}r{]}\textless{}-} text to the syntactically valid reversal
\texttt{)-{[}r{]}-\textgreater{}(} to \texttt{)\textless{}-{[}r{]}-(}.
The focused parser/source/constructor run passed 6/6; the complete
static conformance gate passed 250/250 both normally and with Maven
offline cache, with zero failures, errors, or skips, mutation score
20/20, and proof sync PASS on 2026-08-01.

PA14d discharges selected parser acceptance and AST-kind/projection
agreement for the finite checked corpus. It is not full equality between
Neo4j\textquotesingle s internal AST and the formal raw AST, does not
invoke Neo4j semantic-analysis phases, and does not prove closure for
every constructible plan. Runtime meaning remains conditional on
CY1-\/-CY9 and the recorded real-Neo4j evidence.

\texttt{AdapterAdequate(P,G\_impl,MM,M)} holds iff these observations
are defined for every theorem-reachable argument and agree extensionally
with the canonical observations of \texttt{Phi\_star(MM,M)} after the
identity-preserving object-node map:

\begin{Shaded}
\begin{Highlighting}[]
\NormalTok{objects\_P = objects\_star}
\NormalTok{metamodel\_P = metamodel\_star}
\NormalTok{type\_P = type\_star}
\NormalTok{value\_P = value\_star}
\NormalTok{nav\_P = nav\_star}
\NormalTok{allInstances\_P = allInstances\_star.}
\end{Highlighting}
\end{Shaded}

\textbf{Theorem PA-COMP (AdapterAdequate composition schema).} Fix one
model, metamodel, synchronized graph post-state, renderer version, and
observation interface. Assume (i) exact validation-visible M2
declarations in that same post-state; (ii) PA1-\/-PA8 with all of their
synchronization, admission, and no-external-mutation premises
simultaneously satisfied; and (iii) PA9 for the same plan and renderer.
Then the six equalities above hold and hence
\texttt{AdapterAdequate(P,G\_impl,MM,M)}.

\textbf{Proof.} M2 exactness and PA1 identify the prototype metamodel
observation with the canonical one. PA2/PA3 give object observation
equality under the identity-preserving object-node map. PA4 gives type
equality, PA5 gives scalar value equality, PA6/PA7 give direction- and
qualifier-sensitive navigation equality, and PA8 gives
\texttt{allInstances} equality. PA9 ensures that the renderer accessors
used by the generated plan observe exactly those prototype relations,
rather than a display-name or suffix-based surrogate. Substitution in
the definition yields all six conjuncts. QED under the shared premises.

This theorem is a composition rule, not evidence that the Java
implementation satisfies its premises for every theorem-reachable model.
In particular, PA1-\/-PA9 results recorded on different fixtures or
graph snapshots may not be silently combined, and PA1 key injectivity
alone does not imply exact M2 declarations. The current registry records
PO-21 as partial. The 2026-08-09 checkpoint had a clean
exact-M2/shared-snapshot certificate, but the new codec/model-scope
accessors invalidate that runtime artifact. The composition theorem
remains valid conditionally; the current implementation needs a fresh
certificate before this profile is re-discharged.

This is the exact bridge needed to instantiate Theorem 0 for a
non-canonical layout. The following names must not be conflated without
such a bridge:

{\def\LTcaptype{none} % do not increment counter
\begin{longtable}[]{@{}lll@{}}
\toprule\noalign{}
Canonical proof observation & Prototype-style representation & Required
adapter obligation \\
\midrule\noalign{}
\endhead
\bottomrule\noalign{}
\endlastfoot
injective \texttt{classKey=key\_MM(C)} & exact \texttt{classKey} lookup
& PA1 proves constructor agreement/injectivity; still prove absence of
colliding visible class nodes \\
exact scalar \texttt{attributeKey\_MM(a)} & exact \texttt{attributeKey}
lookup through \texttt{ObjectHasAttribute}; physical slot identity is
unique \texttt{slotKey} & PA5 discharged for primitive scalar attributes
in the checked fixed-M2 synchronization profile; complex attributes are
outside certified admission \\
\texttt{linkAssociation\_G(r)=associationKey\_MM(A)} & exact
relationship \texttt{associationKey} plus unique length-prefixed
\texttt{linkKey} & PA6 discharged for the checked binary-link
synchronization profile; physical \texttt{name} is display/migration
input only \\
ordered \texttt{linkQualifiers\_G(r,d)} & directional
\texttt{sourceQualifiers}/\texttt{targetQualifiers} & PA7 discharged for
the checked admitted primitive-qualifier profile \\
materialized conformance edges & exact \texttt{ObjectInstanceOf} lookup
by \texttt{classKey}; \texttt{COLLECT\ \{\ ...\ RETURN\ DISTINCT\ o\ \}}
& PA4/PA8 discharged for the checked inherited synchronization and
renderer profile \\
raw Cypher AST plus \texttt{renderRaw} & sealed \texttt{RawCypherAst}
and total \texttt{RawCypherRenderer} are reference artifacts; production
\texttt{OclCypherRenderer} constructs text directly & PA13a/b discharge
standalone datatype/printer closure; PA14b/c relate direct text to typed
witnesses and exact canonical trees; PA14d discharges selected Neo4j
2026.06.0 parser acceptance/AST-kind projection for 52 checked queries;
universal/full-AST equality remains open \\
\end{longtable}
}

For the canonical profile, PO-21 retains the historical clean
shared-snapshot checkpoint but is currently partial until recapture.
Even after recapture, Theorems 0-\/-6 do not become an unconditional
implementation theorem: each execution must pass the certificate
boundary, and other graph layouts require their own adapter proof.

\subsection{4.10 Prototype Adapter
Status}\label{410-prototype-adapter-status}

The current prototype was inspected against the preceding definition.
The matrix below records the current static alignment and the premises
that a fresh PO-21 runtime certificate must re-establish for the
canonical profile.

{\def\LTcaptype{none} % do not increment counter
\begin{longtable}[]{@{}lll@{}}
\toprule\noalign{}
Observation & Current prototype evidence & Status against canonical
profile \\
\midrule\noalign{}
\endhead
\bottomrule\noalign{}
\endlastfoot
context/class lookup & \texttt{OclCypherRenderer.renderInvariant}
matches \texttt{cls:UmlClass\ \{modelKey:\$pm,classKey:\$p\}} and
obtains canonical parameters from the bound metamodel context; the
writer and stored schema use the same scope and label & PO-01
vocabulary/syntax is aligned; the changed accessor requires fresh PA
runtime recapture \\
object identity and model isolation & writer merges by the globally
injective model-prefixed \texttt{objectKey}, writes \texttt{modelKey},
refreshes \texttt{use\_id}, and every validation lookup constrains both
fields; the snapshot additionally reports current-key objects declared
in a foreign model & PA2/PA3 and AI-21 are closed statically for the
scoped profile; external mutation and fresh shared-snapshot evidence
remain explicit premises \\
scalar attribute access & writer merges canonical \texttt{slotKey},
stores \texttt{modelKey}, and encodes the payload through
\texttt{CanonicalScalarValueCodec}; renderer follows
\texttt{ObjectHasAttribute}, requires exact
\texttt{(modelKey,attributeKey)}, and decodes according to the bound
static type & PA5/AI-19 are aligned for admitted scalar attributes;
malformed, wrong-tag, collection, and reference payloads are outside
successful certification \\
association lookup & writer merges by canonical \texttt{linkKey} and
stores relationship \texttt{modelKey}; renderer constrains scoped
endpoints/relationship,
\texttt{type(r)\ STARTS\ WITH\ \textquotesingle{}Link\textquotesingle{}},
roles, and exact \texttt{r.associationKey=\$p} & PA6 is aligned for
binary ordinary associations; ternary links are outside the profile and
association-class navigation is now rejected at admission \\
qualifiers & writer encodes ordered end-indexed values with the same
typed wire contract used by scalar storage; renderer selects and encodes
by bound direction/type; graph pull decodes each payload using the
corresponding association-end declaration; bottom cannot match a link &
PA7/AI-20 are aligned statically; fresh selected-runtime evidence is
pending \\
inheritance/allInstances & writer materializes runtime class plus
transitive parents under one \texttt{modelKey}; renderer requires scoped
object/class nodes and uses \texttt{RETURN\ DISTINCT} inside a COLLECT
subquery & PA4/PA8 and AI-18/AI-21 are aligned statically; duplicate
paths preserve finite-set cardinality \\
renderer accessor matrix & typed plan constructors are checked against
exact context/id, scalar attribute, type, navigation, qualifier, and
allInstances templates and canonical parameter values & PA9 discharged
for the checked admitted-constructor profile; raw-AST parse/render
alpha-equivalence remains open \\
bottom representation/admission & semantic set-bottom remains the sole
immutable non-null tagged-map token; stored scalar bottom is the
distinct typed payload \texttt{v1\textbar{}V}; external-parameter guards
and graph scans reject collisions/malformed payloads & PA10/AI-10 are
separated from ordinary strings statically; selected-runtime
\texttt{DISTINCT}/cardinality recapture is pending \\
scalar domain/arithmetic & exact \texttt{BigInteger}/\texttt{BigDecimal}
compile-time evaluation, Int64 bounds, finite canonical Real conversion,
strict static-type decoding, runtime-parameter guards, and finite
observed-domain checks & PA11 and AI-14 are aligned for checked runs;
incomplete dynamic observations remain \texttt{OUT\_OF\_SCOPE} \\
bottom-safe entity receivers & bind \texttt{E.objectKey}, filter
bottom/null before node \texttt{MATCH}, then restore the unique
\texttt{(modelKey,objectKey)} node; the renderer\textquotesingle s alpha
environment keeps \texttt{self}, iterator, and generated aliases
distinct & PA12 plus AI-16/AI-17 are closed statically; fresh runtime
recapture remains pending \\
exact type & \texttt{oclIsTypeOf} is rejected by certified admission
because no proved direct runtime-class accessor is available & correctly
outside \texttt{OCL\_val}; retain a stable negative-admission diagnostic
until a separate accessor and proof are added \\
IR stage boundary & semantic invariants carry \texttt{SEMANTIC}; only
the optimizer creates an opaque current-version \texttt{Artifact};
planner entry points require \texttt{OPTIMIZED} and reject forged/stale
producer versions & AI-23 is closed at the Java API boundary; PO-18
remains the universal optimizer-preservation obligation \\
target syntax & production renders text directly from
\texttt{OclCypherPlan}; an independent typed-token/clause parser
normalization-round-trips all 47 admitted queries; constructor witnesses
cover 17/17 sealed plan constructors; 52 checked complete canonical
token/group trees agree exactly with the reviewed manifest, including a
nested-shadowed OCL fixture and canonical \texttt{UmlClass}; the pinned
Neo4j Cypher 5 parser accepts 52/52 and supplies the selected AST-kind
projection & PA13a/b discharge reference datatype/printer closure;
PA14a/b/c/d discharge the stated finite syntax/parser properties and
PO-14; universal plan closure and full internal-AST equality remain
outside the finite result \\
\end{longtable}
}

\texttt{OclGraphEncodingAdequacyTest} supplies useful regression
evidence for generated query shapes, including \texttt{use\_id},
attribute paths, navigation metadata, \texttt{allInstances}, and one
qualified navigation. It does not by itself establish
\texttt{AdapterAdequate}; the current discharge additionally depends on
the PA1-\/-PA9 matrix, production shared-snapshot certificate, mutation
checks, and the clean real-Neo4j witness. These ingredients must stay
separate in the argument so that a later failure identifies the exact
broken boundary.

The current implementation checkpoint after the AI-05-\/-AI-27 alignment
has 592 discovered \texttt{neo4j-tgg} tests and 27 \texttt{neo4j}
storage/synchronization tests. All ordinary static/unit contracts pass,
the mutation contract kills 20/20 mutants, and the regenerated 52-query
formal-tree manifest agrees with production output. Fifteen opt-in
real-Neo4j tests are skipped by the default suite. On 2026-08-14 the
configured selected runtime was executed directly and all exercised
semantic, representation, case-study, and scale groups passed after the
AI-24-\/-AI-27 repairs. Three freshness guards still report stale
historical evidence because a checked-in certificate requires a clean
common capture commit. Their hashes must not be edited to manufacture a
pass: PA/CY and the 47-query differential evidence must be recaptured
once more from the final clean revision.

On 2026-07-14, the focused Maven run comprising
\texttt{OclCypherRendererTest}, \texttt{OclGraphEncodingAdequacyTest},
and \texttt{OclValidationSemanticsAdequacyTest} passed all 18 tests
(8+5+5). This confirms historical direct-renderer regression evidence. A
current conformance run must be recorded separately after implementation
changes. Neither run executes the formal raw-AST
\texttt{withEntityReceiver} expansion and therefore does not change any
open adapter status above.

On 2026-07-31, after aligning the documented prototype accessors and
rejecting \texttt{oclIsTypeOf} at certified admission, the complete
implementation-conformance script passed 128 tests with no failures,
errors, or skips. This run discharges the corresponding static
admission/regression checks only; it did not opt in to the real-Neo4j
differential or dialect suites and therefore does not discharge
CY1-\/-CY9 or \texttt{AdapterAdequate}.

The same 2026-07-31 vocabulary pass added an executable generated-query
contract that rejects legacy display-name, association-name, and
direction-free qualifier observations. It requires exact
\texttt{associationKey} and role predicates, requires
\texttt{sourceQualifiers} for qualified forward navigation, and requires
\texttt{targetQualifiers} for qualified reverse navigation. The 47-case
static pipeline suite plus six graph-encoding shape tests passed,
including a new reverse-qualified case. The later \texttt{UmlClass}
alignment additionally makes the formal tree, raw AST, stored schema,
and production class target exact, thereby discharging PO-01. R5/R6
graph-level completeness and reflection remain \texttt{AdapterAdequate}
obligations.

The closed-world admission pass was extended on 2026-07-31. The 47
positive \texttt{OCL\_val} cases and 28 negative admission checks
passed, while 104 general compiler regression tests confirmed that
experimental support remains available only through the non-certifying
API. The full static conformance gate then passed 138 tests with no
failures, errors, or skips. These results discharge the executable
admission-boundary obligation, not any graph/runtime premise.

The PA1 pass on 2026-07-31 made the previously implicit canonical-key
domain executable: every normalized component must be non-empty and must
exclude the reserved \texttt{::} delimiter. Ten focused tests passed
(three independent key contract/collision tests and seven graph-encoding
agreement tests). The latter include inherited attribute access and
confirm that binder resolution supplies the defining owner used by the
graph writer\textquotesingle s \texttt{attributeKey}. This discharges
key constructor agreement and injectivity, but not graph-level
uniqueness, completeness, or reflection.

The PA2/PA3 pass on 2026-07-31 extended the independent representation
snapshot with physical-node-token, \texttt{use\_id}, and
\texttt{objectKey} observations. Fifteen focused static
contract/evaluator tests passed. An opt-in production-encoder run then
encoded a three-object qualified/inheritance fixture into a real Neo4j
database: all 11 reported obligations passed, PA2/PA3 had zero missing,
spurious, or detail violations, and duplicate semantic links were zero.
The same run injected duplicate-ID, missing-ID, wrong-key, and
spurious-object counterexamples; both PA2 and PA3 failed as required.
The test used a dedicated timestamped model and cleaned it afterward.
This is conformance evidence for the checked synchronization profile,
not a universal claim about arbitrary graphs or concurrent external
mutation.

The PA4 pass on 2026-07-31 added an explicit bidirectional
type-exactness obligation. Eighteen focused static tests passed,
including complete inherited membership and independent missing/spurious
type mutations. The real Neo4j fixture then passed all 12 baseline
obligations. Removing \texttt{book\_a}\textquotesingle s inherited
\texttt{Publication} edge and adding a false \texttt{book\_b} to
\texttt{Library} edge produced one PA4 missing fact and one PA4 spurious
fact, as required. This evidence exercises the production batch writer
and graph reader; PA4 remains conditional on the documented
synchronization and metamodel premises.

The PA5 pass on 2026-07-31 introduced canonical per-object
\texttt{slotKey} identity, schema migration, and a uniqueness
constraint. Twenty-three focused static tests passed. The production
encoder then passed all 13 obligations on the real Neo4j fixture with
zero PA5 differences or structural violations. A mutation changed the
expected payload and attached an additional unkeyed
\texttt{AttributeValue} for the same \texttt{(object,attributeKey)}; PA5
reported the expected missing/spurious value facts plus missing-key and
duplicate-semantic-slot details. This evidence covers scalar values in
the real fixture and the shared identity path used by complex
attributes. This does not certify their semantic decoding:
collection/reference attribute accesses are now rejected by the
certified bound policy and remain available only through the general
compiler.

The follow-up theory/implementation type audit on 2026-08-09 corrected a
genuine overstatement in the earlier profile: native USE and the
metamodel index type an ordinary multiplicity-one navigation as scalar
\texttt{NODE}, whereas the certified binder had overwritten the property
node with \texttt{Set(NODE)}. The instrumented path now preserves the
source property type and records the singleton/empty finite-set view
only on its consuming collection operation/iterator. The binder consults
USE\textquotesingle s actual result binding so a qualified upper-one
navigation that is nevertheless collection-typed remains a collection
and is normalized to certified \texttt{Set}, rather than being mistaken
for this scalar lift. \texttt{BoundAdm} rejects every other
scalar-as-collection use. All 47 coverage expressions are independently
compiled by the native USE compiler against the same fixture metamodel
and are Boolean; all 47 also pass the certified pipeline. Constructor
checks distinguish scalar attribute access, multiplicity-preserving
binary navigation, its typed to-one collection view, and
\texttt{collect} with an attribute body. Negative checks include
arbitrary scalar-as-collection misuse plus collection/reference
attributes excluded from PA5 and enumeration values excluded from the
four-value-sort certified scalar domain. The separate internal
\texttt{Void} type has only bottom as an inhabitant and does not add an
ordinary scalar value sort.

The Java correction and focused native-USE/bound-tree regression are
executable evidence, not by themselves a proof of the new bridge. The
axiom-free Lean kernel now proves \texttt{LIFT1} by exhaustive
present/absent cases for the independent source, Bound, and VA views,
and the corrected 16-constructor refinement kernel explicitly preserves
\texttt{sourceCollectionType}. This closes the semantic bridge for the
four admitted scalar receivers (\texttt{asSet}, \texttt{size},
\texttt{isEmpty}, and \texttt{notEmpty}). PO-22 is now discharged for
exactly these four consumers after the clean real-Neo4j witness and
provenance were captured on 2026-08-09; no claim is made for direct
scalar iterators or other scalar-as-collection operations.

The same audit found two additional required bridges. \texttt{PO-23}
covers the internal \texttt{Void} type assigned to the literal
\texttt{null}. The rules above make \texttt{Void} conform to every
admitted contextual type, give it only the bottom denotation, and define
its equality, branch, and set-join boundaries. This is a deliberately
small null-bottom type, not the complete OMG OCL invalid/null lattice. A
closed 22-row matrix now checks 13 admitted equality, \texttt{if}/join,
\texttt{let}, mixed-set, scalar, and entity contexts plus nine rejected
Boolean, ordered, arithmetic/property, scalar-collection, Void-only-set,
and heterogeneous-set boundaries against native USE and the Bound
pipeline. This matrix exposed and killed a real mismatch: the binder had
admitted \texttt{Set\{null\}} as \texttt{Set(Void)} although
\texttt{setJoinN} has no inferred element type for a Void-only list. All
13 admitted rows also produce the reviewed native-USE violation IDs on
real Neo4j. PO-23 is now discharged for this 13-admitted/9-rejected
contextual subset after clean committed provenance was captured on
2026-08-09; this evidence still does not establish the complete OMG
null/invalid lattice.

\texttt{PO-24} covers the abstraction from concrete production records
to the proof grammar. It is discharged for that boundary by an 11-row
reflection matrix covering every \texttt{BoundExpression} record and all
12 \texttt{SemanticExpression} targets (the resolved
\texttt{BoundProperty} case splits into attribute and navigation), plus
executable witnesses over all 47 admitted expressions. The axiom-free
Lean theorem \texttt{bound\_va\_abstraction} proves that
\texttt{alpha\_B} preserves recursive evaluation for every shared
payload-parametric primitive algebra, including the
\texttt{sourceCollectionType} payload. This does not formalize the JVM
implementation of the primitive operations or Neo4j; that stronger
instantiation remains inside PO-18.

The 2026-08-09 label audit temporarily reopened PO-21. The production
snapshot reader previously accepted class nodes selected only by key
while the renderer required \texttt{:UmlClass}; removing that label was
therefore a counterexample to the certificate/renderer observation
agreement. The reader now requires \texttt{:UmlClass} for metamodel,
type, and \texttt{allInstances} observations, and a static regression
locks this syntax. A working-revision real shared-snapshot mutation that
removes the label while preserving \texttt{classKey} and
\texttt{ObjectInstanceOf} is killed by both the certificate and renderer
observations. The subsequent clean shared-snapshot run captured that
result and re-discharged PO-21 for the canonical profile; the
counterexample remains a regression guard.

The opt-in real-Neo4j oracle then compared the generated Cypher
violation set with the native USE evaluator on the same model and object
state: 47/47 sets were equal. The initial fixture distribution was 37
all-satisfying, one all-violating, and nine mixed cases. That run was
direct differential evidence for its fixture but had lower mutation
sensitivity in the all-satisfying rows.

The PO-16 non-vacuity pass on 2026-08-01 replaces that weak fixture for
the recorded OCL\_val-47 runtime claim. Before executing Cypher, the
reviewed \texttt{oclval47-nonvacuity.tsv} fixes the context, obligation
class, exact expected violation IDs, and a semantic basis for every
C01-\/-C47 row. A multiplicity-valid fixture supplies both satisfying
and violating objects for every one of the 28 \texttt{MIXED\_REQUIRED}
rows. The independent native USE evaluator matches all fixed sets; after
canonical synchronization, real Neo4j matches the same exact sets for
47/47 queries. The observed distribution is 28 non-vacuous mixed, 19
all-pass, zero all-violate, and zero empty-context cases.

The remaining 19 all-pass cases are not unexplored fixture gaps. Their
profile arguments are: reflexivity/identity for C01, C02, C16, C24;
literal, order-implication, and complementary control for C03, C15, C26;
mandatory \texttt{Company{[}1{]}} reverse navigation and context/type
membership for C09, C11, C33, C34, C40, C41; no-overflow arithmetic from
\texttt{ScalarClosed} for C22 and C23; and non-negative finite-set
cardinality/reflexivity for C30, C31, C35, C37. Thus a violating witness
would contradict a selected model/profile premise rather than improve
test non-vacuity. \texttt{OclVal47NonVacuityContractTest} checks this
classification and every exact USE set;
\texttt{OclValRealNeo4jCoverageTest} checks both USE and Neo4j against
the pre-recorded sets. A freshness-guarded runtime manifest pins the
fixture, expected rows, renderer, encoding, driver dependency, and
runtime profile. These finite witnesses still do not establish
equivalence for arbitrary models or OCL inputs.

The PA6/PA7 pass on 2026-07-31 replaced endpoint/association-only
relationship merging with an injective length-prefixed \texttt{linkKey}
over association, ordered endpoints, and both ordered qualifier lists,
shared with the diff engine\textquotesingle s binary-link identity.
Twenty-seven focused tests passed (ten independent writer/key/schema
contracts and seventeen representation evaluator tests). The production
encoder then passed all 15 obligations on a real Neo4j fixture with two
qualified links, zero missing/spurious facts, and zero duplicate
semantic links. Mutations removing a source witness, adding a spurious
relationship, duplicating a physical \texttt{linkKey}, and changing
direction-selected qualifier payloads made PA6 and PA7 fail with the
expected missing, spurious, and structural reports. This discharges
PA6/PA7 only for binary non-link-object relationships, admitted
primitive qualifiers, complete synchronization/migration, and no
concurrent external graph mutation.

The PA8 pass on 2026-07-31 separated accessor observations from the
evaluator\textquotesingle s type-fact input and exposed a duplicate-path
defect in the renderer\textquotesingle s former pattern comprehension.
The renderer now uses a COLLECT subquery with \texttt{RETURN\ DISTINCT}.
Eighteen evaluator tests and seven graph-encoding tests passed. On the
real inherited fixture, all 16 baseline obligations passed; an extra
physical \texttt{ObjectInstanceOf} edge left
\texttt{Publication.allInstances()-\textgreater{}size()} equal to two,
while independent missing-inherited and spurious-nonconforming mutations
produced the expected PA8 missing and spurious facts. This is checked
implementation-conformance evidence under PA8\textquotesingle s
premises, not a substitute for the universal lemma above.

The PA9 pass on 2026-07-31 strengthened the generated-query verifier
from a navigation-oriented string check into a typed plan/accessor
matrix. One constructor suite covers context/id, inherited scalar
attribute, forward and reverse navigation, directional qualifiers,
\texttt{allInstances}, kind-of, and cast. The verifier also ran over all
47 admitted coverage plans. Twenty targeted renderer/parameter mutants
were killed, including legacy name/suffix access, wrong canonical keys,
missing roles, reversed qualifier properties, schema
\texttt{InstanceOf}, missing local allInstances DISTINCT, and a
corrupted type accessor. The already recorded post-PA8 real-Neo4j
differential remained 47/47. This discharges accessor-template
correspondence, not raw-text alpha-round-trip or all remaining premises
of \texttt{AdapterAdequate}.

The PA10 pass on 2026-07-31 centralized the concrete bottom value in
\texttt{OclBottomToken}, rejected externally supplied occurrences
recursively, and added a model-scoped reserved-marker scan before
document and single-rule validation. Thirteen focused static tests
passed with mutation score 20/20. The opt-in selected-Neo4j probe
passed: two null rows and one integer became a two-element distinct
represented set, both \texttt{COUNT} variants returned two, and an
injected marker changed the checker status
\texttt{PASS\ -\textgreater{}\ FAIL\ -\textgreater{}\ PASS} across
insertion and cleanup. This evidence establishes the checked PA10
premises; it does not by itself discharge PO-12 (now discharged
separately by PA12) or parser-level adapter obligations.

The PA11 pass on 2026-07-31 replaced the earlier literal-divisor
assertion with an executable scalar-domain checker. Boundary tests
detect Int64 overflow, NaN/infinity, zero divisors, inexact
Integer-to-Real conversion, unsupported scalar types, and malformed
Unicode scalar sequences; unobserved dynamic arithmetic reports
\texttt{OUT\_OF\_SCOPE}. A real Neo4j stored-value probe changed status
\texttt{PASS\ -\textgreater{}\ FAIL\ -\textgreater{}\ PASS} around an
injected NaN. The complete 47-case differential fixture was rerun with
per-invariant observed-domain checks and remained 47/47 equivalent, with
every checked \texttt{ScalarClosed} premise passing.

The certified-entry-point integration pass on 2026-08-01 separated the
public general compiler from compiled invariant validation.
\texttt{compileFile} remains an explicitly uncertified experimental API,
whereas document and single-rule validation now call
\texttt{compileFileWithCertifiedContextInvariants}; every supported
\texttt{Class::inv} therefore traverses the same pre-admission,
certified binder, and post-binding admission as the instrumented
T1-\/-T6 path. Before batch or single execution, the service checks
equality with a repeated instrumented result and runs expression-level
\texttt{ScalarClosed} using exact graph \texttt{attributeKey}
observations. The Research Tool now runs graph \texttt{BottomSeparated},
stored scalar domain checks, per-invariant USE-observed
\texttt{ScalarClosed}, and generated-bottom parameter checks before
executing Cypher. The former tautological fixture bottom assertion was
removed. Five focused integration/premise tests passed; the full static
gate passed 247/247 with mutation score 20/20 and proof sync PASS. The
enhanced runtime follow-up then passed on Neo4j Kernel 2026.06.0
Enterprise/Cypher 5/database \texttt{demo}: all 47 USE-\/-Neo4j
violation sets were equivalent with the new graph premise gates, the
then-current five dialect smoke probes passed, and the
bottom/scalar/bottom-safe-receiver probes passed 3/3 without skips. This
discharges the checked entry-point profile for that runtime, not
universal \texttt{AdapterAdequate}. The one-to-one CY1-\/-CY9 matrix is
discharged by the later PO-15 runtime pass recorded below. The selected
finite parser/internal-AST projection is discharged separately by PA14d
below.

The PA14d pass on 2026-08-01 added an actual parser boundary without
changing the production architecture. The test-only bridge loads
\texttt{cypher-parser-factory-2026.06.0.jar}, parses 52/52 formal-tree
queries (47 certified), and deterministically projects
Neo4j\textquotesingle s internal Scala AST. Every constructor-directed
\texttt{AstKind} requirement now needs witnesses from both the
independent closed-fragment parser and the actual Neo4j AST. Direction,
\texttt{COUNT}/\texttt{COLLECT}, syntax rejection, and nested-shadow
scope/alias probes pass. At the PA14d checkpoint the complete static
gate passed 250/250 with mutation score 20/20 and proof sync PASS. This
closes selected finite parser acceptance/AST-kind projection, not full
AST equality, semantic-analysis equivalence, or universal TXT5
implementation closure.

The PO-15 selected-runtime pass on 2026-08-01 executes exactly one
isolated oracle row for each CY1-\/-CY9 assumption through
\texttt{Cypher5ValDialectRealNeo4jTest}. Its transaction-scoped fixture
includes two physical link paths to one semantic target, both navigation
directions, exact \texttt{associationKey}/role/directional-qualifier
predicates, represented bottom, null, duplicates, explicit aliases, and
nested guarded branches. All 9/9 rows passed on Neo4j Kernel 2026.06.0
Enterprise, Cypher 5, database \texttt{demo}, with the Maven driver
dependency pinned to 5.21.0. The checked TSV evidence pins that profile
and SHA-256 values for the renderer, canonical encoding, probe matrix,
and runtime harness. Its static freshness/mutation guard passed 2/2 and
rejects duplicate/failed rows, altered observations, and stale source
hash; the complete static gate then passed offline from Maven cache,
252/252 with mutation score 20/20 and proof sync PASS. This discharges
the selected runtime-evidence obligation, not the axiomatic status of
CY1-\/-CY9 or universal Neo4j semantics.

The PA12 pass on 2026-07-31 exposed a selected-Neo4j failure in the
first prototype attempt that carried a typed node-or-null value through
\texttt{all(...)}. The final lowering carries only
\texttt{receiver.objectKey}, filters null before \texttt{MATCH}, and
re-identifies the node by the PA2 canonical key. Two static tests cover
all three receiver families and nested \texttt{self}/iterator
references. The real fixture executes bottom and entity arms: attribute,
kind, and cast each produce exactly \texttt{\{p0\}}, while the
alias-sensitive nested invariant produces exactly \texttt{\{p2\}}. The
post-change differential regression remains 47/47.

The first PA13 pass on 2026-07-31 added a closed Java raw Cypher AST and
total structural printer matching the Section 6.6 constructor families.
Five tests compare rendered instances with every sealed permitted
subclass, reject syntax injection through typed atoms, verify immutable
children, deterministic capture-free fresh names, and explicit unique
CALL imports in every union arm. Together with six query-model tests the
focused run passed 11/11; the complete implementation-conformance run
passed 231/231 with mutation score 20/20. This discharges standalone raw
datatype/printer totality only. Production retains direct-text rendering
by design. PA14a/b/c/d establish checked normalization,
constructor-witness, finite production-derived canonical
token/group-tree, and selected Neo4j-parser AST-kind/projection
agreement including nested shadowing. After aligning every class target
to \texttt{UmlClass} and regenerating the reviewed 52-tree manifest,
these checks discharge PO-14 for the finite admitted profile. Universal
plan closure and full internal-AST equality remain incomplete; therefore
TXT5 is still not a universal prototype-correctness result.

\begin{center}\rule{0.5\linewidth}{0.5pt}\end{center}

\section{5. Validation Algebra Syntax and
Semantics}\label{5-validation-algebra-syntax-and-semantics}

\subsection{5.1 Core Syntax}\label{51-core-syntax}

The Validation Algebra (VA) contains typed expressions:

\begin{Shaded}
\begin{Highlighting}[]
\NormalTok{v ::= Literal(c)}
\NormalTok{    | SetLiteral([v1,...,vn])}
\NormalTok{    | Self}
\NormalTok{    | Var(x)}
\NormalTok{    | Attribute(v,a)}
\NormalTok{    | NavigationOne(v,A,fromRole,toRole,direction,qbar)}
\NormalTok{    | NavigationMany(v,A,fromRole,toRole,direction,qbar)}
\NormalTok{    | ViewSet(v,D)}
\NormalTok{    | AllInstances(C)}
\NormalTok{    | Not(v)}
\NormalTok{    | And(v,v) | Or(v,v) | Xor(v,v) | Implies(v,v)}
\NormalTok{    | Compare(op,v,v)}
\NormalTok{    | Arith(op,v,v)}
\NormalTok{    | Coerce[tau{-}\textgreater{}upsilon](v)}
\NormalTok{    | If(v,v,v)}
\NormalTok{    | Let(x,tauD,v,v)}
\NormalTok{    | Exists(v,x,tauD,v)}
\NormalTok{    | ForAll(v,x,tauD,v)}
\NormalTok{    | Select(v,x,tauD,v)}
\NormalTok{    | Reject(v,x,tauD,v)}
\NormalTok{    | Collect(v,x,tauD,v)}
\NormalTok{    | Includes(v,v)}
\NormalTok{    | Excludes(v,v)}
\NormalTok{    | IncludesAll(v,v)}
\NormalTok{    | ExcludesAll(v,v)}
\NormalTok{    | Union(v,v) | Intersection(v,v) | AsSet(v)}
\NormalTok{    | IsUnique(v,x,tauD,v)}
\NormalTok{    | Count(v) | Size(v)}
\NormalTok{    | IsEmpty(v) | NotEmpty(v)}
\NormalTok{    | TypeKindOf(v,C)}
\NormalTok{    | Cast(v,C)}
\end{Highlighting}
\end{Shaded}

All collection-valued VA expressions in the theorem denote finite sets.
In later rewrite equations, an abbreviated binder such as
\texttt{Exists(S,x,P)} means that its already-typed \texttt{tauD}
payload is carried unchanged. A specialized navigation plan may omit
that payload only for UML class upcast, whose value embedding is
identity; the body variables still retain their declared static types.

The VA typing judgment is not inferred from the target renderer. Write
\texttt{Prim=\{Boolean,Integer,Real,String\}} and define the exact
resolved-navigation witness

\begin{Shaded}
\begin{Highlighting}[]
\NormalTok{NavWitness\_MM(D,A,fr,tr,dir,E,[kappa\_1,...,kappa\_k],kind)}
\NormalTok{iff exists a source role name r:}
\NormalTok{    resolveNavigation\_MM(D,r,[kappa\_1,...,kappa\_k])}
\NormalTok{      =(A,fr,tr,dir,E,[kappa\_1,...,kappa\_k],kind).}
\end{Highlighting}
\end{Shaded}

The list in this witness is ordered and fixes both arity and declared
types; every \texttt{kappa\_i} must belong to \texttt{Prim}. The VA
typing judgment is the least syntax-directed relation
\texttt{MM;Gamma\ \textbar{}-VA\ v:tau} generated by the following rules
(the resolved metadata in navigation, attribute, and type-operation
constructors is the metadata produced by binding):

\begin{Shaded}
\begin{Highlighting}[]
\NormalTok{Gamma(self)=C                         =\textgreater{} Self:C}
\NormalTok{nearest(Gamma,x)=tau                 =\textgreater{} Var(x):tau}
\NormalTok{typeOfLiteral(c)=tau                 =\textgreater{} Literal(c):tau}

\NormalTok{vi:tau\_i (1\textless{}=i\textless{}=n), n\textgreater{}=1,}
\NormalTok{setJoinN([tau\_1,...,tau\_n])=sigma    =\textgreater{} SetLiteral([v1,...,vn]):Set(sigma)}

\NormalTok{v:D, resolveAttribute\_MM(D,a.name)=(a,tau), tau in Prim}
\NormalTok{                                      =\textgreater{} Attribute(v,a):tau}
\NormalTok{v:D, qbar=[q1,...,qk], qi:kappa\_i, kappa\_i in Prim,}
\NormalTok{NavWitness\_MM(D,A,fr,tr,dir,E,[kappa\_1,...,kappa\_k],ONE)}
\NormalTok{                                      =\textgreater{} NavigationOne(v,A,fr,tr,dir,qbar):E}
\NormalTok{v:D, qbar=[q1,...,qk], qi:kappa\_i, kappa\_i in Prim,}
\NormalTok{NavWitness\_MM(D,A,fr,tr,dir,E,[kappa\_1,...,kappa\_k],MANY)}
\NormalTok{                                      =\textgreater{} NavigationMany(v,A,fr,tr,dir,qbar):Set(E)}
\NormalTok{v:D, D in Class(MM), and v is the admitted direct ONE{-}navigation view}
\NormalTok{                                      =\textgreater{} ViewSet(v,D):Set(D)}
\NormalTok{C in Class(MM)                        =\textgreater{} AllInstances(C):Set(C)}

\NormalTok{p:Boolean                            =\textgreater{} Not(p):Boolean}
\NormalTok{p1:Boolean, p2:Boolean               =\textgreater{} And/Or/Xor/Implies(p1,p2):Boolean}
\NormalTok{e1:tau, e2:sigma, eqJoin(tau,sigma)=upsilon}
\NormalTok{                                      =\textgreater{} Compare(EQ\_OR\_NEQ,e1,e2):Boolean}
\NormalTok{e1:tau, e2:sigma, orderJoin(tau,sigma)=upsilon}
\NormalTok{                                      =\textgreater{} Compare(ORDER,e1,e2):Boolean}
\NormalTok{e1:tau, e2:sigma, arithResult(op,tau,sigma)=upsilon}
\NormalTok{                                      =\textgreater{} Arith(op,e1,e2):upsilon}
\NormalTok{e:tau, upsilon is non{-}set, and iota\_I[tau{-}\textgreater{}upsilon]}
\NormalTok{is the selected canonical embedding}
\NormalTok{                                      =\textgreater{} Coerce[tau{-}\textgreater{}upsilon](e):upsilon}
\NormalTok{c:Boolean, t:tau1, f:tau2, ifJoin(tau1,tau2)=tau}
\NormalTok{                                      =\textgreater{} If(c,t,f):tau}
\NormalTok{init:tau1, tau1\textless{}=tauD, Gamma[x:tauD] |{-}VA body:tau2}
\NormalTok{                                      =\textgreater{} Let(x,tauD,init,body):tau2}

\NormalTok{S:Set(tauE), tauE\textless{}=tauD, Gamma[x:tauD] |{-}VA P:Boolean}
\NormalTok{                                      =\textgreater{} Exists/ForAll(S,x,tauD,P):Boolean}
\NormalTok{                                      =\textgreater{} Select/Reject(S,x,tauD,P):Set(tauE)}
\NormalTok{S:Set(tauE), tauE\textless{}=tauD, Gamma[x:tauD] |{-}VA E:sigma, sigma is non{-}set}
\NormalTok{                                      =\textgreater{} Collect(S,x,tauD,E):Set(sigma)}
\NormalTok{                                      =\textgreater{} IsUnique(S,x,tauD,E):Boolean}
\NormalTok{S:Set(sigma), e:delta, memberJoin(sigma,delta)=upsilon}
\NormalTok{                                      =\textgreater{} Includes/Excludes(S,e):Boolean}
\NormalTok{S:Set(tau), T:Set(sigma), memberJoin(tau,sigma)=upsilon}
\NormalTok{                                      =\textgreater{} IncludesAll/ExcludesAll(S,T):Boolean}
\NormalTok{S:Set(tau), T:Set(sigma), setJoin2(tau,sigma)=upsilon}
\NormalTok{                                      =\textgreater{} Union/Intersection(S,T):Set(upsilon)}
\NormalTok{S:Set(tau)                            =\textgreater{} AsSet(S):Set(tau)}
\NormalTok{S:Set(tau)                            =\textgreater{} Count(S):Integer, Size(S):Integer}
\NormalTok{S:Set(tau)                            =\textgreater{} IsEmpty(S):Boolean, NotEmpty(S):Boolean}
\NormalTok{D,C in Class(MM), e:Void or e:D       =\textgreater{} TypeKindOf(e,C):Boolean}
\NormalTok{D,C in Class(MM), e:Void or e:D       =\textgreater{} Cast(e,C):C}
\end{Highlighting}
\end{Shaded}

No other VA typing rule exists. In particular, an iterator source must
already have static type \texttt{Set(tau)}; \texttt{ViewSet} is
introduced only for the four certified direct consumers \texttt{asSet},
\texttt{size}, \texttt{isEmpty}, and \texttt{notEmpty}, never as an
implicit iterator or collection-property conversion. A VA
\texttt{Coerce} node never has a Set result. The semantic
\texttt{Void-\textgreater{}Set(sigma)} embedding used for a set-valued
\texttt{If} branch or set-equality operand is realized only by the
explicit \texttt{CQEmptySet}/\texttt{setOperand} adapters, not by
\texttt{CQCoerce}. We write \texttt{MM;Gamma\ \textbar{}-\ v:tau} for
\texttt{MM;Gamma\ \textbar{}-VA\ v:tau} after this definition.

\subsubsection{5.1.1 Enforced Semantic/Optimized Stage
Contract}\label{511-enforced-semanticoptimized-stage-contract}

The Java records reuse several payload constructors across stages, so
stage membership is not inferred from a record\textquotesingle s class
name. It is an explicit, checked certificate:

\begin{Shaded}
\begin{Highlighting}[]
\NormalTok{SemanticInvariant = (C,R,v, SEMANTIC, "semantic{-}ir{-}v1")}
\NormalTok{T\_OPT(SemanticInvariant)}
\NormalTok{  = (C,R,nv, OPTIMIZED, optimizerVersion)}

\NormalTok{planInvariant(q) is defined only if}
\NormalTok{  q.stage = OPTIMIZED and q.producerVersion = currentOptimizerVersion.}
\end{Highlighting}
\end{Shaded}

For top-level expression planning, \texttt{T\_OPT} returns an opaque
\texttt{Artifact} whose constructor is not public; the public planner
accepts that artifact rather than an arbitrary \texttt{Expression}. Thus
a caller cannot obtain a production plan by passing a semantic query
directly or by presenting an artifact issued by an older optimizer
version. This is the executable stage boundary corresponding to Theorem
3.

\subsection{5.2 Denotational Semantics}\label{52-denotational-semantics}

Let \texttt{I} be either the object interpretation or the graph
interpretation. The denotation of VA expression \texttt{v} under
environment \texttt{rho} is written:

\begin{Shaded}
\begin{Highlighting}[]
\NormalTok{[[v]]\_I(rho).}
\end{Highlighting}
\end{Shaded}

Within the theorem domain, denotation is total into \texttt{Val}.
Primitive accessors and operations use the following validation-level
completion:

\begin{Shaded}
\begin{Highlighting}[]
\NormalTok{I\_attr(bottom,a) = bottom}
\NormalTok{I\_attr(entity,a) = the typed canonical attribute value}

\NormalTok{asSources(bottom) = empty}
\NormalTok{I\_nav is invoked only on Entity values admitted by T{-}Nav}

\NormalTok{eq\_val\^{}tau(bottom\_tau,bottom\_tau) = true, for non{-}set tau}
\NormalTok{eq\_val\^{}tau(bottom\_tau,v) = false for non{-}set v != bottom\_tau}

\NormalTok{eq\_val\^{}Set(sigma)(A,B)}
\NormalTok{  = true iff finiteSet\_sigma(A)=finiteSet\_sigma(B)}

\NormalTok{an undefined arithmetic operation, including division by zero, yields bottom}
\NormalTok{an ordered comparison with an undefined/incompatible operand yields bottom}

\NormalTok{TypeKindOf(bottom,C) = false}
\NormalTok{Cast(bottom,C) = bottom}
\end{Highlighting}
\end{Shaded}

The equality clauses are a deliberate total typed-equality policy needed
by finite-set membership and its normalization to existential equality.
At Set type, equality is equality of the certified collection
observation, so \texttt{bottom\_Set(sigma)} equals the empty finite set;
this follows the same whole-collection-bottom completion used by every
set operator. They are not the OMG OCL invalid equality rules. Static
typing excludes all other ill-sorted primitive applications.
\texttt{Collect} retains \texttt{bottom} as an element of its finite-set
image. Its admitted body type is non-collection, so neither implicit
flattening nor nested collection realization is claimed.

Core denotations:

\begin{Shaded}
\begin{Highlighting}[]
\NormalTok{[[Literal(c)]]\_I(rho) = c}

\NormalTok{[[SetLiteral\_tau1,...,taun([v1,...,vn])]]\_I(rho)}
\NormalTok{  = \{ iota\_I[tau\_i {-}\textgreater{} upsilon]([[vi]]\_I(rho)) | 1 \textless{}= i \textless{}= n \},}
\NormalTok{    where setJoinN([tau1,...,taun])=upsilon}

\NormalTok{[[Self]]\_I(rho) = rho(self)}

\NormalTok{[[Var(x)]]\_I(rho) = rho(x), or bottom if x is unbound}

\NormalTok{[[Attribute(e,a)]]\_I(rho)}
\NormalTok{  = I\_attr([[e]]\_I(rho), a)}

\NormalTok{targets\_I(e,A,fromRole,toRole,direction,qbar,rho)}
\NormalTok{  = empty, if qualifierValues(qbar,I,rho) is undefined}
\NormalTok{  = union \{ I\_nav(s,A,fromRole,toRole,direction,}
\NormalTok{                  qualifierValues(qbar,I,rho))}
\NormalTok{            | s in asSources([[e]]\_I(rho)) \}, otherwise}

\NormalTok{[[NavigationMany(e,A,fromRole,toRole,direction,qbar)]]\_I(rho)}
\NormalTok{  = targets\_I(e,A,fromRole,toRole,direction,qbar,rho)}

\NormalTok{[[NavigationOne(e,A,fromRole,toRole,direction,qbar)]]\_I(rho)}
\NormalTok{  = oneOrBottom\_I(}
\NormalTok{      targets\_I(e,A,fromRole,toRole,direction,qbar,rho))}

\NormalTok{oneOrBottom\_I(empty)=bottom\_D}
\NormalTok{oneOrBottom\_I(\{d\})=d}
\NormalTok{oneOrBottom\_I(S) is undefined when |S|\textgreater{}1}

\NormalTok{[[ViewSet(e,D)]]\_I(rho)}
\NormalTok{  = empty, when [[e]]\_I(rho)=bottom\_D}
\NormalTok{  = \{d\},   when [[e]]\_I(rho)=d}

\NormalTok{[[AllInstances(C)]]\_I(rho)}
\NormalTok{  = I\_class(C)}
\end{Highlighting}
\end{Shaded}

Qualifier expressions are evaluated before navigation matching:

\begin{Shaded}
\begin{Highlighting}[]
\NormalTok{qualifierValues([],obj,rho) = []}
\NormalTok{qualifierValues([q1,...,qk],obj,rho)}
\NormalTok{  = undefined, if some [[qi]]\_obj(rho)=bottom}
\NormalTok{  = [[[q1]]\_obj(rho),...,[[qk]]\_obj(rho)], otherwise}

\NormalTok{qualifierValues([],graph,eta) = []}
\NormalTok{qualifierValues([q1,...,qk],graph,eta)}
\NormalTok{  = undefined, if some [[qi]]\_graph(eta)=bottom}
\NormalTok{  = encodeQualifierList(}
\NormalTok{      [[[q1]]\_graph(eta),...,[[qk]]\_graph(eta)]), otherwise}
\end{Highlighting}
\end{Shaded}

The object interpretation compares the raw primitive qualifier. The
graph interpretation first evaluates the corresponding graph-side scalar
and then encodes it before comparison with
\texttt{linkQualifiers\_G(r,direction)}. Qualifier preservation is
claimed only for primitive scalar qualifiers covered by the injective
serialization clause of the graph encoding contract.

Finite-set operators:

\begin{Shaded}
\begin{Highlighting}[]
\NormalTok{[[Exists(S,x,tauD,P)]]\_I(rho)}
\NormalTok{  = true iff exists s in finiteSet\_tauE([[S]]\_I(rho)):}
\NormalTok{        bool\_val([[P]]\_I(rho[x {-}\textgreater{} iota\_I[tauE{-}\textgreater{}tauD](s)])) = true}

\NormalTok{[[ForAll(S,x,tauD,P)]]\_I(rho)}
\NormalTok{  = true iff for all s in finiteSet\_tauE([[S]]\_I(rho)):}
\NormalTok{        bool\_val([[P]]\_I(rho[x {-}\textgreater{} iota\_I[tauE{-}\textgreater{}tauD](s)])) = true}

\NormalTok{[[Select(S,x,tauD,P)]]\_I(rho)}
\NormalTok{  = \{ s in finiteSet\_tauE([[S]]\_I(rho))}
\NormalTok{      | bool\_val([[P]]\_I(rho[x {-}\textgreater{} iota\_I[tauE{-}\textgreater{}tauD](s)])) = true \}}

\NormalTok{[[Reject(S,x,tauD,P)]]\_I(rho)}
\NormalTok{  = \{ s in finiteSet\_tauE([[S]]\_I(rho))}
\NormalTok{      | bool\_val([[P]]\_I(rho[x {-}\textgreater{} iota\_I[tauE{-}\textgreater{}tauD](s)])) = false \}}

\NormalTok{[[Collect(S,x,tauD,E)]]\_I(rho)}
\NormalTok{  = \{ [[E]]\_I(rho[x {-}\textgreater{} iota\_I[tauE{-}\textgreater{}tauD](s)])}
\NormalTok{      | s in finiteSet\_tauE([[S]]\_I(rho)) \}}

\NormalTok{[[Includes(S,E)]]\_I(rho)}
\NormalTok{  = true iff}
\NormalTok{      iota\_I[delta {-}\textgreater{} upsilon]([[E]]\_I(rho))}
\NormalTok{      in upSet\_I[sigma {-}\textgreater{} upsilon](finiteSet\_sigma([[S]]\_I(rho))),}
\NormalTok{    where S:Set(sigma), E:delta, memberJoin(sigma,delta)=upsilon}

\NormalTok{[[Excludes(S,E)]]\_I(rho)}
\NormalTok{  = not [[Includes(S,E)]]\_I(rho)}

\NormalTok{[[IncludesAll(S,T)]]\_I(rho)}
\NormalTok{  = true iff}
\NormalTok{      upSet\_I[sigma {-}\textgreater{} upsilon](finiteSet\_sigma([[T]]\_I(rho)))}
\NormalTok{      subseteq}
\NormalTok{      upSet\_I[tau {-}\textgreater{} upsilon](finiteSet\_tau([[S]]\_I(rho))),}
\NormalTok{    where memberJoin(tau,sigma)=upsilon}

\NormalTok{[[ExcludesAll(S,T)]]\_I(rho)}
\NormalTok{  = true iff}
\NormalTok{      upSet\_I[tau {-}\textgreater{} upsilon](finiteSet\_tau([[S]]\_I(rho)))}
\NormalTok{      intersect}
\NormalTok{      upSet\_I[sigma {-}\textgreater{} upsilon](finiteSet\_sigma([[T]]\_I(rho)))}
\NormalTok{      = empty,}
\NormalTok{    where memberJoin(tau,sigma)=upsilon}

\NormalTok{[[Union(S,T)]]\_I(rho)}
\NormalTok{  = upSet\_I[tau {-}\textgreater{} upsilon](finiteSet\_tau([[S]]\_I(rho)))}
\NormalTok{    union upSet\_I[sigma {-}\textgreater{} upsilon](finiteSet\_sigma([[T]]\_I(rho))),}
\NormalTok{    where setJoin2(tau,sigma)=upsilon}

\NormalTok{[[Intersection(S,T)]]\_I(rho)}
\NormalTok{  = upSet\_I[tau {-}\textgreater{} upsilon](finiteSet\_tau([[S]]\_I(rho)))}
\NormalTok{    intersect upSet\_I[sigma {-}\textgreater{} upsilon](finiteSet\_sigma([[T]]\_I(rho))),}
\NormalTok{    where setJoin2(tau,sigma)=upsilon}

\NormalTok{[[AsSet(S)]]\_I(rho) = finiteSet\_tau([[S]]\_I(rho))}

\NormalTok{[[IsUnique(S,x,tauD,E)]]\_I(rho) = true iff}
\NormalTok{  for all s1,s2 in finiteSet\_tauE([[S]]\_I(rho)):}
\NormalTok{    eq\_val\^{}sigma(}
\NormalTok{      [[E]]\_I(rho[x {-}\textgreater{} iota\_I[tauE{-}\textgreater{}tauD](s1)]),}
\NormalTok{      [[E]]\_I(rho[x {-}\textgreater{} iota\_I[tauE{-}\textgreater{}tauD](s2)]))}
\NormalTok{    implies eq\_val\^{}tauE(s1,s2)}

\NormalTok{[[Count(S)]]\_I(rho) = |finiteSet([[S]]\_I(rho))|}

\NormalTok{[[Size(S)]]\_I(rho) = |finiteSet([[S]]\_I(rho))|}

\NormalTok{[[IsEmpty(S)]]\_I(rho) = true iff finiteSet([[S]]\_I(rho)) = empty}

\NormalTok{[[NotEmpty(S)]]\_I(rho) = true iff finiteSet([[S]]\_I(rho)) \textless{}\textgreater{} empty}
\end{Highlighting}
\end{Shaded}

\texttt{Collect} is finite-set image construction with a non-collection
body type only. The theorem does not preserve Bag multiplicity, result
ordering, implicit flattening, or nested collection results of full OMG
OCL \texttt{collect}.

Boolean and scalar operators:

\begin{Shaded}
\begin{Highlighting}[]
\NormalTok{[[Not(E)]]\_I(rho)}
\NormalTok{  = true iff bool\_val([[E]]\_I(rho)) = false}

\NormalTok{[[And(A,B)]]\_I(rho)}
\NormalTok{  = true iff bool\_val([[A]]\_I(rho)) = true}
\NormalTok{         and bool\_val([[B]]\_I(rho)) = true}

\NormalTok{[[Or(A,B)]]\_I(rho)}
\NormalTok{  = true iff bool\_val([[A]]\_I(rho)) = true}
\NormalTok{          or bool\_val([[B]]\_I(rho)) = true}

\NormalTok{[[Xor(A,B)]]\_I(rho)}
\NormalTok{  = true iff bool\_val([[A]]\_I(rho))}
\NormalTok{             != bool\_val([[B]]\_I(rho))}

\NormalTok{[[Implies(A,B)]]\_I(rho)}
\NormalTok{  = true iff bool\_val([[A]]\_I(rho)) = false}
\NormalTok{          or bool\_val([[B]]\_I(rho)) = true}

\NormalTok{[[Compare(op,A,B)]]\_I(rho)}
\NormalTok{  = eq\_val\^{}upsilon(}
\NormalTok{      iota\_I[tau {-}\textgreater{} upsilon]([[A]]\_I(rho)),}
\NormalTok{      iota\_I[sigma {-}\textgreater{} upsilon]([[B]]\_I(rho))) for op = EQ,}
\NormalTok{    where eqJoin(tau,sigma)=upsilon}
\NormalTok{  = Boolean negation of the EQ equation for op = NEQ}
\NormalTok{  = ordered scalar comparison for op in \{LT,LE,GT,GE\} when operands are}
\NormalTok{    compatible ordered scalars; bottom otherwise}

\NormalTok{[[Arith(op,A,B)]]\_I(rho)}
\NormalTok{  = numeric result if operands are numeric and operation is defined;}
\NormalTok{    bottom otherwise}

\NormalTok{[[Coerce[tau{-}\textgreater{}upsilon](E)]]\_I(rho)}
\NormalTok{  = iota\_I[tau{-}\textgreater{}upsilon]([[E]]\_I(rho))}
\end{Highlighting}
\end{Shaded}

Control and binding:

\begin{Shaded}
\begin{Highlighting}[]
\NormalTok{[[If\_tau(C,T:tauT,E:tauE)]]\_I(rho)}
\NormalTok{  = iota\_I[tauT{-}\textgreater{}tau]([[T]]\_I(rho))}
\NormalTok{      if bool\_val([[C]]\_I(rho)) = true}
\NormalTok{  = iota\_I[tauE{-}\textgreater{}tau]([[E]]\_I(rho))}
\NormalTok{      otherwise,}
\NormalTok{    where ifJoin(tauT,tauE)=tau}

\NormalTok{[[Let(x,tauD,E:tau1,B)]]\_I(rho)}
\NormalTok{  = [[B]]\_I(rho[x {-}\textgreater{} iota\_I[tau1{-}\textgreater{}tauD]([[E]]\_I(rho))])}
\end{Highlighting}
\end{Shaded}

Type operations:

\begin{Shaded}
\begin{Highlighting}[]
\NormalTok{[[TypeKindOf(E,C)]]\_I(rho)}
\NormalTok{  = true iff [[E]]\_I(rho) is an Entity conforming to C in interpretation I;}
\NormalTok{    false otherwise}

\NormalTok{[[Cast(E,C)]]\_I(rho)}
\NormalTok{  = [[E]]\_I(rho) if [[E]]\_I(rho) conforms to C; bottom otherwise}
\end{Highlighting}
\end{Shaded}

\begin{center}\rule{0.5\linewidth}{0.5pt}\end{center}

\section{6. Lemma Catalogue}\label{6-lemma-catalogue}

The theorem chain uses the following lemmas.

\subsection{6.1 Representation Lemmas}\label{61-representation-lemmas}

\textbf{R1 Object Injectivity.}

\begin{Shaded}
\begin{Highlighting}[]
\NormalTok{For all o1,o2 in Obj(M):}
\NormalTok{node(o1) = node(o2) iff o1 = o2.}
\end{Highlighting}
\end{Shaded}

Proof. For the forward implication, assume \texttt{node(o1)=node(o2)=n}.
By the object-node clauses, \texttt{prop(n,"use\_id")=id(o1)} and
\texttt{prop(n,"use\_id")=id(o2)}; functionality of a property map gives
\texttt{id(o1)=id(o2)}, and injectivity of \texttt{id} gives
\texttt{o1=o2}. For the reverse implication, substitute \texttt{o2=o1};
functionality of \texttt{node} gives \texttt{node(o1)=node(o1)}. Thus
both implications hold.

\textbf{R2 Object-Node Exactness.}

\begin{Shaded}
\begin{Highlighting}[]
\NormalTok{For every o in Obj(M), there exists a unique n in Node(G) such that:}
\NormalTok{  Rep\_G(o,n).}

\NormalTok{For every n in ObjNode\_val(G), there exists a unique o in Obj(M) such that:}
\NormalTok{  Rep\_G(o,n).}
\end{Highlighting}
\end{Shaded}

The notation \texttt{node(o)} denotes this unique \texttt{n}, and
therefore \texttt{prop(node(o),"use\_id")=id(o)}.

Proof. Fix \texttt{o}. P-Object creates a witness \texttt{n} with
\texttt{Rep\_G(o,n)}. If both \texttt{n1} and \texttt{n2} represent
\texttt{o}, P-Object\textquotesingle s exactly-one clause gives
\texttt{n1=n2}, proving existence and uniqueness. Conversely, fix
\texttt{n\ in\ ObjNode\_val(G)}. P-Exact says that every such node was
introduced by P-Object, so some \texttt{o} satisfies
\texttt{Rep\_G(o,n)}. If \texttt{o1} and \texttt{o2} both do, then
\texttt{n=node(o1)=node(o2)} and R1 gives \texttt{o1=o2}. Therefore
\texttt{Rep\_G} is total and single-valued in both directions, so
\texttt{node} is a bijection between the two displayed domains.

\textbf{R3 Type Preservation.}

\begin{Shaded}
\begin{Highlighting}[]
\NormalTok{o conformsTo C}
\NormalTok{iff}
\NormalTok{node(o) in I\_class\_graph(C).}
\end{Highlighting}
\end{Shaded}

Equivalently:

\begin{Shaded}
\begin{Highlighting}[]
\NormalTok{o conformsTo C}
\NormalTok{iff}
\NormalTok{(node(o)){-}[:ObjectInstanceOf]{-}\textgreater{}(classNode(C)).}
\end{Highlighting}
\end{Shaded}

Proof. If \texttt{o\ conformsTo\ C}, P-Type emits the displayed edge,
hence \texttt{node(o)\ in\ I\_class\_graph(C)}. Conversely, membership
in \texttt{I\_class\_graph(C)} supplies a validation-visible
\texttt{ObjectInstanceOf} edge from \texttt{node(o)} to
\texttt{classNode(C)}. By no-spurious-type-edge/P-Exact, such an edge is
generated only for a true conformance pair, so
\texttt{o\ conformsTo\ C}. These are the two required implications.

\textbf{R4 Attribute Preservation.}

\begin{Shaded}
\begin{Highlighting}[]
\NormalTok{attrVal\_M(o,a) = I\_attr\_graph(node(o),a).}
\end{Highlighting}
\end{Shaded}

Proof. P-Attr creates exactly one attribute-value witness \texttt{av}
for \texttt{(o,a)} and stores \texttt{storeScalar(attrVal\_M(o,a))}. By
definition, \texttt{I\_attr\_graph} applies \texttt{readScalar} to that
unique property. The scalar representation law gives
\texttt{readScalar(storeScalar(v))=v}; substituting
\texttt{v=attrVal\_M(o,a)} proves the equality. P-Exact excludes a
second validation-visible slot that could make the accessor ambiguous.
For an absent/undefined admitted slot, both accessors use the same
declared bottom completion.

The next two lemmas are relative representation results. Their trusted
premises are exposed explicitly rather than hidden inside the phrase "by
the encoding contract":

{\def\LTcaptype{none} % do not increment counter
\begin{longtable}[]{@{}ll@{}}
\toprule\noalign{}
Result & Exact dependencies \\
\midrule\noalign{}
\endhead
\bottomrule\noalign{}
\endlastfoot
R5 association preservation & P-Link completeness;
P-Exact/no-spurious-link reflection; R1 object injectivity; R2
object-node preservation; injectivity of \texttt{associationKey\_MM};
exact source/target role storage; direction-indexed qualifier access;
fixed-arity componentwise injectivity of \texttt{encodeQualifierList} \\
R6 navigation preservation & R5; normalization of direction and role
orientation; equality between evaluated raw qualifiers and the
direction-selected encoding used by \texttt{nav\_graph}; R2 reflection
of graph object nodes; extensional finite-set semantics \\
\end{longtable}
}

\textbf{R5 Association Preservation.}

\begin{Shaded}
\begin{Highlighting}[]
\NormalTok{(o1,sr,sourceQs,o2,tr,targetQs) in links\_M(A)}
\NormalTok{iff}
\NormalTok{exists r in RelLink\_val(G,modelKey\_MM):}
\NormalTok{  src(r)=node(o1)}
\NormalTok{  and trg(r)=node(o2)}
\NormalTok{  and linkAssociation\_G(r)=associationKey\_MM(A)}
\NormalTok{  and linkSourceRole\_G(r)=sr}
\NormalTok{  and linkTargetRole\_G(r)=tr}
\NormalTok{  and decodeQualifierList(r,forward)=sourceQs}
\NormalTok{  and decodeQualifierList(r,reverse)=targetQs.}
\end{Highlighting}
\end{Shaded}

Proof. Forward: take a source link tuple in \texttt{links\_M(A)}. P-Link
creates a relationship with endpoints \texttt{node(o1),node(o2)},
canonical association key, roles, model scope, \texttt{Link*} type, and
both direction-indexed encoded qualifier lists. Strict decode/encode
retraction yields the displayed semantic lists, so the relationship
witnesses the right-hand side. Backward: let \texttt{r} witness the
right-hand side. P-Exact/no-spurious-link implies that \texttt{r} was
generated by P-Link from some source tuple
\texttt{(A\textquotesingle{},o1\textquotesingle{},sr\textquotesingle{},sourceQs\textquotesingle{},o2\textquotesingle{},tr\textquotesingle{},targetQs\textquotesingle{})}.
R1/R2 identify \texttt{o1\textquotesingle{}=o1} and
\texttt{o2\textquotesingle{}=o2}; injectivity of
\texttt{associationKey\_MM} gives \texttt{A\textquotesingle{}=A};
equality of stored roles gives
\texttt{sr\textquotesingle{}=sr,tr\textquotesingle{}=tr}; and strict
typed qualifier decoding at each end\textquotesingle s declared arity
gives \texttt{sourceQs\textquotesingle{}=sourceQs} and
\texttt{targetQs\textquotesingle{}=targetQs}. Hence the required source
tuple exists. Relationship uniqueness is unnecessary because semantic
links and later results use set semantics.

\textbf{R6 Navigation Preservation.}

\begin{Shaded}
\begin{Highlighting}[]
\NormalTok{rawQs = [ [[q1]]\_obj(rho), ..., [[qk]]\_obj(rho) ]}

\NormalTok{if no rawQs component is bottom, encodedQs=encodeQualifierList(rawQs) and:}

\NormalTok{  \{ node(y) | y in navMany\_obj(o,A,fromRole,toRole,direction,rawQs) \}}
\NormalTok{  = navMany\_graph(node(o),A,fromRole,toRole,direction,encodedQs).}

\NormalTok{if some rawQs component is bottom, both displayed navigation results are}
\NormalTok{empty and the graph side uses the distinguished NoMatch qualifier result.}

\NormalTok{For resultKind=ONE, valid multiplicity gives cardinality at most one and:}

\NormalTok{  encodeValue(oneOrBottom\_obj(navMany\_obj(...)))}
\NormalTok{  = oneOrBottom\_graph(navMany\_graph(...)).}
\end{Highlighting}
\end{Shaded}

Proof. In the non-bottom case, prove the \texttt{MANY} equation by two
inclusions. A source target has a link with the resolved association,
direction, roles, and raw qualifiers; R5 gives a relationship in
\texttt{RelLink\_val}, and strict qualifier encoding gives the graph
match. Conversely, R5 reflects each graph match to a source link, while
R2 reflects its endpoint to the unique source object. In the
bottom-qualifier case, both semantic definitions select no link, so both
sets are empty. For \texttt{ONE}, apply the \texttt{MANY} equality. It
maps empty to empty and a singleton \texttt{\{y\}} to
\texttt{\{node(y)\}}; the two equations of \texttt{oneOrBottom} and R1
then give the scalar/bottom encoding equation.

The dependencies above are necessary, not merely convenient. If
\texttt{associationKey\_MM} is not injective, relationships for two
distinct associations can satisfy the same emitted metadata predicate,
invalidating the backward direction of R5. If qualifier serialization is
not injective at the declared arity, two distinct qualified links can
become observationally indistinguishable. If direction or stored roles
are not normalized consistently, R5 may still identify a source link
while R6 returns the opposite endpoint. These countermodels explain why
R5/R6 are claimed only for encodings that discharge the listed premises.

\textbf{R7 allInstances Preservation.}

\begin{Shaded}
\begin{Highlighting}[]
\NormalTok{\{ node(o) | o in allInstances\_obj(C) \}}
\NormalTok{= allInstances\_graph(C).}
\end{Highlighting}
\end{Shaded}

Proof. For arbitrary graph object node \texttt{n}, R2 gives a unique
\texttt{o} with \texttt{n=node(o)}. Then:

\begin{Shaded}
\begin{Highlighting}[]
\NormalTok{n in \{node(o) | o in allInstances\_obj(C)\}}
\NormalTok{iff o conformsTo C}
\NormalTok{iff node(o) in I\_class\_graph(C)                 (R3)}
\NormalTok{iff n in allInstances\_graph(C).}
\end{Highlighting}
\end{Shaded}

Since membership is equivalent for every validation-visible object node,
the two finite sets are extensionally equal.

\subsection{6.2 Transformation Lemmas}\label{62-transformation-lemmas}

\textbf{F1 OCL\_val Representation and Bound Rebinding.} The admission
relation is:

\begin{Shaded}
\begin{Highlighting}[]
\NormalTok{Adm\_MM(astInv)}
\NormalTok{iff}
\NormalTok{RawAdm(astInv)}
\NormalTok{and decodeInvariant\_val(astInv)=I is defined}
\NormalTok{and T\_BIND\_INV(astInv,MM)=bInv is defined}
\NormalTok{and BoundAdm\_MM(bInv).}
\end{Highlighting}
\end{Shaded}

For every source invariant \texttt{I=(cName,R,s)} with
\texttt{s\ in\ OCL\_val}, \texttt{encode\_inv(I)} belongs to
\texttt{ASTInv\_val} and
\texttt{decodeInvariant\_val(encode\_inv(I))=I}. If binding succeeds
with \texttt{resolveClass\_MM(cName)=C} and
\texttt{T\_BIND\_INV(encode\_inv(I),MM)=BInvariant(C,R,v)}, then:

\begin{Shaded}
\begin{Highlighting}[]
\NormalTok{1. encode\_inv(I) conforms to the invariant OCL AST metamodel.}
\NormalTok{2. decodeInvariant\_val(encode\_inv(I)) = I.}
\NormalTok{3. erase\_bound(v) \textasciitilde{}=\_(AST,alpha,loc) predAst(encode\_inv(I)).}
\NormalTok{4. T\_BIND\_EXPR(erase\_bound(v),MM,\{self:C\})=v\textquotesingle{} is defined and}
\NormalTok{   v\textquotesingle{} \textasciitilde{}=\_(B,alpha,meta) v.}
\NormalTok{5. Every construct in v is in the theorem{-}supported fragment.}
\NormalTok{6. No excluded construct, such as any, one, sortedBy, count(element),}
\NormalTok{   Bag/Sequence/OrderedSet{-}specific operation, or fallback{-}only construct,}
\NormalTok{   occurs in v.}
\end{Highlighting}
\end{Shaded}

Operationally, when starting from invariant text \texttt{tInv}, the same
condition is checked in the opposite implementation direction:

\begin{Shaded}
\begin{Highlighting}[]
\NormalTok{astInv = ParseInvariant(tInv)}
\NormalTok{I=(cName,R,s) = decodeInvariant\_val(astInv)}
\NormalTok{resolveClass\_MM(cName)=C}
\NormalTok{BInvariant(C,R,b) = T\_BIND\_INV(astInv,MM)}
\NormalTok{s in OCL\_val}
\NormalTok{b in BoundOCL\_val(MM), with type Boolean under \{self:C\}}
\end{Highlighting}
\end{Shaded}

Proof. The source/AST retraction is M1. For the bound round trip,
proceed by induction on the successful binding derivation. Base bound
constructors erase to the corresponding parser AST constructor. For each
compound constructor, apply the induction hypotheses to its children and
build the parser constructor with the same surface operator and binder
structure. For an attribute, class, declaration type, or navigation
reference, erasure retains the source name/role and ordered qualifier
AST but removes the resolved element; resolver adequacy and
functionality reconstruct that same element when rebinding over the
unchanged \texttt{MM}. Iterator and \texttt{let} cases preserve optional
type names, resolved declaration types, and binder positions, and use
fresh alpha-renaming when necessary, so rebinding reconstructs a
\texttt{\textasciitilde{}=\_(B,alpha,meta)}-equivalent term. The
induction only has cases from the admitted grammar, proving clauses
5-\/-6. Thus \texttt{erase\_bound(v)} is a parser-AST witness and the
round trip has all six stated properties. This does not identify the two
metamodels or assert representability without resolver determinism.

Here \texttt{\textasciitilde{}=\_(AST,alpha,loc)} is AST congruence
modulo capture-avoiding alpha-renaming and non-semantic source
locations/wrapper nodes. \texttt{\textasciitilde{}=\_(B,alpha,meta)} is
the corresponding Bound-term congruence: it permits only
capture-avoiding binder renaming and irrelevant source-location
differences; it does not identify different resolved UML elements,
static types, navigation directions, roles, or qualifier metadata. The
round trip uses resolver functionality from Section 2.3; without
deterministic adequate resolvers, F1 is not claimed.

The F1 induction cases are certified by the following table. "Structural
congruence" means that the child induction hypotheses are reassembled
with the same admitted source operator; it does not erase or identify
resolved semantic metadata.

{\def\LTcaptype{none} % do not increment counter
\begin{longtable}[]{@{}llll@{}}
\toprule\noalign{}
Source constructor family & \texttt{decode\_val(encode\_ast(s))}
obligation & Erase/rebind obligation & Additional premise \\
\midrule\noalign{}
\endhead
\bottomrule\noalign{}
\endlastfoot
\texttt{self}, literal & constructor and payload are preserved directly
& bound base constructor erases to the same AST form & admitted literal
domain \\
variable & variable name and lexical position are preserved & nearest
lookup reconstructs the same binding, modulo alpha-renaming of enclosing
binders & B4 stack discipline \\
attribute call & receiver and surface property name are preserved &
erasure removes resolved attribute; adequate functional resolution
reconstructs that attribute & unique applicable attribute \\
navigation, including qualifiers & receiver, role name, and ordered
qualifier ASTs are preserved & erasure removes
association/roles/direction/types; \texttt{resolveNavigation\_MM}
reconstructs the same tuple & B3; ordered qualifier-type agreement \\
\texttt{C.allInstances()} & class name and operation form are preserved
& class resolver reconstructs the same class element &
injective/adequate class lookup \\
\texttt{oclIsKindOf}, \texttt{oclAsType} & receiver, operation, and
target class name are preserved & target class metadata is reconstructed
by the class resolver & supported type operation and compatible
target \\
\texttt{if}, unary/binary operation, collection operation & operator and
child order are preserved & structural congruence applies to every child
& all children satisfy their F1 induction hypotheses \\
iterator & source, iterator operation, optional type name, binder
position, and body are preserved & \texttt{declType} reconstructs the
same \texttt{tauD}; rebinding under \texttt{Gamma{[}x:tauD{]}}
reconstructs the body modulo fresh alpha-renaming & B4;
\texttt{tauE\ \textless{}=\ tauD}; admitted body type \\
\texttt{let} & initializer, optional type name, binder position, and
body are preserved & \texttt{declType} reconstructs the same
\texttt{tauD} and body; binder names may differ only by
alpha-equivalence & B4; \texttt{tauInit\ \textless{}=\ tauD}; no
variable capture \\
\end{longtable}
}

The table is exhaustive over the source metamodel of \texttt{OCL\_val}.
Because \texttt{decode\_val} is partial, parser AST forms without a row
are rejected before F1 is invoked.

\textbf{B1 Name Resolution Soundness.} If
\texttt{MM;Gamma\ \textbar{}-\ ast\ =\textgreater{}\ b\ :\ tau}, every
variable reference in \texttt{b} denotes the nearest binding in
\texttt{Gamma}, and every class, attribute, navigation, and type
operation reference stored in \texttt{b} is the unique element returned
by the adequate resolver for \texttt{MM}. Therefore each reference
denotes the same semantic entity as the admitted source AST.

Proof. Induct on the final binding rule. For \texttt{self}, the rule
reads the fixed context entry. For a variable, nearest-stack lookup
returns the same lexical declaration used by admitted source semantics.
Attribute/class/type rules call their adequate partial resolver once and
store the returned UML element, so stored and source denotations
coincide. Navigation calls \texttt{resolveNavigation\_MM}; B3 below
gives one tuple containing association, roles, direction, and qualifier
types. Iterator and \texttt{let} rules apply the induction hypothesis
under the same stack extension, while all other compound rules only
assemble already resolved children. There is no successful case for
failed or ambiguous resolution. Hence every stored reference has the
claimed denotation.

The reference-bearing cases used by the B1 induction are:

{\def\LTcaptype{none} % do not increment counter
\begin{longtable}[]{@{}llll@{}}
\toprule\noalign{}
Bound constructor/reference & Resolution source & Stored semantic data &
Why the denotation is preserved \\
\midrule\noalign{}
\endhead
\bottomrule\noalign{}
\endlastfoot
\texttt{BSelf} & fixed invariant context & context class/binding &
source and bound semantics read the same context entry \\
\texttt{BVariable(x)} & nearest stack lookup & selected lexical binding
and type & B4 proves that the selected declaration is the source lexical
declaration \\
\texttt{BAttribute(receiver,a)} & attribute resolver over the static
receiver type & UML attribute identity, declaring class, result type &
resolver adequacy identifies the attribute denoted by the admitted
property call \\
\texttt{BNavigation} & \texttt{resolveNavigation\_MM} & association,
source role, target role, direction, target class, qualifier types & B3
gives a unique tuple and adequacy identifies it with the source
navigation \\
\texttt{BAllInstances(C)} & class resolver & UML class identity/key &
adequate unique lookup preserves the class reference \\
\texttt{BTypeKindOf(receiver,C)} & class resolver plus operation
admission & target UML class & the stored class is exactly the class
named by the admitted type operation \\
\texttt{BCast(receiver,C)} & class resolver plus cast-compatibility
check & target UML class and result type & the stored target is unique
and the type premise records the admitted cast \\
iterator/\texttt{let} declaration and body references &
\texttt{resolveType\_MM} and stack extension \texttt{Gamma{[}x:tauD{]}}
& resolved declaration type and nearest local binding for \texttt{x};
unchanged outer bindings for other names & functionality of
\texttt{declType}; B4 stack equations preserve shadowing and restore the
outer environment after the body \\
\end{longtable}
}

All remaining constructors contain no new name-bearing reference: their
B1 case follows by applying the induction hypotheses to their children.
A failed, ambiguous, or fallback resolver path has no successful binding
derivation and therefore cannot be used as a B1 case.

\textbf{B2 Type Assignment Soundness.} If
\texttt{MM;Gamma\ \textbar{}-\ ast\ =\textgreater{}\ b\ :\ tau}, then
\texttt{MM;Gamma\ \textbar{}-\ b\ :\ tau}. Moreover, every subderivation
assigning \texttt{tau\textquotesingle{}} to a bound subexpression
\texttt{b\textquotesingle{}} satisfies the corresponding typing rule.
For every \texttt{rho\ \textbar{}=\_M\ Gamma} for which that bound
subevaluation is \texttt{EvalClosed\_obj}, whenever evaluation of
\texttt{b\textquotesingle{}} produces a non-\texttt{bottom} value, that
value belongs to the semantic interpretation of
\texttt{tau\textquotesingle{}}.

Proof. Use simultaneous induction on the binding derivation for typing
and semantic membership. Base values belong to their declared tagged
domains by environment satisfaction or literal admission.
Attribute/navigation cases use resolver type adequacy and the declared
result type; navigation returns a finite subset of the resolved target
class. Boolean/comparison cases return a Boolean or admitted bottom, and
arithmetic uses \texttt{ScalarClosed} for the selected result domain.
\texttt{If} uses the join premise and the selected-branch induction
hypothesis. \texttt{Let} and iterators use \texttt{declType} and its
canonical embedding to extend the environment with a value in the
resolved declaration type, so the body induction hypothesis applies.
Collection constructors form finite homogeneous sets, while type
operations use class conformance. Each final bound constructor has
exactly the corresponding typing conclusion, establishing both claims
for every subderivation.

The simultaneous induction is made exhaustive by the following case
matrix. \texttt{Semantic\ membership} concerns only non-\texttt{bottom}
results; admitted bottom propagates according to the validation-level
primitive equations.

{\def\LTcaptype{none} % do not increment counter
\begin{longtable}[]{@{}llll@{}}
\toprule\noalign{}
Successful binder case & Critical premises stored or checked by the rule
& Bound result type & Non-bottom semantic-membership argument \\
\midrule\noalign{}
\endhead
\bottomrule\noalign{}
\endlastfoot
\texttt{self} & context contains \texttt{self:C} & \texttt{C} &
environment satisfaction gives an object conforming to \texttt{C} \\
variable & nearest lookup returns \texttt{x:tau} & \texttt{tau} &
environment satisfaction for the selected stack entry \\
Boolean/Integer/Real/String literal & literal kind and value are
admitted & corresponding primitive type & literal belongs to its tagged
scalar domain \\
\texttt{null} literal & the only \texttt{Void} inhabitant is bottom;
contextual compatibility is checked by the parent rule & \texttt{Void}
before contextual embedding & both Source and Bound denotations are
their tagged bottom value \\
attribute call & receiver type admits resolved attribute \texttt{a:T};
resolver is type-adequate & \texttt{T} & \texttt{attrVal\_M} returns a
declared \texttt{T} value when non-bottom \\
binary navigation & receiver is compatible with the resolved source end;
target class is \texttt{D}; qualifier types match in order & \texttt{D}
when upper multiplicity is one; otherwise certified \texttt{Set(D)} &
multiplicity-one returns one target or bottom; multi-valued navigation
returns a finite endpoint set conforming to \texttt{D} \\
\texttt{C.allInstances()} & class resolver returns the unique class
\texttt{C} & \texttt{Set(C)} & the definition selects exactly finite
objects conforming to \texttt{C} \\
unary \texttt{not} & operand has type \texttt{Boolean} &
\texttt{Boolean} & the admitted truth operation returns Boolean when
non-bottom \\
Boolean binary operation & both operands have type \texttt{Boolean} &
\texttt{Boolean} & the admitted Boolean operation is closed over Boolean
operands \\
equality/inequality & operands have the same admitted type or
deterministic numeric common type & \texttt{Boolean} & typed equality
returns Boolean; entity equality uses identity \\
ordering & operands have compatible ordered scalar types &
\texttt{Boolean} & scalar compatibility gives a Boolean comparison where
defined \\
arithmetic & \texttt{arithResult(op,tau1,tau2)=tau}; scalar-closed
premises hold & \texttt{tau} & \texttt{ScalarClosed} places every
defined result in the selected numeric domain \\
\texttt{if} & condition is Boolean; branch types have admitted join
\texttt{tau} & \texttt{tau} & validation truth selects one branch; its
induction hypothesis gives membership in the joined domain \\
\texttt{let\ x{[}:T?{]}=init\ in\ body} & \texttt{init:tau1};
\texttt{declType\_MM(T?,tau1)=tauD}; body under
\texttt{Gamma{[}x:tauD{]}} & body type \texttt{tau} & canonical
embedding places the initializer in \texttt{tauD}; body IH then yields
\texttt{tau} \\
\texttt{exists}, \texttt{forAll} & source has native type
\texttt{Set(tauE)}; \texttt{declType\_MM(T?,tauE)=tauD}; predicate
Boolean under \texttt{Gamma{[}x:tauD{]}} & \texttt{Boolean} & each
embedded element inhabits \texttt{tauD}; finite witness/counterexample
evaluation returns Boolean \\
\texttt{select}, \texttt{reject} & same declaration premises; predicate
Boolean under \texttt{Gamma{[}x:tauD{]}} & \texttt{Set(tauE)} & the
predicate sees the declared view, while the result remains a finite
subset of original source elements \\
finite-set \texttt{collect} & same declaration premises; body has
non-collection type \texttt{T} & \texttt{Set(T)} & finite image
construction is homogeneous in \texttt{T}; no flattening or bag
semantics is admitted \\
\texttt{includes}, \texttt{excludes} & source is \texttt{Set(sigma)};
argument type \texttt{delta} satisfies \texttt{memberJoin(sigma,delta)}
& \texttt{Boolean} & finite-set membership test returns Boolean \\
\texttt{includesAll}, \texttt{excludesAll} & both operands are
compatible finite sets & \texttt{Boolean} & subset/disjointness
evaluation returns Boolean \\
\texttt{size}, \texttt{isEmpty}, \texttt{notEmpty} & operand has a
finite-set type or a recorded direct to-one collection view &
\texttt{Integer} or \texttt{Boolean} & finite cardinality of the
selected view is representable under \texttt{ScalarClosed}; emptiness
tests are Boolean \\
\texttt{oclIsKindOf(C)} & receiver is entity-or-bottom; class resolver
returns \texttt{C} & \texttt{Boolean} & conformance testing is Boolean;
bottom follows the stated false policy \\
\texttt{oclAsType(C)} & cast compatibility is admitted; class resolver
returns \texttt{C} & \texttt{C} & a successful cast returns an entity
conforming to \texttt{C}; a failed cast is bottom \\
\end{longtable}
}

No row exists for \texttt{any}, \texttt{one}, \texttt{sortedBy},
\texttt{count(element)}, \texttt{oclIsTypeOf}, ordered collections,
bags, or fallback-bound constructs. Consequently an implementation path
that produces one of those constructors cannot appeal to B2 or to the
downstream theorem chain.

\textbf{B3 Navigation Resolution Determinism.} If binding a supported
property call as navigation succeeds, it produces a unique resolved
tuple:

\begin{Shaded}
\begin{Highlighting}[]
\NormalTok{(association, fromRole, toRole, direction, qualifierTypes)}
\end{Highlighting}
\end{Shaded}

compatible with the receiver type and \texttt{MM}.

Proof. Suppose two successful derivations produce tuples \texttt{u1} and
\texttt{u2} for the same metamodel, receiver type, role name, and
ordered qualifier-type list. Both final rules require
\texttt{resolveNavigation\_MM(input)=u1} and
\texttt{resolveNavigation\_MM(input)=u2}. Since the resolver is a
partial function, functionality gives \texttt{u1=u2}, hence
componentwise equality of association, roles, direction, and qualifier
types. If the metamodel offers multiple unresolved candidates, the
resolver is undefined, so there is no admitted countercase.

\textbf{B4 Lexical Scope and Shadowing Preservation.} In every binding
derivation, iterator and \texttt{let} bodies are bound under stack
extension \texttt{Gamma{[}x:tauD{]}}, where
\texttt{declType\_MM(T?,tauActual)=tauD}. Nearest-binding lookup
resolves \texttt{x} to the new binding and leaves all \texttt{y\ !=\ x}
lookups unchanged. On leaving the body, the previous \texttt{Gamma} is
restored. Hence free variables, bound variables, and shadowing agree
with the lexical OCL AST structure.

Proof. The stack equations are:

\begin{Shaded}
\begin{Highlighting}[]
\NormalTok{lookup(Gamma[x:tauD],x) = (x,tauD)}
\NormalTok{lookup(Gamma[x:tauD],y) = lookup(Gamma,y), for y != x.}
\end{Highlighting}
\end{Shaded}

At depth zero there is no local binder. For the induction step, the
displayed equations prove correct lookup inside the new body; apply the
induction hypothesis to nested bodies. Popping the frame restores
exactly \texttt{Gamma}, so no binding escapes. For alpha-renaming to
fresh \texttt{z}, the standard environment renaming bijection maps the
\texttt{x} entry to \texttt{z} and leaves all free-variable entries
unchanged; structural induction on the body then preserves denotation.

Declared-binder conformance sublemma. If
\texttt{declType\_MM(T?,tauActual)=tauD} and \texttt{v} is a non-bottom
value of \texttt{tauActual}, then
\texttt{iota{[}tauActual-\textgreater{}tauD{]}(v)} has type
\texttt{tauD}. The embedding is functional and commutes with the Source,
Bound, VA, and Cypher representations. Proof is by cases on the
generating conformance rule: reflexivity and UML class upcast are
identity; \texttt{Integer\textless{}=Real} uses A8;
\texttt{Void\textless{}=tauD} maps to typed bottom; and finite-Set
conformance follows pointwise plus extensionality. This is the exact
fact used when extending the environment for typed \texttt{let} and
iterator bodies.

\textbf{VA1 Bound-to-VA Type Preservation.} If
\texttt{MM;Gamma\ \textbar{}-\ b\ :\ tau} and \texttt{T\_VA(b)=v}, then
\texttt{MM;Gamma\ \textbar{}-\ v\ :\ tau} in the corresponding VA typing
system.

Proof. Structural induction on \texttt{b}, using the complete VA2 table.
Base rows map to VA constructors with identical declared type.
Attribute, navigation, Boolean, comparison, arithmetic, and
type-operation rows copy the resolved metadata and use the child
induction hypotheses, so the corresponding VA typing rule has the same
premises and conclusion. \texttt{If} preserves branch join; \texttt{Let}
preserves \texttt{variableType=tauD}, initializer conformance, and
checks the body under \texttt{Gamma{[}x:tauD{]}}. Iterator rows preserve
both the finite-set source type and \texttt{iteratorVariableType=tauD},
then apply the body induction hypothesis under
\texttt{Gamma{[}x:tauD{]}}. These rows exhaust Bound constructors
admitted by B2, proving the result type unchanged.

\textbf{VA2 Local Bound-to-VA Simulation.} \texttt{T\_VA} is the total
structural function on successfully bound \texttt{OCL\_val} terms
defined by the following complete table. Let \texttt{V(e)=T\_VA(e)} and
\texttt{Vbar({[}e1,...,ek{]})={[}V(e1),...,V(ek){]}}. Define the typed
collection-boundary translation:

\begin{Shaded}
\begin{Highlighting}[]
\NormalTok{CV(S,Set(sigma)) = V(S),             when type(S)=Set(sigma), for every}
\NormalTok{                                      non{-}collection value type sigma}
\NormalTok{CV(S,Set(D)) = ViewSet(V(S),D),      when D is in Class(MM), type(S)=D, and S}
\NormalTok{                                      is an admitted directly consumed ONE}
\NormalTok{                                      navigation.}
\end{Highlighting}
\end{Shaded}

The first equation includes scalar collections such as
\texttt{Set(Integer)} as well as entity collections. The second equation
is the only scalar-to-set view: it is restricted to a resolved
entity-valued ONE navigation. In particular,
\texttt{Set\{1\}-\textgreater{}exists(x\textbar{}x=1)} uses the first
equation and does not require an entity class \texttt{D}; this keeps
\texttt{T\_VA} total on every iterator source admitted by B2.

{\def\LTcaptype{none} % do not increment counter
\begin{longtable}[]{@{}ll@{}}
\toprule\noalign{}
Bound constructor & VA result \\
\midrule\noalign{}
\endhead
\bottomrule\noalign{}
\endlastfoot
\texttt{BSelf(C)} & \texttt{Self} \\
\texttt{BVar(x,tau)} & \texttt{Var(x)} \\
\texttt{BLiteral(c,tau)} & \texttt{Literal(c)} \\
\texttt{BAttribute(e,a,tau)} & \texttt{Attribute(V(e),a)} \\
\texttt{BNavigation(e,A,fr,tr,dir,qs,ONE,D)} &
\texttt{NavigationOne(V(e),A,fr,tr,dir,Vbar(qs))} \\
\texttt{BNavigation(e,A,fr,tr,dir,qs,MANY,Set(D))} &
\texttt{NavigationMany(V(e),A,fr,tr,dir,Vbar(qs))} \\
\texttt{BAllInstances(C)} & \texttt{AllInstances(C)} \\
\texttt{BNot(e)} & \texttt{Not(V(e))} \\
\texttt{BBinary(and,e1,e2,Boolean)} & \texttt{And(V(e1),V(e2))} \\
\texttt{BBinary(or,e1,e2,Boolean)} & \texttt{Or(V(e1),V(e2))} \\
\texttt{BBinary(xor,e1,e2,Boolean)} & \texttt{Xor(V(e1),V(e2))} \\
\texttt{BBinary(implies,e1,e2,Boolean)} &
\texttt{Implies(V(e1),V(e2))} \\
\texttt{BBinary(op,e1,e2,Boolean)},
\texttt{op\ in\ \{=,\textless{}\textgreater{},\textless{},\textless{}=,\textgreater{},\textgreater{}=\}}
& \texttt{Compare(mapCompare(op),V(e1),V(e2))} \\
\texttt{BBinary(op,e1,e2,tau)}, arithmetic \texttt{op} &
\texttt{Arith(mapArith(op),V(e1),V(e2))} \\
\texttt{BIf(c,t,f,tau)} & \texttt{If(V(c),V(t),V(f))} \\
\texttt{BLet(x,tauD,e,b,tau)} & \texttt{Let(x,tauD,V(e),V(b))} \\
\texttt{BIterator(exists,S,kappa,x,tauD,P,Boolean)} &
\texttt{Exists(CV(S,kappa),x,tauD,V(P))} \\
\texttt{BIterator(forAll,S,kappa,x,tauD,P,Boolean)} &
\texttt{ForAll(CV(S,kappa),x,tauD,V(P))} \\
\texttt{BIterator(select,S,kappa,x,tauD,P,Set(tauE))} &
\texttt{Select(CV(S,kappa),x,tauD,V(P))} \\
\texttt{BIterator(reject,S,kappa,x,tauD,P,Set(tauE))} &
\texttt{Reject(CV(S,kappa),x,tauD,V(P))} \\
\texttt{BIterator(collect,S,kappa,x,tauD,E,Set(sigma))} &
\texttt{Collect(CV(S,kappa),x,tauD,V(E))} \\
\texttt{BCollectionOp(includes,S,kappa,{[}e{]},Boolean)} &
\texttt{Includes(CV(S,kappa),V(e))} \\
\texttt{BCollectionOp(excludes,S,kappa,{[}e{]},Boolean)} &
\texttt{Excludes(CV(S,kappa),V(e))} \\
\texttt{BCollectionOp(includesAll,S,kappa,{[}T{]},Boolean)} &
\texttt{IncludesAll(CV(S,kappa),V(T))} \\
\texttt{BCollectionOp(excludesAll,S,kappa,{[}T{]},Boolean)} &
\texttt{ExcludesAll(CV(S,kappa),V(T))} \\
\texttt{BSetLiteral({[}e1,...,en{]},Set(tau))} &
\texttt{SetLiteral({[}V(e1),...,V(en){]})} \\
\texttt{BIterator(isUnique,S,kappa,x,tauD,E,Boolean)} &
\texttt{IsUnique(CV(S,kappa),x,tauD,V(E))} \\
\texttt{BCollectionOp(union,S,kappa,{[}T{]},Set(tau))} &
\texttt{Union(CV(S,kappa),V(T))} \\
\texttt{BCollectionOp(intersection,S,kappa,{[}T{]},Set(tau))} &
\texttt{Intersection(CV(S,kappa),V(T))} \\
\texttt{BCollectionOp(asSet,S,kappa,{[}{]},Set(tau))} &
\texttt{AsSet(CV(S,kappa))} \\
\texttt{BCollectionOp(size,S,kappa,{[}{]},Integer)} &
\texttt{Count(CV(S,kappa))} \\
\texttt{BCollectionOp(isEmpty,S,kappa,{[}{]},Boolean)} &
\texttt{IsEmpty(CV(S,kappa))} \\
\texttt{BCollectionOp(notEmpty,S,kappa,{[}{]},Boolean)} &
\texttt{NotEmpty(CV(S,kappa))} \\
\texttt{BTypeOp(kindOf,e,C,Boolean)} & \texttt{TypeKindOf(V(e),C)} \\
\texttt{BTypeOp(cast,e,C,C)} & \texttt{Cast(V(e),C)} \\
\end{longtable}
}

\texttt{mapCompare} maps surface comparison tokens to
\texttt{\{EQ,NEQ,LT,LE,GT,GE\}} and \texttt{mapArith} maps arithmetic
tokens to \texttt{\{PLUS,MINUS,TIMES,DIV\}}. \texttt{Count} is the
canonical VA lowering of source \texttt{size()}; \texttt{Size} is a
denotational synonym accepted as VA input but is canonicalized to
\texttt{Count} before count-comparison normalization. Neither
constructor represents OCL \texttt{count(element)}.

\texttt{CV} is a meta-level definition that emits either no constructor
or the actual VA constructor \texttt{ViewSet}; there is no undefined VA
constructor named \texttt{View}. For the only admitted scalar-source
case it is empty on \texttt{bottom\_D} and a singleton on an entity.
This is the exact \texttt{LIFT1} boundary preserved by the Java
\texttt{sourceCollectionType} payload.

The Java \texttt{variableType} and \texttt{iteratorVariableType}
payloads are likewise preserved by \texttt{T\_VA}, optimization, and
planning. At text generation, \texttt{Integer\textless{}=Real} is
rendered with \texttt{toFloat}, finite-Set conformance is lifted
elementwise and re-distincted, and UML class upcast is represented by
identity on the same object node. The optimizer may inline a
\texttt{let} initializer only when its type equals the declared variable
type; otherwise the explicit typed binder is retained. Consequently no
optimization silently removes a required declaration-boundary embedding.

\subsubsection{Lemma LIFT1 (source-\/-Bound-\/-VA to-one collection-view
agreement)}\label{lemma-lift1-source--bound--va-to-one-collection-view-agreement}

Let a native to-one navigation result be represented by
\texttt{t\ in\ Option(D)}, where \texttt{none} denotes an absent target
and \texttt{some(d)} a present target. Define the three views
independently:

\begin{Shaded}
\begin{Highlighting}[]
\NormalTok{View\_src(none) = empty       View\_src(some(d)) = \{d\}}
\NormalTok{View\_B(none)   = empty       View\_B(some(d))   = \{d\}}
\NormalTok{View\_VA(none)  = empty       View\_VA(some(d))  = \{d\}.}
\end{Highlighting}
\end{Shaded}

Let \texttt{Obs(S)} be the tuple consisting of the \texttt{asSet} value,
cardinality, emptiness, and non-emptiness of \texttt{S}. Then:

\begin{Shaded}
\begin{Highlighting}[]
\NormalTok{Obs(View\_src(t)) = Obs(View\_B(t)) = Obs(View\_VA(t)).             (LIFT1)}
\end{Highlighting}
\end{Shaded}

Proof. Case \texttt{t=none}: all three definitions reduce to the empty
collection, so their observations are \texttt{(empty,0,true,false)}.
Case \texttt{t=some(d)}: all three reduce to the same singleton, whose
observations are \texttt{(\{d\},1,false,true)}. These are the only
inhabitants of \texttt{Option(D)}. The Lean theorems
\texttt{lift1\_source\_bound}, \texttt{lift1\_bound\_validation},
\texttt{lift1\_present}, \texttt{lift1\_absent}, and
\texttt{lift1\_consumer\_agreement} check these reductions without
project axioms; the last theorem\textquotesingle s axiom audit is empty.
The result is scoped exactly to the four collection observations
certified for a scalar to-one receiver. An iterator is admitted only
after explicit \texttt{asSet}, at which point its source is already
collection-typed and the ordinary iterator simulation applies. This
lemma does not identify OCL invalid with an absent navigation and does
not establish the full OMG null/invalid lattice.

\textbf{BV0 Bound-Primitive/VA Commutation.} For every admitted Bound
primitive and its corresponding VA primitive, at the canonical result
type and under the local \texttt{EvalClosed\_obj} premise:

\begin{Shaded}
\begin{Highlighting}[]
\NormalTok{erase\_B(op\_B(w1,...,wn))}
\NormalTok{  = op\_VA(erase\_B(w1),...,erase\_B(wn)).}
\end{Highlighting}
\end{Shaded}

This includes typed scalar coercion, \texttt{If} branch embeddings,
\texttt{SetLiteral}, the certified collection view, \texttt{union},
\texttt{intersection}, \texttt{asSet}, \texttt{isUnique}, attribute
access, \texttt{NavigationOne}, \texttt{NavigationMany},
\texttt{allInstances}, kind-of, and cast. In particular:

\begin{Shaded}
\begin{Highlighting}[]
\NormalTok{erase\_B(navMany\_B(w,m,qs))}
\NormalTok{  = navMany\_obj(erase\_B(w),m,map(erase\_B,qs))}

\NormalTok{erase\_B(navOne\_B(w,m,qs))}
\NormalTok{  = oneOrBottom\_obj(}
\NormalTok{      navMany\_obj(erase\_B(w),m,map(erase\_B,qs)))}

\NormalTok{erase\_B(viewSet\_B(w)) = viewSet\_obj(erase\_B(w)).}
\end{Highlighting}
\end{Shaded}

Proof. Unfold the independently defined Bound primitive and the
corresponding VA equation. Scalar cases use the same typed profile.
Finite-set cases follow by extensionality and commutation of the
canonical embeddings. Attribute, navigation, and type cases use the
copied resolved metadata; \texttt{NavigationOne} uses the
empty/singleton cases of \texttt{oneOrBottom}. This lemma relates Bound
and VA semantics and therefore is distinct from SB0, which relates
Source and Bound semantics.

For each table row, if each erased Bound child denotation equals the
corresponding VA child denotation under \texttt{rho}, then erasing the
Bound parent denotation equals the VA parent denotation:

\begin{Shaded}
\begin{Highlighting}[]
\NormalTok{erase\_B([[b\_parent]]\_B(M,tag\_B(rho)))}
\NormalTok{= [[V(b\_parent)]]\_obj(M,rho).}
\end{Highlighting}
\end{Shaded}

Proof of local simulation. For base rows, erasure removes only the Bound
tag, so both sides are the same environment value or literal. For
primitive parent rows, substitute the child equalities and apply BV0.
For iterators, child equality gives the same finite source set, and VA3
below relates every extended body environment, so extensional
quantification, filtering, image, and uniqueness results agree.
\texttt{Let} is the same extension argument for one value. These cases
exhaust the VA2 table and do not invoke Theorem 2, graph encoding, or
\texttt{encodeValue}.

\textbf{VA3 Environment Lookup Preservation.} \texttt{T\_VA} does not
rename free variables, iterator variables, \texttt{let} variables, or
\texttt{self}. Corresponding Bound and VA environments have the same
domain and satisfy \texttt{erase\_B(rho\_B(x))=rho(x)} for every name.
Extension with \texttt{tag\_B(v)} and \texttt{v} uses the same stack
operation, so nearest lookup and shadowing are preserved.

Proof. For the initial environment, define
\texttt{rho(x)=erase\_B(rho\_B(x))} pointwise; domains and all lookups
correspond. \texttt{BoundVar\ -\textgreater{}\ Var} and
\texttt{BoundSelf\ -\textgreater{}\ Self} therefore preserve lookup. For
extension with value \texttt{v}, both stacks add the same key and values
\texttt{tag\_B(v)}/\texttt{v}; erasure gives
\texttt{erase\_B(tag\_B(v))=v} at that key, while older keys are
unchanged. Induction on the number of iterator/\texttt{let} extensions
proves the invariant at every body depth, including shadowing because
both sides use nearest lookup.

\textbf{N0 Coercion and Freshness Adequacy.} Every rewrite is selected
from the typing derivation of its source redex. A generated comparison
records the unique \texttt{memberJoin} result and both canonical
embeddings. Every introduced binder satisfies:

\begin{Shaded}
\begin{Highlighting}[]
\NormalTok{fresh(e1,...,en,Gamma)}
\NormalTok{  notin allNames(e1) union ... union allNames(en) union dom(Gamma).}
\end{Highlighting}
\end{Shaded}

Fix once and for all an infinite ordered supply
\texttt{\_ocl0,\_ocl1,...}; \texttt{fresh(e1,...,en,Gamma)} denotes the
least name in that supply satisfying the displayed condition. Thus
fresh-name selection is a function. Without fixing this supply, all
determinism statements below are understood only modulo capture-avoiding
alpha-equivalence.

Consequently the target is well typed, no free occurrence becomes bound,
and the target coercions are exactly those required by the source
operation.

\textbf{N1 Rewrite Semantic Preservation.} If
\texttt{e\ -\textgreater{}\_bu\ e\textquotesingle{}},
\texttt{rho\ \textbar{}=\_I\ Gamma}, and \texttt{EvalClosed\_I(e,rho)},
then for \texttt{I=obj,graph}:

\begin{Shaded}
\begin{Highlighting}[]
\NormalTok{[[e]]\_I(rho) = [[e\textquotesingle{}]]\_I(rho)}
\NormalTok{and EvalClosed\_I(e\textquotesingle{},rho).}
\end{Highlighting}
\end{Shaded}

Expressions are compared modulo capture-avoiding alpha-equivalence.

\textbf{N2 Rewrite Type Preservation.} If
\texttt{MM;Gamma\ \textbar{}-\ e:tau} and
\texttt{e\ -\textgreater{}\_bu\ e\textquotesingle{}}, then
\texttt{MM;Gamma\ \textbar{}-\ e\textquotesingle{}:tau}. Every explicit
coercion in \texttt{e\textquotesingle{}} is selected by the unique
typing derivation and is total under \texttt{EvalClosed\_I}.

\textbf{N3 Bottom-Up Normalization Termination.} Define
\texttt{-\textgreater{}\_bu} deterministically:

\begin{enumerate}
\def\labelenumi{\arabic{enumi}.}
\tightlist
\item
  rewrite the leftmost child that still has a step;
\item
  apply a root rule only after all children satisfy \texttt{NF\_R};
\item
  choose generated binders with \texttt{fresh}; and
\item
  otherwise stop.
\end{enumerate}

Use the lexicographic measure over \texttt{Nat\ x\ Nat}:

\begin{Shaded}
\begin{Highlighting}[]
\NormalTok{mu(e) = \#ForAll(e)+\#Implies(e)+\#Xor(e)+\#Reject(e)}
\NormalTok{      + \#IsEmpty(e)+\#NotEmpty(e)+\#Includes(e)+\#Excludes(e)}
\NormalTok{      + \#CountCompareOnNavigation(e)}

\NormalTok{measure(e) = (\#Size(e),mu(e)).}
\end{Highlighting}
\end{Shaded}

A contextual child step strictly decreases the corresponding additive
count of the whole term. \texttt{Size\ -\textgreater{}\ Count} strictly
decreases the first component. Every other root rule removes exactly one
counted root redex. \texttt{Xor} duplicates its operands, but the root
rule is enabled only after both operands are in normal form, so neither
duplicate contains a counted redex and \texttt{\#Xor} decreases by one.
Coercion nodes and fresh variables are not redexes. Hence every strategy
step strictly decreases \texttt{measure}.

\textbf{N4 Normal-Form Coverage.} For every well-typed \texttt{e},
\texttt{NF\_R(T\_NORM(e))} holds modulo alpha-equivalence. This is a
syntactic result; it has no environment or runtime-value premise.

Proof. Use well-founded induction on \texttt{measure}, nested with
structural induction for child normalization. N3 proves descent; after
the children are normal, the unique root rule either descends again or
no redex remains. This is relative coverage for the stated rule set, not
a claim about arbitrary OCL rewrites.

\subsection{6.3 Graph and Cypher
Lemmas}\label{63-graph-and-cypher-lemmas}

Attribute and navigation correspondence are intentionally not duplicated
in this group: they are already the primitive representation results R4
and R6. The G-lemmas below lift those primitive observations through
environments, finite sets, scalars, and the validation-level truth
policy.

\textbf{G1 Environment Encoding Preservation.} Index environment
correspondence by its typing context. If
\texttt{rho\ \textasciitilde{}Phi\_Gamma\ eta},
\texttt{v\ in\ ValBottom\_obj(tau)}, and
\texttt{Gamma\textquotesingle{}=Gamma{[}x:tau{]}}, then:

\begin{Shaded}
\begin{Highlighting}[]
\NormalTok{rho[x{-}\textgreater{}v] \textasciitilde{}Phi\_Gamma\textquotesingle{} eta[x{-}\textgreater{}encodeValue\_tau(v)].}
\end{Highlighting}
\end{Shaded}

Proof. By definition, \texttt{rho\ \textasciitilde{}Phi\_Gamma\ eta}
means \texttt{eta(y)=encodeValue\_Gamma(y)(rho(y))} for every visible
binding \texttt{y} in the common typed domain. For the new nearest
binding \texttt{x:tau}, the required equation is the definition of the
graph-side update. Every other visible binding is unchanged and follows
from the premise. This also handles shadowing: the old \texttt{x}
binding is not read inside \texttt{Gamma\textquotesingle{}}, and
restoring \texttt{Gamma} on scope exit restores both old environment
entries.

\textbf{G2 Value-Encoding Injectivity and Finite-Set Homomorphism.} For
each fixed theorem-supported type \texttt{tau},
\texttt{encodeValue\_tau} is injective on \texttt{ValBottom\_obj(tau)}:

\begin{Shaded}
\begin{Highlighting}[]
\NormalTok{encodeValue\_tau(v1) = encodeValue\_tau(v2) iff v1 = v2.}
\end{Highlighting}
\end{Shaded}

No injectivity statement compares values at two different static types;
in particular, typed bottoms may share a backend token without becoming
equal.

The proof is by induction on value structure. Object entities use R1/R2,
supported scalars use identity encoding, \texttt{bottom} and the value
constructors are disjoint tagged sorts, and finite sets use extensional
equality plus the induction hypothesis. Consequently, for finite
\texttt{A,B\ subseteq\ ValBottom\_obj(sigma)}, \texttt{encodeSet\_sigma}
preserves and reflects membership, union, subset, intersection,
emptiness, and cardinality.

More explicitly, membership reflection uses injectivity:

\begin{Shaded}
\begin{Highlighting}[]
\NormalTok{encodeValue\_sigma(x) in encodeSet\_sigma(A)}
\NormalTok{iff exists a in A: encodeValue\_sigma(x)=encodeValue\_sigma(a)}
\NormalTok{iff exists a in A: x=a}
\NormalTok{iff x in A.}
\end{Highlighting}
\end{Shaded}

The union and intersection equations then follow pointwise. Subset
reflection uses membership reflection in both directions. Finally, the
restriction of \texttt{encodeValue\_sigma} to a finite set \texttt{A} is
a bijection from \texttt{A} to \texttt{encodeSet\_sigma(A)}, which
proves emptiness and cardinality preservation. Thus the cardinality
claim does not rely on an unstated global injectivity assumption.

\textbf{G3 Typed Equality and Scalar-Operator Preservation.} For every
fixed theorem-supported type \texttt{tau}, value encoding preserves and
reflects typed equality:

\begin{Shaded}
\begin{Highlighting}[]
\NormalTok{eq\_val\^{}tau(v1,v2)}
\NormalTok{iff}
\NormalTok{eq\_val\^{}tau(encodeValue\_tau(v1),encodeValue\_tau(v2)).}
\end{Highlighting}
\end{Shaded}

For supported compatible scalar types it additionally preserves ordered
comparison, arithmetic where defined, and Boolean truth under
\texttt{bool\_val}. Mixed numeric operations use G3a before applying
this same-type result.

Proof. Proceed by the disjoint tagged value sorts. For entities, typed
equality is object identity on the object side and node identity on the
graph side, so the equivalence is R1. For scalar values,
\texttt{encodeValue} is the canonical scalar injection fixed by
\texttt{ScalarClosed}; its equality, coercion, order, and arithmetic
clauses are exactly the scalar-profile equations in Section 3.4. For
non-set \texttt{bottom}, total typed equality gives true only against
the same typed bottom, whose tag is distinct from every non-bottom
encoding. For sets, equality first applies \texttt{finiteSet} on both
sides; G4 commutes that completion and G2 preserves/reflexes the
resulting extensional equality. Arithmetic is asserted only when the
operation is defined in the closed scalar profile; undefined operations
produce \texttt{bottom} on both sides. These cases exhaust the
theorem-supported typed value domain.

\textbf{G3a Canonical Coercion Commutation.} If a canonical embedding
\texttt{iota\_obj{[}tau-\textgreater{}upsilon{]}} is selected by
\texttt{ifJoin}, \texttt{eqJoin}, \texttt{memberJoin},
\texttt{setJoinN}, \texttt{setJoin2}, \texttt{arithJoin}, or
\texttt{orderJoin}, then for every reachable value on which it is
defined:

\begin{Shaded}
\begin{Highlighting}[]
\NormalTok{encodeValue\_upsilon(iota\_obj[tau{-}\textgreater{}upsilon](v))}
\NormalTok{  = iota\_graph[tau{-}\textgreater{}upsilon](encodeValue\_tau(v)).}
\end{Highlighting}
\end{Shaded}

It is defined on the object side iff it is defined on the related graph
side. Consequently, for every finite set \texttt{S} satisfying
\texttt{S\ subseteq\ dom(iota\_obj{[}tau-\textgreater{}upsilon{]})}
(which follows for reachable sets from \texttt{EvalClosed\_obj}):

\begin{Shaded}
\begin{Highlighting}[]
\NormalTok{encodeSet\_upsilon(upSet\_obj[tau{-}\textgreater{}upsilon](S))}
\NormalTok{  = upSet\_graph[tau{-}\textgreater{}upsilon](encodeSet\_tau(S)).}
\end{Highlighting}
\end{Shaded}

Proof. Cases are identity; \texttt{Void-\textgreater{}upsilon}
(including \texttt{Void-\textgreater{}Set(sigma)}) mapping
\texttt{bottom\_Void} to \texttt{bottom\_upsilon}; typed bottom under a
class upcast; UML class upcast; and exact Integer-to-Real conversion.
The nonnumeric cases are immediate from the typed value encoding and R1.
For Integer-to-Real, A8 gives \texttt{abs(n)\textless{}=2\^{}53}, so
both paths produce the same finite binary64 value. Concrete storage
strings and Cypher\textquotesingle s bottom token do not occur here;
those belong to Theorem 5.

\textbf{G4 Validation-Policy Compatibility.} For theorem-supported
values:

\begin{Shaded}
\begin{Highlighting}[]
\NormalTok{finiteSet\_sigma(encodeValue\_Set(sigma)(v))}
\NormalTok{  = encodeSet\_sigma(finiteSet\_sigma(v))}
\NormalTok{for v in ValBottom\_obj(C):}
\NormalTok{  asSources\_C(encodeValue\_C(v)) = encodeSet\_C(asSources\_C(v))}
\NormalTok{bool\_val(encodeValue\_Boolean(v)) = bool\_val(v).}
\end{Highlighting}
\end{Shaded}

The first equality is used only at well-typed collection positions. The
second is used only at well-typed navigation-source positions. Both
include the stated \texttt{bottom}-to-empty policy and make no claim
about full OMG OCL invalid/null propagation.

Proof. For the first equality, the only admitted collection-position
values are \texttt{Set(S)} and \texttt{bottom}. If \texttt{v=Set(S)},
both sides are \texttt{\{encodeValue(s)\ \textbar{}\ s\ in\ S\}}; if
\texttt{v=bottom}, both sides are empty. For source lifting there are
two admitted cases. An entity \texttt{o} maps to the singleton
\texttt{\{node(o)\}} on both sides, and \texttt{bottom} maps to the
empty set on both sides. Finally, Boolean true is encoded as Boolean
true, while every other supported value remains different from true
because the value sorts and bottom tag are disjoint. Therefore
\texttt{bool\_val} agrees. No case invokes OMG null or invalid
propagation.

\subsection{6.4 Cypher Semantic Interface and Realization
Judgments}\label{64-cypher-semantic-interface-and-realization-judgments}

The proof does not attempt to axiomatize arbitrary Cypher. It uses the
following abstract target fragment, denoted \texttt{Cypher\_val}:

\begin{Shaded}
\begin{Highlighting}[]
\NormalTok{c ::= alias | $parameter | reservedConstant | c.property}
\NormalTok{    | [c1,...,cn] | [x IN c WHERE c | c]}
\NormalTok{    | unary(c) | binary(c,c) | supportedFunction(c1,...,cn) | coalesce(c,c)}
\NormalTok{    | CASE WHEN c THEN c ELSE c END}
\NormalTok{    | EXISTS \{ Q \} | COUNT \{ Q \} | COLLECT \{ Q \}}

\NormalTok{Q ::= MATCH pattern | OPTIONAL MATCH pattern}
\NormalTok{    | Q WHERE c}
\NormalTok{    | Q WITH projection}
\NormalTok{    | Q UNWIND c AS alias}
\NormalTok{    | CALL \{ Q \}}
\NormalTok{    | Q UNION ALL Q}
\NormalTok{    | Q RETURN [DISTINCT] projection}
\end{Highlighting}
\end{Shaded}

Only generated instances of these forms are covered. Patterns are
restricted to labels, relationship types, directions, and property
predicates used by the encoding contract. \texttt{WITH} and
\texttt{UNWIND} are used only when the selected rendering scheme needs
row-wise realization of a finite set, such as \texttt{Select} or
finite-set \texttt{Collect}. Reserved constants are limited to Boolean
constants, the unit projection \texttt{1}, Cypher null, and the
collision-free \texttt{BOTTOM\_TOKEN}.

Three finite, type-indexed expression builders are used below:
\texttt{decodeScalar\_C{[}tau{]}},
\texttt{coerce\_C{[}tau-\textgreater{}upsilon{]}}, and
\texttt{encScalar\_C{[}tau{]}}/\texttt{qualList\_C}. They are
target-syntax builders, not semantic functions and not metavariables
that may be filled arbitrarily. For each certified primitive type their
output must be a closed term of the displayed raw \texttt{CExpr} grammar
and must satisfy BR4, BR9, or BR10, respectively, under CY1. We write
\texttt{BuilderClosure\_C} for this finite family of syntax-totality and
commutation obligations. The equations below specify their observable
behavior; a line-by-line expansion into \texttt{Case} and the selected
\texttt{FunctionId}s is implementation evidence for
\texttt{BuilderClosure\_C}, not something proved merely by naming the
builder. Consequently TXT3/TXT5 and Theorem 5 are conditional on these
CY1/BR obligations. This prevents a prose macro from being mistaken for
an already derived raw-AST theorem.

For each static type \texttt{tau}, let \texttt{ValBottom\_G(tau)} be the
theorem-supported graph-value domain and let \texttt{CyVal\_C(tau)} be
the Cypher runtime values generated at that type by parameters,
canonical graph lookups, and raw-AST evaluation in
\texttt{CYPHER5\_val}. Values outside the indexed domain are not
silently assigned an OCL meaning. The static index is essential: one
runtime null/token can represent different typed bottoms in different
typing derivations without identifying those semantic values across
types. The representation interface is:

\begin{Shaded}
\begin{Highlighting}[]
\NormalTok{repr\_C\^{}tau : ValBottom\_G(tau) {-}\textgreater{} CyVal\_C(tau)}
\NormalTok{abs\_C\^{}tau  : CyVal\_C(tau) {-}\textgreater{} ValBottom\_G(tau)}
\NormalTok{abs\_C?\^{}tau : CyVal {-}\textgreater{} ValBottom\_G(tau) + \{unsupported\}}

\NormalTok{abs\_C\^{}tau(repr\_C\^{}tau(v)) = v}
\NormalTok{abs\_C\^{}tau(CypherNull) = bottom\_tau}
\NormalTok{repr\_C\^{}tau(bottom\_tau) = BOTTOM\_TOKEN}
\NormalTok{abs\_C?\^{}tau(c) = abs\_C\^{}tau(c), if c in CyVal\_C(tau)}
\NormalTok{abs\_C?\^{}tau(c) = unsupported, otherwise}
\end{Highlighting}
\end{Shaded}

Here \texttt{CypherNull} and \texttt{BOTTOM\_TOKEN} are members of each
applicable indexed \texttt{CyVal\_C(tau)}, but they are distinct: null
denotes an undefined scalar result at an expression boundary, whereas
the non-null token represents semantic \texttt{bottom} when it must
survive as an element of a finite set. Every realization judgment below
implicitly quantifies only over executions whose observed values lie in
\texttt{CyVal\_C}; this is part of target-profile closure, not a claim
that arbitrary Neo4j values can be abstracted.

The pair and token must satisfy all of the following representation
laws:

\begin{Shaded}
\begin{Highlighting}[]
\NormalTok{BR1 Retraction:       abs\_C\^{}tau(repr\_C\^{}tau(v)) = v for every supported v:tau.}
\NormalTok{BR2 Injectivity:      repr\_C\^{}tau(v1)=repr\_C\^{}tau(v2) iff v1=v2, for one fixed tau.}
\NormalTok{BR3 Non{-}collision:    BOTTOM\_TOKEN is unequal to every represented scalar,}
\NormalTok{                      entity, class key, qualifier payload, and object id.}
\NormalTok{BR4 Lookup/codec stability: parameter lookup and a directly represented}
\NormalTok{                      property preserve repr\_C\^{}tau(v). A canonical tagged}
\NormalTok{                      AttributeValue.value is not directly represented. For}
\NormalTok{                      \textasciigrave{}c:CExpr(String)\textasciigrave{} whose evaluation is an admitted wire}
\NormalTok{                      payload, its generated decoder satisfies}
\NormalTok{                      abs\_C\^{}tau(EvalExpr\_G,pi(}
\NormalTok{                        decodeScalar\_C[tau](c),row))}
\NormalTok{                      = decScalar\_tau(abs\_C\^{}String(}
\NormalTok{                          EvalExpr\_G,pi(c,row))); instantiated at}
\NormalTok{                      the canonical slot property, this is \textasciigrave{}value\_G(n,a)\textasciigrave{}.}
\NormalTok{BR5 Set stability:    DISTINCT treats BOTTOM\_TOKEN as one ordinary non{-}null}
\NormalTok{                      represented value; it is neither dropped nor merged with}
\NormalTok{                      any non{-}bottom value.}
\NormalTok{BR6 Count stability:  COUNT(DISTINCT x) counts BOTTOM\_TOKEN exactly once when}
\NormalTok{                      it occurs and ignores CypherNull according to CY5.}
\NormalTok{BR7 Boundary rule:    scalar undefinedness may be CypherNull and abstracts to}
\NormalTok{                      bottom; a semantic bottom projected as a set element is}
\NormalTok{                      converted to BOTTOM\_TOKEN before DISTINCT or COUNT.}
\NormalTok{BR8 Set representation: for every finite A subseteq ValBottom\_G(sigma),}
\NormalTok{                      repr\_C\^{}Set(sigma)(A) is a finite Cypher list L with}
\NormalTok{                      \{abs\_C\^{}sigma(c) | c occurs in L\}=A; a whole{-}set bottom}
\NormalTok{                      may be CypherNull/BOTTOM\_TOKEN and is completed to [],}
\NormalTok{                      while an element bottom occurs as BOTTOM\_TOKEN in L.}
\NormalTok{                      Conversely, for every admitted finite runtime list L,}
\NormalTok{                      finiteSet\_sigma(abs\_C\^{}Set(sigma)(L))}
\NormalTok{                      = \{abs\_C\^{}sigma(c) | c occurs in L\}; list order and}
\NormalTok{                      multiplicity are observationally irrelevant.}
\NormalTok{BR9 Coercion commutation: whenever the selected canonical embedding}
\NormalTok{                      iota\_graph[tau{-}\textgreater{}upsilon] is defined and c:tau,}
\NormalTok{                      abs\_C\^{}upsilon(EvalExpr\_G,pi(}
\NormalTok{                        coerce\_C[tau{-}\textgreater{}upsilon](c),row))}
\NormalTok{                      = iota\_graph[tau{-}\textgreater{}upsilon](}
\NormalTok{                          abs\_C\^{}tau(EvalExpr\_G,pi(c,row))).}
\NormalTok{BR10 Qualifier codec commutation: for every \textasciigrave{}c:CExpr(tau)\textasciigrave{} and row such that}
\NormalTok{                      abs\_C\^{}tau(EvalExpr\_G,pi(c,row))=v is a supported}
\NormalTok{                      non{-}bottom qualifier,}
\NormalTok{                      EvalExpr\_G,pi(encScalar\_C[tau](c),row)}
\NormalTok{                      = repr\_C\^{}String(encScalar\_tau(v)). For an ordered typed}
\NormalTok{                      tuple, qualList\_C evaluates componentwise in the same}
\NormalTok{                      order to encodeQualifierList; if any component is}
\NormalTok{                      bottom, it evaluates to CypherNull and matches no link.}
\end{Highlighting}
\end{Shaded}

Define \texttt{BottomSeparated(MM,M,G,e,pi)} to mean that the concrete
value bound to \texttt{\$oclBottom} is outside the ranges of every
theorem-visible scalar encoding, model literal, stored attribute value,
object id, metamodel key, and encoded qualifier reachable from
\texttt{e}. In the selected concrete profile, \texttt{BOTTOM\_TOKEN} is
supplied through this reserved parameter as the tagged map
\texttt{\{\_\_oclBottom:\ true\}} and is never persisted as an ordinary
UML value. Maps are outside every admitted Boolean, numeric, String,
entity, identifier, and metamodel-key domain, so separation follows
structurally instead of from an unusual string. The adapter must reject
this tagged map in stored scalar and qualifier positions. An
implementation using another representation must re-establish
BR1-\/-BR10.

Let a Cypher row be a finite alias-to-Cypher-value map. Let
\texttt{alpha} be an injective variable-to-alias map whose domain is the
typing context. Define the type-directed observation of a runtime cell
and a graph-semantic value by:

\begin{Shaded}
\begin{Highlighting}[]
\NormalTok{cellObs\_tau(c) = abs\_C\^{}tau(c),                         if tau is non{-}set}
\NormalTok{cellObs\_Set(sigma)(c) = finiteSet\_sigma(abs\_C\^{}Set(sigma)(c))}

\NormalTok{valueObs\_tau(v) = v,                                  if tau is non{-}set}
\NormalTok{valueObs\_Set(sigma)(v) = finiteSet\_sigma(v)}

\NormalTok{RowCorr\_Gamma(alpha,eta,row)}
\NormalTok{iff}
\NormalTok{  dom(alpha)=dom(eta)=dom(Gamma)}
\NormalTok{  and for every x in dom(Gamma):}
\NormalTok{        cellObs\_Gamma(x)(row(alpha(x)))}
\NormalTok{        = valueObs\_Gamma(x)(eta(x)).}
\end{Highlighting}
\end{Shaded}

The observational clause for \texttt{Set(sigma)} is deliberate: OCL\_val
identifies a whole-collection bottom with the empty set at every
certified collection boundary. Non-set aliases still correspond by exact
typed value equality.

Let \texttt{EvalExpr\_G,pi(c,row)} denote expression evaluation and let
\texttt{EvalRows\_G,pi(Q,R)} denote the output row bag obtained by
executing \texttt{Q} from input row bag \texttt{R} under parameter map
\texttt{pi}. These functions are used only under the Cypher axioms
listed in Theorem 5. Define semantic projection as a set:

\begin{Shaded}
\begin{Highlighting}[]
\NormalTok{projectSemanticSet\_sigma(a,B)}
\NormalTok{  = \{ abs\_C\^{}sigma(row(a)) | row occurs in row bag B \}.}

\NormalTok{SemSet\_G,pi,sigma(SetPlan(Q,a),row)}
\NormalTok{  = projectSemanticSet\_sigma(a,EvalRows\_G,pi(Q,\{row\})).}
\end{Highlighting}
\end{Shaded}

\texttt{SemSet} is the only one-argument abbreviation for observing a
\texttt{SetPlan}; it always retains the graph, parameter map, element
type, and correlated input row shown in its subscript/argument.

For a finite alias set \texttt{K}, define preservation across a scalar
prelude:

\begin{Shaded}
\begin{Highlighting}[]
\NormalTok{Preserves(K,row,row\textquotesingle{})}
\NormalTok{iff}
\NormalTok{for every a in K:}
\NormalTok{  a in dom(row) intersect dom(row\textquotesingle{})}
\NormalTok{  and row\textquotesingle{}(a) = row(a).}
\end{Highlighting}
\end{Shaded}

\texttt{ExecPrelude\_G,pi(L,row)} abbreviates
\texttt{EvalRows\_G,pi(Seq(L),\{row\})}, with an empty clause list
acting as identity. The scalar-plan well-formedness invariant is
continuation-sensitive:

\begin{Shaded}
\begin{Highlighting}[]
\NormalTok{WFScalar\_G,pi(K, ScalarPlan(L,c), row)}
\NormalTok{iff}
\NormalTok{  InputFV(L) subseteq dom(row)}
\NormalTok{  and}
\NormalTok{  ExecPrelude\_G,pi(L,row) is a singleton bag \{row\textquotesingle{}\}}
\NormalTok{  and Preserves(K,row,row\textquotesingle{})}
\NormalTok{  and FV(c) subseteq dom(row\textquotesingle{}).}
\end{Highlighting}
\end{Shaded}

\texttt{InputFV(L)} contains exactly aliases read by the sequential
prelude before that prelude binds them; an alias created by a
\texttt{WITH\ ...\ AS\ a} in \texttt{L} is not an input requirement
merely because \texttt{c} later reads \texttt{a}. Here \texttt{K} is the
set of aliases already live on entry and required after the prelude. It
contains the relevant part of \texttt{range(alpha)}, but never a future
result alias created by \texttt{L}; such an alias is added to the live
set only for a subsequent child after materialization. This is a
pointwise property, not a claim about arbitrary ill-typed backend rows.
In a realization judgment it is required for every admitted input row
quantified by \texttt{AdmRow}; this is stronger than mere freshness and
prevents a generated \texttt{WITH} from dropping \texttt{self}, iterator
aliases, or correlated outer aliases.

For \texttt{MM;Gamma\ \textbar{}-\ v:tau}, define
\texttt{AdmRow\_Gamma(v,alpha,eta,row)} to mean
\texttt{eta\ \textbar{}=\_G\ Gamma}, \texttt{EvalClosed\_graph(v,eta)},
and \texttt{RowCorr\_Gamma(alpha,eta,row)}. The realization judgments
quantify only over these reachable, well-typed, scalar-closed
environments-\/-not over arbitrary alias maps or backend values. The
renderer may choose an expression, a projected subquery, or both,
according to the static result type.

\begin{Shaded}
\begin{Highlighting}[]
\NormalTok{Gamma;alpha,pi |{-} c realizes\_scalar v : tau}
\NormalTok{iff}
\NormalTok{for every eta,row with AdmRow\_Gamma(v,alpha,eta,row):}
\NormalTok{  abs\_C\^{}tau(EvalExpr\_G,pi(c,row)) = [[v]]\_graph\^{}tau(G,eta).}

\NormalTok{Gamma;alpha,pi;K |{-} ScalarPlan(L,c) realizes\_scalar v : tau}
\NormalTok{iff}
\NormalTok{for every eta,row with AdmRow\_Gamma(v,alpha,eta,row):}
\NormalTok{  WFScalar\_G,pi(K,ScalarPlan(L,c),row)}
\NormalTok{  and ExecPrelude\_G,pi(L,row)=\{row\textquotesingle{}\} for one row\textquotesingle{}}
\NormalTok{  and RowCorr\_Gamma(alpha,eta,row\textquotesingle{})}
\NormalTok{  and abs\_C\^{}tau(EvalExpr\_G,pi(c,row\textquotesingle{})) = [[v]]\_graph\^{}tau(G,eta).}

\NormalTok{Gamma;alpha,pi |{-} c realizes\_pred v}
\NormalTok{iff}
\NormalTok{for every eta,row with AdmRow\_Gamma(v,alpha,eta,row):}
\NormalTok{  bool\_val(abs\_C\^{}Boolean(EvalExpr\_G,pi(c,row)))}
\NormalTok{  = bool\_val([[v]]\_graph\^{}Boolean(G,eta)).}

\NormalTok{Gamma;alpha,pi;K |{-} ScalarPlan(L,c) realizes\_pred v}
\NormalTok{iff}
\NormalTok{for every eta,row with AdmRow\_Gamma(v,alpha,eta,row):}
\NormalTok{  WFScalar\_G,pi(K,ScalarPlan(L,c),row)}
\NormalTok{  and ExecPrelude\_G,pi(L,row)=\{row\textquotesingle{}\} for one row\textquotesingle{}}
\NormalTok{  and RowCorr\_Gamma(alpha,eta,row\textquotesingle{})}
\NormalTok{  and}
\NormalTok{  bool\_val(abs\_C\^{}Boolean(EvalExpr\_G,pi(c,row\textquotesingle{})))}
\NormalTok{  = bool\_val([[v]]\_graph\^{}Boolean(G,eta)).}

\NormalTok{Gamma;alpha,pi |{-} Q =\textgreater{} a realizes\_set v : Set(sigma)}
\NormalTok{iff}
\NormalTok{for every eta,row with AdmRow\_Gamma(v,alpha,eta,row):}
\NormalTok{  projectSemanticSet\_sigma(a, EvalRows\_G,pi(Q,\{row\}))}
\NormalTok{  = finiteSet\_sigma([[v]]\_graph\^{}Set(sigma)(G,eta)).}
\end{Highlighting}
\end{Shaded}

The expression-only judgments are the special case \texttt{L={[}{]}}.
General lowering uses the scalar-plan judgments. The predicate judgment
is separate because validation correctness requires preservation of
\texttt{bool\_val}, not equality with full OMG four-valued truth. The
set judgment explicitly forgets row multiplicity.

For a context class \texttt{C}, define the invariant-query judgment:

\begin{Shaded}
\begin{Highlighting}[]
\NormalTok{pi |{-} q realizes\_inv(C,nva)}
\NormalTok{iff}
\NormalTok{returnedIds(q,G,pi)}
\NormalTok{= \{ id(o)}
\NormalTok{    | o in Obj(M)}
\NormalTok{      and classOf(o) conformsTo C}
\NormalTok{      and bool\_val([[nva]]\_graph\^{}Boolean}
\NormalTok{                     (G,self {-}\textgreater{} node(o))) = false \}.}
\end{Highlighting}
\end{Shaded}

This judgment compares identifiers only. It never treats a result table
as a VA value and never assumes that the query returns graph nodes.

\textbf{C1 Alias and Live-Projection Preservation.} Suppose
\texttt{RowCorr\_Gamma(alpha,eta,row)}, \texttt{a} is fresh,
\texttt{w:tau}, and the materialized cell \texttt{c} satisfies
\texttt{cellObs\_tau(c)=valueObs\_tau(w)}. Then:

\begin{Shaded}
\begin{Highlighting}[]
\NormalTok{RowCorr\_Gamma[x:tau](alpha[x{-}\textgreater{}a],eta[x{-}\textgreater{}w],row[a{-}\textgreater{}c]).}
\end{Highlighting}
\end{Shaded}

For non-set \texttt{tau}, choose \texttt{c=repr\_C\^{}tau(w)} and use
BR1. For \texttt{tau=Set(sigma)}, \texttt{bindValue} may store either
the BR8 canonical list for \texttt{finiteSet\_sigma(w)} or a
whole-set-bottom representation later normalized by C5a; both satisfy
the displayed observational equation. This is the set-valued
\texttt{let} case that exact untyped row equality could not express.

For any \texttt{K\ subseteq\ dom(row)}, projecting every alias in
\texttt{K} by identity and a fresh expression \texttt{e} as \texttt{a}
yields a row that satisfies
\texttt{Preserves(K,row,row\textquotesingle{})}. Conversely, omitting an
alias in \texttt{K} violates the live-projection premise and is not a
valid translation derivation. Lexical restoration on leaving a generated
subquery preserves any shadowed outer binding.

\textbf{C2 Parameter Soundness.} If the renderer associates parameter
\texttt{p} with a supported literal/scalar value \texttt{s:tau}, then
\texttt{pi(p)=repr\_C\^{}tau(encodeValue\_tau(s))}. Consequently, for
every row,
\texttt{abs\_C\^{}tau(EvalExpr\_G,pi(Param(p),row))=encodeValue\_tau(s)}.

\textbf{C3 Primitive Graph Access Realization.} Under
\texttt{AdmRow\_Gamma(v,alpha,eta,row)}, the generated aliases, property
lookups, type-membership patterns, attribute accessors, and directional
binary-navigation patterns satisfy the appropriate
\texttt{realizes\_scalar}, \texttt{realizes\_pred}, or
\texttt{realizes\_set} judgment for \texttt{Self}, \texttt{Var},
\texttt{Literal}, \texttt{Attribute}, \texttt{AllInstances}, and
\texttt{Navigation}. The navigation case includes ordered qualifier-list
evaluation, matching, and the preserved \texttt{ONE}/\texttt{MANY}
result kind; \texttt{ViewSet} is discharged by LIFT1. For entity
primitives this lemma is restricted to represented non-bottom receivers;
C3a lifts the result to the admitted bottom case without assuming
node-pattern safety.

\textbf{C3a Bottom-Safe Entity Receiver.} Let \texttt{P} realize a value
of static entity type with possible semantic bottom, and let \texttt{F}
be a singleton entity-body builder realizing primitive operation
\texttt{f} for every represented entity while preserving live set
\texttt{K}. If \texttt{b} represents \texttt{f(bottom)}, then:

\begin{Shaded}
\begin{Highlighting}[]
\NormalTok{withEntityReceiver(P,K,b,F)}
\end{Highlighting}
\end{Shaded}

is a well-formed scalar plan realizing \texttt{f} on the complete
admitted receiver domain. Proof: if the receiver abstracts to bottom,
total \texttt{isBottom} enables only the pattern-free bottom arm.
Otherwise static typing/profile closure make it a node and enable only
\texttt{F}. CY3 prevents clauses after the false guard from seeing a
row; CY9 preserves imported bindings and combines only the enabled
singleton result. Therefore exactly one output row exists and C1 gives
\texttt{WFScalar\_G,pi(K,...,row)} for every admitted input row.
Attribute, \texttt{TypeKindOf}, and \texttt{Cast} instantiate this lemma
with R4, R3, and bottom results null, false, and null, respectively.

\textbf{C4 Scalar-Plan, Predicate, and Existential Composition.} If
child plans satisfy their realization judgments and every generated
\texttt{WITH} projects the computed live-out set, sequential prelude
composition remains singleton and preserves \texttt{RowCorr\_Gamma}.
Generated Boolean/comparison expressions preserve \texttt{bool\_val};
\texttt{WHERE\ c} retains exactly rows for which the realized predicate
is validation-true; and \texttt{EXISTS\ \{Q\}} is validation-true
exactly when the realized set/witness query has at least one row. This
establishes compositional realization for \texttt{Not}, \texttt{And},
\texttt{Or}, comparisons, \texttt{If}, \texttt{Let}, \texttt{Exists},
and the normalized absence-of-counterexample forms.

\textbf{C5 Finite-Set and Cardinality Composition.} If source and body
fragments satisfy their realization judgments, generated filtering,
image construction, membership, subset, disjointness, emptiness, and
cardinality fragments satisfy the corresponding set/scalar judgments.
Every semantic-element projection uses set projection, and every
cardinality/uniqueness observation consumes \texttt{closeSet}, whose
final \texttt{RETURN\ DISTINCT\ setOut(...)} selects exactly one row per
abstract element. Thus Cypher bag multiplicity cannot affect VA
finite-set cardinality or \texttt{isUnique}. If a set contains semantic
\texttt{bottom}, the renderer projects \texttt{BOTTOM\_TOKEN}, not
Cypher null, before \texttt{DISTINCT} or \texttt{COUNT}.

\textbf{C5a Collection-Bottom Boundary.} Whole-collection bottom and a
bottom set element use different lowerings:

\begin{Shaded}
\begin{Highlighting}[]
\NormalTok{collection\_C(c) = CASE WHEN isBottom(c) THEN [] ELSE c END}
\NormalTok{viewOne\_C(c)    = CASE WHEN isBottom(c) THEN [] ELSE [c] END}

\NormalTok{collectionRows(ScalarPlan(L,c),z)}
\NormalTok{  = SetPlan(Seq(L ++ [Unwind(collection\_C(c),z),}
\NormalTok{                      Return(true,[Projection(setOut(Alias(z)),z)])]),z).}
\end{Highlighting}
\end{Shaded}

For the internal \texttt{P=CellPlan(Set(sigma))}, represented as
\texttt{ScalarPlan(L,c)}, and the unique
\texttt{ExecPrelude\_G,pi(L,row)=\{row\textquotesingle{}\}}:

\begin{Shaded}
\begin{Highlighting}[]
\NormalTok{SemSet\_G,pi,sigma(collectionRows(P,z),row)}
\NormalTok{  = finiteSet\_sigma(abs\_C\^{}Set(sigma)(EvalExpr\_G,pi(c,row\textquotesingle{}))).}
\end{Highlighting}
\end{Shaded}

Proof. If the whole value is \texttt{bottom\_Set(tau)},
\texttt{collection\_C} produces \texttt{{[}{]}} and therefore zero rows.
If it is a finite set, \texttt{UNWIND} produces its members,
\texttt{setOut} retains an element-level bottom as
\texttt{BOTTOM\_TOKEN}, and \texttt{DISTINCT} quotients exactly by typed
equality using BR2/BR5/BR8. Every static collection receiver is passed
through this boundary before size, quantification, filter/image,
union/intersection, or \texttt{asSet}.

Production-discharge note. \texttt{renderFiniteCollectionValue}
implements the concrete boundary

\begin{Shaded}
\begin{Highlighting}[]
\NormalTok{CASE WHEN receiver IS NULL}
\NormalTok{       OR coalesce(receiver = $oclBottom,false)}
\NormalTok{     THEN [] ELSE receiver END.}
\end{Highlighting}
\end{Shaded}

\texttt{renderCollectionView} applies it to every already
collection-typed primary receiver. The same boundary is applied
independently to every collection-valued operand or embedded source
before subset/disjointness, union/intersection, extensional equality,
flattening, collection definition checks, entity-source \texttt{UNWIND},
and collection attribute projection. This is strictly stronger than
\texttt{coalesce(receiver,{[}{]})}: the latter does not convert the
non-null \texttt{BOTTOM\_TOKEN} representation. The admitted null-valued
primary/RHS regressions and selected-runtime discriminators pass.
Additional token-shaped tests exercise direct bound/plan renderer
totality, but their present producers are outside frozen OCL\_val
admission and therefore are not theorem evidence. These finite
executions do not by themselves prove universal \texttt{PlanAdequacy},
and the cases still require inclusion in the common clean evidence
capture before they can support a released production certificate.

\textbf{C5b Extensional Set Equality.} If \texttt{R} and \texttt{T}
realize \texttt{Set(tau)}, use the explicit two-counterexample raw
builder \texttt{extSetEq\_tau(R,T,imports)} defined in Section 6.6. Let
\texttt{ScalarPlan({[}{]},e\_eq)=extSetEq\_tau(R,T,imports)}. For every
admitted correlated input row:

\begin{Shaded}
\begin{Highlighting}[]
\NormalTok{bool\_val(abs\_C\^{}Boolean(EvalExpr\_G,pi(e\_eq,row))) = true}
\NormalTok{iff SemSet\_G,pi,tau(R,row)=SemSet\_G,pi,tau(T,row).}
\end{Highlighting}
\end{Shaded}

\texttt{NEQ} is its Boolean negation. Proof. Each subset direction is
the absence of a left element with no typed-equal right witness. C5,
BR2, BR8, BR9, and C5a identify those witnesses with semantic
membership, so the conjunction is exactly mutual inclusion.

Define:

\begin{Shaded}
\begin{Highlighting}[]
\NormalTok{ContextIds(C,M)}
\NormalTok{  = \{ id(o) | o in Obj(M), classOf(o) conformsTo C \}}

\NormalTok{ViolIds\_\{M,G\}(C,nva)}
\NormalTok{  = \{ id(o) | o in Obj(M), classOf(o) conformsTo C,}
\NormalTok{               bool\_val([[nva]]\_graph\^{}Boolean}
\NormalTok{                 (G,self{-}\textgreater{}node(o)))=false \}.}
\end{Highlighting}
\end{Shaded}

\textbf{C6 Invariant Wrapper and No-Ghost Realization.} Let
\texttt{Gamma0=\{self:C\}},
\texttt{alpha0=\{self-\textgreater{}a\_self\}}, \texttt{K=\{a\_self\}},
and
\texttt{Gamma0;alpha0,pi;K\ \textbar{}-\ ScalarPlan(L,c)\ realizes\_pred\ nva},
where \texttt{pi} already contains the reserved model parameter
\texttt{pM}. Choose fresh \texttt{pCtx} and define:

\begin{Shaded}
\begin{Highlighting}[]
\NormalTok{pi\_ctx = pi[pCtx {-}\textgreater{} repr\_C\^{}String(key\_MM(C))].}
\end{Highlighting}
\end{Shaded}

Choose fresh \texttt{r\_ctx,cls\_ctx}. Let \texttt{q} be exactly the
model-scoped context match for \texttt{C} (both \texttt{Object} and
\texttt{UmlClass} constrain \texttt{\$pM=modelKey\_MM}, and the class
uses \texttt{\$pCtx=key\_MM(C)}), followed by \texttt{L}, followed by
the violation filter and identifier projection:

\begin{Shaded}
\begin{Highlighting}[]
\NormalTok{q = Seq([Match(false,[}
\NormalTok{   Rel(Node(a\_self,Object,[(modelKey,Param(pM))]),}
\NormalTok{       r\_ctx,OUT,ObjectInstanceOf,}
\NormalTok{       Node(cls\_ctx,UmlClass,[(modelKey,Param(pM)),}
\NormalTok{                              (classKey,Param(pCtx))]))])] ++ L ++}
\NormalTok{[Where(Unary(NOT,truth(c))),}
\NormalTok{ Return(true,[Projection(Property(Alias(a\_self),use\_id),useId)])]).}
\end{Highlighting}
\end{Shaded}

Then:

\begin{Shaded}
\begin{Highlighting}[]
\NormalTok{returnedIds(q,G,pi\_ctx)=ViolIds\_\{M,G\}(C,nva)}
\NormalTok{and returnedIds(q,G,pi\_ctx) subseteq ContextIds(C,M)}
\NormalTok{and ContextIds(C,M) subseteq \{id(o) | o in Obj(M)\}.}
\end{Highlighting}
\end{Shaded}

Proof. R2/R3 enumerate exactly the canonical context nodes.
\texttt{WFScalar\_G,pi} makes \texttt{L} singleton for each input and
preserves \texttt{a\_self}; C4 gives the exact keep/drop decision.
Stable \texttt{use\_id} yields \texttt{id(o)}, and \texttt{DISTINCT}
changes only multiplicity. Conversely, every returned row descends from
such a context row, so no identifier outside the two displayed supersets
can appear.

\subsection{6.5 Syntax-Directed T\_TEXT Translation
Rules}\label{65-syntax-directed-t_text-translation-rules}

The renderer is specified by two mutually recursive, syntax-directed,
plan-valued judgments:

\begin{Shaded}
\begin{Highlighting}[]
\NormalTok{alpha ; pi ; K |{-} v =\textgreater{}E P : tau ; pi\textquotesingle{}       where P:ScalarPlan}
\NormalTok{alpha ; pi ; K |{-} v =\textgreater{}S R : Set(tau) ; pi\textquotesingle{}  where R:SetPlan}
\end{Highlighting}
\end{Shaded}

The first emits a complete \texttt{ScalarPlan(L,c)}, including every
prelude needed to evaluate its final cell; the second emits a complete
\texttt{SetPlan(Q,a)}. Formally:

\begin{Shaded}
\begin{Highlighting}[]
\NormalTok{alpha;pi;K |{-} v =\textgreater{}E P:tau;pi\textquotesingle{}}
\NormalTok{iff BuildCQ(v)=V, type(v)=tau is non{-}set,}
\NormalTok{    and ExpandE(V,alpha,pi,K)=(P,pi\textquotesingle{})}

\NormalTok{alpha;pi;K |{-} v =\textgreater{}S R:Set(tau);pi\textquotesingle{}}
\NormalTok{iff BuildCQ(v)=V, type(v)=Set(tau),}
\NormalTok{    and ExpandS(V,alpha,pi,K)=(R,pi\textquotesingle{}).}
\end{Highlighting}
\end{Shaded}

The parameter state is monotone
(\texttt{pi\ subseteq\ pi\textquotesingle{}}), and every new
parameter/alias is fresh. A successful derivation has exactly one
query-model constructor selected by the outer VA constructor and static
result kind. The tables below are the constructor-coverage index; the
raw-AST equations in Section 6.6 are the defining lowerings.
Consequently a displayed final cell is never detached from its
\texttt{ScalarPlan} prelude. \texttt{T\_TEXT} is defined by the
derivation root plus TXT-Inv below.

The selected concrete target is the \textbf{Neo4j Cypher 5 validation
profile} \texttt{CYPHER5\_val}. The project uses Java driver 5.21.0, but
driver version is not a server-semantics premise. Before claiming
execution evidence, the experiment must record
\texttt{CALL\ dbms.components()\ YIELD\ versions\ RETURN\ versions{[}0{]}}
and run the dialect probes listed below on that server. The proof covers
a server only when those probes validate the following generated forms.

\texttt{\$oclBottom}, \texttt{\$pa}, and \texttt{\$pC} below are
generated parameters. \texttt{classKey} is the injective key from
Section 4.1. \texttt{attributeKey(a)} is the canonical injective key
used by the selected attribute-value encoding. Native operators appear
only under the scalar compatibility premises of Section 3.3.

Entity-consuming operations use the bottom-safe scalar-plan combinator
\texttt{withEntityReceiver(P,K,b,F)}. It first materializes receiver
plan \texttt{P} exactly once as fresh alias \texttt{recv}, then emits
one correlated \texttt{CALL} with two complementary \texttt{UNION\ ALL}
arms:

\begin{Shaded}
\begin{Highlighting}[]
\NormalTok{bottom arm:}
\NormalTok{  WITH imported aliases}
\NormalTok{  WHERE isBottom(recv)}
\NormalTok{  RETURN b AS entityOut}

\NormalTok{entity arm:}
\NormalTok{  WITH imported aliases}
\NormalTok{  WHERE NOT isBottom(recv)}
\NormalTok{  \textless{}F(recv), which may now use recv in node positions\textgreater{}}
\NormalTok{  RETURN the entity{-}case result AS entityOut}
\end{Highlighting}
\end{Shaded}

The static receiver type and target-profile closure imply that every
non-bottom \texttt{recv} reaching the entity arm is represented by a
graph node. CY3 removes the bottom row before any pattern in \texttt{F}
is evaluated. The two guards are complementary because \texttt{isBottom}
is total Boolean. The bottom arm produces one row, and \texttt{F} is
required to produce exactly one row for each entity input; therefore the
correlated call is singleton and preserves the live aliases in
\texttt{K}. This construction does not rely on \texttt{CASE}
short-circuiting around a graph pattern.

{\def\LTcaptype{none} % do not increment counter
\begin{longtable}[]{@{}lllll@{}}
\toprule\noalign{}
VA constructor/combinator & Concrete \texttt{CYPHER5\_val} expansion &
Premises & Result type & Null/bottom policy \\
\midrule\noalign{}
\endhead
\bottomrule\noalign{}
\endlastfoot
\texttt{isBottom(c)} &
\texttt{c\ IS\ NULL\ OR\ coalesce(c\ =\ \$oclBottom,false)} & BR1-\/-BR4
& Boolean & total Boolean; true for \texttt{CypherNull} and represented
bottom \\
\texttt{truth(c)} & \texttt{coalesce(c\ =\ true,false)} &
\texttt{c:Boolean\ or\ bottom} & Boolean & every non-true value is
validation-false \\
\texttt{setOut(c)} &
\texttt{CASE\ WHEN\ c\ IS\ NULL\ THEN\ \$oclBottom\ ELSE\ c\ END} &
BR3-\/-BR7 & represented element & converts null-like bottom before set
projection \\
\texttt{collection\_C(c)} &
\texttt{CASE\ WHEN\ isBottom(c)\ THEN\ {[}{]}\ ELSE\ c\ END} & C5a &
collection/list receiver & whole-collection bottom becomes zero
elements; element bottom is retained later by \texttt{setOut} \\
\texttt{viewOne\_C(c)} &
\texttt{CASE\ WHEN\ isBottom(c)\ THEN\ {[}{]}\ ELSE\ {[}c{]}\ END} &
LIFT1 & native ONE navigation at an admitted collection boundary &
absent/bottom becomes empty; present entity becomes singleton \\
\texttt{coerce\_C{[}tau-\textgreater{}upsilon{]}(c)} & identity/class
upcast preserves the representation; Integer-to-Real performs the exact
admitted conversion; bottom remains bottom & BR9, M3a, G3,
\texttt{ScalarClosed} & represented \texttt{upsilon} value & coercion
occurs before \texttt{setOut}, equality, membership, or deduplication \\
\texttt{semEq(c,d)} &
\texttt{CASE\ WHEN\ isBottom(c)\ AND\ isBottom(d)\ THEN\ true\ WHEN\ isBottom(c)\ OR\ isBottom(d)\ THEN\ false\ ELSE\ c\ =\ d\ END}
& same admitted type or deterministic numeric common type & Boolean &
bottom equals bottom only under \texttt{eq\_val} \\
\texttt{setEq\_tau(c,d)} & \texttt{semEq\_tau(setOut(c),setOut(d))} &
binder fixes one comparable element type \texttt{tau} & Boolean & total
typed equality after bottom tokenization \\
\texttt{semOrd(op,c,d)} &
\texttt{CASE\ WHEN\ isBottom(c)\ OR\ isBottom(d)\ THEN\ null\ ELSE\ c\ op\ d\ END}
& compatible ordered scalar types; \texttt{op} is \texttt{\textless{}},
\texttt{\textless{}=}, \texttt{\textgreater{}}, or
\texttt{\textgreater{}=} & Boolean or null & null abstracts to bottom \\
\texttt{semArith(op,c,d)} &
\texttt{CASE\ WHEN\ isBottom(c)\ OR\ isBottom(d)\ THEN\ null\ ELSE\ c\ op\ d\ END}
& \texttt{arithResult} defined; common representable range; non-zero
divisor for \texttt{/} & numeric or null & null abstracts to bottom;
excluded backend errors are not theorem inputs \\
\texttt{attr\_C(P,a)} & bottom-safe re-identification as
\texttt{(recv:Object\ \{modelKey:\$pm,objectKey:key\})}; exact
\texttt{AttributeValue.attributeKey=\$pa\ AND\ AttributeValue.modelKey=\$pm};
decode \texttt{value} by \texttt{decScalar\_type(a)} & exactly one
canonical slot or none; strict codec; accessor equals \texttt{value\_G}
& \texttt{ScalarPlan} of declared type & missing or `v1 \\
\texttt{qualList\_C({[}c1,...,ck{]})} & componentwise
\texttt{encScalar\_tau}; also require every qualifier expression to
satisfy \texttt{NOT\ isBottom(ci)} before comparing the
direction-selected list & declared type/order/arity; typed injective
serialization & ordered tagged-text list or no match & either concrete
bottom representation matches no navigation row \\
\texttt{kind\_C(P,C)} & bottom-safe re-identification;
\texttt{EXISTS\ \{\ MATCH\ (recv:Object\ \{modelKey:\$pm,...\})-{[}:ObjectInstanceOf{]}-\textgreater{}(:UmlClass\ \{modelKey:\$pm,classKey:\$pC\})\ \}}
& materialized conformance and model isolation & Boolean
\texttt{ScalarPlan} & bottom returns false without a graph pattern \\
\texttt{cast\_C(P,C)} & same scoped lookup as \texttt{kind\_C}; return
the re-identified receiver iff conformance holds & same premises as
\texttt{kind\_C} & entity-or-bottom \texttt{ScalarPlan} & bottom and
failed cast return bottom \\
\texttt{AllInstances(C)} &
\texttt{MATCH\ (n:Object\ \{modelKey:\$pm\})-{[}:ObjectInstanceOf{]}-\textgreater{}(:UmlClass\ \{modelKey:\$pm,classKey:\$pC\})\ RETURN\ DISTINCT\ n}
& R3/R7 and model isolation & \texttt{Set(C)} rows & no null or
foreign-model rows \\
forward \texttt{Navigation} & bottom-safe owners;
\texttt{MATCH\ (s:Object\ \{modelKey:\$pm\})-{[}r{]}-\textgreater{}(t:Object\ \{modelKey:\$pm\})\ WHERE\ type(r)\ STARTS\ WITH\ \textquotesingle{}Link\textquotesingle{}\ AND\ r.modelKey=\$pm\ AND\ r.associationKey=\$pA\ AND\ r.sourceRole=\$pFr\ AND\ r.targetRole=\$pTr\ AND\ r.sourceQualifiers=qualList\_C(qbar)\ RETURN\ DISTINCT\ t}
& exact scoped key/roles/source qualifiers; R5/R6 & target set rows &
bottom qualifier or foreign row retains no result \\
reverse \texttt{Navigation} & symmetric incoming pattern with scoped
endpoints/relationship and direction-selected
\texttt{r.targetQualifiers} & exact scoped key/roles/target qualifiers;
R5/R6 & target set rows & same as forward navigation \\
\texttt{Count(S)} &
\texttt{size(COLLECT\ \{\ Qs\ RETURN\ DISTINCT\ setOut(s)\ AS\ value\ \})}
& \texttt{Qs,s} realizes an extensional set; COLLECT-subquery probe
passes & Integer & duplicate physical paths count once \\
set-valued \texttt{If} & correlated
\texttt{CALL\ \{\ WITH\ imported\ aliases\ ...\ UNION\ ALL\ ...\ \}},
with complementary \texttt{WHERE\ truth(cc)} and
\texttt{WHERE\ NOT\ truth(cc)} branches, followed by
\texttt{RETURN\ DISTINCT\ setOut(value)} & both branches project the
same alias and compatible type; correlated-call probe passes &
finite-set rows & only selected branch contributes; set projection
preserves bottom \\
\end{longtable}
}

\texttt{Link*} in the graph contract is therefore meta-notation for the
concrete predicate
\texttt{type(r)\ STARTS\ WITH\ \textquotesingle{}Link\textquotesingle{}};
it is never emitted as a relationship type token. A deployment may
instead standardize one relationship type such as \texttt{:UmlLink}, but
then both the encoding contract and these two navigation rows must be
changed together and R5/R6 re-established.

The set-valued \texttt{If} rule is defined formally by the raw-AST
equation for \texttt{CQIfSet} in Section 6.6. Its recursively
constructed branch queries are alpha-renamed, import the sorted aliases
of their free variables through an explicit \texttt{Call(imports,Q)}
node, use complementary \texttt{Where} guards, and project one common
fresh output column. No literal \texttt{IMPORT} keyword or angle-bracket
placeholder is emitted. Both \texttt{UnionAll} arms therefore have the
same one-column schema and joined static element type.

Each row has the local realization obligation:

\begin{Shaded}
\begin{Highlighting}[]
\NormalTok{abs\_C\^{}tau(EvalExpr\_G,pi(expansion,row)) = required graph denotation of type tau}
\end{Highlighting}
\end{Shaded}

or, for set-producing rows, equality of \texttt{projectSemanticSet} with
the required finite set. \texttt{isBottom}, \texttt{truth},
\texttt{setOut}, equality, ordering, arithmetic, and typed coercion
follow by case analysis using BR1-\/-BR10, the scalar contract, and CY7.
Attribute and type rows additionally use R3/R4; navigation uses R5/R6;
cardinality uses C5. Thus the combinators are definitions in
\texttt{CYPHER5\_val}, not unexpanded correctness assumptions.

The mandatory dialect probe suite parses and executes one isolated
oracle row for each CY1-\/-CY9 assumption. Its single executable source
is \texttt{Cypher5ValAssumptionMatrix}: it owns the query text,
parameters, transaction-scoped fixture, and expected observations, while
\texttt{Cypher5ValDialectRealNeo4jTest} verifies the selected server
profile and runs all rows.
\texttt{md/research/cypher5-dialect-probes.cypher} is only a manual
launcher and version smoke query, so it cannot drift into a second
normative probe copy. A failed row places that server outside Theorem 5
until the template is replaced by an extensionally equivalent supported
form.

Non-normative execution evidence recorded on 2026-08-01 shows all 9/9
rows passing on Neo4j Kernel 2026.06.0 Enterprise, Cypher component 5,
database \texttt{demo}, with Java driver dependency 5.21.0. Exact
observations are recorded in
\texttt{md/research/cypher5-dialect-probe-results.md} and the
machine-checked manifest
\texttt{md/research/evidence/cypher5-val-runtime-2026-08-01.tsv}. The
manifest guard checks one-to-one row agreement, pinned profile, and
current renderer/encoding/ probe/harness hashes. This evidence validates
the selected representatives on that runtime; it does not replace the
local realization proof or extend Theorem 5 to other versions.

Let \texttt{SRC\_alpha,pi,K(e)=(Q\_s,s,pi\textquotesingle{})} be derived
as follows:

\begin{Shaded}
\begin{Highlighting}[]
\NormalTok{for the certified navigation receiver e:D, derive}
\NormalTok{alpha;pi;K |{-} e =\textgreater{}E ScalarPlan(L,c):D;pi\textquotesingle{} and emit}
\NormalTok{  Q\_s = Seq(L ++ [withBind(K union \{s\},c,s),}
\NormalTok{                  Where(NOT isBottom(Alias(s))),}
\NormalTok{                  retD(Alias(s),s)])}
\end{Highlighting}
\end{Shaded}

Thus \texttt{SRC} realizes exactly the entity/bottom cases of
\texttt{asSources}. There is no \texttt{Set(D)} receiver case because
implicit collection-property navigation is not admitted. The notation
\texttt{BODY(Q\_s,s,x,F)} means that \texttt{F} is derived under
\texttt{alpha{[}x-\textgreater{}s{]}}; freshness and row correspondence
are provided by C1.

For an equality whose \texttt{eqJoin} result is \texttt{Set(sigma)},
define the total type-directed adapter:

\begin{Shaded}
\begin{Highlighting}[]
\NormalTok{setOperand\_sigma(E) = BuildCQSet(E),       if type(E)=Set(sigma)}
\NormalTok{setOperand\_sigma(E) = CQEmptySet(sigma),  if type(E)=Void.}
\end{Highlighting}
\end{Shaded}

The second case realizes
\texttt{finiteSet\_sigma(iota{[}Void-\textgreater{}Set(sigma){]}(bottom\_Void))=empty};
it is not a general scalar-to-set conversion. No other source type has a
set-operand case.

\subsubsection{Expression-Producing
Rules}\label{expression-producing-rules}

{\def\LTcaptype{none} % do not increment counter
\begin{longtable}[]{@{}lll@{}}
\toprule\noalign{}
Rule family & Static premise & Query-model result; defining lowering \\
\midrule\noalign{}
\endhead
\bottomrule\noalign{}
\endlastfoot
TXT-Self / TXT-Var & \texttt{alpha} contains the nearest bound alias &
\texttt{CQAlias}; \texttt{ExpandE} returns
\texttt{ScalarPlan({[}{]},Alias(...))} \\
TXT-Lit & admitted literal of type \texttt{tau} &
\texttt{CQLiteral(c,tau)}; \texttt{alloc\_tau} returns the parameter
together with its extended state \\
TXT-Attr & scalar entity receiver and primitive resolved attribute &
\texttt{CQAttribute}; the Section 6.6 rule expands the whole receiver
plan before \texttt{withEntityReceiver} \\
TXT-Not & Boolean child plan & \texttt{CQNot}; the child prelude is
retained and only its final cell is wrapped \\
TXT-And / TXT-Or / TXT-Xor / TXT-Implies & two Boolean child plans &
\texttt{CQBoolean}; \texttt{compose2} evaluates both plans left-to-right
and applies \texttt{boolAst} to their materialized cells \\
TXT-Compare & non-set operands and the selected
\texttt{eqJoin}/\texttt{orderJoin} & \texttt{CQCompare};
\texttt{compose2} retains both preludes and applies the typed coercions
and \texttt{semEq}/\texttt{semOrd} \\
TXT-SetEq & result of \texttt{eqJoin} is \texttt{Set(tau)}; each operand
passes \texttt{setOperand\_tau} & \texttt{CQSetEquality};
\texttt{extSetEq\_tau} performs the two typed counterexample tests \\
TXT-Arith & \texttt{arithResult} is defined & \texttt{CQArith};
\texttt{compose2} plus typed coercion and \texttt{semArith} \\
TXT-Coerce & selected non-set canonical embedding & \texttt{CQCoerce};
retain the child \texttt{ScalarPlan} prelude and wrap its final cell
with \texttt{coerce\_C} \\
TXT-Nav-One & exact primitive qualifier signature and
\texttt{resultKind=ONE} & \texttt{CQNavigationOne}; lower the complete
many-target plan, close it, and apply \texttt{HEAD(COLLECT\ \{...\})} \\
TXT-If-E & Boolean condition and non-set joined result &
\texttt{CQIfExpr}; materialize the condition once and use
\texttt{ifScalar} so only the selected branch prelude executes \\
TXT-Let-E & initializer is a \texttt{CQValue} of
\texttt{tau1\textless{}=tauD}; body has non-set type under
\texttt{x:tauD} & \texttt{CQLetExpr}; \texttt{ExpandV} evaluates and
embeds the initializer once, then expands the body under the extended
alias/live set \\
TXT-Exists / TXT-ForAll & native finite-set source \texttt{Set(tauE)};
\texttt{tauE\textless{}=tauD}; Boolean body under \texttt{x:tauD} &
\texttt{CQExists} / \texttt{CQForAll}; \texttt{existsRows} embeds each
source element and executes the complete body prelude \\
TXT-Includes / TXT-Excludes & scalar member type \texttt{delta};
\texttt{memberJoin(sigma,delta)=upsilon} & \texttt{CQMembership};
\texttt{memberIn} coerces both represented values before typed
equality \\
TXT-IncludesAll / TXT-ExcludesAll & two set sources and defined
\texttt{memberJoin} & \texttt{CQSetRelation}; absence of a typed
counterexample realizes subset/disjointness \\
TXT-Count & finite-set source & \texttt{CQCount};
\texttt{CountExpr(closeSet(...))} counts the distinct semantic
projection \\
TXT-IsEmpty / TXT-NotEmpty & finite-set source & \texttt{CQEmpty}; apply
\texttt{ExistsExpr} to the complete closed set plan and negate exactly
for \texttt{EMPTY} \\
TXT-KindOf / TXT-Cast & scalar entity-or-bottom receiver &
\texttt{CQKindOf} / \texttt{CQCast}; class parameters are allocated
before \texttt{withEntityReceiver} executes the complete plan \\
\end{longtable}
}

\texttt{Size(S)} uses TXT-Count after the canonical
\texttt{Size\ -\textgreater{}\ Count} phase. \texttt{Select},
\texttt{Reject}, and \texttt{Collect} are set-producing and are listed
below. An expression rule requiring a child expression is applicable
only when the child has a non-set static type.

\subsubsection{Set-Producing Rules}\label{set-producing-rules}

{\def\LTcaptype{none} % do not increment counter
\begin{longtable}[]{@{}lll@{}}
\toprule\noalign{}
Rule family & Static premise & Query-model result; defining lowering \\
\midrule\noalign{}
\endhead
\bottomrule\noalign{}
\endlastfoot
TXT-Var-S & bound variable has \texttt{Set(tau)} & \texttt{CQSetAlias};
\texttt{collectionRows} normalizes whole-set bottom before
\texttt{UNWIND} \\
TXT-SetLit & nonempty elements and defined \texttt{setJoinN} &
\texttt{CQSetLiteral}; \texttt{materializeList} retains every child
prelude before typed coercion and projection \\
TXT-AllInst & resolved class & \texttt{CQAllInstances}; allocate the
class parameter, then emit the model-scoped type pattern \\
TXT-Nav-Many & exact primitive qualifier signature and
\texttt{resultKind=MANY} & \texttt{CQNavigation}; \texttt{ExpandSrc} and
\texttt{materializeList} retain the receiver and qualifier preludes
before the scoped link pattern \\
TXT-ViewSet & directly consumed admitted ONE navigation &
\texttt{CQViewSet}; execute the whole scalar plan, then unwind
\texttt{viewOne\_C} \\
TXT-Select / TXT-Reject & native \texttt{Set(tauE)},
\texttt{tauE\textless{}=tauD}, and Boolean body under \texttt{x:tauD} &
\texttt{CQSelect} / \texttt{CQReject}; \texttt{filterRows} embeds the
iterator view but returns original source elements \\
TXT-Collect & native \texttt{Set(tauE)}, \texttt{tauE\textless{}=tauD},
and noncollection body & \texttt{CQCollect}; \texttt{mapRows} executes
the body under the embedded iterator view and projects
\texttt{setOut} \\
TXT-IsUnique & finite \texttt{Set(tauE)}, \texttt{tauE\textless{}=tauD},
and noncollection projection & \texttt{CQIsUnique}; \texttt{uniquePlan}
compares the distinct source and image cardinalities \\
TXT-Union / TXT-Intersection & \texttt{setJoin2(tau,sigma)=upsilon} &
\texttt{CQUnion} / \texttt{CQIntersection}; both arm types and the
joined type are retained for BR9 coercions \\
TXT-AsSet & finite-set source & \texttt{CQAsSet}; distinct semantic
reprojection \\
TXT-If-S & joined result \texttt{Set(tau)}; a \texttt{Void} branch is
adapted to \texttt{CQEmptySet(tau)} & \texttt{CQIfSet}; materialize the
condition once and use guarded correlated branches through
\texttt{ifSet} \\
TXT-Let-S & initializer is a \texttt{CQValue}; body has set type &
\texttt{CQLetSet}; \texttt{ExpandV} binds once and \texttt{prefixSet}
executes the complete body under the extended live set \\
\end{longtable}
}

The \texttt{TXT-If-S} branch aliases are renamed to one fresh output
alias before the union. Both branches project exactly the same
one-column schema and compatible static type. Their first subquery
clause imports precisely
\texttt{free(cond)\ union\ free(t)\ union\ free(f)} through
\texttt{WITH}; no undeclared outer alias is visible. \texttt{UNION\ ALL}
preserves all branch rows and the outer \texttt{DISTINCT} restores
finite-set semantics. Since \texttt{truth(cc)} is total Boolean and the
guards are syntactic complements, exactly one branch is enabled for each
corresponding input row. This establishes branch exclusivity, schema
compatibility, and alias scope independently of row multiplicity.

The two \texttt{TXT-Let} rules do not expand by syntactic substitution.
The helper \texttt{bindValue} evaluates the initializer exactly once,
materializes its typed representation under a fresh alias \texttt{a},
and establishes:

\begin{Shaded}
\begin{Highlighting}[]
\NormalTok{RowCorr\_Gamma(alpha,eta,row)}
\NormalTok{implies}
\NormalTok{RowCorr\_Gamma[x:tauD](alpha[x{-}\textgreater{}a],}
\NormalTok{  eta[x{-}\textgreater{}iota\_graph[tau1{-}\textgreater{}tauD]([[init]]\_graph\^{}tau1(eta))],row\textquotesingle{}).}
\end{Highlighting}
\end{Shaded}

Here \texttt{tau1\textless{}=tauD} is the declaration payload carried by
the plan. Integer-to-Real uses \texttt{toFloat}; UML class upcast is
identity; a Set embedding maps the element embedding and reapplies
distinctness. For a scalar/entity initializer \texttt{bindValue}
materializes a \texttt{ScalarPlan}; for a set initializer it
materializes a canonical list through C5a and exposes it as a set source
when \texttt{x} is used. C1 and the initializer realization judgment
prove the displayed relation. Recursive calls are on the two proper AST
children \texttt{init} and \texttt{body}, so ordinary structural
recursion proves termination and no duplicated evaluation or
substitution-size argument is needed.

\subsubsection{Invariant Wrapper Rule}\label{invariant-wrapper-rule}

\begin{Shaded}
\begin{Highlighting}[]
\NormalTok{(TXT{-}Inv)}
\NormalTok{alpha0 = \{ self {-}\textgreater{} s \}}
\NormalTok{alpha0 ; pi; \{s\} |{-} nva =\textgreater{}E ScalarPlan(L,c) : Boolean ; pi\textquotesingle{}}
\NormalTok{pi already contains pM{-}\textgreater{}repr\_C\^{}String(modelKey\_MM)}
\NormalTok{fresh parameter pCtx with}
\NormalTok{pi\textquotesingle{}\textquotesingle{}=pi\textquotesingle{}[pCtx{-}\textgreater{}repr\_C\^{}String(key\_MM(C))]}
\NormalTok{fresh aliases r\_ctx,cls\_ctx}
\NormalTok{Qinv = Seq([}
\NormalTok{  Match(false,[Rel(Node(s,Object,[(modelKey,Param(pM))]),}
\NormalTok{                   r\_ctx,OUT,ObjectInstanceOf,}
\NormalTok{                   Node(cls\_ctx,UmlClass,[(modelKey,Param(pM)),}
\NormalTok{                                          (classKey,Param(pCtx))]))])]}
\NormalTok{  ++ L ++}
\NormalTok{  [Where(Unary(NOT,truth(c))),}
\NormalTok{   Return(true,[Projection(Property(Alias(s),use\_id),useId)])])}
\NormalTok{{-}{-}{-}{-}{-}{-}{-}{-}{-}{-}{-}{-}{-}{-}{-}{-}{-}{-}{-}{-}{-}{-}{-}{-}{-}{-}{-}{-}{-}{-}{-}{-}{-}{-}{-}{-}{-}{-}{-}{-}{-}{-}{-}{-}{-}{-}{-}{-}{-}{-}{-}{-}{-}{-}{-}{-}{-}{-}{-}{-}{-}{-}{-}{-}{-}{-}{-}{-}{-}}
\NormalTok{T\_TEXT\^{}spec(C,nva) = (renderRaw(Qinv),pi\textquotesingle{}\textquotesingle{}).}
\end{Highlighting}
\end{Shaded}

\subsubsection{Translation Determinism and
Coverage}\label{translation-determinism-and-coverage}

\textbf{TXT1 Determinism.} Fix the deterministic fresh-name supply,
resolver metadata, and canonical combinator expansions. If two
derivations have the same judgment input, their outputs are
alpha-equivalent and their parameter maps agree up to fresh parameter
renaming. Proof: induction on the unique outer constructor and static
result kind.

\textbf{TXT2 Type/Kind Preservation.} An \texttt{=\textgreater{}E}
derivation exists only for non-set result types and produces one
represented semantic value; an \texttt{=\textgreater{}S} derivation
exists only for \texttt{Set(tau)} and projects one semantic element
alias. Proof: induction on the translation derivation using the VA
typing rules.

\textbf{TXT3 Fragment Closure.} Every emitted template expands only to
\texttt{Cypher\_val}. \texttt{TXT-If-S} additionally uses guarded
\texttt{CALL} and \texttt{UNION\ ALL}, whose behavior is stated in CY9
below. No arbitrary user Cypher is introduced.

\textbf{TXT4 Realization.} Every successful \texttt{=\textgreater{}E} or
\texttt{=\textgreater{}S} derivation satisfies the corresponding
realization judgment in Section 6.4. Proof: induction on the translation
derivation. Base rules use C2/C3; composition rules use C1/C4/C5; set
rules use typed coercion, semantic projection, and \texttt{setOut}; type
rules use R3. The complete non-immediate induction cases are D1-\/-D13
after the total \texttt{Expand} equations in Section 6.6; TXT4 is a
derived meta-theorem, not an assumed compiler contract.

\subsection{6.6 Raw Cypher AST and Total
Expansion}\label{66-raw-cypher-ast-and-total-expansion}

Section 6.5 is a readable presentation of the translation derivation.
Its symbols \texttt{Qs}, \texttt{Qt}, and \texttt{Qf} range over results
of recursive derivations; they are not pieces of text and are never
constructors of the target language. The certified target is the
following parameterized raw Cypher AST.

\begin{Shaded}
\begin{Highlighting}[]
\NormalTok{CExpr ::=}
\NormalTok{    Alias(AliasId) | Param(ParamId) | Null | Bool(Boolean) | Int(Integer)}
\NormalTok{  | Property(CExpr,PropertyKey) | List([CExpr])}
\NormalTok{  | Unary(UnaryOp,CExpr) | Binary(BinaryOp,CExpr,CExpr)}
\NormalTok{  | Case([(CExpr,CExpr)],CExpr) | Function(FunctionId,[CExpr])}
\NormalTok{  | ListComp(AliasId,CExpr,Optional(CExpr),CExpr)}
\NormalTok{  | ExistsExpr(CQuery) | CountExpr(CQuery) | CollectExpr(CQuery)}

\NormalTok{CNodePattern ::= Node(AliasId,Optional(Label),[(PropertyKey,CExpr)])}
\NormalTok{CRelPattern  ::= Rel(CNodePattern,AliasId,Direction,Optional(RelType),CNodePattern)}
\NormalTok{CPattern     ::= CNodePattern | CRelPattern}
\NormalTok{CProjection ::= Projection(CExpr,AliasId)}

\NormalTok{CClause ::=}
\NormalTok{    Match(Boolean optional,[CPattern]) | Where(CExpr) | Unwind(CExpr,AliasId)}
\NormalTok{  | With(Boolean distinct,[CProjection])}
\NormalTok{  | Return(Boolean distinct,[CProjection])}
\NormalTok{  | Call([AliasId],CQuery)}

\NormalTok{CQuery ::= Seq([CClause]) | UnionAll(CQuery,CQuery).}
\end{Highlighting}
\end{Shaded}

\texttt{UnaryOp}, \texttt{BinaryOp}, \texttt{FunctionId},
\texttt{PropertyKey}, \texttt{Label}, and \texttt{RelType} are finite
certified enumerations. Identifiers are AST atoms, not interpolated
strings, and model values occur only below \texttt{Param}.
\texttt{renderRaw} is a total, parenthesizing pretty-printer over this
grammar; it performs no semantic choice and introduces no identifier or
literal.

Because a scalar expression can require preceding \texttt{WITH},
\texttt{MATCH}, or \texttt{CALL} clauses, expansion uses two closed
host-language plan datatypes:

\begin{Shaded}
\begin{Highlighting}[]
\NormalTok{ScalarPlan ::= ScalarPlan([CClause] prelude,CExpr result)}
\NormalTok{SetPlan    ::= SetPlan(CQuery rows,AliasId exported)}
\NormalTok{ValuePlan  ::= ValuePlan([CClause] prelude,AliasId bound,Type type)}

\NormalTok{prelude(ScalarPlan(L,c))=L       result(ScalarPlan(L,c))=c}
\NormalTok{rows(SetPlan(Q,a))=Q             exported(SetPlan(Q,a))=a}
\NormalTok{prelude(ValuePlan(L,a,tau))=L    bound(ValuePlan(L,a,tau))=a}
\NormalTok{type(ValuePlan(L,a,tau))=tau}
\end{Highlighting}
\end{Shaded}

\texttt{ScalarPlan} is the implementation-facing historical name for a
one-cell plan, not an assertion that its cell has a primitive scalar
type. Formally it is indexed by the static type carried by the
construction derivation; write \texttt{CellPlan(tau)} for that fibre.
The public \texttt{=\textgreater{}E} judgment uses only
\texttt{CellPlan(tau)} with non-set \texttt{tau}. The internal
C5a/CQSetAlias bridge is the single additional use
\texttt{CellPlan(Set(sigma))}: its cell is a materialized list
representation consumed immediately by \texttt{collectionRows}, never
returned by an \texttt{=\textgreater{}E} derivation. This removes any
coercion between a Set value and a scalar value while retaining the
production datatype name.

The well-formedness invariant is that a \texttt{ScalarPlan} prelude
leaves all free aliases of \texttt{result} in scope, and a
\texttt{SetPlan} query returns exactly one column named
\texttt{exported}. Row multiplicity is not part of this invariant. These
plans are construction datatypes, not additional target syntax.
\texttt{closeScalar(P,z)} appends
\texttt{Return(false,{[}Projection(P.result,z){]})} to
\texttt{P.prelude}. The set closure is the mandatory extensional
boundary:

\begin{Shaded}
\begin{Highlighting}[]
\NormalTok{closeSet(SetPlan(Q,a))}
\NormalTok{  = Seq([Call(sortAlias(FV(Q)),Q),}
\NormalTok{         Return(true,[Projection(setOut(Alias(a)),a)])]).}
\end{Highlighting}
\end{Shaded}

Thus \texttt{closeSet} returns one represented row per abstract set
element even when an internal \texttt{Q} has duplicate physical rows; it
also preserves an element bottom as \texttt{BOTTOM\_TOKEN}. Both
closures return a raw \texttt{CQuery}.

The certified query model is:

\begin{Shaded}
\begin{Highlighting}[]
\NormalTok{CQExpr ::= CQAlias(x) | CQLiteral(v,tau) | CQAttribute(P,a)}
\NormalTok{         | CQNot(P) | CQBoolean(op,P,P)}
\NormalTok{         | CQCompare(op,P,tau,P,sigma,upsilon)}
\NormalTok{         | CQSetEquality(EQ\_OR\_NEQ,R,R,tau)}
\NormalTok{         | CQArith(op,P,tau,P,sigma,upsilon) | CQCoerce(tau,upsilon,P)}
\NormalTok{         | CQIfExpr(P,P,P) | CQLetExpr(x,V,P)}
\NormalTok{         | CQNavigationOne(P,A,fromRole,toRole,direction,[CQExpr])}
\NormalTok{         | CQExists(R,x,P) | CQForAll(R,x,P)}
\NormalTok{         | CQIsUnique(R,x,P)}
\NormalTok{         | CQMembership(mode,R,sigma,P,delta,upsilon)}
\NormalTok{         | CQSetRelation(mode,R,tau,R,sigma,upsilon)}
\NormalTok{         | CQCount(R) | CQEmpty(mode,R) | CQKindOf(P,C) | CQCast(P,C)}

\NormalTok{CQSource ::= CQSourceExpr(CQExpr) | CQSourceSet(CQSet)}

\NormalTok{CQSet  ::= CQEmptySet(tau) | CQSetAlias(x,tau)}
\NormalTok{         | CQSetLiteral([CQExpr],[tau],upsilon)}
\NormalTok{         | CQAllInstances(C)}
\NormalTok{         | CQNavigation(CQSource,A,fromRole,toRole,direction,[CQExpr])}
\NormalTok{         | CQViewSet(CQExpr,D)}
\NormalTok{         | CQSelect(R,x,P) | CQReject(R,x,P) | CQCollect(R,x,P)}
\NormalTok{         | CQUnion(R,tau,R,sigma,upsilon)}
\NormalTok{         | CQIntersection(R,tau,R,sigma,upsilon) | CQAsSet(R)}
\NormalTok{         | CQIfSet(P,R,R) | CQLetSet(x,V,R)}

\NormalTok{CQValue ::= CQValueExpr(CQExpr,tau) | CQValueSet(CQSet,Set(sigma))}

\NormalTok{CQInv  ::= CQInvariant(Class,CQExpr).}
\end{Highlighting}
\end{Shaded}

\texttt{BuildCQ} is Section 6.5 read as a function. TXT1 makes its
result unique up to fresh-name alpha-equivalence. Its comparison,
arithmetic, membership, and set combination constructors copy all
operand/common-result type indices from the unique VA typing derivation;
lowering never guesses a coercion from a runtime value. The two
specialised names used later are restrictions/constructors, not
additional unspecified translations:

\begin{Shaded}
\begin{Highlighting}[]
\NormalTok{BuildCQSet(v) = BuildCQ(v),}
\NormalTok{  when MM;Gamma |{-}VA v:Set(sigma) and BuildCQ(v):CQSet;}

\NormalTok{BuildCQInvariant(C,v) = CQInvariant(C,BuildCQ(v)),}
\NormalTok{  when C in Class(MM), MM;\{self:C\} |{-}VA v:Boolean,}
\NormalTok{       and BuildCQ(v):CQExpr.}
\end{Highlighting}
\end{Shaded}

For the fixed canonical encoding profile of Sections 4 and 6.5,
expansion has the signatures:

\begin{Shaded}
\begin{Highlighting}[]
\NormalTok{ExpandE : CQExpr x AliasEnv x ParamState x LiveSet}
\NormalTok{          {-}\textgreater{} ScalarPlan x ParamState}
\NormalTok{ExpandS : CQSet  x AliasEnv x ParamState x LiveSet}
\NormalTok{          {-}\textgreater{} SetPlan x ParamState}
\NormalTok{ExpandV : CQValue x AliasEnv x ParamState x LiveSet}
\NormalTok{          {-}\textgreater{} ValuePlan x ParamState}
\NormalTok{ExpandSrc : CQSource x AliasEnv x ParamState x LiveSet}
\NormalTok{            {-}\textgreater{} SetPlan x ParamState}
\NormalTok{ExpandI : CQInv  x ParamState            {-}\textgreater{} CQuery    x ParamState.}
\end{Highlighting}
\end{Shaded}

The fourth \texttt{ExpandE} argument is the continuation live-out set
\texttt{K}. Every call includes \texttt{range(alpha)} and any fresh
result/free aliases required by unexpanded siblings and the enclosing
continuation. Every returned plan is syntactically scoped; TXT4 proves
\texttt{WFScalar\_G,pi\textquotesingle{}(K,P,row)} for each admitted
correlated input row of its realization judgment. The fourth
\texttt{ExpandS}/\texttt{ExpandSrc} argument is the same correlated live
set; every generated \texttt{WITH}, nested \texttt{CALL}, qualifier
prelude, or branch must retain it until the set\textquotesingle s
exported alias has been consumed.

\texttt{ExpandV} is total by cases. For \texttt{CQValueExpr(P,tau)}, it
expands \texttt{P} to a scalar plan and materializes its result once
under a fresh alias. For \texttt{CQValueSet(R,Set(sigma))}, it expands
\texttt{R=SetPlan(Q,a)} and materializes
\texttt{CollectExpr(closeSet(R))} once under a fresh alias. CY5 and BR8
make that list cell observationally equal to
\texttt{finiteSet\_sigma({[}{[}R{]}{]})}, including the
whole-set-bottom-to-empty completion. Thus scalar- and set-valued
\texttt{let} initializers both have a proper \texttt{CQValue} child and
neither is encoded by substitution.

\texttt{ExpandSrc(CQSourceSet(R),alpha,pi,K)=ExpandS(R,alpha,pi,K)}. For
\texttt{CQSourceExpr(P)}, let
\texttt{(ScalarPlan(L,e),pi\textquotesingle{})=ExpandE(P,alpha,pi,K)}
and choose fresh \texttt{s}; return the pair
\texttt{(SetPlan(Seq(L\ ++\ {[}withBind(K\ union\ \{s\},e,s),\ Where(Unary(NOT,isBottom(Alias(s)))),retD(Alias(s),s){]}),s),pi\textquotesingle{})}.
This is the raw-AST definition of the singleton/bottom cases of
\texttt{asSources}. \texttt{CQSourceSet} remains an internal source
adapter for non-navigation set combinators, but a certified
\texttt{CQNavigation} is constructed only with \texttt{CQSourceExpr};
otherwise BuildCQ would reintroduce the excluded implicit
\texttt{collection.property} rule.

The following total raw-AST builders are definitions, not semantic
premises. Here \texttt{++} is clause-list concatenation,
\texttt{retD(e,a)} abbreviates
\texttt{Return(true,{[}Projection(e,a){]})}, and all aliases returned by
\texttt{fresh} are new.

\begin{Shaded}
\begin{Highlighting}[]
\NormalTok{alloc\_tau(pi,v) =}
\NormalTok{  let p be the least parameter name outside dom(pi) and the reserved names;}
\NormalTok{  (p,pi[p{-}\textgreater{}repr\_C\^{}tau(v)])}

\NormalTok{allocClass(pi,C) = alloc\_String(pi,key\_MM(C))}
\NormalTok{allocAttribute(pi,a) = alloc\_String(pi,attributeKey\_MM(a))}

\NormalTok{allocNavigation(pi,A,fr,tr) =}
\NormalTok{  let (pA,pi1)=alloc\_String(pi,associationKey\_MM(A));}
\NormalTok{      (pFr,pi2)=alloc\_String(pi1,fr);}
\NormalTok{      (pTr,pi3)=alloc\_String(pi2,tr);}
\NormalTok{  (pA,pFr,pTr,pi3)}

\NormalTok{keep(K) = [Projection(Alias(k),k) | k in sortAlias(K)]}

\NormalTok{withBind(K,e,a) =}
\NormalTok{  With(false,keep(K {-} \{a\}) ++ [Projection(e,a)])}

\NormalTok{bindScalar(ScalarPlan(L,e),a,K) =}
\NormalTok{  ScalarPlan(L ++ [withBind(K union \{a\},e,a)],Alias(a))}

\NormalTok{appendClauses(ScalarPlan(L,e),K) = ScalarPlan(L ++ K,e)}

\NormalTok{callSet(SetPlan(Q,a),imports,K) =}
\NormalTok{  [Call(imports,Q)] ++ K}

\NormalTok{existsRows(SetPlan(Q,a),imports,L,p) =}
\NormalTok{  ExistsExpr(Seq([Call(imports,Q)] ++ L ++ [Where(p),}
\NormalTok{                  Return(false,[Projection(Int(1),fresh(unit))])]))}

\NormalTok{typeStartsWithLink(r) =}
\NormalTok{  Binary(STARTS\_WITH,Function(TYPE,[Alias(r)]),Param(pLinkPrefix)).}

\NormalTok{isBottom(c) =}
\NormalTok{  Binary(OR,Unary(IS\_NULL,c),Binary(EQ,c,Param(pBottom)))}

\NormalTok{truth(c) = Function(COALESCE,[Binary(EQ,c,Bool(true)),Bool(false)])}
\NormalTok{setOut(c) = Case([(Unary(IS\_NULL,c),Param(pBottom))],c)}
\NormalTok{collection\_C(c) = Case([(isBottom(c),List([]))],c)}
\NormalTok{viewOne\_C(c) = Case([(isBottom(c),List([]))],List([c]))}

\NormalTok{boolAst(AND,x,y) = Binary(AND,truth(x),truth(y))}
\NormalTok{boolAst(OR,x,y)  = Binary(OR,truth(x),truth(y))}
\NormalTok{boolAst(XOR,x,y) =}
\NormalTok{  Binary(OR,}
\NormalTok{    Binary(AND,truth(x),Unary(NOT,truth(y))),}
\NormalTok{    Binary(AND,Unary(NOT,truth(x)),truth(y)))}
\NormalTok{boolAst(IMPLIES,x,y) = Binary(OR,Unary(NOT,truth(x)),truth(y))}

\NormalTok{semEq\_upsilon(x,y) =}
\NormalTok{  Case([(Binary(AND,isBottom(x),isBottom(y)),Bool(true)),}
\NormalTok{        (Binary(OR,isBottom(x),isBottom(y)),Bool(false))],}
\NormalTok{       Binary(EQ,x,y))}

\NormalTok{semOrd\_upsilon(op,x,y) =}
\NormalTok{  Case([(Binary(OR,isBottom(x),isBottom(y)),Null)],Binary(op,x,y))}

\NormalTok{semArith\_upsilon(op,x,y) =}
\NormalTok{  Case([(Binary(OR,isBottom(x),isBottom(y)),Null)],Binary(op,x,y))}
\end{Highlighting}
\end{Shaded}

The remaining builders are defined compositionally:

\begin{Shaded}
\begin{Highlighting}[]
\NormalTok{eqProp(r,k,p) = Binary(EQ,Property(Alias(r),k),Param(p))}
\NormalTok{andAll([]) = Bool(true)}
\NormalTok{andAll([e1,...,ek]) = Binary(AND,e1,andAll([e2,...,ek]))}

\NormalTok{qualifierProperty(OUT) = sourceQualifiers}
\NormalTok{qualifierProperty(IN)  = targetQualifiers}

\NormalTok{sourceRoleParam(OUT,pFr,pTr)=pFr}
\NormalTok{sourceRoleParam(IN,pFr,pTr)=pTr}
\NormalTok{targetRoleParam(OUT,pFr,pTr)=pTr}
\NormalTok{targetRoleParam(IN,pFr,pTr)=pFr}

\NormalTok{linkPredicate(r,pM,pA,pFr,pTr,d,qexpr) =}
\NormalTok{  andAll([typeStartsWithLink(r),}
\NormalTok{          eqProp(r,modelKey,pM),}
\NormalTok{          eqProp(r,associationKey,pA),}
\NormalTok{          eqProp(r,sourceRole,sourceRoleParam(d,pFr,pTr)),}
\NormalTok{          eqProp(r,targetRole,targetRoleParam(d,pFr,pTr)),}
\NormalTok{          Binary(EQ,Property(Alias(r),qualifierProperty(d)),qexpr)])}

\NormalTok{escapeString\_C(c) =}
\NormalTok{  Function(REPLACE,[}
\NormalTok{    Function(REPLACE,[c,Param(pPercent),Param(pEscPercent)]),}
\NormalTok{    Param(pBar),Param(pEscBar)])}

\NormalTok{scalarText\_C[Boolean](c) =}
\NormalTok{  Case([(Binary(EQ,c,Bool(true)),Param(pTrueText))],Param(pFalseText))}
\NormalTok{scalarText\_C[Integer](c) = Function(TO\_STRING,[c])}
\NormalTok{scalarText\_C[Real](c)    = Function(TO\_STRING,[c])}
\NormalTok{scalarText\_C[String](c)  = escapeString\_C(c)}

\NormalTok{encScalar\_C[tau](c) =}
\NormalTok{  Case([(isBottom(c),Param(pVoidWire))],}
\NormalTok{       Binary(CONCAT,Param(pPrefix\_tau),scalarText\_C[tau](c)))}

\NormalTok{anyBottom\_C([]) = Bool(false)}
\NormalTok{anyBottom\_C([(c,tau)] ++ qs) =}
\NormalTok{  Binary(OR,isBottom(c),anyBottom\_C(qs))}

\NormalTok{qualList\_C(qs=[(c1,tau1),...,(ck,tauk)]) =}
\NormalTok{  Case([(anyBottom\_C(qs),Null)],}
\NormalTok{       List([encScalar\_C[tau1](c1),...,encScalar\_C[tauk](ck)]))}

\NormalTok{memberIn[sigma,delta{-}\textgreater{}upsilon](SetPlan(Q,a),imports,L,e) =}
\NormalTok{  existsRows(}
\NormalTok{    SetPlan(Q,a),imports,L,}
\NormalTok{    setEq\_upsilon(}
\NormalTok{      coerce\_C[sigma{-}\textgreater{}upsilon](Alias(a)),}
\NormalTok{      coerce\_C[delta{-}\textgreater{}upsilon](e))).}

\NormalTok{notSubset\_tau(SetPlan(Q,a),SetPlan(R,b),imports) =}
\NormalTok{  existsRows(}
\NormalTok{    SetPlan(Q,a),imports,[],}
\NormalTok{    Unary(NOT,}
\NormalTok{      memberIn[tau,tau{-}\textgreater{}tau](SetPlan(R,b),imports union \{a\},[],Alias(a))))}

\NormalTok{extSetEq\_tau(R,T,imports) =}
\NormalTok{  ScalarPlan([],}
\NormalTok{    Binary(AND,}
\NormalTok{      Unary(NOT,notSubset\_tau(R,T,imports)),}
\NormalTok{      Unary(NOT,notSubset\_tau(T,R,imports))))}

\NormalTok{filterRows(SetPlan(Q,a),imports,L,p,out) =}
\NormalTok{  SetPlan(Seq([Call(imports,Q)] ++ L ++}
\NormalTok{              [Where(p),retD(setOut(Alias(a)),out)]),out)}

\NormalTok{mapRows(SetPlan(Q,a),imports,L,e,out) =}
\NormalTok{  SetPlan(Seq([Call(imports,Q)] ++ L ++}
\NormalTok{              [With(false,[Projection(setOut(e),out)]),}
\NormalTok{               retD(Alias(out),out)]),out)}

\NormalTok{importClauses([]) = []}
\NormalTok{importClauses(as) = [With(false,keep(as))] when as != []}

\NormalTok{prefixScalar(L,ScalarPlan(K,e)) = ScalarPlan(L ++ K,e)}
\NormalTok{prefixSet(L,SetPlan(Q,a),imports,z) =}
\NormalTok{  SetPlan(Seq(L ++ [Call(imports,Q),retD(Alias(a),z)]),z)}

\NormalTok{materialize(P,a,K) = bindScalar(P,a,K)}

\NormalTok{ExpandV(CQValueExpr(P,tau),alpha,pi,K) =}
\NormalTok{  let a=fresh;}
\NormalTok{      (ScalarPlan(L,e),pi\textquotesingle{})=ExpandE(P,alpha,pi,K);}
\NormalTok{  (ValuePlan(L ++ [withBind(K union \{a\},e,a)],a,tau),pi\textquotesingle{})}

\NormalTok{ExpandV(CQValueSet(R,Set(sigma)),alpha,pi,K) =}
\NormalTok{  let a=fresh;}
\NormalTok{      (SetPlan(Q,z),pi\textquotesingle{})=ExpandS(R,alpha,pi,K);}
\NormalTok{  (ValuePlan([withBind(K union \{a\},}
\NormalTok{                  CollectExpr(closeSet(SetPlan(Q,z))),a)],a,Set(sigma)),pi\textquotesingle{})}

\NormalTok{compose2(f,P1,P2,alpha,pi,K) =}
\NormalTok{  let a1,a2 be fresh;}
\NormalTok{      (P1\textquotesingle{},pi1)=ExpandE(P1,alpha,pi,K);}
\NormalTok{      X1=materialize(P1\textquotesingle{},a1,K);}
\NormalTok{      (P2\textquotesingle{},pi2)=ExpandE(P2,alpha,pi1,K union \{a1\});}
\NormalTok{      X2=materialize(P2\textquotesingle{},a2,K union \{a1\});}
\NormalTok{  in (ScalarPlan(X1.prelude ++ X2.prelude,}
\NormalTok{                 f(Alias(a1),Alias(a2))),pi2)}

\NormalTok{materializeList([],alpha,pi,K) = ([],[],K,pi)}
\NormalTok{materializeList([P1,...,Pn],alpha,pi0,K0) =}
\NormalTok{  for i=1..n from left to right, choose fresh ai,}
\NormalTok{    (Pi\textquotesingle{},pi\_i)=ExpandE(Pi,alpha,pi\_(i{-}1),K\_(i{-}1)),}
\NormalTok{    Xi=materialize(Pi\textquotesingle{},ai,K\_(i{-}1)),}
\NormalTok{    K\_i=K\_(i{-}1) union \{ai\};}
\NormalTok{  return (X1.prelude ++ ... ++ Xn.prelude,}
\NormalTok{          [Alias(a1),...,Alias(an)],K\_n,pi\_n)}

\NormalTok{setLiteralPlan(as,taus,upsilon) =}
\NormalTok{  let z=fresh;}
\NormalTok{      es=[setOut(coerce\_C[tau\_i{-}\textgreater{}upsilon](Alias(ai))) | i=1..n];}
\NormalTok{  SetPlan(Seq([Unwind(List(es),z),retD(Alias(z),z)]),z)}

\NormalTok{unionArm(SetPlan(Q,a),tau,upsilon,imports,z) =}
\NormalTok{  Seq([Call(imports,Q),}
\NormalTok{       Return(false,[Projection(}
\NormalTok{         setOut(coerce\_C[tau{-}\textgreater{}upsilon](Alias(a))),z)])])}

\NormalTok{unionPlan(S:tau,T:sigma,upsilon,imports) =}
\NormalTok{  let z=fresh;}
\NormalTok{  SetPlan(Seq([}
\NormalTok{    Call(imports,}
\NormalTok{      UnionAll(unionArm(S,tau,upsilon,imports,z),}
\NormalTok{               unionArm(T,sigma,upsilon,imports,z))),}
\NormalTok{    retD(Alias(z),z)]),z)}

\NormalTok{intersectionWitness(SetPlan(R,b),sigma,a,tau,upsilon,imports) =}
\NormalTok{  let unit=fresh;}
\NormalTok{  Seq([Call(imports union \{a\},R),}
\NormalTok{       Where(setEq\_upsilon(}
\NormalTok{         coerce\_C[tau{-}\textgreater{}upsilon](Alias(a)),}
\NormalTok{         coerce\_C[sigma{-}\textgreater{}upsilon](Alias(b)))),}
\NormalTok{       Return(false,[Projection(Int(1),unit)])])}

\NormalTok{intersectPlan(S=SetPlan(Q,a):tau,T:Set(sigma),upsilon,imports) =}
\NormalTok{  let z=fresh;}
\NormalTok{  SetPlan(Seq([Call(imports,Q),}
\NormalTok{    Where(ExistsExpr(intersectionWitness(T,sigma,a,tau,upsilon,imports))),}
\NormalTok{    retD(setOut(coerce\_C[tau{-}\textgreater{}upsilon](Alias(a))),z)]),z)}

\NormalTok{asSetPlan(SetPlan(Q,a),imports) =}
\NormalTok{  let z=fresh;}
\NormalTok{  SetPlan(Seq([Call(imports,Q),retD(setOut(Alias(a)),z)]),z)}

\NormalTok{uniquePlan(S=SetPlan(Q,a),imports,ScalarPlan(L,e)) =}
\NormalTok{  let z=fresh;}
\NormalTok{      sourceQ=injectImports(imports,closeSet(S));}
\NormalTok{      imageQ=Seq([Call(imports,Q)] ++ L ++}
\NormalTok{                 [retD(setOut(e),z)]);}
\NormalTok{  ScalarPlan([],Binary(EQ,CountExpr(sourceQ),CountExpr(imageQ)))}

\NormalTok{scalarArm(imports,g,ScalarPlan(L,e),out) =}
\NormalTok{  Seq(importClauses(imports union FV(g)) ++ [Where(g)] ++ L ++}
\NormalTok{      [Return(false,[Projection(e,out)])])}

\NormalTok{setArm(imports,g,SetPlan(Q,a),out) =}
\NormalTok{  Seq(importClauses(imports union FV(g)) ++ [Where(g),Call(imports,Q),}
\NormalTok{      Return(false,[Projection(setOut(Alias(a)),out)])])}

\NormalTok{ifScalar(imports,cc,T,F) =}
\NormalTok{  let out=fresh(imports union \{cc\} union aliases(T) union aliases(F));}
\NormalTok{  ScalarPlan([Call(imports,}
\NormalTok{    UnionAll(scalarArm(imports,truth(Alias(cc)),T,out),}
\NormalTok{             scalarArm(imports,Unary(NOT,truth(Alias(cc))),F,out)))],}
\NormalTok{    Alias(out))}

\NormalTok{ifSet(imports,cc,T,F,z) =}
\NormalTok{  let out=fresh(imports union \{cc,z\} union aliases(T) union aliases(F));}
\NormalTok{  SetPlan(Seq([Call(imports,}
\NormalTok{    UnionAll(setArm(imports,truth(Alias(cc)),T,out),}
\NormalTok{             setArm(imports,Unary(NOT,truth(Alias(cc))),F,out))),}
\NormalTok{    retD(Alias(out),z)]),z).}
\end{Highlighting}
\end{Shaded}

The finite codec-parameter family used by \texttt{encScalar\_C} is
reserved in \texttt{pi0(MM)} and is disjoint from the fresh allocator;
its values are exactly \texttt{v1\textbar{}V}, the type prefixes
\texttt{v1\textbar{}B\textbar{}}, \texttt{v1\textbar{}I\textbar{}},
\texttt{v1\textbar{}R\textbar{}}, \texttt{v1\textbar{}S\textbar{}}, the
Boolean texts, and the four escape strings as applicable. CY1 requires
\texttt{TO\_STRING} on admitted Int64/finite Real64 and the two
\texttt{REPLACE} calls to agree with the canonical Section 4.4 text
equations. If any qualifier is bottom, \texttt{qualList\_C} returns
\texttt{Null}; equality with a stored qualifier list is then not Cypher
true, and CY3 removes the row. Hence this helper expands to a finite raw
\texttt{CExpr} tree.

Entity-consuming scalar operations use the following total builder. A
supplied entity builder \texttt{F(recv,out)} is a clause list ending in
exactly one projection named \texttt{out} for every represented entity
receiver and may place \texttt{recv} in node positions.

\begin{Shaded}
\begin{Highlighting}[]
\NormalTok{bottomArm(recv,b,out) =}
\NormalTok{  Seq([Where(isBottom(Alias(recv))),}
\NormalTok{       Return(false,[Projection(b,out)])])}

\NormalTok{entityArm(recv,F,out) =}
\NormalTok{  Seq([Where(Unary(NOT,isBottom(Alias(recv))))] ++ F(recv,out))}

\NormalTok{withEntityReceiver(P,K,b,F) =}
\NormalTok{  let recv,out be fresh;}
\NormalTok{      X = materialize(P,recv,K);}
\NormalTok{      B = UnionAll(bottomArm(recv,b,out),entityArm(recv,F,out));}
\NormalTok{  in ScalarPlan(X.prelude ++ [Call(K union \{recv\},B)],Alias(out))}

\NormalTok{kindExpr(recv,pC) =}
\NormalTok{  {-}{-} rt, cls, and unit are fresh; pM/pC denote modelKey\_MM/key\_MM(C)}
\NormalTok{  ExistsExpr(Seq([}
\NormalTok{    Match(false,[Rel(Node(recv,Object,[(modelKey,Param(pM))]),}
\NormalTok{                     rt,OUT,ObjectInstanceOf,}
\NormalTok{                     Node(cls,UmlClass,[(modelKey,Param(pM)),}
\NormalTok{                                        (classKey,Param(pC))]))]),}
\NormalTok{    Return(false,[Projection(Int(1),unit)])]))}

\NormalTok{attrBody(a,pa)(recv,out) =}
\NormalTok{  {-}{-} ra and av are fresh; pa is bound to attributeKey\_MM(a)}
\NormalTok{  [Match(true,[Rel(Node(recv,Object,[(modelKey,Param(pM))]),}
\NormalTok{                   ra,OUT,ObjectHasAttribute,}
\NormalTok{                   Node(av,AttributeValue,[(modelKey,Param(pM)),}
\NormalTok{                                           (attributeKey,Param(pa))]))]),}
\NormalTok{   Return(false,[Projection(}
\NormalTok{     decodeScalar\_C[type(a)](Property(Alias(av),value)),out)])]}

\NormalTok{kindBody(pC)(recv,out) =}
\NormalTok{  [Return(false,[Projection(kindExpr(recv,pC),out)])]}

\NormalTok{castBody(pC)(recv,out) =}
\NormalTok{  [Return(false,[Projection(}
\NormalTok{     Case([(kindExpr(recv,pC),Alias(recv))],Null),out)])]}
\end{Highlighting}
\end{Shaded}

Under \texttt{BuilderClosure\_C}, \texttt{decodeScalar\_C{[}tau{]}(e)}
is a total raw-\texttt{CExpr} builder, not a semantic cast. It expands
the inverse wire equations of Section 4.4 into \texttt{Case}, equality,
substring/split, and the selected typed conversion functions. Its first
arm is \texttt{e\ IS\ NULL\ -\textgreater{}\ Null}, covering an absent
\texttt{OPTIONAL\ MATCH} slot; \texttt{v1\textbar{}V} also maps to
\texttt{Null}; the unique tag for \texttt{tau} maps to its canonical
Boolean/Int64/finite Real64/String value. Its syntactic default is
\texttt{Null}, so the builder remains a total \texttt{CExpr}; however
every other non-null payload violates the strict-codec part of
\texttt{GraphAdequate} and is outside the realization theorem rather
than being assigned a certified OCL value. Under \texttt{GraphAdequate},
only the missing, \texttt{v1\textbar{}V}, and correctly tagged arms are
reachable and BR4 gives:

\begin{Shaded}
\begin{Highlighting}[]
\NormalTok{abs\_C\^{}tau(EvalExpr\_G,pi(}
\NormalTok{  decodeScalar\_C[tau](Property(Alias(av),value)),row))}
\NormalTok{  = value\_G(row(recv),a),}
\end{Highlighting}
\end{Shaded}

for every admitted entity-arm row after the \texttt{OPTIONAL\ MATCH};
when the slot is absent, \texttt{row(av)=CypherNull} and both sides are
\texttt{bottom\_tau}.

The finite \texttt{FunctionId} enumeration includes exactly the
conversion/string functions used by this builder. Adding a new scalar
type requires extending the codec equations, this builder,
\texttt{BuilderClosure\_C}, CY1 evidence, and R4 together.

\texttt{Call} injects the listed imports into both union arms. If the
receiver is bottom, CY3 removes the entity arm\textquotesingle s row
before \texttt{F} is evaluated; therefore no node pattern receives
\texttt{CypherNull} or \texttt{BOTTOM\_TOKEN}. If it is not bottom, the
well-typed receiver and target-profile closure make it a represented
graph node, the bottom arm is empty, and \texttt{F} returns exactly one
row. Hence exactly one arm contributes one row, \texttt{out} is defined,
and aliases in \texttt{K} remain live.

\texttt{imports} is the sorted list of aliases denoting the free
variables of the embedded plan. The parameter state always contains
\texttt{pi(pLinkPrefix)="Link"} and
\texttt{pi(pM)=repr\_C\^{}String(modelKey\_MM)}. Each class, attribute,
and association builder adds its typed key parameters.
\texttt{qualList\_C} is the fully defined \texttt{Case/List} AST in the
concrete combinator table, not a text macro. The complete constructor
equations are as follows; \texttt{E\_K(P)} abbreviates the plan
component of \texttt{ExpandE(P,alpha,pi,K)} and \texttt{S\_K(R)}
abbreviates the plan component of \texttt{ExpandS(R,alpha,pi,K)} under
the current state. In every table row, recursive calls are evaluated
left-to-right and the parameter state returned by one call is the input
state of the next; the row returns the final state. Thus a notation such
as \texttt{S\_K(R),S\_K(T)} is a state-threaded sequence, never two
allocations from the same state. The completed recursive results are not
holes in an output AST.

{\def\LTcaptype{none} % do not increment counter
\begin{longtable}[]{@{}ll@{}}
\toprule\noalign{}
Query-model constructor & Defining raw-AST equation \\
\midrule\noalign{}
\endhead
\bottomrule\noalign{}
\endlastfoot
\texttt{CQAlias(x)} & \texttt{ScalarPlan({[}{]},Alias(alpha(x)))};
\texttt{alpha(x)} is in \texttt{K} whenever required later \\
\texttt{CQLiteral(v,tau)} & let
\texttt{(p,pi\textquotesingle{})=alloc\_tau(pi,v)}; return
\texttt{(ScalarPlan({[}{]},Param(p)),pi\textquotesingle{})} \\
\texttt{CQSetLiteral(Ps,taus,upsilon)} & let
\texttt{(L,aliases,K\textquotesingle{},pi\textquotesingle{})=materializeList(Ps,alpha,pi,K)}
and \texttt{R=setLiteralPlan(aliases,taus,upsilon)}; return
\texttt{(prefixSet(L,R,K\textquotesingle{},exported(R)),pi\textquotesingle{})},
so every element alias is bound before \texttt{UNWIND} \\
\texttt{CQAttribute(P,a)} & let \texttt{(PP,pi1)=ExpandE(P,alpha,pi,K)}
and \texttt{(pa,pi2)=allocAttribute(pi1,a)}; return
\texttt{(withEntityReceiver(PP,K,Null,attrBody(a,pa)),pi2)}; the bottom
arm contains no pattern and the entity arm uses the unique canonical
attribute lookup \\
\texttt{CQNot(P)} & if \texttt{E\_K(P)=ScalarPlan(L,e)}, return
\texttt{ScalarPlan(L,Unary(NOT,truth(e)))} \\
\texttt{CQBoolean(op,P1,P2)} &
\texttt{compose2(lambda\ x,y.\ boolAst(op,x,y),P1,P2,alpha,pi,K)} \\
\texttt{CQCompare(op,P1,tau,P2,sigma,upsilon)} & use \texttt{compose2}
with \texttt{semEq\_upsilon} or \texttt{semOrd\_upsilon} after
\texttt{coerce\_C{[}tau-\textgreater{}upsilon{]}} and
\texttt{coerce\_C{[}sigma-\textgreater{}upsilon{]}} on its two
materialized arguments \\
\texttt{CQSetEquality(op,R,T,tau)} & expand \texttt{S\_K(R)} and then
\texttt{S\_K(T)} with the returned state; return
\texttt{extSetEq\_tau(S\_K(R),S\_K(T),imports)} and wrap its result
expression in Boolean negation for \texttt{NEQ} \\
\texttt{CQArith(op,P1,tau,P2,sigma,upsilon)} & use \texttt{compose2}
with \texttt{semArith\_upsilon(op,...)} after the same two canonical
coercions \\
\texttt{CQCoerce(tau,upsilon,P)} & if \texttt{E\_K(P)=ScalarPlan(L,e)},
return
\texttt{ScalarPlan(L,coerce\_C{[}tau-\textgreater{}upsilon{]}(e))} \\
\texttt{CQIfExpr(C,T,F)} & choose fresh \texttt{cc}; let
\texttt{XC=materialize(E\_K(C),cc,K)}, \texttt{I=K\ union\ \{cc\}},
\texttt{PT=E\_I(T)}, and \texttt{PF=E\_I(F)}; return
\texttt{prefixScalar(XC.prelude,ifScalar(I,cc,PT,PF))} \\
\texttt{CQLetExpr(x,V,B)} & let
\texttt{(ValuePlan(Lx,a,tau),pi1)=ExpandV(V,alpha,pi,K)}; expand
\texttt{B} under \texttt{alpha{[}x-\textgreater{}a{]}}, \texttt{pi1},
and live set \texttt{K\ union\ \{a\}}; prefix its scalar plan with
\texttt{Lx}, return its final state, then restore the outer
\texttt{alpha} \\
\texttt{CQExists(R,x,P)} & let \texttt{S\_K(R)=SetPlan(Q,a)}; under
\texttt{alpha\textquotesingle{}=alpha{[}x-\textgreater{}a{]}} use live
set \texttt{K\textquotesingle{}=K\ union\ \{a\}} and let
\texttt{E\_K\textquotesingle{}(P)=ScalarPlan(L,p)}; return
\texttt{ScalarPlan({[}{]},existsRows(S\_K(R),imports,L,truth(p)))} \\
\texttt{CQForAll(R,x,P)} & use the same extended alias map/live set and
return
\texttt{ScalarPlan({[}{]},Unary(NOT,existsRows(S\_K(R),imports,L,Unary(NOT,truth(p)))))} \\
\texttt{CQMembership(IN,R,sigma,Y,delta,upsilon)} & first expand
\texttt{S\_K(R)}; choose fresh \texttt{y}, expand/materialize \texttt{Y}
with the returned state and live set \texttt{K}, and return
\texttt{ScalarPlan(XY.prelude,memberIn{[}sigma,delta-\textgreater{}upsilon{]}(S\_K(R),imports\ union\ \{y\},{[}{]},Alias(y)))};
\texttt{NOT\_IN} wraps the result in \texttt{Unary(NOT,...)} \\
\texttt{CQSetRelation(SUBSET,R,tau,T,sigma,upsilon)} & expand
\texttt{S\_K(R)} then \texttt{S\_K(T)}; for each exported
\texttt{t:sigma} of \texttt{T}, call
\texttt{memberIn{[}tau,sigma-\textgreater{}upsilon{]}(S\_K(R),imports,{[}{]},Alias(t))};
return the negation of the \texttt{existsRows} counterexample in which
that membership is false \\
\texttt{CQSetRelation(DISJOINT,R,tau,T,sigma,upsilon)} & use the same
typed \texttt{memberIn} and return the negation of the
\texttt{existsRows} witness in which a right element is also a left
member \\
\texttt{CQIsUnique(R,x,E)} & let \texttt{S\_K(R)=SetPlan(Q,a)}; under
\texttt{alpha{[}x-\textgreater{}a{]}} and \texttt{K\ union\ \{a\}},
expand the body with the returned state to \texttt{ScalarPlan(L,e)} and
return \texttt{uniquePlan(S\_K(R),imports,ScalarPlan(L,e))} \\
\texttt{CQCount(R)} & return
\texttt{ScalarPlan({[}{]},CountExpr(closeSet(S\_K(R))))}; the set plan
already returns \texttt{DISTINCT\ setOut(element)} \\
\texttt{CQEmpty(EMPTY,R)} & return
\texttt{ScalarPlan({[}{]},Unary(NOT,ExistsExpr(closeSet(S\_K(R)))))};
\texttt{NONEMPTY} returns
\texttt{ScalarPlan({[}{]},ExistsExpr(closeSet(S\_K(R))))} \\
\texttt{CQKindOf(P,C)} & let \texttt{(PP,pi1)=ExpandE(P,alpha,pi,K)} and
\texttt{(pC,pi2)=allocClass(pi1,C)}; return
\texttt{(withEntityReceiver(PP,K,Bool(false),kindBody(pC)),pi2)} \\
\texttt{CQCast(P,C)} & let \texttt{(PP,pi1)=ExpandE(P,alpha,pi,K)} and
\texttt{(pC,pi2)=allocClass(pi1,C)}; return
\texttt{(withEntityReceiver(PP,K,Null,castBody(pC)),pi2)}; the receiver
is materialized exactly once \\
\texttt{CQNavigationOne(P,A,fr,tr,dir,qs)} & form
\texttt{RN=CQNavigation(CQSourceExpr(P),A,fr,tr,dir,qs)}, compute
\texttt{(R,pi\textquotesingle{})=ExpandS(RN,alpha,pi,K)}, and return
\texttt{(ScalarPlan({[}{]},Function(HEAD,{[}CollectExpr(closeSet(R)){]})),pi\textquotesingle{})};
valid multiplicity gives at most one semantic target and CY5 maps empty
to null/bottom \\
\texttt{CQEmptySet(tau)} & for fresh \texttt{z}, return a set plan that
unwinds \texttt{{[}{]}} and exports \texttt{z}, hence has zero rows \\
\texttt{CQSetAlias(x,tau)} & for fresh \texttt{z}, return
\texttt{collectionRows(ScalarPlan({[}{]},Alias(alpha(x))),z)} \\
\texttt{CQAllInstances(C)} & let \texttt{(pC,pi1)=allocClass(pi,C)}; for
fresh \texttt{n,rt,cls}, return with \texttt{pi1} a \texttt{SetPlan}
matching
\texttt{Rel(Node(n,Object,{[}(modelKey,Param(pM)){]}),rt,OUT,ObjectInstanceOf,Node(cls,UmlClass,{[}(modelKey,Param(pM)),(classKey,Param(pC)){]}))},
followed by \texttt{retD(Alias(n),n)} \\
\texttt{CQNavigation(src,A,fr,tr,dir,qs)} & let
\texttt{(SetPlan(Q,s),pi1)=ExpandSrc(src,alpha,pi,K)} and
\texttt{K0=K\ union\ \{s\}}; let
\texttt{materializeList(qs,alpha,pi1,K0)=(Lq,qexprs,Kq,pi2)} and
\texttt{(pA,pFr,pTr,pi3)=allocNavigation(pi2,A,fr,tr)}; call \texttt{Q},
append \texttt{Lq} (which preserves \texttt{K0}), then match scoped
\texttt{Node(s,Object,...)} and \texttt{Node(t,Object,...)} in the
resolved direction; append
\texttt{Where(linkPredicate(r,pM,pA,pFr,pTr,dir,qualList\_C(zip(qexprs,qualifierTypes))))};
finish with \texttt{retD(setOut(Alias(t)),t)} and return \texttt{pi3} \\
\texttt{CQViewSet(P,D)} & materialize \texttt{E(P)=ScalarPlan(L,c)} and,
for fresh \texttt{z}, unwind \texttt{viewOne\_C(c)} after \texttt{L};
LIFT1 maps scalar bottom to zero rows and an entity to one row \\
\texttt{CQUnion(P,tau,Q,sigma,upsilon)} & expand \texttt{S\_K(P)} and
then \texttt{S\_K(Q)}; return
\texttt{unionPlan(S\_K(P):tau,S\_K(Q):sigma,upsilon,imports)} \\
\texttt{CQIntersection(P,tau,Q,sigma,upsilon)} & expand \texttt{S\_K(P)}
and then \texttt{S\_K(Q)}; return
\texttt{intersectPlan(S\_K(P):tau,S\_K(Q):sigma,upsilon,imports)} \\
\texttt{CQAsSet(P)} & return \texttt{asSetPlan(S\_K(P),imports)} \\
\texttt{CQSelect(R,x,P)} & if \texttt{S\_K(R)=SetPlan(Q,a)}, then under
\texttt{alpha{[}x-\textgreater{}a{]}}, live set
\texttt{K\ union\ \{a\}}, and the returned state expand \texttt{P} to
\texttt{ScalarPlan(L,p)}; return
\texttt{filterRows(S\_K(R),imports,L,truth(p),a)} \\
\texttt{CQReject(R,x,P)} & under the same bindings, return
\texttt{filterRows(S\_K(R),imports,L,Unary(NOT,truth(p)),a)} \\
\texttt{CQCollect(R,x,P)} & under the same bindings expand
\texttt{P=ScalarPlan(L,p)} and, for fresh \texttt{z}, return
\texttt{mapRows(S\_K(R),imports,L,p,z)} \\
\texttt{CQIfSet(C,T,F)} & choose fresh \texttt{cc,z}; let
\texttt{XC=materialize(E\_K(C),cc,K)}; with \texttt{I=K\ union\ \{cc\}},
expand \texttt{S\_I(T)} and then \texttt{S\_I(F)}; return
\texttt{prefixSet(XC.prelude,ifSet(I,cc,S\_I(T),S\_I(F),z),I,z)} \\
\texttt{CQLetSet(x,V,B)} & let
\texttt{(ValuePlan(Lx,a,tau),pi1)=ExpandV(V,alpha,pi,K)} and
\texttt{(SetPlan(Q,z),pi2)=ExpandS(B,alpha{[}x-\textgreater{}a{]},pi1,K\ union\ \{a\})};
return \texttt{(prefixSet(Lx,SetPlan(Q,z),K\ union\ \{a\},z),pi2)}, then
restore the outer \texttt{alpha} \\
\texttt{CQInvariant(C,P)} & under
\texttt{alpha0=\{self-\textgreater{}self\}}, \texttt{K=\{self\}}, and an
input state containing \texttt{pM}, let
\texttt{(ScalarPlan(L,p),pi1)=ExpandE(P,alpha0,pi,K)} and
\texttt{(pCtx,pi2)=allocClass(pi1,C)}; emit with \texttt{pi2} a scoped
context match
\texttt{(self:Object\ \{modelKey:\$pM\})-{[}:ObjectInstanceOf{]}-\textgreater{}(cls:UmlClass\ \{modelKey:\$pM,classKey:\$pCtx\})},
followed by all of \texttt{L}, then \texttt{Where(Unary(NOT,truth(p)))}
and
\texttt{Return(true,{[}Projection(Property(Alias(self),use\_id),useId){]})} \\
\end{longtable}
}

The apparent names \texttt{Q}, \texttt{a}, \texttt{p}, \texttt{s}, and
\texttt{t} in this table are mathematical names bound by \texttt{let} in
the defining equation. Every right-hand side consists only of raw AST
constructors, total builders defined above, recursive results, and
deterministic \texttt{fresh}/parameter allocation. In particular, branch
imports, branch aliases, attribute keys, and result aliases are
constructed AST atoms; no uninterpreted surface placeholder belongs to
the codomain of \texttt{Expand}.

\subsubsection{TXT4 Hard-Case Derivation
Certificate}\label{txt4-hard-case-derivation-certificate}

TXT4 is discharged by induction on the query-model constructor, using
the raw equations above. The following cases record the non-immediate
derivations; the remaining scalar homomorphism and base cases are direct
applications of C1-\/-C5.

\textbf{D1 Qualified navigation.} Assume the source plan realizes
\texttt{S} and each qualifier plan realizes \texttt{qi}.
\texttt{materializeList} executes qualifier preludes left-to-right and,
by repeated C1, leaves a row corresponding to the environment with every
qualifier alias live. Therefore:

\begin{Shaded}
\begin{Highlighting}[]
\NormalTok{[abs\_C\^{}tau1(EvalExpr\_G,pi(qexpr1,row)),...,}
\NormalTok{ abs\_C\^{}tauk(EvalExpr\_G,pi(qexprk,row))]}
\NormalTok{= [[[q1]]\_graph,...,[[qk]]\_graph].}
\end{Highlighting}
\end{Shaded}

Admission gives \texttt{k=arity(A,dir)} and fixes component \texttt{i}
to the declared type of qualifier \texttt{i} at the direction-selected
end. \texttt{qualList\_C} applies the fixed injective serialization
componentwise without reordering. The raw expansion uses
\texttt{qualifierProperty(OUT)=sourceQualifiers} and
\texttt{qualifierProperty(IN)=targetQualifiers}, exactly the concrete
realization of \texttt{linkQualifiers\_G(r,dir)}. Likewise
\texttt{sourceRoleParam/targetRoleParam} use \texttt{(fr,tr)} for
\texttt{OUT} and \texttt{(tr,fr)} for \texttt{IN}, exactly matching the
two role equations in \texttt{nav\_graph}; reverse navigation therefore
does not test the roles in forward order. The source plan either
contributes no row for bottom or a represented node. R6 and CY2 then
identify the matched targets with \texttt{nav\_graph};
\texttt{retD(setOut(Alias(t)),t)} and C5 quotient duplicate relationship
rows. Hence the projected semantic set is exactly the graph denotation
of the qualified navigation. Injectivity and length preservation also
give the reflection direction: a row cannot match a different qualifier
tuple or a tuple of different arity.

\textbf{D2 Nested \texttt{IncludesAll}.} Let \texttt{RS=SetPlan(Q\_R,a)}
realize \texttt{R:Set(tau)} and \texttt{TS=SetPlan(Q\_T,t)} realize
\texttt{T:Set(sigma)}, with \texttt{memberJoin(tau,sigma)=upsilon}. For
each right-side execution row \texttt{row\_T}, BR9 and C5 make the
evaluation of
\texttt{memberIn{[}tau,sigma-\textgreater{}upsilon{]}(RS,...,Alias(t))}
true exactly when
\texttt{iota\_graph{[}sigma-\textgreater{}upsilon{]}(abs\_C\^{}sigma(row\_T(t)))}
belongs to
\texttt{upSet\_graph{[}tau-\textgreater{}upsilon{]}(finiteSet\_tau({[}{[}R{]}{]}\_graph))}.
Consequently the outer counterexample query has a row iff:

\begin{Shaded}
\begin{Highlighting}[]
\NormalTok{exists t in finiteSet\_sigma([[T]]\_graph):}
\NormalTok{  iota\_graph[sigma{-}\textgreater{}upsilon](t)}
\NormalTok{    notin upSet\_graph[tau{-}\textgreater{}upsilon](finiteSet\_tau([[R]]\_graph)).}
\end{Highlighting}
\end{Shaded}

Negating that existential is precisely subset semantics. Explicit
imports keep the outer \texttt{t} correlated with the inner membership
query. Repeating the same argument without the inner negation proves
\texttt{ExcludesAll} as disjointness.

\textbf{D3 Iteration over scalar sets.} Let an element \texttt{s} of the
source denote a scalar. C5 gives an output row with
\texttt{abs\_C\^{}tauE(row(a))=s}. Under
\texttt{alpha{[}x-\textgreater{}a{]}} and live set
\texttt{K\ union\ \{a\}}, the declared-binder embedding and C1 give
\texttt{RowCorr\_Gamma{[}x:tauD{]}} for
\texttt{eta{[}x-\textgreater{}iota\_graph{[}tauE-\textgreater{}tauD{]}(s){]}}.
The body induction hypothesis therefore applies exactly as it does for
an object element, without invoking \texttt{node(s)}. This discharges
\texttt{Exists}, \texttt{Select}, \texttt{Reject}, and \texttt{Collect}
over both object and scalar finite sets.

\textbf{D4 Set-valued \texttt{If}.} The condition is materialized once
as \texttt{cc}. Since \texttt{truth(cc)} is total Boolean, exactly one
of \texttt{truth(cc)} and \texttt{NOT\ truth(cc)} holds. CY3 eliminates
the unselected arm before its branch query executes; CY9 concatenates
the selected rows only. Both arms return the same freshly generated
output alias and joined element type. C5 applies \texttt{DISTINCT} after
\texttt{setOut}, so the result is exactly the selected finite-set
denotation, independently of Cypher row multiplicity. When the selected
source branch has type \texttt{Void}, its
\texttt{Void\ -\textgreater{}\ Set(tau)} embedding is
\texttt{bottom\_Set(tau)} and C5a\textquotesingle s
\texttt{CQEmptySet(tau)} realization contributes zero rows, as required.

\textbf{D5 \texttt{Collect} containing bottom.} Let the body expansion
be \texttt{ScalarPlan(L,p)} and let
\texttt{ExecPrelude\_G,pi(L,row\_s)=\{row\textquotesingle{}\_s\}} for a
source element \texttt{s}. The body induction hypothesis yields either
\texttt{repr\_C\^{}sigma(v)} or \texttt{CypherNull} for semantic bottom.
\texttt{setOut} maps the latter to the non-null \texttt{BOTTOM\_TOKEN}.
BR1-\/-BR9 imply:

\begin{Shaded}
\begin{Highlighting}[]
\NormalTok{abs\_C\^{}sigma(EvalExpr\_G,pi(setOut(p),row\textquotesingle{}\_s))}
\NormalTok{  = [[body]]\_graph\^{}sigma(eta[x{-}\textgreater{}s]).}
\end{Highlighting}
\end{Shaded}

\texttt{DISTINCT} removes only duplicate semantic image values and
neither drops nor merges the token with a non-bottom value. Thus
\texttt{mapRows} realizes finite-set image construction including bottom
as one possible image element.

\textbf{D6 Typed or inferred \texttt{Let}.} Let the initializer have
type \texttt{tau1} and the stored declaration type be \texttt{tauD},
with \texttt{tau1\textless{}=tauD}. Expand the initializer once, apply
the canonical declaration-boundary embedding, and materialize the result
as fresh alias \texttt{a}. C1 gives:

\begin{Shaded}
\begin{Highlighting}[]
\NormalTok{RowCorr\_Gamma[x:tauD](alpha[x{-}\textgreater{}a],}
\NormalTok{        eta[x{-}\textgreater{}iota\_graph[tau1{-}\textgreater{}tauD]([[init]]\_graph\^{}tau1(eta))],}
\NormalTok{        row\textquotesingle{}).}
\end{Highlighting}
\end{Shaded}

The induction hypothesis for the proper body child therefore yields
exactly the denotation of the body under the let-extended environment.
Freshness and live projection prevent alias capture and preserve outer
bindings. Both recursive calls are on proper constructor children, so
structural recursion proves termination.

\textbf{D7 Attribute, kind-of, and cast with bottom.} Let the receiver
plan realize \texttt{E}. If \texttt{{[}{[}E{]}{]}\_graph=bottom},
BR1-\/-BR4 make \texttt{isBottom(recv)=true}; CY3 removes the entity arm
before its graph pattern and the bottom arm returns respectively null,
false, or null. These abstract to the primitive graph denotations of
attribute, \texttt{TypeKindOf}, and \texttt{Cast} on bottom. If the
receiver is an entity, target-profile closure represents it by a node,
only the entity arm executes, and R4 or R3 proves the returned value.
Complementary guards and singleton entity bodies prove
\texttt{WFScalar}; C1 preserves every alias in \texttt{K} across the
correlated call.

\textbf{D8 Set literal.} \texttt{materializeList} evaluates element
plans left-to-right and C1 preserves all live aliases. M3a and BR9 make
each \texttt{coerce\_C{[}tau\_i-\textgreater{}upsilon{]}} denote the
unique canonical lifting to the \texttt{setJoinN} result.
\texttt{setOut} retains bottom as the reserved non-null token, and
BR2/BR5 make \texttt{DISTINCT} quotient exactly semantic typed equality.
The exported column is therefore the extensional finite set denoted by
the literal.

\textbf{D9 Union and intersection.} For union, CY9 concatenates the two
row bags after BR9 identifies both emitted coercions with their semantic
liftings to \texttt{upsilon}; BR5 then removes exactly semantic
duplicates. For intersection, a coerced left value survives exactly when
CY4 finds a coerced right witness satisfying total typed
\texttt{setEq\_upsilon}. BR1-\/-BR4 make this condition equivalent to
membership in both finite sets, including bottom.

\textbf{D10 \texttt{asSet}.} The induction hypothesis already gives an
extensional finite set. Reprojecting its column through \texttt{setOut}
and \texttt{DISTINCT} neither introduces nor removes a semantic value by
BR5, so \texttt{asSetPlan} realizes the identity of the Set-only
fragment.

\textbf{D11 \texttt{isUnique}.} The closed source plan has one row per
semantic source element. The body plan is evaluated exactly once for
each such row. CY5 and BR6 make equality between the source count and
\texttt{COUNT(DISTINCT\ setOut(body))} hold exactly when no two distinct
source elements have equal typed body values. The empty source yields
\texttt{0=0}, as required.

\textbf{D12 Whole-collection bottom and \texttt{ViewSet}.} C5a proves
that a collection-typed bottom yields zero rows, whereas
\texttt{viewOne\_C} proves the separate scalar-navigation
empty/singleton cases of LIFT1. In both builders, an element-level
bottom is passed to \texttt{setOut} only after membership in an actual
finite set has been established.

\textbf{D13 Extensional set equality.} Expand both operands as set
plans. C5b proves that the two generated subset checks are jointly true
exactly when their projected semantic sets are equal. Boolean negation
gives \texttt{NEQ}; neither case uses element equality as if it were
whole-set equality.

Together D1-\/-D13 cover every nontrivial constructor family. Thus TXT4
is not used as an assumption: it is the conclusion of the constructor
induction from C1-\/-C5b, R3-\/-R6, BR1-\/-BR10, and CY1-\/-CY9.

\subsubsection{\texorpdfstring{Total
\texttt{renderRaw}}{Total renderRaw}}\label{total-renderraw}

Alias identifiers are generated from
\texttt{{[}A-Za-z\_{]}{[}A-Za-z0-9\_{]}*}; labels, property keys,
relationship types, operators, and function names come from the finite
certified enums and are escaped by their enum-specific printers. Define
explicit import insertion before concrete rendering:

\begin{Shaded}
\begin{Highlighting}[]
\NormalTok{injectImports([],Q) = Q}
\NormalTok{injectImports(K,Seq(cs)) = Seq(importClauses(K) ++ cs)}
\NormalTok{injectImports(K,UnionAll(Q1,Q2))}
\NormalTok{  = UnionAll(injectImports(K,Q1),injectImports(K,Q2)).}
\end{Highlighting}
\end{Shaded}

The pretty-printer is the total structural function below. \texttt{join}
is ordinary delimiter insertion and \texttt{paren} always emits
parentheses, so precedence is not an implicit premise.

{\def\LTcaptype{none} % do not increment counter
\begin{longtable}[]{@{}ll@{}}
\toprule\noalign{}
Raw constructor & \texttt{renderRaw} equation \\
\midrule\noalign{}
\endhead
\bottomrule\noalign{}
\endlastfoot
\texttt{Alias(a)}, \texttt{Param(p)} & \texttt{printAlias(a)},
respectively \texttt{\$} followed by \texttt{printParam(p)} \\
\texttt{Null}, \texttt{Bool(b)}, \texttt{Int(i)} & \texttt{null},
\texttt{true/false}, and canonical decimal integer text \\
\texttt{Property(e,k)} &
\texttt{paren(render(e))\ +\ "."\ +\ printKey(k)} \\
\texttt{List(es)} &
\texttt{"{[}"\ +\ join(",",map(render,es))\ +\ "{]}"} \\
\texttt{Unary(op,e)} & \texttt{paren(printUnary(op)+"\ "+render(e))} \\
\texttt{Binary(op,l,r)} &
\texttt{paren(render(l)+"\ "+printBinary(op)+"\ "+render(r))} \\
\texttt{Case(bs,d)} & \texttt{CASE} followed by each
\texttt{WHEN\ render(c)\ THEN\ render(v)}, then
\texttt{ELSE\ render(d)\ END} \\
\texttt{Function(f,args)} &
\texttt{printFunction(f)+"("+join(",",map(render,args))+")"} \\
\texttt{ListComp(a,s,p,e)} & \texttt{{[}a\ IN\ render(s)} plus optional
\texttt{WHERE\ render(p)}, then \texttt{\textbar{}\ render(e){]}} \\
\texttt{ExistsExpr(Q)} & \texttt{EXISTS\ \{\ renderRaw(Q)\ \}} \\
\texttt{CountExpr(Q)} & \texttt{COUNT\ \{\ renderRaw(Q)\ \}} \\
\texttt{CollectExpr(Q)} & \texttt{COLLECT\ \{\ renderRaw(Q)\ \}} \\
\texttt{Node(a,label,props)} & parenthesized alias, optional escaped
label, and rendered parameterized property map \\
\texttt{Rel(n1,r,dir,type,n2)} & rendered endpoint nodes and
relationship variable/type with the selected arrow direction \\
\texttt{Match(false,ps)} & \texttt{MATCH\ } plus joined rendered
patterns \\
\texttt{Match(true,ps)} & \texttt{OPTIONAL\ MATCH\ } plus joined
rendered patterns \\
\texttt{Where(e)} & \texttt{WHERE\ render(e)} \\
\texttt{Unwind(e,a)} & \texttt{UNWIND\ render(e)\ AS\ printAlias(a)} \\
\texttt{With(d,ps)} & \texttt{WITH\ } plus optional \texttt{DISTINCT}
and rendered projections \\
\texttt{Return(d,ps)} & \texttt{RETURN\ } plus optional
\texttt{DISTINCT} and rendered projections \\
\texttt{Call(K,Q)} &
\texttt{CALL\ \{\ renderRaw(injectImports(K,Q))\ \}} \\
\texttt{Seq(cs)} & newline-join \texttt{renderRaw} of clauses in
order \\
\texttt{UnionAll(Q1,Q2)} &
\texttt{renderRaw(Q1)\ +\ newline\ +\ "UNION\ ALL"\ +\ newline\ +\ renderRaw(Q2)} \\
\end{longtable}
}

Projection rendering is \texttt{render(e)\ AS\ printAlias(a)}. No model
value is printed as syntax: every such value remains a \texttt{Param}
and is returned in the parameter map. This table defines concrete
\texttt{CALL} import lowering for the selected Cypher 5 profile and is
total for both \texttt{Seq} and \texttt{UnionAll} subqueries.

\texttt{Expand} is total by simultaneous well-founded recursion on the
lexicographic measure
\texttt{(\#\ of\ not-yet-lowered\ CQNavigationOne\ roots,\ total\ CQ\ syntax\ size)}.
Ordinary recursive calls, including both children of \texttt{Let}, go to
proper children. The only non-child call is the \texttt{CQNavigationOne}
adapter: replacing its root by
\texttt{CQNavigation(CQSourceExpr(P),...)} removes exactly one
not-yet-lowered \texttt{CQNavigationOne} root, so the first measure
component decreases even if the adapter term is larger. Each equation
preserves the plan well-formedness invariant by freshness and explicit
alias import. Consequently \texttt{closeScalar}, \texttt{closeSet}, or
\texttt{ExpandI} always yields a closed raw \texttt{CQuery} with no
holes.

Choose pairwise distinct reserved names \texttt{pBottom},
\texttt{pLinkPrefix}, and \texttt{pM}, disjoint from the fresh
allocator, and define the exact initial state:

\begin{Shaded}
\begin{Highlighting}[]
\NormalTok{codecParams = \{}
\NormalTok{  pVoidWire  {-}\textgreater{} repr\_C\^{}String("v1|V"),}
\NormalTok{  pPrefix\_B  {-}\textgreater{} repr\_C\^{}String("v1|B|"),}
\NormalTok{  pPrefix\_I  {-}\textgreater{} repr\_C\^{}String("v1|I|"),}
\NormalTok{  pPrefix\_R  {-}\textgreater{} repr\_C\^{}String("v1|R|"),}
\NormalTok{  pPrefix\_S  {-}\textgreater{} repr\_C\^{}String("v1|S|"),}
\NormalTok{  pTrueText  {-}\textgreater{} repr\_C\^{}String("true"),}
\NormalTok{  pFalseText {-}\textgreater{} repr\_C\^{}String("false"),}
\NormalTok{  pPercent   {-}\textgreater{} repr\_C\^{}String("\%"),}
\NormalTok{  pEscPercent{-}\textgreater{} repr\_C\^{}String("\%25"),}
\NormalTok{  pBar       {-}\textgreater{} repr\_C\^{}String("|"),}
\NormalTok{  pEscBar    {-}\textgreater{} repr\_C\^{}String("\%7C")}
\NormalTok{\}}

\NormalTok{pi0(MM) = \{}
\NormalTok{  pBottom     {-}\textgreater{} BOTTOM\_TOKEN,}
\NormalTok{  pLinkPrefix {-}\textgreater{} repr\_C\^{}String("Link"),}
\NormalTok{  pM          {-}\textgreater{} repr\_C\^{}String(modelKey\_MM)}
\NormalTok{\} union codecParams.}
\end{Highlighting}
\end{Shaded}

Every class/attribute/association/role/qualifier parameter, including
the context-class parameter \texttt{pCtx}, is then allocated once from
the threaded state. The formal specification of the final pipeline stage
is factored as:

\begin{Shaded}
\begin{Highlighting}[]
\NormalTok{(Qspec,pispec) = ExpandI(BuildCQInvariant(C,nva),pi0(MM))}
\NormalTok{T\_TEXT\^{}spec(C,nva) = (renderRaw(Qspec),pispec).}
\end{Highlighting}
\end{Shaded}

This equation specifies the reference transformation.
\texttt{BuildCQInvariant} is the definitional presentation of
\texttt{T\_CQM\^{}spec}, and \texttt{renderRaw\ o\ ExpandI} is the
definitional presentation of \texttt{T\_TEXT\^{}spec}; Raw Cypher AST is
not a separately trusted normative layer. A Java implementation that
renders from another plan is supporting evidence only until a separate
implementation-refinement theorem is supplied.

Define the structural reference relation:

\begin{Shaded}
\begin{Highlighting}[]
\NormalTok{SpecPlanAdequacy(MM,C,nva,plan,q,pi,G)}
\NormalTok{iff ValidNVA(nva)}
\NormalTok{  and WF\_CQM(plan)}
\NormalTok{  and SpecPlanSim(MM,C,nva,plan)}
\NormalTok{  and ParamCorr(plan,q,pi)}
\NormalTok{  and AliasCorr(nva,plan,q)}
\NormalTok{  and ScopeCorr(nva,plan,q)}
\NormalTok{  and MultiplicityCorr(nva,plan)}
\NormalTok{  and CollectionShapeCorr(nva,plan)}
\NormalTok{  and q=parse\_Cypher(T\_TEXT\^{}spec(plan))}
\NormalTok{  and AdmissibleExec(MM,C,nva,plan,q,pi,G).}
\end{Highlighting}
\end{Shaded}

No evaluation equality occurs in this definition. \texttt{SpecPlanSim}
has exactly one rule for each of the 26 concrete reachable NVA
constructors and carries the parameter, alias, scope, multiplicity, and
collection-shape conditions. The base rules are \texttt{NvaVariable} and
\texttt{NvaLiteral}; list and optional children use the corresponding
pointwise relation; binder rules extend the related environments with
the same declared type.

\textbf{SpecPlanSim\_sound.} If every local NVA/CQM constructor pair
satisfies its typed primitive evaluation equation under CY1-\/-CY9, then
structural induction over \texttt{SpecPlanSim} gives:

\begin{Shaded}
\begin{Highlighting}[]
\NormalTok{SpecPlanAdequacy(MM,C,nva,plan,q,pi,G)}
\NormalTok{implies}
\NormalTok{ExecCypher(q,G,pi) = EvalGraph(nva,G)}
\end{Highlighting}
\end{Shaded}

and the invariant wrapper/C6 specialization gives:

\begin{Shaded}
\begin{Highlighting}[]
\NormalTok{returnedIds(q,G,pi)=ViolIds\_\{M,G\}(C,nva).}
\end{Highlighting}
\end{Shaded}

The induction composition is mechanized by
\texttt{SpecificationPlanRefinement.spec\_plan\_sim\_sound}; the local
primitive equations are exactly the \texttt{LR-*}, \texttt{C*},
\texttt{BR*}, and \texttt{CY*} obligations enumerated in the
constructor-coverage certificate. Thus equality is a theorem conclusion,
not an adequacy premise.

\subsubsection{Non-normative Java optimizer/planner
bridge}\label{non-normative-java-optimizerplanner-bridge}

The remainder of this bridge subsection records the optional Java
\texttt{T\_OPT} refinement obligation. It is not consumed by Theorems
5-\/-6.

Make that bridge an explicit predicate. Let:

\begin{Shaded}
\begin{Highlighting}[]
\NormalTok{(Qspec,pispec) = ExpandI(BuildCQInvariant(C,nva),pi0(MM))}

\NormalTok{AdmissibleExec(MM,C,nva,Qspec,pispec,Qprod,piprod,G)}
\NormalTok{iff Qspec and Qprod belong to CYPHER5\_val}
\NormalTok{  and Params(Qspec) subseteq dom(pispec)}
\NormalTok{  and Params(Qprod) subseteq dom(piprod)}
\NormalTok{  and CY1{-}{-}CY9 hold for both evaluations on G}
\NormalTok{  and every reached runtime value belongs to its indexed CyVal\_C fibre}
\NormalTok{  and both public useId columns, when evaluated, contain only ObjectId values.}

\NormalTok{PlanAdequacy(MM,C,nva,opt,Qprod,piprod,G)}
\NormalTok{iff exists an interface{-}alias bijection hLive:}
\NormalTok{  ProdPlanSim(MM,C,nva,opt,Qspec,hLive,Qprod)}
\NormalTok{  and ParamCorr(Qspec,pispec,Qprod,piprod)}
\NormalTok{  and AdmissibleExec(MM,C,nva,Qspec,pispec,Qprod,piprod,G).}
\end{Highlighting}
\end{Shaded}

The semantic result equality is deliberately not part of
\texttt{PlanAdequacy}. It is the conclusion of the separate
production-refinement theorem below. This separation prevents the
adequacy predicate from assuming the exact observation that the
compiler-correctness argument is required to derive.

Define the constructor-universal obligation:

\begin{Shaded}
\begin{Highlighting}[]
\NormalTok{ProdPlanSim\_sound(MM,C,nva,opt,Qspec,pispec,Qprod,piprod,G)}
\NormalTok{iff}
\NormalTok{  PlanAdequacy(MM,C,nva,opt,Qprod,piprod,G)}
\NormalTok{  implies returnedIds(Qprod,G,piprod)}
\NormalTok{      = returnedIds(Qspec,G,pispec).}
\end{Highlighting}
\end{Shaded}

\texttt{ProdPlanSim\_sound} is not a definitional expansion of
\texttt{PlanAdequacy}. Its proof must proceed by induction over the
certified \texttt{ProdPlanSim} constructor rules, using the local typed
realization judgment, \texttt{hLive}, \texttt{ParamCorr}, and the
wrapper/no-ghost conditions.

The induction domain is now fixed explicitly rather than being left as
an informal phrase. Let \texttt{K\_CQM} be the certified constructors of
the production plan (the general aggregate constructor is excluded by
\texttt{CertifiedCQM}):

\begin{Shaded}
\begin{Highlighting}[]
\NormalTok{K\_CQM = \{}
\NormalTok{  VariablePlan, LiteralPlan, SetLiteralPlan, NotPlan, BinaryPlan,}
\NormalTok{  IfPlan, LetPlan, AttributeAccessPlan, NavigationAccessPlan,}
\NormalTok{  MethodCallPlan, CollectionOperationPlan, IteratorOperationPlan,}
\NormalTok{  ExistsSubqueryPlan, NotExistsSubqueryPlan,}
\NormalTok{  CountSubqueryComparisonPlan, NavigationUniquenessPlan}
\NormalTok{\}}
\end{Highlighting}
\end{Shaded}

For each \texttt{K} in \texttt{K\_CQM}, the proof obligation is one
local constructor rule:

\begin{Shaded}
\begin{Highlighting}[]
\NormalTok{PS{-}K:}
\NormalTok{  WF\_CQM(q\_K) and Ref\_K(v\_K,q\_K) and}
\NormalTok{  (forall child i. ProdPlanSim(child\_i, child\textquotesingle{}\_i)) and}
\NormalTok{  ParamCorr\_K and ScopeCorr\_K and AliasFresh\_K}
\NormalTok{  {-}{-}{-}{-}{-}{-}{-}{-}{-}{-}{-}{-}{-}{-}{-}{-}{-}{-}{-}{-}{-}{-}{-}{-}{-}{-}{-}{-}{-}{-}{-}{-}{-}{-}{-}{-}{-}{-}{-}{-}{-}{-}{-}{-}{-}{-}{-}{-}}
\NormalTok{  ProdPlanSim(K(v\_i), q\_K)}

\NormalTok{SIM{-}K:}
\NormalTok{  ProdPlanSim(K(v\_i), q\_K) and}
\NormalTok{  (forall child i. EvalSim(child\_i, child\textquotesingle{}\_i)) and}
\NormalTok{  PrimitiveAgreement\_K and}
\NormalTok{  CollectionBoundary\_K and NoGhost\_K}
\NormalTok{  {-}{-}{-}{-}{-}{-}{-}{-}{-}{-}{-}{-}{-}{-}{-}{-}{-}{-}{-}{-}{-}{-}{-}{-}{-}{-}{-}{-}{-}{-}{-}{-}{-}{-}{-}{-}{-}{-}{-}{-}{-}{-}{-}{-}{-}{-}{-}{-}}
\NormalTok{  Eval\_G(q\_K) = Eval\_OCL(v\_K)}
\end{Highlighting}
\end{Shaded}

\texttt{Ref\_K} checks the complete payload relation (field values,
resolved references, multiplicities and explicit erasures);
\texttt{ScopeCorr\_K} checks free and bound variables;
\texttt{AliasFresh\_K} prevents capture; and
\texttt{CollectionBoundary\_K} is required for
\texttt{sourceCollectionType}, iterator sources and bottom-to-empty set
observation. The base cases \texttt{VariablePlan} and
\texttt{LiteralPlan} discharge without recursive premises. The recursive
cases use the induction hypotheses for every child and list element; the
subquery cases additionally use the \texttt{NavigationMatchPlan} rule
and its owner/target alias correspondence. The wrapper rule is separate
and proves that only \texttt{ObjectId} values are returned.

The constructor induction theorem is therefore the following conditional
contract (not a claim that the current Java implementation has already
discharged it):

\begin{Shaded}
\begin{Highlighting}[]
\NormalTok{LOCAL\_SIM(K) for every K in K\_CQM}
\NormalTok{and all WF/Ref/Param/Scope/Alias/Primitive/Boundary/NoGhost premises}
\NormalTok{{-}{-}{-}{-}{-}{-}{-}{-}{-}{-}{-}{-}{-}{-}{-}{-}{-}{-}{-}{-}{-}{-}{-}{-}{-}{-}{-}{-}{-}{-}{-}{-}{-}{-}{-}{-}{-}{-}{-}{-}{-}{-}{-}{-}{-}{-}{-}{-}{-}{-}{-}{-}{-}{-}{-}{-}{-}{-}{-}{-}{-}{-}{-}}
\NormalTok{ProdPlanSim\_sound(MM,C,nva,opt,Qspec,pispec,Qprod,piprod,G)}
\end{Highlighting}
\end{Shaded}

The checked \texttt{cqm\_ast\_lowering\_rules.csv} is an executable
structural witness for the next lowering boundary, but it does not
discharge \texttt{PrimitiveAgreement} or \texttt{EvalSim}. Those remain
explicit proof obligations in PO-18 until a constructor-wise semantic
proof or a separately verified backend evaluator is added.

\texttt{hLive} relates only aliases visible at corresponding constructor
interfaces (free-variable imports, exported set columns, scalar results,
and \texttt{self} at the wrapper). Private fresh aliases are quantified
locally by the corresponding \texttt{ProdPlanSim} rule; different
correct raw shapes therefore need not have the same number of helper
aliases. \texttt{ParamCorr} ranges only over \texttt{Params(Qspec)} and
\texttt{Params(Qprod)}, requires every syntactically used lookup to be
defined, and checks the exact typed literal/metadata value required by
the local simulation rule. It ignores unused entries in either parameter
map and allows shape-specific private parameters only when their values
are justified by that rule. In particular, eagerly reserved but unused
codec entries in \texttt{pispec} do not need production counterparts.

For a hypothetical Java-production theorem, the arguments would be fixed
by \texttt{opt=T\_OPT(va)},
\texttt{(tprod,piprod)=T\_TEXT\^{}prod(C,R,opt)}, and
\texttt{Qprod=parse\_Cypher(tprod)}, together with the still-optional
premise \texttt{OptRefines(opt,T\_NORM(va))}. This paragraph is not used
by the reference theorem. Consequently the implementation predicate has
no free \texttt{MM}, resolver, source, or optimizer variables; changing
any input requires a new adequacy derivation.

\texttt{PlanAdequacy} is therefore graph-specific and contains no free
execution variable. A future universal production theorem must prove
this predicate for every \texttt{G} satisfying \texttt{AdmissibleExec};
a finite test suite establishes it only for the executions actually
checked.

\texttt{ProdPlanSim} names the required constructor-directed certificate
over the sealed optimized IR and production plans. A complete discharge
must define an inductive rule for every sealed constructor and prove
\texttt{ProdPlanSim\_sound}; that universal certificate is PO-18 and is
not yet available. The intended rules relate each production constructor
to the corresponding formal constructor, preserves static type, resolved
metadata, lexical scope, live imports, collection boundaries, parameter
values, and the invariant wrapper, and requires both lowerings to
satisfy the same local typed realization judgment. It deliberately
permits different \texttt{CASE}, list-comprehension, re-identification,
or correlated-subquery shapes. The parser round trip is required only
between \texttt{tprod} and the production plan\textquotesingle s own raw
AST \texttt{Qprod}; it is not required to be alpha-equivalent to
\texttt{Qspec}. Thus \texttt{PlanAdequacy} records the structural
bridge, while \texttt{ProdPlanSim\_sound} remains the explicit universal
\texttt{T\_OPT}-to-\texttt{T\_NORM} refinement obligation. Once
\texttt{ProdPlanSim\_sound} is discharged, it yields:

\begin{Shaded}
\begin{Highlighting}[]
\NormalTok{PlanAdequacy(MM,C,nva,opt,Qprod,piprod,G)}
\NormalTok{implies}
\NormalTok{returnedIds(Qprod,G,piprod)=returnedIds(Qspec,G,pispec).}
\end{Highlighting}
\end{Shaded}

The core parametric constructor induction is mechanized as
\texttt{JavaIrRefinement.prod\_plan\_sim\_sound}; it proves composition
under the local \texttt{AlgebraAgreement} premise for the complete
constructor family. This is not yet an instantiation of
\texttt{ProdPlanSim} with the Java CQM planner, production raw-AST
lowering, Neo4j evaluator, \texttt{hLive}, \texttt{ParamCorr}, and C6
wrapper. That instantiation is the remaining PO-18 proof required for a
universal production claim. PA14-style finite tree/parser tests are
evidence for selected inputs, not that universal proof. Until the
adapter-specific premises are discharged, Theorems 5-\/-6 consume
\texttt{PlanAdequacy} only as the explicit structural bridge premise and
must list \texttt{ProdPlanSim\_sound} as a separate conditional premise.

\textbf{TXT5 Raw-AST Closure and Totality.} Every admitted normalized
invariant has one \texttt{BuildCQInvariant} result, one total
\texttt{ExpandI} result up to fresh-name alpha-equivalence, and one
parameterized raw Cypher text. This part follows from TXT1 and
simultaneous structural/well-founded induction on the equations. Define
the separate selected-dialect predicate:

\begin{Shaded}
\begin{Highlighting}[]
\NormalTok{ParserRoundTrip(Q)}
\NormalTok{iff parse\_Cypher(renderRaw(Q)) is defined}
\NormalTok{    and alpha{-}equivalent to Q.}
\end{Highlighting}
\end{Shaded}

TXT5 does not prove \texttt{ParserRoundTrip(Q)} universally from printer
totality. A runtime/text theorem must assume it for the concrete
\texttt{Q} (A6 does so for the production query). The 52 generated
parser cases are finite evidence, not a proof for every raw AST produced
by every future \texttt{Expand} derivation.

\begin{center}\rule{0.5\linewidth}{0.5pt}\end{center}

\section{7. Theorem 0: Object-Graph
Representation}\label{7-theorem-0-object-graph-representation}

\subsection{Statement}\label{statement}

For every valid UML metamodel/object-model pair \texttt{(MM,M)}, if
\texttt{GraphAdequate(MM,M,G)} holds (in particular, when an encoder
\texttt{Phi} has been proved to establish it), then \texttt{G} preserves
all validation-relevant observations used by \texttt{OCL\_val}:

\begin{Shaded}
\begin{Highlighting}[]
\NormalTok{object identity}
\NormalTok{class conformance}
\NormalTok{attribute values}
\NormalTok{binary association links}
\NormalTok{navigation results}
\NormalTok{allInstances results}
\end{Highlighting}
\end{Shaded}

Formally:

\begin{Shaded}
\begin{Highlighting}[]
\NormalTok{lift\_node(Obs\_obj(MM,M)) = Obs\_graph(MM,G).}
\end{Highlighting}
\end{Shaded}

Equivalently, \texttt{(MM,M)\ ==\_val\ G}. This is equality of the
observation structure from Section 4.7, not syntactic equality of model
files, comments, diagram layout, or tool metadata.

\subsection{Required Lemmas}\label{required-lemmas}

\begin{Shaded}
\begin{Highlighting}[]
\NormalTok{R1 Object Injectivity}
\NormalTok{R2 Object{-}Node Exactness}
\NormalTok{R3 Type Preservation}
\NormalTok{R4 Attribute Preservation}
\NormalTok{R5 Association Preservation}
\NormalTok{R6 Navigation Preservation}
\NormalTok{R7 allInstances Preservation}
\end{Highlighting}
\end{Shaded}

\subsection{Proof}\label{proof}

Define \texttt{Psi\_val} by reading \texttt{Obs\_graph} through the
canonical accessors:

\begin{Shaded}
\begin{Highlighting}[]
\NormalTok{objects       from object nodes and use\_id}
\NormalTok{types         from ObjectInstanceOf edges}
\NormalTok{attributes    from value\_G}
\NormalTok{links         from Link* relationships through linkAssociation\_G,}
\NormalTok{              linkSourceRole\_G, linkTargetRole\_G, and linkQualifiers\_G}
\NormalTok{navigation    from nav\_graph}
\NormalTok{allInstances  from ObjectInstanceOf membership to classNode(C)}
\end{Highlighting}
\end{Shaded}

Object-node preservation and reflection follow from R2: \texttt{node} is
a bijection between source objects and validation-visible object nodes.
Object injectivity also follows from R1 and injectivity of \texttt{id}.
Thus the object component read by \texttt{Psi\_val} contains neither
missing nor spurious objects.

Type preservation follows directly from materialized type membership.
The forward direction uses the rule that every conformance fact creates
an \texttt{ObjectInstanceOf} edge. The backward direction uses
no-spurious-type-edge.

Attribute preservation follows from R4, i.e.,
\texttt{I\_attr\_graph(node(o),a)=attrVal\_M(o,a)}.

Association preservation follows from R5 and is proved by two
inclusions. If a UML link exists, \texttt{GraphAdequate} guarantees at
least one \texttt{Link*} relationship with the same exact
\texttt{associationKey}, roles, endpoints, and direction-indexed encoded
\texttt{sourceQualifiers}/\texttt{targetQualifiers}. Conversely, any
relevant \texttt{Link*} relationship corresponds to a UML association
link by no-spurious-link. Relationship uniqueness is not required;
object injectivity maps endpoints back uniquely and later set
realization removes duplicates.

Navigation preservation follows from R6 as a corollary of association
preservation. \texttt{MANY} is equality of target sets after applying
\texttt{node}, including the no-match bottom-qualifier case. For
\texttt{ONE}, valid multiplicity restricts that set to empty or
singleton, and R6 proves that \texttt{oneOrBottom} commutes with the
object-node encoding. Thus neither result kind is silently coerced into
the other.

\texttt{allInstances} preservation follows from R7 and type
preservation:

\begin{Shaded}
\begin{Highlighting}[]
\NormalTok{o in allInstances\_obj(C)}
\NormalTok{iff o conformsTo C}
\NormalTok{iff (node(o)){-}[:ObjectInstanceOf]{-}\textgreater{}(classNode(C))}
\NormalTok{iff node(o) in allInstances\_graph(C).}
\end{Highlighting}
\end{Shaded}

Each component of \texttt{lift\_node(Obs\_obj(MM,M))} therefore equals
its corresponding component of \texttt{Obs\_graph(MM,G)}. By extensional
tuple equality, the two observation structures are equal, so the given
\texttt{G} is adequate for evaluating the supported validation fragment.
As a separate corollary, if \texttt{G=Phi(MM,M)} and \texttt{Phi} has
been proved to establish \texttt{GraphAdequate}, then that particular
encoder output is adequate. No property of an arbitrary \texttt{Phi} is
derived from adequacy of an unrelated \texttt{G}.

\begin{center}\rule{0.5\linewidth}{0.5pt}\end{center}

\section{8. Theorem 1: Binder
Soundness}\label{8-theorem-1-binder-soundness}

\subsection{Statement}\label{statement-1}

If:

\begin{Shaded}
\begin{Highlighting}[]
\NormalTok{ParseInvariant(tInv)=astInv}
\NormalTok{decodeInvariant\_val(astInv)=I=(cName,R,e)}
\NormalTok{resolveClass\_MM(cName)=C}
\NormalTok{T\_BIND\_INV(astInv,MM)=BInvariant(C,R,b)}
\NormalTok{Adm\_MM(astInv)}
\end{Highlighting}
\end{Shaded}

with \texttt{e\ in\ OCL\_val}, \texttt{b\ in\ BoundOCL\_val(MM)}, and
\texttt{MM;\{self:C\}\ \textbar{}-\ b:Boolean}, then the equations below
hold for every valid model \texttt{M} conforming to \texttt{MM} and
every \texttt{rho} such that \texttt{rho\ \textbar{}=\_M\ \{self:C\}}
and \texttt{EvalClosed\_obj(e,rho)}. For invariant-root environments
this closure premise is supplied by A8; it is necessary because the
canonical Integer-to-Real embedding is partial outside the exact
binary64 range.

Let \texttt{tag\_S(rho)} and \texttt{tag\_B(rho)} be the corresponding
source and bound tagged environments from Section 2.5. Then:

\begin{Shaded}
\begin{Highlighting}[]
\NormalTok{[[e]]\_S(M,tag\_S(rho)) \textasciitilde{}=SB [[b]]\_B(M,tag\_B(rho)).}
\end{Highlighting}
\end{Shaded}

Consequently, by SB-Erase and the derived public denotations:

\begin{Shaded}
\begin{Highlighting}[]
\NormalTok{erase\_S([[e]]\_S(M,tag\_S(rho)))}
\NormalTok{= erase\_B([[b]]\_B(M,tag\_B(rho))),}

\NormalTok{that is,}

\NormalTok{[[e]]\_OCLval(M,rho) = [[b]]\_Bound(M,rho).}
\end{Highlighting}
\end{Shaded}

No denotation is assigned to \texttt{tInv} or \texttt{astInv}.

\subsection{Required Lemmas}\label{required-lemmas-1}

\begin{Shaded}
\begin{Highlighting}[]
\NormalTok{M2 Bound Subexpression Closure and Structural Induction}
\NormalTok{M3 Canonical Type Uniqueness}
\NormalTok{M3a Set Join Functionality}
\NormalTok{M4 Admission and Binding Soundness}
\NormalTok{B0 Binding Functionality}
\NormalTok{B1 Name Resolution Soundness}
\NormalTok{B2 Type Assignment Soundness}
\NormalTok{B3 Navigation Resolution Determinism}
\NormalTok{B4 Lexical Scope and Shadowing Preservation}
\NormalTok{SB0 Independent{-}Operator Correspondence}
\end{Highlighting}
\end{Shaded}

\subsection{Proof}\label{proof-1}

M4 places the result in the finite typed theorem domain, M2 justifies
the induction, M3 fixes the canonical type used by each local equation,
and M3a fixes the element type of set literals, union, and intersection.
The proof is by structural induction on the unique successful binding
derivation given by B0. The complete case split is the table in Section
2.5: each binder rule has exactly one denotational row, so no admitted
source constructor is omitted.

Base cases are direct under \texttt{\textasciitilde{}=SB}. \texttt{self}
maps to \texttt{BoundSelf}, variables map to related tagged environment
bindings, and literals map to corresponding tagged scalar values.

For attribute and navigation property calls, the induction hypothesis
preserves the receiver. Name resolution identifies whether the property
denotes an attribute or an association role. Type assignment ensures
applicability. For a navigation, B3 supplies the unique association,
roles, direction, result kind, and qualifier types used by both
interpretations. \texttt{MANY} follows by finite-set extensionality;
\texttt{ONE} additionally applies the corresponding empty/singleton
\texttt{oneOrBottom} equation.

For Boolean, comparison, arithmetic, and type operations, the
syntax-directed binding rule preserves the operator and B2 supplies the
exact typing premise from Section 2.3. The induction hypotheses relate
operands; SB0 then relates the independently defined source and bound
operator results. The proof does not identify the two evaluators or
appeal to VA semantics.

For \texttt{if}, first apply the induction hypothesis to the condition.
SB0 makes the two independent truth tests select the same branch. Parent
\texttt{EvalClosed\_obj} then supplies closure for that selected branch,
so the branch induction hypothesis is applied only there; no closure or
denotation premise is required for the unselected branch.

For a typed or inferred \texttt{let}, apply the initializer hypothesis
and the componentwise embedding-correspondence case of SB0. The
\texttt{declType} premise ensures that both environments extend
\texttt{x} at the same resolved \texttt{tauD}; then apply the body
hypothesis. Thus the annotation changes static checking but does not
introduce a new runtime computation beyond the canonical embedding.

For iterator expressions, the induction hypothesis preserves the source
set. Iterator scope preservation B4 states that the body is evaluated
under corresponding extended environments
\texttt{rho{[}x\ -\textgreater{}\ iota{[}tauE-\textgreater{}tauD{]}(s){]}}.
The iterator typing premise and B2 imply the extension satisfies
\texttt{Gamma{[}x:tauD{]}}. Applying the induction hypothesis to the
body yields equal predicate/body results for each element. Set literals,
the certified scalar-to-set view, \texttt{isUnique}, \texttt{union},
\texttt{intersection}, and \texttt{asSet} are separate induction cases
in Section 2.5; SB0 discharges their typed lifting and extensional
finite-set equations.

Unsupported source constructs have no \texttt{decodeInvariant\_val}
derivation, while ambiguous or ill-typed admitted constructs are
rejected by \texttt{T\_BIND\_INV}; neither enters the successful theorem
case. F1 separately establishes that every theorem-covered
\texttt{OCL\_val} expression has a corresponding \texttt{AST\_val}
representation; it is a representability/domain lemma, not a premise
needed by the structural induction itself.

\begin{center}\rule{0.5\linewidth}{0.5pt}\end{center}

\section{9. Theorem 2: Validation Algebra
Abstraction}\label{9-theorem-2-validation-algebra-abstraction}

\subsection{Statement}\label{statement-2}

If \texttt{MM;Gamma\ \textbar{}-\ b:tau}, \texttt{T\_VA(b)=va},
\texttt{rho\ \textbar{}=\_M\ Gamma}, and
\texttt{EvalClosed\_obj(va,rho)}, then:

\begin{Shaded}
\begin{Highlighting}[]
\NormalTok{MM;Gamma |{-} va:tau}

\NormalTok{erase\_B([[b]]\_B(M,tag\_B(rho))) = [[va]]\_obj\^{}tau(M,rho).}

\NormalTok{Equivalently, by the derived notation:}

\NormalTok{[[b]]\_Bound(M,rho) = [[va]]\_obj(M,rho).}
\end{Highlighting}
\end{Shaded}

\subsection{Required Lemmas}\label{required-lemmas-2}

\begin{Shaded}
\begin{Highlighting}[]
\NormalTok{M2 Bound Subexpression Closure and Structural Induction}
\NormalTok{M3 Canonical Type Uniqueness}
\NormalTok{M3a Set Join Functionality}
\NormalTok{VA1 Bound{-}to{-}VA Type Preservation}
\NormalTok{VA2 Local Bound{-}to{-}VA Simulation}
\NormalTok{VA3 Environment Lookup Preservation}
\NormalTok{BV0 Bound{-}Primitive/VA Commutation}
\NormalTok{B4 Lexical Scope and Shadowing Preservation}
\end{Highlighting}
\end{Shaded}

\subsection{Proof}\label{proof-2}

M2 justifies structural induction over the finite Bound OCL expression,
and M3 fixes the canonical type at each constructor.

Base cases are immediate:

\begin{Shaded}
\begin{Highlighting}[]
\NormalTok{BoundSelf       {-}\textgreater{} Self}
\NormalTok{BoundVariable   {-}\textgreater{} Var}
\NormalTok{BoundLiteral    {-}\textgreater{} Literal}
\end{Highlighting}
\end{Shaded}

For \texttt{BoundVariable} and \texttt{BoundSelf}, VA3 ensures that
erasing the Bound-tagged environment value yields the corresponding VA
environment value. This property is preserved by iterator and
\texttt{let} environment extension, including shadowing.

For attributes and navigations, the induction hypothesis preserves the
receiver. The VA2 translation equation copies the resolved metadata and
result kind unchanged. BV0 gives the \texttt{NavigationMany} equation;
for \texttt{NavigationOne} it also commutes \texttt{oneOrBottom}. A
certified \texttt{ViewSet} uses its two bottom/entity equations.

For Boolean, comparison, arithmetic, and type operations, apply the
reached child induction hypotheses and BV0 at the primitive parent. For
\texttt{if}, first relate the condition, use BV0 to select the same
branch, derive local closure for that selected branch from parent
\texttt{EvalClosed\_obj}, and apply the induction hypothesis only to it.
For \texttt{let}, relate the initializer and then apply the body
hypothesis in the reached extended environment supplied by VA3. This is
not an assumption of whole-expression equality: each local result
follows by unfolding one pair of constructor definitions. No graph
encoding or \texttt{encodeValue} operation is involved; both sides use
the object interpretation.

For iterators, the induction hypothesis preserves the source finite set.
VA3 and B4 give identical stack extension for each element; the body
induction hypothesis therefore applies pointwise. Unfolding the local
VA2 equation yields the same existential, universal, filtering, or
finite-set image result. \texttt{SetLiteral}, \texttt{union},
\texttt{intersection}, and \texttt{asSet} use M3a, BV0, and typed
finite-set extensionality. \texttt{IsUnique} uses the body induction
hypothesis for each pair of source elements. Other collection operations
follow directly from their VA2 equation and the child induction
hypotheses. These cases exhaust the Bound syntax.

Type preservation follows compositionally from the typing rules of Bound
OCL and VA.

\begin{center}\rule{0.5\linewidth}{0.5pt}\end{center}

\section{10. Theorem 3: Normalization
Preservation}\label{10-theorem-3-normalization-preservation}

\subsection{Statement}\label{statement-3}

Let:

\begin{Shaded}
\begin{Highlighting}[]
\NormalTok{MM;Gamma |{-} va:tau}
\NormalTok{rho |=\_I Gamma}
\NormalTok{EvalClosed\_I(va,rho)}
\NormalTok{nva = T\_NORM(va)}
\end{Highlighting}
\end{Shaded}

where \texttt{T\_NORM} uses the deterministic bottom-up relation
\texttt{-\textgreater{}\_bu}. Then:

\begin{Shaded}
\begin{Highlighting}[]
\NormalTok{1. T\_NORM terminates independently of rho.}
\NormalTok{2. MM;Gamma |{-} nva:tau.}
\NormalTok{3. NF\_R(nva), modulo alpha{-}equivalence.}
\NormalTok{4. [[va]]\_I(rho) = [[nva]]\_I(rho).}
\NormalTok{5. EvalClosed\_I(nva,rho).}
\end{Highlighting}
\end{Shaded}

Clauses 4-\/-5 hold for \texttt{I=obj,graph} whenever the corresponding
local closure premise holds. For an admitted invariant, A8 supplies it
for every reachable object-side environment.

\subsection{Rewrite Rules}\label{rewrite-rules}

Let \texttt{z=fresh(S,y,Gamma)} for membership rules and
\texttt{u=fresh(S,Gamma)} for binder-free emptiness/count rules. The
reference normalizer uses:

\begin{Shaded}
\begin{Highlighting}[]
\NormalTok{Size(S)                 {-}\textgreater{} Count(S)}
\NormalTok{ForAll(S,x,P)       {-}\textgreater{} Not(Exists(S,x,Not(P)))}
\NormalTok{NotEmpty(S)         {-}\textgreater{} Exists(S,u,true)}
\NormalTok{IsEmpty(S)          {-}\textgreater{} Not(Exists(S,u,true))}
\NormalTok{Reject(S,x,P)       {-}\textgreater{} Select(S,x,Not(P))}
\NormalTok{Implies(A,B)        {-}\textgreater{} Or(Not(A),B)}
\NormalTok{Xor(A,B)            {-}\textgreater{} Or(And(A,Not(B)),And(Not(A),B))}

\NormalTok{Includes(S:Set(sigma),y:delta)}
\NormalTok{  {-}\textgreater{} Exists(S,z,}
\NormalTok{       Compare(EQ,}
\NormalTok{         Coerce[sigma{-}\textgreater{}upsilon](Var(z)),}
\NormalTok{         Coerce[delta{-}\textgreater{}upsilon](y)))}
\NormalTok{  where memberJoin(sigma,delta)=upsilon}

\NormalTok{Excludes(S:Set(sigma),y:delta)}
\NormalTok{  {-}\textgreater{} Not(Exists(S,z,}
\NormalTok{       Compare(EQ,}
\NormalTok{         Coerce[sigma{-}\textgreater{}upsilon](Var(z)),}
\NormalTok{         Coerce[delta{-}\textgreater{}upsilon](y))))}
\NormalTok{  where memberJoin(sigma,delta)=upsilon}

\NormalTok{Count(NavigationMany(...)) \textgreater{} 0 {-}\textgreater{} Exists(NavigationMany(...),u,true)}
\NormalTok{Count(NavigationMany(...)) = 0 {-}\textgreater{} Not(Exists(NavigationMany(...),u,true))}
\end{Highlighting}
\end{Shaded}

\texttt{Coerce{[}Void-\textgreater{}upsilon{]}(bottom\_Void)=bottom\_upsilon};
class upcasts preserve identity. The typing derivation emits
\texttt{Coerce{[}Integer-\textgreater{}Real{]}} independently of runtime
values. Exactness is required only by \texttt{EvalClosed\_I} when N1
evaluates that node. A rule whose required static join is undefined has
no certified normalization derivation; a reached inexact conversion
falsifies N1\textquotesingle s closure premise but does not change
\texttt{T\_NORM} or N3/N4.

Implementation boundary: this list defines the logical reference
transform \texttt{T\_NORM}; it is not a claim that Java must construct
the same syntax tree. \texttt{OclIrOptimizer} implements a
graph-oriented \texttt{T\_OPT}: it expands \texttt{Implies} and
\texttt{Xor} directly, specializes navigation
\texttt{ForAll}/emptiness/filter/size cases to
\texttt{NavigationPredicateCheck} or \texttt{NavigationCountComparison},
and can retain general collection operations for the planner. Applying
Theorem 3 to that implementation therefore requires semantic refinement
\texttt{T\_OPT(va)\ \textasciitilde{}=\ T\_NORM(va)}, not syntactic
equality. The finite rewrite and differential suites are scoped evidence
for this refinement. Iterator fusion now additionally checks
\texttt{canRenameWithoutCapture}; an unsafe target binder causes a
conservative fallback to the general iterator IR. The Lean kernel proves
type-index preservation for the four abstract
\texttt{Bool}/\texttt{Nat}/\texttt{Elem}/\texttt{Set} indices and a de
Bruijn binder-boundary renaming theorem. It now also proves lexical
named-to-de-Bruijn semantic correspondence and soundness/semantic
preservation of the exact three-disjunct
\texttt{canRenameWithoutCapture} guard over the Boolean binder kernel.
Metamodel-specific OCL subtyping/coercion, exhaustive
production-\texttt{OclIr} constructor mapping to that kernel, and
universal \texttt{T\_OPT} semantic refinement remain open.

\subsection{Normal-Form Predicate}\label{normal-form-predicate}

Let \texttt{Redex\_R(t)} hold when \texttt{t} matches the left-hand side
of one of the stated normalization rules, including \texttt{Size}
canonicalization and the two count-comparison patterns. Define:

\begin{Shaded}
\begin{Highlighting}[]
\NormalTok{NF\_R(e)}
\NormalTok{iff}
\NormalTok{for every subterm t of e: not Redex\_R(t).}
\end{Highlighting}
\end{Shaded}

This is a relative normal form for the finite rule set \texttt{R}; it
does not claim a canonical form for arbitrary OCL. Fresh-variable side
conditions are part of the \texttt{Includes}, \texttt{Excludes},
emptiness, and count macro-rules. Rewriting is capture-avoiding and
expressions are identified up to alpha-equivalence.

\subsection{Measure}\label{measure}

Use the N3 measure:

\begin{Shaded}
\begin{Highlighting}[]
\NormalTok{mu(e) =}
\NormalTok{  \#ForAll}
\NormalTok{  + \#Implies}
\NormalTok{  + \#Xor}
\NormalTok{  + \#Reject}
\NormalTok{  + \#IsEmpty}
\NormalTok{  + \#NotEmpty}
\NormalTok{  + \#Includes}
\NormalTok{  + \#Excludes}
\NormalTok{  + \#CountCompareOnNavigation}

\NormalTok{measure(e) = (\#Size(e), mu(e)), ordered lexicographically.}
\end{Highlighting}
\end{Shaded}

The bottom-up enabling condition is part of the rewrite relation, not
merely an implementation convention. It is required for \texttt{Xor}:
its expansion duplicates both operands, so \texttt{mu} decreases only
because those operands are already in \texttt{NF\_R} and contain no
counted redex. \texttt{Excludes} is one macro-step, not a first step
that introduces \texttt{Includes}. N3 therefore proves that every
\texttt{-\textgreater{}\_bu} step strictly decreases \texttt{measure}.

\subsection{Per-Rule Proof-Obligation
Table}\label{per-rule-proof-obligation-table}

Let \texttt{D\_I(e,rho)={[}{[}e{]}{]}\_I(rho)}. The following table
discharges N1-\/-N3 locally. All semantic equations hold for both
\texttt{I=obj} and \texttt{I=graph}. Child expressions are assumed well
typed, as provided by the enclosing derivation.

{\def\LTcaptype{none} % do not increment counter
\begin{longtable}[]{@{}lllll@{}}
\toprule\noalign{}
Rule & Typing and freshness premises & Semantic equality & Type
preservation & Decrease/no regeneration \\
\midrule\noalign{}
\endhead
\bottomrule\noalign{}
\endlastfoot
N-Size & \texttt{S:Set(tau)} &
\texttt{D\_I(Size(S),rho)=card(finiteSet(D\_I(S,rho)))=D\_I(Count(S),rho)}
& \texttt{Integer\ -\textgreater{}\ Integer} & \texttt{\#Size} decreases
by 1; target contains no \texttt{Size} \\
N-ForAll & \texttt{S:Set(tau)}; \texttt{P:Boolean} under
\texttt{Gamma{[}x:tau{]}} &
\texttt{for\ all\ s\ in\ finiteSet(S):\ bool\_val(P\_s)=true} iff no
\texttt{s} makes \texttt{Not(P\_s)} validation-true &
\texttt{Boolean\ -\textgreater{}\ Boolean} & \texttt{\#ForAll} decreases
by 1; target contains no counted source constructor \\
N-NotEmpty & \texttt{S:Set(tau)}; \texttt{\_x} fresh &
\texttt{finiteSet(S)\ !=\ empty} iff
\texttt{exists\ s\ in\ finiteSet(S):\ true} &
\texttt{Boolean\ -\textgreater{}\ Boolean} & \texttt{\#NotEmpty}
decreases by 1; fresh binder prevents capture \\
N-IsEmpty & \texttt{S:Set(tau)}; \texttt{\_x} fresh &
\texttt{finiteSet(S)=empty} iff no \texttt{s} witnesses true &
\texttt{Boolean\ -\textgreater{}\ Boolean} & \texttt{\#IsEmpty}
decreases by 1; no source pattern regenerated \\
N-Reject & \texttt{S:Set(tau)}; \texttt{P:Boolean} under
\texttt{Gamma{[}x:tau{]}} & the subset with
\texttt{bool\_val(P\_s)=false} equals the subset with
\texttt{bool\_val(Not(P\_s))=true} &
\texttt{Set(tau)\ -\textgreater{}\ Set(tau)} & \texttt{\#Reject}
decreases by 1 \\
N-Implies & \texttt{A:Boolean}; \texttt{B:Boolean} & both sides are true
iff \texttt{bool\_val(A)=false} or \texttt{bool\_val(B)=true} &
\texttt{Boolean\ -\textgreater{}\ Boolean} & \texttt{\#Implies}
decreases by 1 \\
N-Xor & \texttt{A:Boolean}; \texttt{B:Boolean}; both operands in
\texttt{NF\_R} & both sides are true iff exactly one operand is
validation-true & \texttt{Boolean\ -\textgreater{}\ Boolean} &
duplicated operands have no counted redex; \texttt{\#Xor} decreases by
1 \\
N-Includes & \texttt{S:Set(sigma)}; \texttt{y:delta};
\texttt{memberJoin(sigma,delta)=upsilon}; \texttt{z=fresh(S,y,Gamma)};
N1 additionally assumes \texttt{EvalClosed\_I} for reached coercions &
lifted membership equals existence of \texttt{s} satisfying typed
equality after both embeddings to \texttt{upsilon} &
\texttt{Boolean\ -\textgreater{}\ Boolean}; both comparison operands
have type \texttt{upsilon} & \texttt{\#Includes} decreases by 1;
coercions introduce no redex \\
N-Excludes & same join/coercion/freshness premises & lifted
non-membership equals absence of the same typed-equality witness &
\texttt{Boolean\ -\textgreater{}\ Boolean} & target contains neither
\texttt{Excludes} nor \texttt{Includes}; \texttt{mu} decreases by 1 \\
N-CountPos & \texttt{N=Navigation(...):Set(E)}; \texttt{\_x} fresh &
\texttt{card(finiteSet(D\_I(N,rho)))\textgreater{}0} iff an element
witnesses true & \texttt{Boolean\ -\textgreater{}\ Boolean} & one
\texttt{CountCompareOnNavigation} is removed \\
N-CountZero & same & \texttt{card(finiteSet(D\_I(N,rho)))=0} iff no
element witnesses true & \texttt{Boolean\ -\textgreater{}\ Boolean} &
one \texttt{CountCompareOnNavigation} is removed \\
\end{longtable}
}

For N-ForAll, N-Reject, N-Includes, and N-Excludes, the target binder is
either the original binder or a fresh name. The capture-avoiding
substitution lemma is:

\begin{Shaded}
\begin{Highlighting}[]
\NormalTok{if y is fresh for e, then}
\NormalTok{[[e[x:=y]]] \_I (rho[y{-}\textgreater{}v]) = [[e]]\_I(rho[x{-}\textgreater{}v]).}
\end{Highlighting}
\end{Shaded}

It is proved by structural induction on \texttt{e} and justifies
alpha-renaming in the table. Only \texttt{Xor} duplicates subterms, and
its bottom-up enabling premise guarantees that neither duplicate
contains a redex counted by \texttt{mu}.

\subsection{Proof}\label{proof-3}

Semantic preservation is by the equations in the proof-obligation table.
The following paragraphs expand the nontrivial cases.

\texttt{Size(S)\ -\textgreater{}\ Count(S)}: both constructors denote
\texttt{card(finiteSet({[}{[}S{]}{]}\_I(rho)))}; the rewrite only
selects the canonical VA cardinality constructor.

\texttt{ForAll(S,x,P)\ -\textgreater{}\ Not(Exists(S,x,Not(P)))}: under
finite-set semantics, \texttt{forAll} is true iff there is no element
for which \texttt{P} is not validation-true. The right-hand side
expresses exactly the absence of such a counterexample. Both sides are
Boolean.

\texttt{NotEmpty(S)\ -\textgreater{}\ Exists(S,\_x,true)}: a finite set
is non-empty iff there exists an element in it. Both sides are Boolean
and \texttt{\_x} is fresh.

\texttt{IsEmpty(S)\ -\textgreater{}\ Not(Exists(S,\_x,true))}: a finite
set is empty iff no element exists in it. Both sides are Boolean.

\texttt{Reject(S,x,P)\ -\textgreater{}\ Select(S,x,Not(P))}: rejection
keeps exactly elements for which the predicate is not validation-true.
Selection with negated predicate returns the same finite set and
preserves element type.

\texttt{Implies(A,B)\ -\textgreater{}\ Or(Not(A),B)}: by
validation-level Boolean semantics, implication is equivalent to
disjunction of the negated antecedent and the consequent. Both sides are
Boolean.

For \texttt{Includes}, let \texttt{S:Set(sigma)}, \texttt{y:delta}, and
\texttt{memberJoin(sigma,delta)=upsilon},
\texttt{SY=upSet\_I{[}sigma-\textgreater{}upsilon{]}(finiteSet\_sigma({[}{[}S{]}{]}\_I(rho)))},
and
\texttt{y\textquotesingle{}=iota\_I{[}delta-\textgreater{}upsilon{]}({[}{[}y{]}{]}\_I(rho))}.
Then:

\begin{Shaded}
\begin{Highlighting}[]
\NormalTok{[[Includes(S,y)]]\_I(rho)=true}
\NormalTok{iff y\textquotesingle{} in SY}
\NormalTok{iff exists s in finiteSet\_sigma([[S]]\_I(rho)):}
\NormalTok{      eq\_val\^{}upsilon(iota\_I[sigma{-}\textgreater{}upsilon](s),y\textquotesingle{})=true}
\NormalTok{iff [[Exists(S,z,}
\NormalTok{      Compare(EQ,}
\NormalTok{        Coerce[sigma{-}\textgreater{}upsilon](Var(z)),}
\NormalTok{        Coerce[delta{-}\textgreater{}upsilon](y)))]]\_I(rho)=true.}
\end{Highlighting}
\end{Shaded}

The middle equivalence is extensional membership at one static type.
\texttt{z} is fresh, so extending the environment cannot capture a free
occurrence in \texttt{y}. \texttt{Excludes} is the Boolean negation of
the same witness equivalence. Neither rule introduces a coercion absent
from the source operation\textquotesingle s typed membership semantics,
so reachable closure is preserved.

\texttt{Count(Nav(...))\ \textgreater{}\ 0\ -\textgreater{}\ Exists(Nav(...),\_x,true)}:
a finite navigation result has cardinality greater than zero iff it
contains at least one element.

\texttt{Count(Nav(...))\ =\ 0\ -\textgreater{}\ Not(Exists(Nav(...),\_x,true))}:
a finite navigation result has cardinality zero iff it contains no
element.

Type preservation follows in each case from the derivations recorded in
the table. Termination follows because each one-step strategy rewrite
decreases the lexicographic \texttt{measure}, expression trees are
finite, and the normalizer does not apply inverse rules. N4 proves
\texttt{NF\_R(nva)} by well-founded/structural induction over the
deterministic bottom-up traversal.

Therefore \texttt{T\_NORM} terminates, preserves type, denotation, and
reachable evaluation closure, and produces \texttt{NF\_R} modulo
alpha-equivalence.

\begin{center}\rule{0.5\linewidth}{0.5pt}\end{center}

\section{11. Theorem 4: Validation
Preservation}\label{11-theorem-4-validation-preservation}

\subsection{Statement}\label{statement-4}

Let:

\begin{Shaded}
\begin{Highlighting}[]
\NormalTok{MM;Gamma |{-} nva:tau}
\NormalTok{NF\_R(nva)}
\NormalTok{rho |=\_M Gamma}
\NormalTok{eta |=\_G Gamma}
\NormalTok{rho \textasciitilde{}Phi\_Gamma eta}
\NormalTok{EvalClosed\_obj(nva,rho).}
\end{Highlighting}
\end{Shaded}

Assume \texttt{GraphAdequate(MM,M,G)}, G1-\/-G4 (including G3a), and
that each reached iterator/\texttt{let} environment inherits the local
closure premise. Then:

\begin{Shaded}
\begin{Highlighting}[]
\NormalTok{encodeValue\_tau([[nva]]\_obj\^{}tau(M,rho))}
\NormalTok{= [[nva]]\_graph\^{}tau(G,eta)}
\NormalTok{and EvalClosed\_graph(nva,eta).}
\end{Highlighting}
\end{Shaded}

For invariant predicates:

\begin{Shaded}
\begin{Highlighting}[]
\NormalTok{bool\_val([[nva]]\_obj(M,self {-}\textgreater{} o))}
\NormalTok{= bool\_val([[nva]]\_graph\^{}Boolean(G,self {-}\textgreater{} node(o))).}
\end{Highlighting}
\end{Shaded}

\subsection{Required Lemmas}\label{required-lemmas-3}

\begin{Shaded}
\begin{Highlighting}[]
\NormalTok{Theorem 0 Object{-}Graph Representation}
\NormalTok{G1 Environment Encoding Preservation}
\NormalTok{G2 Finite{-}Set Homomorphism}
\NormalTok{G3 Typed Equality and Scalar{-}Operator Preservation}
\NormalTok{G3a Canonical Coercion Commutation}
\NormalTok{G4 Validation{-}Policy Compatibility}
\end{Highlighting}
\end{Shaded}

\subsection{Constructor Exhaustiveness
Matrix}\label{constructor-exhaustiveness-matrix}

The induction is over the normalized VA typing derivation. The matrix
below is exhaustive with respect to the VA grammar in Section 5.1.
``Unreachable in NVA'' means N4 excludes that outer constructor after
\texttt{T\_NORM}; the general VA preservation argument is nevertheless
recorded for reuse and auditability.

{\def\LTcaptype{none} % do not increment counter
\begin{longtable}[]{@{}lllll@{}}
\toprule\noalign{}
Constructor(s) & Induction hypotheses & Principal result used & Value
case & Normalized status \\
\midrule\noalign{}
\endhead
\bottomrule\noalign{}
\endlastfoot
\texttt{Literal}, \texttt{Self}, \texttt{Var} & none & environment
relation and scalar encoding & scalar/entity/bottom & reachable \\
\texttt{SetLiteral} & every element & M3a, G2, G3a typed finite-set
image & finite set, including bottom & reachable \\
\texttt{Attribute(e,a)} & receiver & R4 Attribute Preservation &
scalar/entity/bottom & reachable \\
\texttt{NavigationMany(...)} & receiver and every qualifier & R6-Many,
G4 \texttt{asSources} & finite entity set & reachable \\
\texttt{NavigationOne(...)} & receiver and every qualifier & R6-One,
\texttt{oneOrBottom} commutation & entity/bottom & reachable only at its
certified parent boundary \\
\texttt{ViewSet(e,D)} & receiver & G4 and empty/singleton equations &
finite entity set & reachable \\
\texttt{AllInstances(C)} & none & R3 and R7 & finite entity set &
reachable \\
\texttt{Not}, \texttt{And}, \texttt{Or} & operand(s) & G3/G4 validation
truth & Boolean/bottom policy & reachable \\
\texttt{Implies} & both operands & validation Boolean equation & Boolean
& unreachable in NVA by N4 \\
\texttt{Xor} & both operands & validation Boolean equation & Boolean &
unreachable in NVA by N4 \\
\texttt{Compare(EQ/NEQ,...)} & both operands & G3 and G3a typed equality
preservation/reflection & scalar/entity/set & reachable \\
ordered \texttt{Compare} & both operands & A8, G3, G3a & scalar &
reachable \\
\texttt{Arith} & both operands & A8, G3, G3a & scalar & reachable \\
\texttt{Coerce{[}tau-\textgreater{}upsilon{]}} & operand & G3a &
scalar/entity/bottom & reachable when introduced by normalization \\
scalar/entity \texttt{If} & condition and selected branch & G4 truth,
branch IH & scalar/entity/bottom & reachable \\
set-valued \texttt{If} & condition and selected branch & G4 truth, set
IH & finite set/bottom & reachable \\
\texttt{Let} & initializer and body under extension & G1 & every
admitted type & reachable \\
\texttt{Exists} & source and body under
\texttt{x-\textgreater{}encodeValue(s)} & G1, G2, G4 & object or scalar
finite set & reachable \\
\texttt{ForAll} & source and body under extension & counterexample
equivalence & object or scalar finite set & unreachable in NVA by N4 \\
\texttt{Select} & source and predicate under extension & G1/G2 & finite
set & reachable \\
\texttt{Reject} & source and predicate under extension & G1/G2 & finite
set & unreachable in NVA by N4 \\
\texttt{Collect} & source and body under extension & G1/G2, set-image
definition & finite set, including represented bottom & reachable \\
\texttt{IsUnique} & source and projection body & G1, G2/G3 equality
reflection & finite set to Boolean & reachable \\
\texttt{Includes}/\texttt{Excludes} & set and element & G2 typed
membership & finite set & unreachable in NVA by N4 \\
\texttt{IncludesAll}/\texttt{ExcludesAll} & both sets & G2 and G3a
subset/disjointness & finite sets & reachable \\
\texttt{Union}/\texttt{Intersection}/\texttt{AsSet} & operand set(s) &
M3a, G2 and G3a typed lifting and extensional set laws & finite sets &
reachable \\
\texttt{Count} & source & G2 cardinality & finite set to Integer &
reachable unless count-navigation redex is eliminated \\
\texttt{Size} & source & same denotation as \texttt{Count} & finite set
to Integer & unreachable in NVA by N4 \\
\texttt{IsEmpty}/\texttt{NotEmpty} & source & G2 emptiness & finite set
to Boolean & unreachable in NVA by N4 \\
\texttt{TypeKindOf} & receiver & R3 materialized conformance &
entity/bottom & reachable \\
\texttt{Cast} & receiver & R3 and matching failure-to-bottom policy &
entity/bottom & reachable \\
\end{longtable}
}

Every reachable normalized constructor has exactly one row. Together
with N4, this discharges constructor coverage rather than leaving
exhaustiveness to an informal comparison between syntax and proof prose.

\subsection{Proof}\label{proof-4}

Proceed by induction on the normalized VA typing derivation,
strengthened to every related environment pair reachable from the
initial pair and satisfying the inherited local closure premise.

Base cases:

\begin{Shaded}
\begin{Highlighting}[]
\NormalTok{Literal: scalars encode to themselves.}
\NormalTok{Self:    rho(self)=o corresponds to eta(self)=node(o).}
\NormalTok{Var:     follows from rho \textasciitilde{}Phi\_Gamma eta.}
\end{Highlighting}
\end{Shaded}

For a set literal, the induction hypotheses relate every element. M3a
selects one canonical result element type, and G3a makes each canonical
lifting commute with \texttt{encodeValue}. G2 then identifies the two
extensional finite images, including one bottom element when present.

Attribute case. For \texttt{Attribute(e,a)}, the induction hypothesis
gives corresponding receiver values. There are two exhaustive well-typed
runtime cases. If the receiver is \texttt{bottom}, the validation-level
primitive equations give \texttt{I\_attr\_obj(bottom,a)=bottom} and
\texttt{I\_attr\_graph(bottom,a)=bottom}, so the result is preserved
immediately. If the receiver is an entity \texttt{o} and the graph
receiver is \texttt{node(o)}, Theorem 0 gives:

\begin{Shaded}
\begin{Highlighting}[]
\NormalTok{attrVal\_M(o,a) = value\_G(node(o),a).}
\end{Highlighting}
\end{Shaded}

Navigation cases. For \texttt{NavigationMany}, the receiver and
qualifier induction hypotheses supply corresponding values. If a
qualifier is bottom, both sides produce the R6 \texttt{NoMatch} empty
result. Otherwise strict componentwise encoding and R6-Many give
corresponding target sets. \texttt{NavigationOne} applies R6-One and the
empty/singleton \texttt{oneOrBottom} equations to that same result.
\texttt{ViewSet} then maps corresponding bottoms to empty sets and
corresponding entities to corresponding singletons.

\texttt{AllInstances(C)} case. By materialized inheritance and Theorem
0:

\begin{Shaded}
\begin{Highlighting}[]
\NormalTok{encodeSet(allInstances\_obj(C)) = allInstances\_graph(C).}
\end{Highlighting}
\end{Shaded}

Boolean, comparison, arithmetic, and coercion cases. The induction
hypotheses preserve operands. G3a first commutes every selected
canonical embedding; G3 then preserves and reflects same-type equality,
order, and arithmetic. \texttt{EvalClosed} ensures the reached operation
is defined on one side exactly when it is defined on the other. G4 makes
validation truth, \texttt{finiteSet}, and \texttt{asSources} commute
with value encoding, including the stated bottom-to-empty policies.

For \texttt{Union} and \texttt{Intersection}, M3a and G3a first give
corresponding operands at the same \texttt{setJoin2} result type. G2
then preserves union, typed membership, and intersection in both
directions. \texttt{AsSet} applies \texttt{finiteSet}, so its
whole-set-bottom case commutes by G4 rather than by an identity law. For
\texttt{IsUnique}, G1 pairs source elements under the extended
environments and G3 reflects equality of projection values; hence a
collision exists on one side exactly when the corresponding collision
exists on the other side.

\texttt{Exists(S,x,tauD,P)} case. Let the source element type be
\texttt{tauE}. The induction hypothesis for \texttt{S} gives
corresponding finite source sets through \texttt{encodeValue}. For each
object or scalar element \texttt{s}, the corresponding graph-side value
is \texttt{encodeValue(s)}: if \texttt{s} is an object then
\texttt{encodeValue(s)=node(s)}, and if \texttt{s} is a scalar then
\texttt{encodeValue\_tauE(s)=s} at scalar identity types. The
declared-binder embedding commutes with encoding, so environment
extension preserves \texttt{\textasciitilde{}Phi\_Gamma}:

\begin{Shaded}
\begin{Highlighting}[]
\NormalTok{rho[x {-}\textgreater{} iota\_obj[tauE{-}\textgreater{}tauD](s)] \textasciitilde{}Phi\_Gamma[x:tauD]}
\NormalTok{  eta[x {-}\textgreater{} iota\_graph[tauE{-}\textgreater{}tauD](encodeValue\_tauE(s))].}
\end{Highlighting}
\end{Shaded}

Applying the induction hypothesis to \texttt{P} gives equal validation
truth for each witness. Therefore an object-side witness exists iff a
graph-side witness exists.

\texttt{ForAll(S,x,tauD,P)} case. The same environment-extension
argument applies. A counterexample exists on the object side iff the
corresponding graph counterexample exists. Equivalently, use the
normalized \texttt{not\ exists\ not} form.

\texttt{Select} and \texttt{Reject} cases. The induction hypothesis for
the predicate gives the same keep/drop decision for each corresponding
element under \texttt{encodeValue}. Hence selected or rejected finite
sets correspond under \texttt{encodeValue(Set(...))}.

\texttt{Collect} case. The source elements correspond by induction
through \texttt{encodeValue}. Applying the induction hypothesis to the
body under extended environments gives corresponding collected values.
Set semantics removes duplicate multiplicities.

\texttt{Includes}, \texttt{Excludes}, \texttt{IncludesAll}, and
\texttt{ExcludesAll} first use G3a to lift both operands to the unique
\texttt{memberJoin} type and then use G2. \texttt{Count}, \texttt{Size},
\texttt{IsEmpty}, and \texttt{NotEmpty} follow from the total
\texttt{finiteSet} homomorphism, including
\texttt{bottom\_Set\ -\textgreater{}\ empty}.

\texttt{If} case. The induction hypothesis preserves validation truth of
the condition, so both interpretations select corresponding branches.
Apply the induction hypothesis only to the selected, hence reachable,
branch, then G3a to the branch-to-\texttt{ifJoin} embedding (notably
\texttt{Void\ -\textgreater{}\ Set(sigma)}).

\texttt{Let} case. The induction hypothesis preserves the bound value.
Environment extension preserves the context-indexed
\texttt{\textasciitilde{}Phi}; the hereditary closure premise applies to
the reached body environment, so the body follows by induction.

Type-operation cases. For \texttt{TypeKindOf(E,C)}, a \texttt{bottom}
receiver evaluates to \texttt{false} in both interpretations by the
validation-level primitive equation; an entity receiver is preserved by
R3. For \texttt{Cast(E,C)}, a \texttt{bottom} receiver evaluates to
\texttt{bottom} in both interpretations. For an entity receiver,
conformance success is preserved by R3 and returns corresponding
entities, while conformance failure yields \texttt{bottom} on both
sides. \texttt{oclIsTypeOf} has no case because it is outside the
certified grammar.

All constructors in the normalized fragment preserve the strengthened
induction invariant. Hence the typed commutation equation holds;
applying G4 to a Boolean result gives the invariant-truth corollary.

\begin{center}\rule{0.5\linewidth}{0.5pt}\end{center}

\section{12. Theorem 5: Reference Cypher Realization under Cypher
Assumptions}\label{12-theorem-5-reference-cypher-realization-under-cypher-assumptions}

\subsection{Statement}\label{statement-5}

Theorem 5 has two layers and uses the realization judgments of Section
6.4, thereby avoiding a type mismatch between VA values and Cypher
result tables.

\textbf{Expression realization.} Suppose \texttt{b} is a Bound
subexpression of an admitted invariant body \texttt{bInv},
\texttt{BoundAdm\_MM(bInv)}, \texttt{MM;Gamma\ \textbar{}-\ b:tau},
\texttt{va=T\_VA(b)}, and \texttt{nva=T\_NORM(va)}. Thus, by M2 and
Theorems 2 and 3, \texttt{MM;Gamma\ \textbar{}-VA\ nva:tau},
\texttt{NF\_R(nva)}, and \texttt{nva} belongs to the certified range
whose attributes and navigation qualifiers satisfy the exact primitive
metadata premises above. Let the renderer be initialized with an
injective alias map \texttt{alpha} with \texttt{dom(alpha)=dom(Gamma)},
and let \texttt{pi} be a sound parameter state. Realization is
quantified only over \texttt{AdmRow\_Gamma(nva,alpha,eta,row)}, which
includes \texttt{eta\ \textbar{}=\_G\ Gamma} and
\texttt{EvalClosed\_graph(nva,eta)}. If \texttt{tau} is non-set,
\texttt{BuildCQ} yields a unique expression-kind query-model derivation
and, for every valid live set \texttt{K}:

\begin{Shaded}
\begin{Highlighting}[]
\NormalTok{(ScalarPlan(L,c),pi\textquotesingle{})=ExpandE(BuildCQ(nva),alpha,pi,K),}
\NormalTok{pi subseteq pi\textquotesingle{},}
\NormalTok{Gamma;alpha,pi\textquotesingle{};K |{-} ScalarPlan(L,c) realizes\_scalar nva:tau}
\end{Highlighting}
\end{Shaded}

(or \texttt{realizes\_pred} when only validation truth is observed). If
\texttt{tau=Set(sigma)}, then for every valid correlated live set
\texttt{K}:

\begin{Shaded}
\begin{Highlighting}[]
\NormalTok{(SetPlan(Q,a),pi\textquotesingle{})=ExpandS(BuildCQ(nva),alpha,pi,K),}
\NormalTok{pi subseteq pi\textquotesingle{},}
\NormalTok{Gamma;alpha,pi\textquotesingle{} |{-} Q =\textgreater{} a realizes\_set nva:Set(sigma).}
\end{Highlighting}
\end{Shaded}

Every newly allocated entry of \texttt{pi\textquotesingle{}} satisfies
C2, so the extension is sound. TXT5 closes either result to a raw Cypher
AST. By TXT4 the displayed realization judgments hold, and by TXT1/TXT5
the result is unique up to fresh alias and parameter renaming.

\textbf{Invariant query realization.} For this second layer, specialize
the expression result to the admitted invariant body itself:
\texttt{BInvariant(C,R,bInv)}, \texttt{b=bInv},
\texttt{Gamma=\{self:C\}}, \texttt{tau=Boolean},
\texttt{va=T\_VA(bInv)}, and \texttt{nva=T\_NORM(va)}. Additionally fix
\texttt{M} and \texttt{G} such that \texttt{GraphAdequate(MM,M,G)}, and
use the invariant-root closure/representation premises A8-\/-A9. Let the
formal plan be
\texttt{(Qspec,pispec)=ExpandI(BuildCQInvariant(C,nva),pi0(MM))}. TXT4,
TXT5, C5a/C5b, and C6 give:

\begin{Shaded}
\begin{Highlighting}[]
\NormalTok{returnedIds(Qspec,G,pispec)=ViolIds\_\{M,G\}(C,nva).}
\end{Highlighting}
\end{Shaded}

For the reference query, let \texttt{plan=T\_CQM\^{}spec(C,R,nva)},
\texttt{(tq,pi)=T\_TEXT\^{}spec(plan)}, and
\texttt{q=parse\_Cypher(tq)}. If
\texttt{SpecPlanAdequacy(MM,C,nva,plan,q,pi,G)} holds, then
\texttt{SpecPlanSim\_sound}, derived by constructor induction rather
than assumed as result equality, gives:

\begin{Shaded}
\begin{Highlighting}[]
\NormalTok{returnedIds(q,G,pi)}
\NormalTok{= ViolIds\_\{M,G\}(C,nva)}
\NormalTok{= \{ id(o)}
\NormalTok{    | o in Obj(M)}
\NormalTok{      and classOf(o) conformsTo C}
\NormalTok{      and bool\_val([[nva]]\_graph\^{}Boolean}
\NormalTok{                     (G,self {-}\textgreater{} node(o))) = false \}.}
\end{Highlighting}
\end{Shaded}

Consequently:

\begin{Shaded}
\begin{Highlighting}[]
\NormalTok{for every o in Obj(M):}

\NormalTok{id(o) in returnedIds(q,G,pi)}
\NormalTok{iff}
\NormalTok{classOf(o) conformsTo C}
\NormalTok{and bool\_val([[nva]]\_graph\^{}Boolean(G,self {-}\textgreater{} node(o))) = false;}

\NormalTok{and for every u in returnedIds(q,G,pi),}
\NormalTok{  there exists a unique o in Obj(M) with u=id(o).}
\end{Highlighting}
\end{Shaded}

The pointwise equivalence is a corollary of the set equality and
no-ghost property; it is not used as a weaker replacement for them.

\subsection{Selected Dialect and Cypher
Assumptions}\label{selected-dialect-and-cypher-assumptions}

The theorem is realization under the selected \texttt{CYPHER5\_val}
profile, not a backend-independent theorem. It relies on the following
extensional behavioral axioms for only the generated subset. The
concrete server version is an experiment parameter and must pass the
Section 6.5 dialect probes; Java driver 5.21.0 alone does not discharge
this premise.

\begin{Shaded}
\begin{Highlighting}[]
\NormalTok{CY1 Representation/lookup: repr\_C\^{}tau and abs\_C\^{}tau satisfy BR1{-}{-}BR10 and}
\NormalTok{    BuilderClosure\_C holds;}
\NormalTok{    parameter and property lookup return the represented stored value, and an}
\NormalTok{    absent property yields CypherNull, abstracted as typed bottom; the emitted}
\NormalTok{    scalar decoder obeys the typed codec equation in BR4, the emitted coercer}
\NormalTok{    obeys BR9, and the emitted qualifier encoder obeys BR10.}
\NormalTok{CY2 Pattern: MATCH extends each input row with every graph assignment satisfying}
\NormalTok{    the emitted label, direction, type(r) STARTS WITH \textquotesingle{}Link\textquotesingle{}, and property}
\NormalTok{    predicates, with ordinary Cypher bag multiplicity. OPTIONAL MATCH has the}
\NormalTok{    same result when at least one assignment exists; when none exists, it emits}
\NormalTok{    exactly one extension of the input row with every newly introduced pattern}
\NormalTok{    alias bound to CypherNull.}
\NormalTok{CY3 Filter: WHERE retains exactly rows whose condition evaluates to Cypher true.}
\NormalTok{CY4 Existential: EXISTS \{Q\} is true exactly when}
\NormalTok{    EvalRows\_G,pi(Q,\{row\}) is non{-}empty for the current graph, parameter map,}
\NormalTok{    and correlated input row.}
\NormalTok{CY5 Aggregation/collection subquery: \textasciigrave{}COLLECT \{Q RETURN DISTINCT x\}\textasciigrave{} returns}
\NormalTok{    the projected extensional list of}
\NormalTok{    \textasciigrave{}EvalRows\_G,pi(Q,\{row\})\textasciigrave{} for the current correlated input \textasciigrave{}row\textasciigrave{};}
\NormalTok{    \textasciigrave{}COUNT \{Q RETURN DISTINCT x\}\textasciigrave{} is the cardinality of that same distinct}
\NormalTok{    projection, \textasciigrave{}size(...)\textasciigrave{} returns the finite list cardinality, and}
\NormalTok{    collection{-}subquery order is never observed. On the ONE{-}navigation lists emitted by the compiler,}
\NormalTok{    \textasciigrave{}head([])=CypherNull\textasciigrave{} and \textasciigrave{}head([x])=x\textasciigrave{}; the compiler never applies \textasciigrave{}head\textasciigrave{}}
\NormalTok{    to a list with more than one semantic element. Reference raw plans may use}
\NormalTok{    an extensionally equivalent \textasciigrave{}COUNT(DISTINCT x)\textasciigrave{} form. Semantic bottom}
\NormalTok{    inside a VA set is emitted as non{-}null BOTTOM\_TOKEN.}
\NormalTok{CY6 Projection: RETURN/WITH project the stated expressions; DISTINCT replaces}
\NormalTok{    the projected row bag by one representative of each equal projected row.}
\NormalTok{CY7 Control/truth: Boolean, comparison, arithmetic, CASE, and coalesce have the}
\NormalTok{    standard behavior used by the emitted expressions; the renderer maps every}
\NormalTok{    graph{-}side non{-}validation{-}true result to false before invariant filtering.}
\NormalTok{CY8 Row expansion: UNWIND emits one row per member of the finite realized list}
\NormalTok{    used by the renderer; no order property is used by the proof.}
\NormalTok{CY9 Branch composition: inside a correlated CALL, \textasciigrave{}Q\textasciigrave{} can see exactly the}
\NormalTok{    explicitly named imported aliases. After the call, each returned subquery}
\NormalTok{    row is joined with the entire unchanged input row (subject only to fresh,}
\NormalTok{    non{-}colliding output aliases); UNION ALL concatenates branch row bags. Guarded}
\NormalTok{    complementary branches therefore realize validation{-}level conditional sets.}
\end{Highlighting}
\end{Shaded}

The assumptions form three different parts of the trusted semantic
boundary:

{\def\LTcaptype{none} % do not increment counter
\begin{longtable}[]{@{}llll@{}}
\toprule\noalign{}
Class & Assumptions & What is trusted & Required discharge/evidence \\
\midrule\noalign{}
\endhead
\bottomrule\noalign{}
\endlastfoot
Representation interface & CY1 together with BR1-\/-BR10 and
\texttt{BuilderClosure\_C} & representation/abstraction retraction,
raw-builder closure, lookup/codec/coercion behavior, bottom-token
separation, qualifier serialization, and stability through
projection/counting & mathematical proof for typed
\texttt{repr\_C}/\texttt{abs\_C}, explicit finite builder equations plus
codec/coercion commutation, qualifier encoding, and token non-collision;
implementation conformance tests for parameter/property access \\
Core generated-query semantics & CY2-\/-CY8 & extensional behavior of
pattern matching, filtering, existence, aggregation, projection, scalar
control/truth, and finite row expansion & selected Cypher semantic
specification plus isolated executable probes for every emitted form
used by a proof case \\
Version-sensitive composition & CY9 & correlated imports, branch scope,
row joining, and \texttt{UNION\ ALL} behavior & mandatory server-version
probe covering imported aliases, nested correlation, complementary
branches, and common branch schemas \\
\end{longtable}
}

The corresponding evidence obligations are tracked more finely below.

{\def\LTcaptype{none} % do not increment counter
\begin{longtable}[]{@{}llll@{}}
\toprule\noalign{}
Assumption & Used to justify & Minimum evidence before claiming runtime
coverage & What the evidence does not prove \\
\midrule\noalign{}
\endhead
\bottomrule\noalign{}
\endlastfoot
CY1 & literal/parameter access, property/codec access, typed coercion
and qualifier serialization & representation unit tests over every
admitted scalar kind, object identifier, missing property,
\texttt{BOTTOM\_TOKEN}, every qualifier type/escape, and mixed numeric
coercion & arbitrary driver conversions or values outside
\texttt{CyVal\_C} \\
CY2 & context matching, navigation, \texttt{allInstances}, optional
attribute slot & forward/reverse relationship probes with duplicate
paths, metadata predicates, swapped reverse roles, qualifiers, and both
match/no-match \texttt{OPTIONAL\ MATCH} rows & arbitrary patterns or
planner optimizations outside generated ASTs \\
CY3 & predicate filtering and complementary guards & tests for true,
false, Cypher null, represented bottom, and guarded entity receivers &
full OCL four-valued logic \\
CY4 & \texttt{Exists}, membership, emptiness, subset/disjointness
counterexamples & empty/non-empty and correlated witness probes &
arbitrary user-written subqueries \\
CY5 & \texttt{COLLECT}, \texttt{COUNT}, \texttt{head}, and finite
cardinality & correlated empty/singleton collection,
\texttt{head({[}{]})}, \texttt{head({[}x{]})}, duplicate, null, distinct
scalar/object, and \texttt{BOTTOM\_TOKEN} probes & Bag multiplicity,
order, or \texttt{head} on lists outside the ONE invariant \\
CY6 & live projection, set quotienting, returned IDs &
\texttt{WITH}/\texttt{RETURN} alias-scope and \texttt{DISTINCT} probes &
preservation of columns not explicitly projected \\
CY7 & Boolean/scalar operations, \texttt{truth}, \texttt{CASE},
\texttt{coalesce} & boundary tests for every admitted scalar operator
and validation-false policy & overflow, NaN/infinity, locale collation,
or unsupported coercions excluded by \texttt{ScalarClosed} \\
CY8 & renderer-owned finite list expansion & empty, singleton,
duplicate, and bottom-token list probes & order-sensitive semantics \\
CY9 & set conditionals and bottom-safe receiver branching & correlated
\texttt{CALL}, explicit imports, \texttt{UNION\ ALL}, nested
iterator/self, and one-enabled-arm probes on the selected server & other
Neo4j/Cypher versions that have not passed the same suite \\
\end{longtable}
}

The historical selected-runtime evidence manifest dated 2026-08-01
contains one row for each CY label and records \texttt{PASS} for those
then-current probes. The strengthened OPTIONAL-MATCH, \texttt{head},
typed coercion, reverse-role, and qualifier-codec subclauses above
require corresponding probe observations and a fresh manifest hash
before they are treated as discharged. In the historical run, CY2
observes two physical forward and reverse paths but one distinct target;
CY3 separates Cypher null from the non-null represented-bottom guard;
CY5 and CY8 make row, null, duplicate, and semantic-set cardinalities
observable; CY9 exercises both complementary arms with nested explicit
imports. \texttt{Cypher5ValRuntimeEvidenceManifestTest} makes the
artifact stale when the selected runtime constants, driver dependency,
renderer, canonical encoding, probe matrix, or runtime harness no longer
match.

These are assumptions about Neo4j behavior, not results proved in this
document. Specification references explain the intended semantics, while
dialect probes and conformance tests provide runtime evidence; neither
converts CY1-\/-CY9 into internally derived lemmas. No correctness claim
is made for arbitrary Cypher, untested server versions, or target
constructs outside \texttt{Cypher\_val}.

\subsection{Distinctness Discipline}\label{distinctness-discipline}

Cypher \texttt{MATCH} may produce duplicate rows. Therefore:

\begin{Shaded}
\begin{Highlighting}[]
\NormalTok{1. Violation results are compared as sets of returnedIds.}
\NormalTok{2. Invariant queries use RETURN DISTINCT self.use\_id.}
\NormalTok{3. Cardinality over VA sets uses \textasciigrave{}size(COLLECT \{ ... RETURN DISTINCT x \})\textasciigrave{}}
\NormalTok{   (or a proved equivalent \textasciigrave{}COUNT(DISTINCT x)\textasciigrave{} reference form) whenever graph}
\NormalTok{   patterns can introduce duplicate rows.}
\NormalTok{4. Intermediate expression realization must preserve VA set semantics, not}
\NormalTok{   Cypher bag semantics.}
\end{Highlighting}
\end{Shaded}

\subsection{Required Lemmas}\label{required-lemmas-4}

\begin{Shaded}
\begin{Highlighting}[]
\NormalTok{C1 Alias and Live{-}Projection Preservation}
\NormalTok{C2 Parameter Soundness}
\NormalTok{C3 Primitive Graph Access Realization}
\NormalTok{C3a Bottom{-}Safe Entity Receiver}
\NormalTok{C4 Scalar{-}Plan, Predicate, and Existential Composition}
\NormalTok{C5 Finite{-}Set and Cardinality Realization}
\NormalTok{C5a Collection{-}Bottom Boundary}
\NormalTok{C5b Extensional Set Equality}
\NormalTok{C6 Invariant Wrapper and No{-}Ghost Realization}
\NormalTok{SpecPlanAdequacy and SpecPlanSim\_sound for the reference plan/text}
\NormalTok{TXT1{-}{-}TXT4 Translation Determinism, Kind, Closure, and Realization}
\NormalTok{TXT5 Raw{-}AST Closure and Totality}
\NormalTok{Neo4j subset assumptions}
\end{Highlighting}
\end{Shaded}

\subsection{Constructor-Coverage
Certificate}\label{constructor-coverage-certificate}

The following matrix makes the structural induction obligation explicit.
A row is \texttt{Covered} only when all four items are available: a
syntax-directed \texttt{BuildCQ} rule, a total raw-AST expansion, a
local semantic argument, and the listed premises. Constructors removed
by normalization are marked \texttt{Unreachable} rather than silently
omitted.

{\def\LTcaptype{none} % do not increment counter
\begin{longtable}[]{@{}lllll@{}}
\toprule\noalign{}
Normalized VA constructor family & Query-model/raw-AST realization &
Local proof obligation & Main dependencies & Status \\
\midrule\noalign{}
\endhead
\bottomrule\noalign{}
\endlastfoot
\texttt{Self}, scalar/entity variable, set variable & \texttt{CQAlias},
\texttt{CQSetAlias} & alias denotes the corresponding environment value
and remains live; set use passes through C5a & C1, C5a,
\texttt{RowCorr\_Gamma} & Covered \\
literal & \texttt{CQLiteral} with fresh parameter & parameter
abstraction equals the admitted literal & C2, CY1 & Covered \\
\texttt{SetLiteral} & \texttt{CQSetLiteral}, \texttt{setLiteralPlan} &
canonical lifting and distinct projection equal the semantic finite set
& D8, C5, BR2, BR5 & Covered \\
attribute & \texttt{CQAttribute}, \texttt{withEntityReceiver},
\texttt{attrBody} & bottom receiver is not used in a node pattern;
entity receiver returns \texttt{I\_attr\_graph} & LR-TypeAccess, C3a,
R4, CY1, CY3, CY9 & Covered \\
\texttt{AllInstances} & \texttt{CQAllInstances} & projected nodes equal
the graph instance set & C3, R3/R7, CY2, CY6 & Covered \\
directional/qualified navigation (\texttt{ONE}/\texttt{MANY}) and
\texttt{ViewSet} & \texttt{CQNavigationOne}, \texttt{CQNavigation},
\texttt{CQViewSet} & matched targets equal R6; one/empty and set
boundary preserve result kind & D1, D12, R5/R6, BR10, C5a, CY1-\/-CY3,
CY6 & Covered \\
\texttt{Not}, Boolean connective & \texttt{CQNot}, \texttt{CQBoolean} &
target truth expression agrees with validation truth & C4, CY7 &
Covered \\
scalar equality/order/arithmetic/coercion & \texttt{CQCompare},
\texttt{CQArith}, \texttt{CQCoerce} & represented scalar operation
abstracts to the graph VA operation & BR1-\/-BR9, scalar compatibility,
G3a, CY7 & Covered \\
finite-set equality/inequality & \texttt{CQSetEquality},
\texttt{extSetEq} & mutual inclusion equals extensional set equality &
D13, C5b & Covered \\
scalar \texttt{If} & \texttt{CQIfExpr}, \texttt{ifScalar} & exactly the
validation-selected branch contributes one scalar result &
LR-Conditional, C1, CY3, CY9 & Covered \\
set-valued \texttt{If} & \texttt{CQIfSet}, \texttt{ifSet} & exactly the
selected branch contributes its semantic set & D4, LR-Conditional, C5,
CY3, CY9 & Covered \\
scalar/set typed or inferred \texttt{Let} & \texttt{CQLetExpr},
\texttt{CQLetSet}, \texttt{bindValue} & initializer is evaluated once,
embedded to \texttt{tauD}, and body sees \texttt{x:tauD} & D6, LR-Let,
C1, A8 & Covered \\
typed or inferred \texttt{Exists}, direct \texttt{ForAll} &
\texttt{CQExists}, \texttt{CQForAll} & each \texttt{tauE} row is
embedded to \texttt{tauD}; witness/counterexample agrees with VA & D3,
LR-Quantifier, C1, C4, C5, CY4 & Covered \\
typed or inferred \texttt{Select}, \texttt{Reject} & \texttt{CQSelect},
\texttt{CQReject}, \texttt{filterRows} & body sees the embedded iterator
view while projected result retains original elements & D3,
LR-SelectReject, C1, C5, CY3, CY6 & Covered \\
typed or inferred finite-set \texttt{Collect} & \texttt{CQCollect},
\texttt{mapRows} & body sees the embedded iterator view; result is the
finite image, including bottom & D3, D5, LR-Collect, C1, C5, BR1-\/-BR9,
CY6 & Covered \\
\texttt{Includes}, \texttt{Excludes} & \texttt{CQMembership},
\texttt{memberIn} & membership test agrees with semantic finite-set
membership & LR-Membership, C5, CY4, CY7 & Covered \\
\texttt{IncludesAll}, \texttt{ExcludesAll} & \texttt{CQSetRelation} &
absence of a counterexample realizes subset/disjointness & D2,
LR-SetRelation, C5, CY4 & Covered \\
\texttt{Union}, \texttt{Intersection}, \texttt{AsSet} &
\texttt{CQUnion}, \texttt{CQIntersection}, \texttt{CQAsSet} & canonical
lifting followed by extensional set operation & D9-\/-D10, C5, G3a &
Covered \\
\texttt{IsUnique} & \texttt{CQIsUnique}, \texttt{uniquePlan} &
projection is injective exactly when source and distinct-image
cardinalities agree & D11, C5, BR5-\/-BR6 & Covered \\
\texttt{Count}, normalized \texttt{Size} & \texttt{CQCount},
\texttt{CountExpr(closeSet(...))} & count equals semantic finite-set
cardinality & LR-Cardinality, C5, BR5-\/-BR7, CY5, CY6 & Covered \\
\texttt{IsEmpty}, \texttt{NotEmpty} & \texttt{CQEmpty},
\texttt{ExistsExpr(closeSet(...))} & row absence/presence equals
finite-set emptiness/non-emptiness & LR-Cardinality, C5, CY4 &
Covered \\
\texttt{TypeKindOf}, \texttt{Cast} & \texttt{CQKindOf}, \texttt{CQCast},
\texttt{withEntityReceiver} & bottom and entity branches agree with
graph type semantics & D7, LR-TypeAccess, C3a, R3, CY3, CY9 & Covered \\
invariant root & \texttt{CQInvariant}, \texttt{ExpandI} & context match,
validation-false filter, and distinct identifier projection realize
\texttt{realizes\_inv} & C6, R2/R3, CY2, CY3, CY6 & Covered \\
\texttt{Implies}, \texttt{Xor}, \texttt{Reject} when eliminated,
\texttt{Size}, \texttt{IsEmpty}, \texttt{NotEmpty}, \texttt{ForAll} when
rewritten & normalization rules & no independent target case is required
for a constructor absent from the reached normal form & N1-\/-N4 &
Unreachable when eliminated; otherwise covered by the direct row
above \\
\texttt{any}, \texttt{one}, \texttt{sortedBy}, \texttt{count(element)},
\texttt{oclIsTypeOf}, ordered/bag constructs & no admitted rule &
outside \texttt{OCL\_val}/normalized VA codomain & admission boundary &
Excluded \\
\end{longtable}
}

The matrix is exhaustive over the normalized VA grammar. Therefore
Theorem 5 does not rely on an implicit default renderer case. Adding a
constructor to the grammar invalidates the exhaustiveness claim until a
new row, expansion, local realization lemma, and proof case are
supplied.

\subsection{Local Realization Lemmas}\label{local-realization-lemmas}

The following lemmas expose the semantic steps that were previously
compressed inside the prose induction. In each statement, the source and
body plans are assumed to satisfy their induction hypotheses under
\texttt{RowCorr\_Gamma}, and all projections use
\texttt{projectSemanticSet} so Cypher bag multiplicity is not confused
with VA finite-set semantics.

\textbf{LR-Quantifier (existential and universal realization).} If
\texttt{RS=SetPlan(Q,a)} realizes \texttt{S} and
\texttt{ScalarPlan(L,p)} realizes \texttt{P} under
\texttt{alpha{[}x\ -\textgreater{}\ a{]}}, then:

\begin{Shaded}
\begin{Highlighting}[]
\NormalTok{bool\_val(abs\_C\^{}Boolean(EvalExpr\_G,pi(}
\NormalTok{  existsRows(RS,imports,L,truth(p)),row))) = true}
\NormalTok{iff}
\NormalTok{exists s in finiteSet\_tauE([[S]]\_graph):}
\NormalTok{  bool\_val([[P]]\_graph(}
\NormalTok{    eta[x {-}\textgreater{} iota\_graph[tauE{-}\textgreater{}tauD](s)])) = true.}
\end{Highlighting}
\end{Shaded}

Proof. C5 gives one semantic representative row for each element of
\texttt{finiteSet\_tauE({[}{[}S{]}{]}\_graph)}. The declared-binder
embedding and C1 turn extension by alias \texttt{a} into environment
extension at \texttt{x:tauD}, including scalar elements as established
in D3. The body induction hypothesis makes \texttt{truth(p)} true
exactly for semantic witnesses. CY3 removes non-witness rows and CY4
equates subquery non-emptiness with \texttt{EXISTS}. Negating the same
construction with \texttt{NOT\ truth(p)} proves the direct
\texttt{ForAll} rule by absence of a counterexample.

\textbf{LR-SelectReject (finite-set filtering).} Under the same
source/body premises:

\begin{Shaded}
\begin{Highlighting}[]
\NormalTok{SemSet\_G,pi,tauE(filterRows(RS,imports,L,truth(p),a),row)}
\NormalTok{= \{ s in finiteSet\_tauE([[S]]\_graph)}
\NormalTok{    | bool\_val([[P]]\_graph(}
\NormalTok{        eta[x {-}\textgreater{} iota\_graph[tauE{-}\textgreater{}tauD](s)])) = true \}.}
\end{Highlighting}
\end{Shaded}

Replacing \texttt{truth(p)} by \texttt{NOT\ truth(p)} yields the
corresponding \texttt{Reject} equality.

Proof. For each semantic source element, declared-binder embedding plus
C1 establishes the extended \texttt{RowCorr\_Gamma{[}x:tauD{]}}; the
body induction hypothesis gives the same validation truth as the VA
predicate. CY3 keeps exactly the rows satisfying the selected guard. CY6
and C5 project and quotient these rows by semantic equality.
Consequently no satisfying element is lost, no non-satisfying element is
introduced, and duplicate Cypher rows do not alter the result set.

\textbf{LR-Collect (finite-set image realization).} If
\texttt{RS=SetPlan(Q,a)} realizes \texttt{S} and the body scalar plan
realizes \texttt{P} under \texttt{alpha{[}x\ -\textgreater{}\ a{]}},
then:

\begin{Shaded}
\begin{Highlighting}[]
\NormalTok{SemSet\_G,pi,sigma(mapRows(RS,imports,L,p,z),row)}
\NormalTok{= \{ [[P]]\_graph(eta[x {-}\textgreater{} iota\_graph[tauE{-}\textgreater{}tauD](s)])}
\NormalTok{    | s in finiteSet\_tauE([[S]]\_graph) \}.}
\end{Highlighting}
\end{Shaded}

Proof. Source-row correspondence and C1 reduce each output row to the
body induction hypothesis. \texttt{setOut} converts \texttt{CypherNull}
representing semantic bottom to the reserved non-null
\texttt{BOTTOM\_TOKEN}; BR1-\/-BR9 make abstraction commute with that
conversion. CY6 \texttt{DISTINCT} removes exactly duplicate image
values. D5 shows that bottom is retained as one semantic set element, so
the equality also covers bottom-producing bodies. No multiplicity,
order, or flattening property is claimed.

\textbf{LR-Membership (membership and non-membership realization).}
Suppose \texttt{R:Set(sigma)}, \texttt{Y:delta},
\texttt{memberJoin(sigma,delta)=upsilon}, the set plan realizes
\texttt{R}; write \texttt{RS=SetPlan(Q,a)} for that plan. Let
\texttt{PY=ScalarPlan(LY,eY)} realize \texttt{Y}, let \texttt{y} be
fresh, and put \texttt{XY=materialize(PY,y,K)}. For every admitted input
row, let \texttt{ExecPrelude\_G,pi(XY.prelude,row)=\{rowY\}}; hence
\texttt{Alias(y)} is defined in \texttt{rowY} and C1 preserves the
correlated outer bindings. Then:

\begin{Shaded}
\begin{Highlighting}[]
\NormalTok{bool\_val(abs\_C\^{}Boolean(EvalExpr\_G,pi(}
\NormalTok{  memberIn[sigma,delta{-}\textgreater{}upsilon](}
\NormalTok{    RS,imports union \{y\},[],Alias(y)),rowY))) = true}
\NormalTok{iff}
\NormalTok{iota\_graph[delta{-}\textgreater{}upsilon]([[Y]]\_graph\^{}delta(G,eta))}
\NormalTok{  in upSet\_graph[sigma{-}\textgreater{}upsilon](}
\NormalTok{       finiteSet\_sigma([[R]]\_graph\^{}Set(sigma)(G,eta))).}
\end{Highlighting}
\end{Shaded}

Proof. The scalar realization of \texttt{PY} and C1 identify
\texttt{rowY(y)} with the semantic value of \texttt{Y}; no expression is
evaluated against the pre-materialized row. C5 identifies the projected
source rows with the semantic set. BR9 makes both generated coercions
commute with the corresponding canonical embeddings, and typed
\texttt{setEq\_upsilon} is true exactly when the two coerced values
abstract to the same \texttt{upsilon} value. CY4 therefore equates
existence of an equal row with membership in the lifted left set.
Boolean negation proves \texttt{Excludes}.

\textbf{LR-SetRelation (subset and disjointness realization).}

For \texttt{R:Set(tau)}, \texttt{T:Set(sigma)}, and
\texttt{memberJoin(tau,sigma)=upsilon}, put:

\begin{Shaded}
\begin{Highlighting}[]
\NormalTok{R\_up = upSet\_graph[tau{-}\textgreater{}upsilon](finiteSet\_tau([[R]]\_graph)),}
\NormalTok{T\_up = upSet\_graph[sigma{-}\textgreater{}upsilon](finiteSet\_sigma([[T]]\_graph)).}

\NormalTok{SUBSET(R,T) = true       iff T\_up subseteq R\_up,}
\NormalTok{DISJOINT(R,T) = true     iff R\_up intersect T\_up = empty.}
\end{Highlighting}
\end{Shaded}

Proof. For \texttt{SUBSET}, D2 and LR-Membership identify the generated
nested subquery with a right-side counterexample absent from the left
side. CY4 and outer negation state that no such counterexample exists.
For \texttt{DISJOINT}, the same argument searches for a right-side
element that is also in the left side and negates its existence.
Explicit imports preserve correlation in both nested queries.

\textbf{LR-Cardinality (count and emptiness realization).} If
\texttt{RS=SetPlan(Q,a)} realizes \texttt{S}, then:

\begin{Shaded}
\begin{Highlighting}[]
\NormalTok{abs\_C\^{}Integer(EvalExpr\_G,pi(CountExpr(closeSet(RS)),row))}
\NormalTok{= |finiteSet([[S]]\_graph)|,}

\NormalTok{bool\_val(abs\_C\^{}Boolean(EvalExpr\_G,pi(}
\NormalTok{  ExistsExpr(closeSet(RS)),row))) = true}
\NormalTok{iff}
\NormalTok{finiteSet([[S]]\_graph) != empty.}
\end{Highlighting}
\end{Shaded}

Proof. \texttt{closeSet} projects \texttt{DISTINCT\ setOut(a)}. C5 and
BR5 identify these non-null projected values one-to-one with semantic
set elements, including at most one bottom token. CY5 counts precisely
those distinct rows. CY4 gives the emptiness equivalence, from which
\texttt{IsEmpty} and \texttt{NotEmpty} follow by Boolean negation or
identity.

\textbf{LR-TypeAccess (attribute, kind-of, and cast realization).} Let a
scalar plan realize an attribute receiver \texttt{E:D}, or let
\texttt{K:kappa} be a type-operation receiver with
\texttt{kappa\ in\ \{Void,D\}}. For each expression
\texttt{X\ in\ \{Attribute(E,a),TypeKindOf(K,C),Cast(K,C)\}}, destruct
its expansion as
\texttt{(ScalarPlan(L\_X,c\_X),pi\_X)=ExpandE(BuildCQ(X),alpha,pi,K\_live)}
and let \texttt{ExecPrelude\_G,pi\_X(L\_X,row)=\{row\_X\}}.
\texttt{withEntityReceiver} produces exactly this one result row and:

\begin{Shaded}
\begin{Highlighting}[]
\NormalTok{abs\_C\^{}tau(EvalExpr\_G,pi\_Attribute(c\_Attribute,row\_Attribute))}
\NormalTok{  = I\_attr\_graph\^{}tau([[E]]\_graph\^{}D,a),}
\NormalTok{abs\_C\^{}Boolean(EvalExpr\_G,pi\_TypeKindOf(c\_TypeKindOf,row\_TypeKindOf))}
\NormalTok{  = TypeKindOf\_graph([[K]]\_graph\^{}kappa,C),}
\NormalTok{abs\_C\^{}C(EvalExpr\_G,pi\_Cast(c\_Cast,row\_Cast))}
\NormalTok{  = Cast\_graph\^{}C([[K]]\_graph\^{}kappa,C).}
\end{Highlighting}
\end{Shaded}

Proof. If the receiver is bottom, BR1-\/-BR4 make only the bottom arm
survive CY3, and no graph pattern receives the bottom representation. If
it is an entity, target-profile closure supplies its node representation
and only the entity arm survives. R4 proves the attribute result; R3
proves conformance and the successful/failed cast result. CY9 preserves
imported aliases and combines the complementary singleton arms. This is
the local result used by D7.

\textbf{LR-Conditional (scalar and set conditional realization).} If the
condition and both branches are realized, \texttt{ifScalar} returns the
semantic value of exactly the branch selected by
\texttt{bool\_val(condition)}, and \texttt{ifSet} returns exactly its
semantic finite set.

Proof. \texttt{truth(condition)} is total Boolean by BR1-\/-BR4 and CY7.
The guards \texttt{truth(cc)} and \texttt{NOT\ truth(cc)} are
complementary, so CY3 enables one branch only. CY9 preserves correlation
and concatenates only the enabled branch rows. For sets, CY6/C5 remove
duplicate rows after \texttt{setOut}; for scalars, the singleton branch
obligation establishes \texttt{WFScalar}. D4 is the set-valued hard
case.

\textbf{LR-Let (environment-binding realization).} For both scalar and
set results, the expansion of \texttt{Let(x,tauD,V:tau1,B)} realizes:

\begin{Shaded}
\begin{Highlighting}[]
\NormalTok{[[B]]\_graph(eta[x {-}\textgreater{} iota\_graph[tau1{-}\textgreater{}tauD]([[V]]\_graph(eta))]).}
\end{Highlighting}
\end{Shaded}

Proof. The initializer induction hypothesis and \texttt{bindValue}
materialize the canonically embedded \texttt{{[}{[}V{]}{]}\_graph(eta)}
exactly once under a fresh alias. C1 establishes row correspondence for
the displayed \texttt{tauD} environment; the body induction hypothesis
then gives the displayed result. Restoring the alias map after the body
implements lexical scope. Both recursive calls are on proper children,
which proves termination. This is the local obligation used by D6.

\subsection{Proof}\label{proof-5}

Expression realization is proved by induction on the syntax-directed
\texttt{BuildCQ} derivation from Section 6.5 and its constructor-matched
expansion in Section 6.6. The induction invariant is the applicable
realization judgment from Section 6.4 under every corresponding input
row. TXT2 guarantees that the derivation kind agrees with the static
result type, TXT3 keeps every target fragment inside
\texttt{Cypher\_val}, and TXT5 ensures that every case closes to a raw
AST without holes.

Variables and \texttt{self} are realized by aliases assigned by the
planner. C1 and \texttt{RowCorr\_Gamma} ensure that each alias denotes
the corresponding VA environment binding and is not captured by nested
subqueries.

Literals are realized by parameters. C2 establishes the scalar
realization judgment.

Attributes are realized using \texttt{withEntityReceiver} and the same
canonical accessor policy as \texttt{value\_G}. The bottom branch
returns null without evaluating a node pattern; the entity branch is
admitted only after the non-bottom guard and returns
\texttt{I\_attr\_graph} by C3/R4. The D7 singleton argument preserves
\texttt{WFScalar} and every live alias.

Navigation is realized by directional \texttt{MATCH} patterns with an
unrestricted relationship variable, the concrete predicate
\texttt{type(r)\ STARTS\ WITH\ \textquotesingle{}Link\textquotesingle{}},
and metadata predicates on:

\begin{Shaded}
\begin{Highlighting}[]
\NormalTok{associationKey}
\NormalTok{sourceRole}
\NormalTok{targetRole}
\NormalTok{sourceQualifiers for forward navigation}
\NormalTok{targetQualifiers for reverse navigation}
\end{Highlighting}
\end{Shaded}

By Neo4j \texttt{MATCH} and \texttt{WHERE} semantics, and by the graph
encoding contract, the primitive pattern enumerates exactly
\texttt{nav\_graph} by C3, modulo duplicate rows. C5 restores VA set
semantics.

\texttt{Exists} is realized by \texttt{EXISTS\ \{\ ...\ \}}. Neo4j
returns true iff the subquery has at least one semantic witness by C4.
By induction on the source and predicate realization, this is exactly VA
existential semantics.

\texttt{ForAll}, when present after normalization, is realized either
directly or as absence of a counterexample. Correctness follows from
\texttt{Exists} realization and Boolean realization.

\texttt{Select} is realized by filtering the realized finite source set
with the realized predicate. \texttt{Reject} is realized either directly
as the complementary filter or through its normalized
\texttt{Select(...,\ Not(P))} form. In both cases, the induction
hypotheses for the source and predicate give the same keep/drop decision
as VA semantics.

\texttt{Collect} is realized as finite-set image construction over the
realized source set. The body expression is realized under the iterator
alias. \texttt{DISTINCT} is used after \texttt{setOut} according to C5,
so a bottom image remains one semantic set element. This does not claim
Bag multiplicity, ordering, or implicit flattening.

A set literal evaluates and materializes every element exactly once,
coerces each result to the functional \texttt{setJoinN} type, applies
\texttt{setOut}, and projects distinct values. D8 therefore gives
exactly its extensional typed set. \texttt{Union} uses CY9 followed by
distinct projection; \texttt{Intersection} retains a left row precisely
when a typed-equal right witness exists. \texttt{AsSet} only reprojects
an already realized finite set. These are D9-\/-D10. \texttt{IsUnique}
compares the source cardinality with the distinct, bottom-tokenized
projection cardinality; D11 proves that this is equivalent to
injectivity of the VA body image.

The reference cardinality lowering uses
\texttt{size(COLLECT\ \{\ ...\ RETURN\ DISTINCT\ setOut(x)\ \})} over
the realized set. The reference raw-AST proof may use
\texttt{COUNT(DISTINCT\ setOut(x))} only under the CY5 equivalence
stated above. In either form, C5 requires quotienting duplicate graph
paths before cardinality is observed.

\texttt{IsEmpty} and \texttt{NotEmpty} are realized by absence/presence
of rows or by cardinality comparison over the realized finite set.
Correctness follows from the cardinality realization and finite-set
semantics.

\texttt{Includes} and \texttt{Excludes} are realized as membership and
negated membership over the realized finite set. \texttt{IncludesAll} is
realized as a subset check: there is no element of the right-hand
realized set missing from the left-hand realized set.
\texttt{ExcludesAll} is realized as disjointness: there is no element
that belongs to both realized sets. When these sets are obtained from
\texttt{MATCH}, the distinctness discipline is applied before
membership, subset, or disjointness checks.

Every collection-typed receiver first passes through C5a. Thus a
whole-set bottom contributes zero rows, while bottom produced as an
actual set element survives as \texttt{BOTTOM\_TOKEN}. Finite-set
\texttt{EQ}/\texttt{NEQ} uses C5b\textquotesingle s two inclusions; it
never applies element-level \texttt{setEq} directly to two whole
collections.

Boolean, comparison, and arithmetic operations are realized by
\texttt{boolAst}, \texttt{semEq}, \texttt{semOrd}, and
\texttt{semArith}, restricted to compatible operands. Their shared total
\texttt{isBottom} test recognizes both Cypher null and
\texttt{BOTTOM\_TOKEN}: equality implements the typed bottom table,
while ordering and arithmetic produce Cypher null whenever either
operand is bottom. Validation truth is implemented under C4 by:

\begin{Shaded}
\begin{Highlighting}[]
\NormalTok{truth(predicate) = coalesce(predicate = true,false)}
\end{Highlighting}
\end{Shaded}

so the expression is always Boolean and null-like Cypher results are not
validation-true. In particular, a represented bottom token is compared
with \texttt{true}; it is never passed through \texttt{coalesce} as a
map-valued predicate.

\texttt{TypeKindOf(E,C)} and \texttt{Cast(E,C)} use the same bottom-safe
receiver split. Their bottom arms return false and null, respectively,
without a node pattern. Their entity arms check
\texttt{ObjectInstanceOf} membership to \texttt{classNode(C)}; R3
preserves conformance, and cast returns the receiver exactly on success.
Validation truth follows the graph interpretation\textquotesingle s
\texttt{bool\_val} policy.

After structural induction, the predicate is a
\texttt{ScalarPlan(L,predicateExpr)}. \texttt{ExpandI} appends the
context match, then the entire prelude \texttt{L}, then
\texttt{Where(NOT\ truth(predicateExpr))} and
\texttt{Return(DISTINCT,self.use\_id\ AS\ useId)}. C6 therefore gives
the exact formal set equality and no-ghost property. TXT4 composes
C1-\/-C5b; TXT5 proves closure and rendering of the formal raw AST.

Finally, \texttt{SpecPlanSim\_sound} derives \texttt{returnedIds}
equality for \texttt{q=parse\_Cypher(tq)} and
\texttt{(tq,pi)=T\_TEXT\^{}spec(plan)} from the structural
\texttt{SpecPlanAdequacy} certificate and the local realization lemmas.
The equality is not projected from adequacy. Finite parser-tree fixtures
or successful Java examples remain implementation evidence and are not
premises of this reference theorem.

\begin{center}\rule{0.5\linewidth}{0.5pt}\end{center}

\section{13. Theorem 6: End-to-End Validation
Equivalence}\label{13-theorem-6-end-to-end-validation-equivalence}

\subsection{Statement}\label{statement-6}

Assume A1-\/-A9 and the successful invariant-rooted chain:

\begin{Shaded}
\begin{Highlighting}[]
\NormalTok{ParseInvariant(tInv)=astInv}
\NormalTok{decodeInvariant\_val(astInv)=I=(cName,R,e)}
\NormalTok{resolveClass\_MM(cName)=C}
\NormalTok{T\_BIND\_INV(astInv,MM)=BInvariant(C,R,b)}
\NormalTok{va = T\_VA(b)}
\NormalTok{nva = T\_NORM(va)}
\NormalTok{plan = T\_CQM\^{}spec(C,R,nva)}
\NormalTok{(tq,pi) = T\_TEXT\^{}spec(plan)}
\NormalTok{q = parse\_Cypher(tq)}
\NormalTok{GraphAdequate(MM,M,G)}
\NormalTok{SpecPlanAdequacy(MM,C,nva,plan,q,pi,G)}
\end{Highlighting}
\end{Shaded}

Then:

\begin{Shaded}
\begin{Highlighting}[]
\NormalTok{returnedIds(q,G,pi)}
\NormalTok{= id[Viol\_OCL(e,C,M)]}
\NormalTok{= \{ id(o) | o in Viol\_OCL(e,C,M) \}.}
\end{Highlighting}
\end{Shaded}

By injectivity of \texttt{id}, the following are corollaries:

\begin{Shaded}
\begin{Highlighting}[]
\NormalTok{Viol\_OCL(e,C,M)}
\NormalTok{= \{ o in Obj(M) | id(o) in returnedIds(q,G,pi) \};}

\NormalTok{o in Viol\_OCL(e,C,M)}
\NormalTok{iff}
\NormalTok{id(o) in returnedIds(q,G,pi);}

\NormalTok{returnedIds(q,G,pi) subseteq \{id(o) | o in Obj(M)\}.}
\end{Highlighting}
\end{Shaded}

The first displayed ID-set equality is the principal conclusion. The
object-set/pointwise forms alone would not exclude an extra returned
identifier that denotes no source object.

\subsection{Required Results}\label{required-results}

\begin{Shaded}
\begin{Highlighting}[]
\NormalTok{Theorem 0 Object{-}Graph Representation}
\NormalTok{Theorem 1 Binder Soundness}
\NormalTok{Theorem 2 Validation Algebra Abstraction}
\NormalTok{Theorem 3 Normalization Preservation}
\NormalTok{Theorem 4 Validation Preservation}
\NormalTok{Theorem 5 Cypher Realization}
\NormalTok{object{-}id injectivity}
\NormalTok{set semantics for returnedIds}
\end{Highlighting}
\end{Shaded}

F1 establishes source/AST representation and bound rebinding, but it is
not a separate premise once A1-\/-A2 provide the successful invariant
chain.

The proof is intentionally compositional. The following trace records
the exact equality or membership fact supplied at each boundary; Theorem
6 does not silently repeat constructor proofs already discharged in the
component theorems.

{\def\LTcaptype{none} % do not increment counter
\begin{longtable}[]{@{}ll@{}}
\toprule\noalign{}
Boundary & Result consumed by Theorem 6 \\
\midrule\noalign{}
\endhead
\bottomrule\noalign{}
\endlastfoot
admitted source to Bound OCL & Theorem 1 and SB-Erase identify source
evaluation with bound evaluation \\
Bound OCL to VA & Theorem 2 identifies bound denotation with object-side
VA denotation \\
VA to normalized VA & Theorem 3 preserves object-side denotation and
typing \\
object-side VA to graph-side VA & Theorem 4 preserves encoded values
and, in particular, validation truth \\
graph-side normalized VA to generated query & Theorem 5 identifies
validation-false context objects with \texttt{returnedIds} \\
object model to graph representation & Theorem 0 supplies context
membership, object-node correspondence, and identifier preservation \\
returned identifier to source object & object-ID injectivity gives a
unique object representative \\
\end{longtable}
}

\subsection{Proof}\label{proof-6}

For every context object \texttt{o}, Theorems 1 and 2 give:

\begin{Shaded}
\begin{Highlighting}[]
\NormalTok{[[e]]\_OCLval(M,self{-}\textgreater{}o)}
\NormalTok{= [[b]]\_Bound(M,self{-}\textgreater{}o)}
\NormalTok{= [[va]]\_obj(M,self{-}\textgreater{}o).}
\end{Highlighting}
\end{Shaded}

Theorem 3 and A8 then give:

\begin{Shaded}
\begin{Highlighting}[]
\NormalTok{bool\_val([[va]]\_obj(M,self{-}\textgreater{}o))}
\NormalTok{= bool\_val([[nva]]\_obj(M,self{-}\textgreater{}o)).}
\end{Highlighting}
\end{Shaded}

Theorem 4, using Theorem 0\textquotesingle s observation adequacy,
gives:

\begin{Shaded}
\begin{Highlighting}[]
\NormalTok{bool\_val([[nva]]\_obj(M,self{-}\textgreater{}o))}
\NormalTok{= bool\_val([[nva]]\_graph\^{}Boolean(G,self{-}\textgreater{}node(o))).}
\end{Highlighting}
\end{Shaded}

Hence the object set

\begin{Shaded}
\begin{Highlighting}[]
\NormalTok{Vg = \{ o in Obj(M)}
\NormalTok{       | classOf(o) conformsTo C}
\NormalTok{         and bool\_val([[nva]]\_graph\^{}Boolean}
\NormalTok{                (G,self{-}\textgreater{}node(o)))=false \}}
\end{Highlighting}
\end{Shaded}

is extensionally equal to \texttt{Viol\_OCL(e,C,M)}. Theorem 0 also
preserves the context observation and maps each such object to its
unique canonical graph node with property \texttt{use\_id=id(o)}.

Theorem 5, including C6 and \texttt{SpecPlanSim\_sound}, supplies the
exact target equality:

\begin{Shaded}
\begin{Highlighting}[]
\NormalTok{returnedIds(q,G,pi)=\{id(o) | o in Vg\}.}
\end{Highlighting}
\end{Shaded}

Substituting \texttt{Vg=Viol\_OCL(e,C,M)} yields:

\begin{Shaded}
\begin{Highlighting}[]
\NormalTok{returnedIds(q,G,pi)=id[Viol\_OCL(e,C,M)].}
\end{Highlighting}
\end{Shaded}

C6 already proves the no-ghost inclusion. Finally, injectivity of
\texttt{id} turns the ID-set equality into the object-set and pointwise
corollaries stated above.

\subsection{Applicability and Failure
Boundary}\label{applicability-and-failure-boundary}

Theorem 6 is internally complete relative to its premises, but it is not
applicable when any required boundary is undischargeable. In particular:

\begin{Shaded}
\begin{Highlighting}[]
\NormalTok{failed admission/binding}
\NormalTok{or failed graph{-}encoding conformance}
\NormalTok{or a normalized constructor without a Theorem 5 coverage row}
\NormalTok{or SpecPlanAdequacy is not discharged for the reference plan/text}
\NormalTok{or a collection{-}valued bottom can cross a collection boundary without}
\NormalTok{   normalization to zero rows/[]}
\NormalTok{or a selected Neo4j runtime that does not satisfy the mandatory dialect probes}
\NormalTok{or a violation of ScalarClosed/BottomSeparated}
\NormalTok{=\textgreater{} no Theorem 6 runtime claim.}
\end{Highlighting}
\end{Shaded}

Such a failure does not establish that the generated query is incorrect;
it means only that this theorem chain supplies no correctness result for
that input, encoding, renderer, or runtime. Conversely, passing tests
alone does not replace Theorems 1-\/-5: tests discharge
implementation/runtime conformance obligations, while the theorem chain
supplies the universal conditional argument for admitted models and
expressions.

\begin{center}\rule{0.5\linewidth}{0.5pt}\end{center}

\section{14. Reviewer-Facing
Boundary}\label{14-reviewer-facing-boundary}

The proof establishes conditional preservation, not universal compiler
correctness. The exact claim is:

\begin{Shaded}
\begin{Highlighting}[]
\NormalTok{For well{-}typed expressions in OCL\_val, under finite{-}set validation semantics,}
\NormalTok{if parsing, binding, algebra construction, normalization, Cypher realization,}
\NormalTok{and graph encoding all satisfy their stated contracts, then the generated}
\NormalTok{invariant query returns exactly the same violating objects as object{-}based OCL}
\NormalTok{validation.}
\end{Highlighting}
\end{Shaded}

The proof does not establish:

\begin{Shaded}
\begin{Highlighting}[]
\NormalTok{correctness for full OMG OCL}
\NormalTok{correctness for arbitrary Cypher}
\NormalTok{correctness of the unproved Java T\_OPT optimizer or arbitrary Java planning paths}
\NormalTok{unbounded Neo4j behavior outside the selected CY1{-}{-}CY9 runtime profile}
\NormalTok{preservation of Bag/Sequence/OrderedSet multiplicities or order}
\NormalTok{ordered collection operators such as sortedBy}
\NormalTok{deterministic value equality for any without a choice policy}
\NormalTok{one and count(element) without dedicated rewrites/operators and proof cases}
\NormalTok{correctness of fallback{-}only constructs}
\NormalTok{correctness of graph encodings that omit inherited ObjectInstanceOf edges}
\end{Highlighting}
\end{Shaded}

\begin{center}\rule{0.5\linewidth}{0.5pt}\end{center}

\section{15. Safe Paper Wording}\label{15-safe-paper-wording}

Use:

\begin{Shaded}
\begin{Highlighting}[]
\NormalTok{We prove conditional end{-}to{-}end preservation of violation sets for a finite{-}set}
\NormalTok{OCL validation fragment over an explicit UML{-}to{-}property{-}graph encoding.}
\end{Highlighting}
\end{Shaded}

Use:

\begin{Shaded}
\begin{Highlighting}[]
\NormalTok{The reference NVA{-}to{-}CQM{-}to{-}Cypher query realizes the graph interpretation of}
\NormalTok{the normalized validation algebra for the supported Cypher subset.}
\end{Highlighting}
\end{Shaded}

Avoid:

\begin{Shaded}
\begin{Highlighting}[]
\NormalTok{We prove correctness of OCL{-}to{-}Cypher.}
\NormalTok{We support full OCL semantics.}
\NormalTok{The generated Cypher is correct for all queries.}
\NormalTok{The graph representation is always faithful.}
\end{Highlighting}
\end{Shaded}

\begin{center}\rule{0.5\linewidth}{0.5pt}\end{center}

\section{16. Checklist of Remaining Proof
Obligations}\label{16-checklist-of-remaining-proof-obligations}

The current obligation status is summarized normatively in Section 1.2.
Most items below are maintenance conditions for already discharged
profile-specific evidence; they are not all open proof obligations. The
principal remaining correctness work is PO-18, while PO-20 is
publication-only and out of scope.

\begin{Shaded}
\begin{Highlighting}[]
\NormalTok{1. PA4 verifies that the checked synchronization profile creates}
\NormalTok{   ObjectInstanceOf edges to every conforming supertype; retain inherited{-}type}
\NormalTok{   and missing/spurious mutation coverage.}
\NormalTok{2. No{-}spurious type/link reflection is discharged by PA4/PA6 for their checked}
\NormalTok{   synchronization profiles; extend it to any future link representation.}
\NormalTok{3. PA1/PA6 discharge shared injective \textasciigrave{}associationKey\textasciigrave{} construction and logical}
\NormalTok{   association reflection for the checked binary{-}link profile; reopen this item}
\NormalTok{   for ternary or association{-}class representations.}
\NormalTok{4. PA5 discharges exact scalar \textasciigrave{}attributeKey\textasciigrave{} slot existence, uniqueness,}
\NormalTok{   no{-}spurious reflection, and equality with \textasciigrave{}value\_G\textasciigrave{} for the checked fixed{-}M2}
\NormalTok{   synchronization profile.}
\NormalTok{5. PA9 verifies invariant queries use RETURN DISTINCT self.use\_id; retain the}
\NormalTok{   generated{-}query mutation and interpret returnedIds as a set.}
\NormalTok{6. Retain the production cardinality shape}
\NormalTok{   \textasciigrave{}size(COLLECT \{ ... RETURN DISTINCT setOut(x) \})\textasciigrave{}; any alternative COUNT}
\NormalTok{   lowering must first prove the same bottom{-}token and duplicate behavior.}
\NormalTok{7. Keep oclIsTypeOf outside the certified OCL\_val profile. Reintroducing it}
\NormalTok{   requires a direct runtime{-}class accessor, typing/VA rules, and realization}
\NormalTok{   proof cases.}
\NormalTok{8. Keep any, sortedBy, Bag, Sequence, OrderedSet, and fallback{-}only constructs}
\NormalTok{   outside Theorem 6 unless additional proof cases are added.}
\NormalTok{   Keep collection{-} and reference{-}valued attributes outside the certified}
\NormalTok{   \textasciigrave{}VAttribute\textasciigrave{} case until their value domains and accessors are proved.}
\NormalTok{   Keep enumeration values outside until \textasciigrave{}Enum(E)\textasciigrave{}, typed storage/equality, and}
\NormalTok{   PA5 are extended together.}
\NormalTok{9. Retain the exact CY1{-}{-}CY9 assumptions for MATCH, WHERE, EXISTS, COLLECT,}
\NormalTok{   size/cardinality, WITH, UNWIND, CALL, UNION ALL, RETURN,}
\NormalTok{   property/parameter lookup, CASE, coalesce, and DISTINCT.}
\NormalTok{10. Back Theorem 0 and Theorem 5 with implementation conformance tests.}
\NormalTok{11. Verify \textasciigrave{}one\textasciigrave{} is either excluded or fully rewritten and proved.}
\NormalTok{12. Verify \textasciigrave{}count(element)\textasciigrave{} is either excluded or represented by \textasciigrave{}CountElem\textasciigrave{}.}
\NormalTok{13. Keep \textasciigrave{}collect\textasciigrave{} as finite{-}set image with a non{-}collection body type only.}
\NormalTok{    Retain executable post{-}bind checks that navigation and \textasciigrave{}collect\textasciigrave{} results are}
\NormalTok{    \textasciigrave{}Set\textasciigrave{}, and that every collection operator has a collection{-}typed source.}
\NormalTok{14. Verify \textasciigrave{}asSources\textasciigrave{}, \textasciigrave{}finiteSet\textasciigrave{}, and \textasciigrave{}encodeValue\textasciigrave{} are defined consistently.}
\NormalTok{15. PA7 verifies stored qualifier direction/arity/order and the typed qualifier}
\NormalTok{    codec; retain the independent expected{-}payload oracle and verify that every}
\NormalTok{    qualified{-}navigation lowering evaluates expressions before matching.}
\NormalTok{16. Reverse navigation storage/accessor direction is discharged by PA7 for the}
\NormalTok{    checked binary{-}link profile; retain binder/renderer regression coverage.}
\NormalTok{17. Record the exact Neo4j server version and execute every mandatory}
\NormalTok{    \textasciigrave{}CYPHER5\_val\textasciigrave{} dialect probe before reporting runtime realization evidence.}
\NormalTok{18. Retain graph{-}level PA5 regression evidence for completeness and reflection}
\NormalTok{    of the exact scalar \textasciigrave{}attributeKey\textasciigrave{} lookup against the same \textasciigrave{}value\_G\textasciigrave{} accessor.}
\NormalTok{19. Retain BR1{-}{-}BR10, \textasciigrave{}repr\_C\textasciigrave{}/\textasciigrave{}abs\_C\textasciigrave{}, and the non{-}null \textasciigrave{}BOTTOM\_TOKEN\textasciigrave{}}
\NormalTok{    discipline whenever semantic bottom can occur as a finite{-}set element.}
\NormalTok{20. PO{-}18 now mechanizes predicate{-}set algebra, flat typed{-}value}
\NormalTok{    \textasciigrave{}encodeValue\textasciigrave{} injectivity, all eleven normalization rewrite{-}family}
\NormalTok{    equations after validation{-}truth projection, typed N2 preservation for}
\NormalTok{    four abstract result indices, lexical named{-}to{-}de{-}Bruijn semantic}
\NormalTok{    correspondence, soundness and scoped{-}semantic preservation of the exact}
\NormalTok{    Java capture guard over the Boolean binder kernel, relative N4 normal form,}
\NormalTok{    idempotence and lexicographic root decrease, the}
\NormalTok{    structural Boolean core of Theorem 4, and both abstract Theorem 6}
\NormalTok{    inclusions in Lean 4.32.2. Complete metamodel{-}specific OCL}
\NormalTok{    subtyping/coercion and concrete semantic agreement for every Java \textasciigrave{}T\_OPT\textasciigrave{}}
\NormalTok{    rule; instantiate the 16{-}constructor relational kernel with the actual}
\NormalTok{    object/graph evaluators; and complete the full typed \textasciigrave{}OCL\_val\textasciigrave{} induction}
\NormalTok{    before calling the theorem chain fully mechanized. PO{-}18 remains classified}
\NormalTok{    as recommended, but the active claim configuration deliberately treats that}
\NormalTok{    classification as blocking.}
\NormalTok{21. Because production renders text directly, retain PA14a{-}{-}d checks from}
\NormalTok{    generated text to the independent canonical tree and selected Neo4j parser}
\NormalTok{    projection. Any optional \textasciigrave{}parse\_Cypher(renderRaw(Q)) =alpha Q\textasciigrave{} result applies}
\NormalTok{    only to the reference raw AST unless production is explicitly migrated to it.}
\NormalTok{22. Retain and rerun the PO{-}21 \textasciigrave{}AdapterAdequate\textasciigrave{} certificate whenever graph}
\NormalTok{    vocabulary, synchronization, snapshot observation, codec, or renderer}
\NormalTok{    accessors change. The current static contracts include \textasciigrave{}modelKey\textasciigrave{}, but the}
\NormalTok{    prior runtime discharge is stale until the shared{-}snapshot mutations are}
\NormalTok{    recaptured on the canonical profile.}
\NormalTok{23. PA10 enforces \textasciigrave{}BottomSeparated\textasciigrave{} with two non{-}colliding representations:}
\NormalTok{    the non{-}null tagged{-}map set token and typed stored payload \textasciigrave{}v1|V\textasciigrave{}.}
\NormalTok{    Retain malformed/wrong{-}tag/string{-}\textasciigrave{}Undefined\textasciigrave{} collision tests and selected{-}}
\NormalTok{    runtime \textasciigrave{}DISTINCT\textasciigrave{}/cardinality probes whenever representation changes.}
\NormalTok{24. PA12 discharges bottom{-}safe attribute/kind/cast receivers and live{-}alias}
\NormalTok{    preservation for the checked canonical{-}objectKey profile. Retain generated}
\NormalTok{    guard{-}order tests and the nested real fixture whenever receiver lowering,}
\NormalTok{    object identity, or expression{-}subquery scoping changes.}
\NormalTok{25. PA8 discharges inherited \textasciigrave{}allInstances\textasciigrave{} for the checked profile; retain the}
\NormalTok{    independent accessor observation, missing/spurious mutations, and duplicate{-}}
\NormalTok{    membership cardinality test whenever the renderer template changes.}
\NormalTok{26. PA11 enforces \textasciigrave{}ScalarClosed\textasciigrave{} for checked runs with complete finite operand}
\NormalTok{    observations. Retain Int64/Real64, zero{-}divisor, exact{-}coercion, Unicode,}
\NormalTok{    graph{-}boundary, and OUT\_OF\_SCOPE regression tests whenever scalar lowering}
\NormalTok{    or the admitted scalar profile changes.}
\NormalTok{27. PA13a/b provide a sealed reference raw{-}AST datatype and total structural}
\NormalTok{    printer. Production intentionally retains direct{-}string rendering.}
\NormalTok{    PA14a/b/c/d establish normalization stability, 17/17 constructor witnesses,}
\NormalTok{    exact equality with the production{-}derived canonical token/group{-}tree, and}
\NormalTok{    selected Neo4j 2026.06.0 parser acceptance/AST{-}kind projection for 52}
\NormalTok{    checked cases including nested shadowing. Complete universal plan closure}
\NormalTok{    and full internal{-}AST equality before treating TXT5 as implementation}
\NormalTok{    closure.}
\NormalTok{28. PO{-}01/PO{-}14 class{-}node syntax now uses \textasciigrave{}UmlClass\textasciigrave{} plus the joint}
\NormalTok{    \textasciigrave{}(modelKey,classKey)\textasciigrave{} scope in the canonical graph, formal contract,}
\NormalTok{    production renderer, and all 52 reviewed canonical trees. Retain wrong{-}}
\NormalTok{    label, wrong{-}model, and cross{-}model{-}edge mutations and rerun runtime}
\NormalTok{    evidence whenever either predicate changes.}
\NormalTok{29. Keep association{-}class navigation outside certified admission. Supporting}
\NormalTok{    it later requires an explicit link{-}object/spoke graph contract, binder}
\NormalTok{    metadata, renderer rules, and new PA6/PA7/Theorem 4{-}{-}5 cases.}
\NormalTok{30. Preserve the \textasciigrave{}SEMANTIC {-}\textgreater{} OPTIMIZED\textasciigrave{} IR stage transition. No public planner}
\NormalTok{    entry may accept an unversioned expression; any new optimizer producer}
\NormalTok{    version must invalidate stale artifacts and rerun preservation evidence.}
\end{Highlighting}
\end{Shaded}

\begin{center}\rule{0.5\linewidth}{0.5pt}\end{center}

\section{17. What Was Fixed in This
Revision}\label{17-what-was-fixed-in-this-revision}

This revision tightens the proof without widening the claim.

\begin{Shaded}
\begin{Highlighting}[]
\NormalTok{1. Removed one from the main OCL\_val grammar and placed it outside Theorem 6}
\NormalTok{   unless rewritten to Count(Select(...)) = 1 and proved separately.}
\NormalTok{2. Removed count(element) from the main OCL\_val grammar; Count(S) is the}
\NormalTok{   internal VA finite{-}set cardinality lowering of source size().}
\NormalTok{3. Defined encodeValue uniformly for objects, scalars, bottom, and finite sets.}
\NormalTok{4. Defined \textasciigrave{}asSources\textasciigrave{} for singleton/set navigation receivers and kept it}
\NormalTok{   distinct from \textasciigrave{}finiteSet\textasciigrave{} at collection{-}operator positions.}
\NormalTok{5. Clarified collect as finite{-}set image construction only, with no Bag}
\NormalTok{   multiplicity, order, or implicit flattening claim.}
\NormalTok{6. Clarified qualified navigation: qualifier expressions are evaluated before}
\NormalTok{   matching the encoded qualifier payload, and only scalar primitive qualifiers}
\NormalTok{   with fixed serialization are covered.}
\NormalTok{7. Clarified reverse navigation through normalized binder metadata:}
\NormalTok{   association, roles, direction, and optional qualifier.}
\NormalTok{8. Completed missing Theorem 5 realization cases for Select, Reject, Collect,}
\NormalTok{   IncludesAll, ExcludesAll, IsEmpty, NotEmpty, TypeKindOf, and Cast; TypeOf is}
\NormalTok{   explicitly outside the certified fragment.}
\NormalTok{9. Kept the theorem chain and numbering unchanged.}
\NormalTok{10. Kept the claim conditional and restricted to the finite{-}set validation}
\NormalTok{    fragment and supported Cypher subset.}
\end{Highlighting}
\end{Shaded}

\begin{center}\rule{0.5\linewidth}{0.5pt}\end{center}

\section{18. What Was Tightened}\label{18-what-was-tightened}

This final tightening pass made the existing proof internally more
consistent without changing the pipeline, theorem numbering, or claim.

\begin{Shaded}
\begin{Highlighting}[]
\NormalTok{1. Added e{-}\textgreater{}collect(x | e) to the OCL\_val grammar because Collect is present}
\NormalTok{   in VA syntax, semantics, Theorem 4, and Theorem 5.}
\NormalTok{2. Kept collect scoped to finite{-}set image construction with a non{-}collection}
\NormalTok{   body type; no Bag multiplicity, ordering, flattening, or nested collection}
\NormalTok{   semantics are claimed.}
\NormalTok{3. Replaced object{-}only iterator correspondence in Theorem 4 with}
\NormalTok{   encodeValue(s), covering both object elements and scalar elements.}
\NormalTok{4. Added finiteSet(v) for collection operators and specified that}
\NormalTok{   finiteSet(bottom)=empty as a validation{-}level policy.}
\NormalTok{5. Updated Exists, ForAll, Select, Reject, Collect, Includes, Excludes,}
\NormalTok{   IncludesAll, ExcludesAll, Count, Size, IsEmpty, and NotEmpty to use}
\NormalTok{   finiteSet([[S]]).}
\NormalTok{6. Clarified qualified navigation evaluation as an ordered list:}
\NormalTok{   rawQs = [[[q1]]\_I(rho),...,[[qk]]\_I(rho)] and}
\NormalTok{   encodedQs = encodeQualifierList(rawQs), followed by payload comparison.}
\NormalTok{7. Kept qualifier preservation limited to scalar primitive qualifiers with}
\NormalTok{   fixed, componentwise injective serialization and matching declared arity.}
\NormalTok{8. Refined F1 as OCL\_val Representability and AST Admission. OCL\_val is now}
\NormalTok{   related to the parser AST by partial T\_BIND admission and a forgetful erase}
\NormalTok{   projection, rather than being written as a literal AST subset.}
\NormalTok{9. Split Theorem 0 representation support into seven lemmas: object}
\NormalTok{   injectivity, object{-}node preservation, type preservation, attribute}
\NormalTok{   preservation, association preservation, navigation preservation, and}
\NormalTok{   allInstances preservation.}
\NormalTok{10. Defined Rep\_G(o,n) before node(o), so R2 proves total and unique object}
\NormalTok{    representation rather than the nearly tautological existence of}
\NormalTok{    n=node(o).}
\NormalTok{11. Split binding support into B1{-}{-}B4 by separating deterministic navigation}
\NormalTok{    resolution from iterator/let scope preservation.}
\NormalTok{12. Added A3 environment lookup preservation and made A2 explicitly}
\NormalTok{    object{-}side, without encodeValue.}
\NormalTok{13. Added G4 to state compatibility of encodeValue with finiteSet and bool\_val.}
\NormalTok{14. Replaced the overly broad Cypher pattern lemma with C3 primitive graph}
\NormalTok{    access realization, then separated filters/existentials, set/cardinality,}
\NormalTok{    and the invariant wrapper into C4{-}{-}C6.}
\NormalTok{15. Added typed object/graph value domains and well{-}typed environment}
\NormalTok{    satisfaction, including equality of environment domains.}
\NormalTok{16. Fixed the result typing of arithmetic, annotated let bindings, and the}
\NormalTok{    uniform finite{-}set result policy for binary navigation.}
\NormalTok{17. Type{-}indexed Denote\_graph so primitive scalar access is not incorrectly}
\NormalTok{    assigned a set codomain.}
\NormalTok{18. Added an explicit Cypher\_val grammar, row/environment correspondence,}
\NormalTok{    scalar/predicate/set/invariant realization judgments, and CY1{-}{-}CY9}
\NormalTok{    behavioral assumptions.}
\NormalTok{19. Removed the node{-}returning formulation of Theorem 5. The theorem and the}
\NormalTok{    reverse direction of Theorem 6 now reason only about returnedIds.}
\NormalTok{20. Restricted F1 representability to derivable, admitted OCL\_val terms rather}
\NormalTok{    than quantifying over arbitrary metamodel instances.}
\NormalTok{21. Added repr\_C/abs\_C and a reserved non{-}null BOTTOM\_TOKEN so set cardinality}
\NormalTok{    remains correct when finite{-}set collect contains semantic bottom.}
\NormalTok{22. Added a scalar{-}compatibility contract and excluded unproved overflow, NaN,}
\NormalTok{    infinity, locale{-}sensitive collation, and backend{-}specific behavior.}
\NormalTok{23. Replaced representative binder prose with a complete syntax{-}directed}
\NormalTok{    construction table for every admitted OCL\_val AST form.}
\NormalTok{24. Replaced representative A2 equations with a complete Bound{-}to{-}VA mapping.}
\NormalTok{25. Added mutually recursive =\textgreater{}E/=\textgreater{}S T\_TEXT derivations, complete translation}
\NormalTok{    tables, invariant{-}root rule, and TXT1{-}{-}TXT5 meta{-}properties.}
\NormalTok{26. Added Size{-}to{-}Count canonicalization and a per{-}rule normalization proof}
\NormalTok{    table with typing, semantics, freshness, no{-}regeneration, and the}
\NormalTok{    lexicographic termination measure M.}
\NormalTok{27. Added a complete Source OCL/Bound OCL denotation correspondence table for}
\NormalTok{    every successful binder construction rule.}
\NormalTok{28. Instantiated the scalar compatibility premise with explicit policies for}
\NormalTok{    numeric coercion, division, representable range, overflow, String equality,}
\NormalTok{    NaN/infinity, and entity identity.}
\NormalTok{29. Strengthened the Cypher representation interface with BR1{-}{-}BR7, including}
\NormalTok{    token non{-}collision and stability through lookup, DISTINCT, and COUNT.}
\NormalTok{30. Added concrete \textasciigrave{}CYPHER5\_val\textasciigrave{} expansions and local realization obligations}
\NormalTok{    for target combinators, class/type access, directional navigation,}
\NormalTok{    cardinality, and set{-}valued conditionals.}
\NormalTok{31. Replaced emitted \textasciigrave{}Link*\textasciigrave{} notation by an unrestricted relationship variable}
\NormalTok{    plus \textasciigrave{}type(r) STARTS WITH \textquotesingle{}Link\textquotesingle{}\textasciigrave{}; \textasciigrave{}Link*\textasciigrave{} remains meta{-}notation only.}
\NormalTok{32. Added mandatory server{-}version and dialect probes and clarified that driver}
\NormalTok{    5.21.0 does not determine the Neo4j server semantics.}
\NormalTok{33. Added a complete Theorem 4 constructor exhaustiveness matrix, including}
\NormalTok{    normal{-}form reachability status for constructors eliminated by N4.}
\NormalTok{34. Made class lookup namespace{-}safe through injective \textasciigrave{}classKey\textasciigrave{}, strengthened}
\NormalTok{    the Let substitution/size argument, and fixed schema/scope obligations for}
\NormalTok{    set{-}valued If.}
\NormalTok{35. Added a constructive least{-}graph presentation \textasciigrave{}Phi\_star\textasciigrave{} while retaining}
\NormalTok{    Theorem 0\textquotesingle{}s paper{-}safe status as adequacy under the encoding contract.}
\NormalTok{36. Separated Source OCL and Bound OCL into disjoint tagged value domains and}
\NormalTok{    proved binder preservation through the logical relation \textasciigrave{}\textasciitilde{}=SB\textasciigrave{}.}
\NormalTok{37. Made raw Cypher AST the official target, introduced scalar/set construction}
\NormalTok{    plans, and factored T\_TEXT into total BuildCQ, Expand, and renderRaw stages.}
\NormalTok{38. Removed textual holes from Expand: branch imports, aliases, attribute}
\NormalTok{    lookup, navigation predicates, subqueries, and unions are raw constructors}
\NormalTok{    produced by deterministic pattern matching.}
\NormalTok{39. Made every \textasciigrave{}ExpandE\textasciigrave{} equation codomain{-}correct: set relations, count, and}
\NormalTok{    emptiness now return \textasciigrave{}ScalarPlan\textasciigrave{}, never a bare \textasciigrave{}CExpr\textasciigrave{}.}
\NormalTok{40. Added \textasciigrave{}withEntityReceiver\textasciigrave{}, whose complementary correlated{-}call branches}
\NormalTok{    prevent bottom/null/token values from entering graph{-}node pattern positions.}
\NormalTok{41. Re{-}expressed Attribute, TypeKindOf, and Cast lowering through the}
\NormalTok{    bottom{-}safe receiver combinator and proved singleton/live{-}alias preservation.}
\NormalTok{42. Added D1{-}{-}D11 derivations for qualified navigation, nested set relations,}
\NormalTok{    scalar{-}set iteration, set{-}valued If, bottom{-}valued Collect, Let, and}
\NormalTok{    bottom{-}sensitive entity operations, making TXT4 a derived result.}
\NormalTok{43. Audited \textasciigrave{}AdapterAdequate\textasciigrave{} against the concrete renderer and separated}
\NormalTok{    partial query{-}shape evidence from the still{-}open observation{-}equivalence}
\NormalTok{    obligations.}
\NormalTok{44. Synchronized the ScalarPlan, bottom{-}safe lowering, TXT4 cases, adapter}
\NormalTok{    boundary, and updated C1 name with the LaTeX proof artifact.}
\NormalTok{45. Assigned stable assumption, theorem{-}contract, and semantic{-}function IDs:}
\NormalTok{    \textasciigrave{}A1{-}{-}A9\textasciigrave{}, \textasciigrave{}PC{-}T0{-}{-}PC{-}T6\textasciigrave{}, and \textasciigrave{}SF{-}01{-}{-}SF{-}29\textasciigrave{}.}
\NormalTok{46. Replaced the scattered semantic signature declarations by the canonical}
\NormalTok{    registry in Section 3.2 and made all later notation an abbreviation of it.}
\NormalTok{47. Expanded R1{-}{-}R7, F1, B1{-}{-}B4, A1{-}{-}A3, and G1{-}{-}G4 into explicit direction,}
\NormalTok{    constructor, or typed{-}value case arguments instead of one{-}line sketches.}
\NormalTok{48. Added a versioned machine{-}readable proof registry, a content digest in}
\NormalTok{    both publication artifacts, and an executable synchronization check.}
\NormalTok{49. Added an exhaustive Theorem 5 constructor{-}coverage certificate linking}
\NormalTok{    each normalized VA family to its expansion, local proof obligation,}
\NormalTok{    dependencies, and coverage status.}
\NormalTok{50. Factored quantifier, filter, collect, membership, set{-}relation,}
\NormalTok{    cardinality, type{-}access, conditional, and let reasoning into explicit}
\NormalTok{    \textasciigrave{}LR{-}*\textasciigrave{} local realization lemmas.}
\NormalTok{51. Expanded B2 with a complete binder typing and non{-}bottom}
\NormalTok{    semantic{-}membership case matrix.}
\NormalTok{52. Classified CY1{-}{-}CY9 into representation, core{-}query, and version{-}sensitive}
\NormalTok{    trusted boundaries, with minimum runtime evidence and explicit}
\NormalTok{    non{-}coverage for each assumption.}
\NormalTok{53. Exposed the exact dependency certificates of R5/R6 and added countermodels}
\NormalTok{    showing why association{-}key, qualifier, role, and direction premises are}
\NormalTok{    necessary.}
\NormalTok{54. Added exhaustive F1 and B1 case certificates for source/AST round trips}
\NormalTok{    and all reference{-}bearing bound constructors.}
\NormalTok{55. Added a Theorem 6 dependency trace and an explicit applicability/failure}
\NormalTok{    boundary separating theorem reasoning from runtime conformance evidence.}
\NormalTok{56. Added a one{-}to{-}one CY1{-}{-}CY9 selected{-}runtime matrix, duplicate{-}path fixture,}
\NormalTok{    exact expected/observed evidence rows, and profile/source{-}hash freshness}
\NormalTok{    guard while retaining every CY item as an explicit trusted assumption.}
\NormalTok{57. Replaced the weak OCL\_val{-}47 seed by a multiplicity{-}valid non{-}vacuity}
\NormalTok{    fixture, fixed exact expected IDs before Cypher execution, classified 28}
\NormalTok{    countermodel{-}capable rows and proved the profile basis of 19 tautological}
\NormalTok{    rows, then recorded 47/47 USE{-}{-}Neo4j exact{-}ID agreement.}
\NormalTok{58. Upgraded the proof registry to schema 2 with explicit PO status/evidence,}
\NormalTok{    runtime profile, 47 admitted constructors, and eight implementation}
\NormalTok{    vocabulary terms; exact generated Markdown/LaTeX contract blocks and six}
\NormalTok{    killed proof{-}sync mutations now protect the normative projection.}
\NormalTok{59. Added a Lean 4.32.2 project linked to the proof{-}contract version and}
\NormalTok{    registry digest. The Lean kernel first checked 15 named theorems for finite{-}set}
\NormalTok{    transport, flat value{-}encoding injectivity, selected normalization laws,}
\NormalTok{    structural Boolean preservation, and the two abstract Theorem 6}
\NormalTok{    inclusions. Static guards reject hidden \textasciigrave{}sorry\textasciigrave{}/\textasciigrave{}admit\textasciigrave{}/\textasciigrave{}axiom\textasciigrave{}/\textasciigrave{}opaque\textasciigrave{},}
\NormalTok{    and six killed mutations demonstrate hash, theorem{-}list, project{-}axiom,}
\NormalTok{    placeholder, kernel{-}type{-}error, and unapproved{-}core{-}axiom sensitivity.}
\NormalTok{    The audit records Lean\textquotesingle{}s core \textasciigrave{}propext\textasciigrave{} dependency explicitly. At this}
\NormalTok{    initial checkpoint PO{-}18 remained partial; item 60 records the subsequent}
\NormalTok{    all{-}rule normalization extension.}
\NormalTok{60. Extended the Lean normalization kernel to all eleven formal rule heads.}
\NormalTok{    A typed semantic certificate covers Size, ForAll, NotEmpty, IsEmpty,}
\NormalTok{    Reject, Implies, Xor, Includes, Excludes, Count{-}positive and Count{-}zero;}
\NormalTok{    the structural normalizer is total, reaches the stated relative normal}
\NormalTok{    form, is idempotent, and every bottom{-}up root step decreases the formal}
\NormalTok{    lexicographic measure. The required theorem inventory at that checkpoint}
\NormalTok{    was 19. This does not equate Java \textasciigrave{}T\_OPT\textasciigrave{} syntax with formal \textasciigrave{}T\_NORM\textasciigrave{}.}
\NormalTok{61. Added intrinsic \textasciigrave{}Bool\textasciigrave{}/\textasciigrave{}Nat\textasciigrave{}/\textasciigrave{}Elem\textasciigrave{}/\textasciigrave{}Set\textasciigrave{} indices for the normalization}
\NormalTok{    rule heads, erasure/type inference, and the kernel theorem}
\NormalTok{    \textasciigrave{}typed\_rewrite\_preserves\_type\textasciigrave{}. Added de Bruijn{-}scoped expressions and}
\NormalTok{    \textasciigrave{}scoped\_rename\_preserves\_binder\_boundary\textasciigrave{}. The Java optimizer now refuses}
\NormalTok{    iterator fusion when renaming would cross a target{-}named binder, with a}
\NormalTok{    regression preserving the original free iterator reference. The required}
\NormalTok{    theorem inventory at that checkpoint was 21; the full metamodel{-}specific type lattice,}
\NormalTok{    named{-}Java/de{-}Bruijn semantic correspondence, and universal \textasciigrave{}T\_OPT\textasciigrave{}}
\NormalTok{    refinement were still open at that checkpoint.}
\NormalTok{62. Strengthened the PO{-}19 publication gate. \textasciigrave{}check{-}proof{-}sync.ps1\textasciigrave{} can now}
\NormalTok{    retain a non{-}empty compiled PDF and emit a JSON manifest binding it to the}
\NormalTok{    proof{-}contract ID, registry hash, LaTeX{-}source hash, PDF hash, byte count,}
\NormalTok{    compiler version, and \textasciigrave{}SOURCE\_DATE\_EPOCH\textasciigrave{}. CI uploads that pair as a named}
\NormalTok{    artifact. The artifact branch and all hash checks pass under an isolated}
\NormalTok{    compiler smoke test. At that checkpoint PO{-}19 remained open pending a clean}
\NormalTok{    CI run; the 261{-}test checkpoint was discharged on 2026{-}08{-}04, while paper}
\NormalTok{    compilation moved to PO{-}20. PO{-}19 was reopened after the registry, Lean and}
\NormalTok{    Java evidence changed for the 263{-}test payload{-}refinement revision.}
\NormalTok{63. Added a lexical named{-}expression kernel, partial name resolution to}
\NormalTok{    de Bruijn indices, and independent named/scoped evaluators. Lean proves}
\NormalTok{    \textasciigrave{}named\_to\_scoped\_semantic\_correspondence\textasciigrave{}, proves the Boolean decision}
\NormalTok{    \textasciigrave{}JavaCaptureGuard\textasciigrave{} sound for the same three disjuncts used by}
\NormalTok{    \textasciigrave{}canRenameWithoutCapture\textasciigrave{}, and proves an accepted guarded rename preserves}
\NormalTok{    compiled scoped semantics. The inventory is now 24. This closes the}
\NormalTok{    abstract Boolean{-}binder bridge for guarded fusion, not the exhaustive map}
\NormalTok{    from every production \textasciigrave{}OclIr\textasciigrave{} constructor or the remaining optimizer rules.}
\NormalTok{64. Replaced the payload{-}free MK{-}T4 syntax by a payload{-}parametric 16{-}constructor}
\NormalTok{    kernel. The Java matrix now enumerates every record component and possible}
\NormalTok{    target plan constructor, and executable witnesses validate all 16 source}
\NormalTok{    constructors plus payload/lowering mutations. Concrete primitive semantics}
\NormalTok{    and optimizer{-}rule equivalence remain open, so PO{-}18 stays partial.}
\NormalTok{65. Recorded the successful clean machine{-}verification run for commit}
\NormalTok{    \textasciigrave{}24ce4f26\textasciigrave{} and discharged PO{-}19 for that 261{-}test checkpoint. PO{-}19 was}
\NormalTok{    reopened when the registry, Lean and Java evidence changed for 263 tests.}
\NormalTok{    Paper publication remains PO{-}20 and is not correctness evidence.}
\NormalTok{66. Applied the local{-}only publication policy: no paper{-}source mirror is tracked}
\NormalTok{    under \textasciigrave{}verification/\textasciigrave{}, and GitHub Actions neither compiles nor uploads the}
\NormalTok{    paper. The ignored sources under \textasciigrave{}md/\textasciigrave{} remain the local publication source.}
\NormalTok{67. Separated invariant text, parser AST, unresolved source OCL, resolved Bound}
\NormalTok{    OCL, and VA terms. The current refinement admits both inferred and explicit}
\NormalTok{    \textasciigrave{}let\textasciigrave{}/iterator declaration types through the functional \textasciigrave{}declType\textasciigrave{} rule,}
\NormalTok{    and preserves the resolved declaration type through BoundOCL, VA, and plan.}
\NormalTok{68. Closed the typed semantic domains through \textasciigrave{}Void\textasciigrave{}, scalar/entity finite}
\NormalTok{    Sets, indexed bottom values, \textasciigrave{}repr\_C\^{}tau\textasciigrave{}/\textasciigrave{}abs\_C\^{}tau\textasciigrave{}, BR1{-}{-}BR10, and the}
\NormalTok{    complete source{-}{-}Bound{-}{-}VA constructor and typing tables.}
\NormalTok{69. Added total collection{-}boundary, scalar{-}bottom, signed{-}zero, alias{-}scope,}
\NormalTok{    validation{-}truth, and \textasciigrave{}Void\textasciigrave{}{-}{-}\textasciigrave{}Set\textasciigrave{} equality repairs with admitted}
\NormalTok{    counterexamples and selected{-}runtime checks; the audit now records}
\NormalTok{    AI{-}01{-}{-}AI{-}34 rather than treating the earlier ten rows as exhaustive.}
\NormalTok{70. Made \textasciigrave{}PlanAdequacy\textasciigrave{} graph{-}specific, separated its open universal}
\NormalTok{    production certificate from finite evidence, repaired \textasciigrave{}WFScalar\textasciigrave{}, \textasciigrave{}CV\textasciigrave{},}
\NormalTok{    the direct{-}ONE \textasciigrave{}collView\textasciigrave{} boundary, and exact ID/no{-}ghost theorem codomains;}
\NormalTok{    the English LaTeX report is generated from this Markdown and guarded by}
\NormalTok{    the canonical source SHA{-}256.}
\end{Highlighting}
\end{Shaded}

\begin{center}\rule{0.5\linewidth}{0.5pt}\end{center}

\section{19. Proof-Artifact Synchronization
Discipline}\label{19-proof-artifact-synchronization-discipline}

The normative source is this document:

\begin{Shaded}
\begin{Highlighting}[]
\NormalTok{md/research/formal{-}theorems{-}and{-}proofs.md}
\end{Highlighting}
\end{Shaded}

Its canonical JSON block fixes the contract version and registry schema,
stable assumptions, scope lemmas \texttt{M1-\/-M5} (including
\texttt{M3a}), theorem statement kernels and exact dependencies,
semantic functions, PO-01-\/-PO-24 classification/status/evidence, the
selected runtime profile, the 47 admitted constructor rows, the 17 plan
constructors, the 16 Java IR constructors, and implementation
vocabulary. The derived registry is:

\begin{Shaded}
\begin{Highlighting}[]
\NormalTok{verification/contract/proof{-}contract{-}registry.json}
\end{Highlighting}
\end{Shaded}

The tracked Lean artifact carries the SHA-256 digest of that byte-stable
projection. After changing the formal contract block, run:

\begin{Shaded}
\begin{Highlighting}[]
\NormalTok{powershell }\OperatorTok{{-}}\NormalTok{NoProfile }\OperatorTok{{-}}\NormalTok{ExecutionPolicy Bypass }\OperatorTok{{-}}\NormalTok{File verification}\OperatorTok{/}\NormalTok{scripts}\OperatorTok{/}\NormalTok{sync{-}proof{-}contract{-}from{-}formal}\OperatorTok{.}\FunctionTok{ps1} \OperatorTok{{-}}\NormalTok{UpdateDerived}
\NormalTok{powershell }\OperatorTok{{-}}\NormalTok{NoProfile }\OperatorTok{{-}}\NormalTok{ExecutionPolicy Bypass }\OperatorTok{{-}}\NormalTok{File }\FunctionTok{md}\OperatorTok{/}\NormalTok{research}\OperatorTok{/}\NormalTok{check{-}proof{-}sync}\OperatorTok{.}\FunctionTok{ps1} \OperatorTok{{-}}\NormalTok{UpdateProjections }\OperatorTok{{-}}\NormalTok{SkipArtifactBuild}
\NormalTok{powershell }\OperatorTok{{-}}\NormalTok{NoProfile }\OperatorTok{{-}}\NormalTok{ExecutionPolicy Bypass }\OperatorTok{{-}}\NormalTok{File verification}\OperatorTok{/}\NormalTok{scripts}\OperatorTok{/}\NormalTok{check{-}verification{-}contract}\OperatorTok{.}\FunctionTok{ps1}
\NormalTok{powershell }\OperatorTok{{-}}\NormalTok{NoProfile }\OperatorTok{{-}}\NormalTok{ExecutionPolicy Bypass }\OperatorTok{{-}}\NormalTok{File verification}\OperatorTok{/}\NormalTok{scripts}\OperatorTok{/}\NormalTok{check{-}verification{-}contract{-}mutations}\OperatorTok{.}\FunctionTok{ps1}
\end{Highlighting}
\end{Shaded}

The check first rejects any derived registry that differs from the
embedded canonical block, then rejects stale hashes, missing/extra
theorem premises or dependencies, unknown/missing evidence,
registry-\/-CSV-\/-admission constructor drift, exact vocabulary drift,
and an active prototype-correctness claim while a configured blocking PO
is open or partial. The current synchronization/contract mutation gates
demonstrate seven rejected drifts, including an attempted edit of only
the derived registry. failures. The historical clean record is commit
\texttt{7b7bdd25}: 311/311 selected Java tests, Lean 32/32, both
six-mutant gates, and uploaded proof/build artifacts. PO-19 is
discharged only for that checkpoint. It does not certify the present
working copy: the current 335/335 conformance gate passes, but
renderer/runtime-source hashes changed and the three common runtime
manifests deliberately remain stale pending a new clean capture. PO-20
is explicitly local-only and outside the machine-correctness contract:
GitHub contains no paper mirror, build job, PDF, or publication hash.
The remaining drift/conformance checks are not a proof assistant:
mathematical validity of each equation remains a proof obligation.

\subsection{19.1 Selected Lean Mechanization
Boundary}\label{191-selected-lean-mechanization-boundary}

The checked proof-assistant artifact is:

\begin{Shaded}
\begin{Highlighting}[]
\NormalTok{verification/lean/Ocl2CypherProof.lean}
\NormalTok{Lean 4.32.2}
\NormalTok{proof contract PC{-}2026{-}07{-}22.3}
\end{Highlighting}
\end{Shaded}

The registry fixes the toolchain, source/checker paths, 32 required
theorem names, coverage status, and open scope.
\texttt{check-mechanized-proof.ps1} compiles the source with the Lean
kernel and rejects unapproved theorem-list drift plus \texttt{sorry},
\texttt{admit}, project-declared \texttt{axiom}, or \texttt{opaque}. The
axiom audit permits Lean core dependencies explicitly: \texttt{propext}
occurs in twelve registered proofs and two dependent-scope proofs also
use \texttt{Quot.sound}. The all-rule semantic certificate and
registered pointwise T6 theorem must report no axiom dependency. Its
mutation companion must kill stale registry-hash, placeholder,
project-axiom, missing-theorem, real kernel type-error, and
unapproved-core- axiom mutations.

The mechanized statements establish the following mathematical kernels:

\begin{Shaded}
\begin{Highlighting}[]
\NormalTok{MK{-}FINITE{-}SET  predicate{-}set image/subset/exists/forall transport}
\NormalTok{MK{-}LIFT1       to{-}one absent/empty and present/singleton collection views}
\NormalTok{MK{-}ENCODE      injectivity for flat bottom/scalar/entity/finite{-}set values}
\NormalTok{MK{-}NORMALIZE   11 semantics; types; named/de{-}Bruijn guard bridge; NF/idempotence/decrease}
\NormalTok{MK{-}T4          payload{-}parametric 16{-}constructor relational lifting}
\NormalTok{MK{-}BOUND{-}VA    11 Bound records to 12 semantic{-}IR targets under shared primitives}
\NormalTok{MK{-}PA{-}COMP     AdapterAdequate composition from exact M2 and PA1{-}{-}PA9}
\NormalTok{MK{-}T6          forward/backward/pointwise ID inclusions under stated agreement}
\end{Highlighting}
\end{Shaded}

PO-18 as a whole remains deliberately \texttt{partial}. The
normalization artifact covers all listed rule families after
\texttt{bool\_val}, four abstract typed N2 result indices, lexical
named-to-de-Bruijn semantic correspondence, and guarded-renaming
preservation over its Boolean binder kernel. It does not cover the
metamodel-specific OCL subtype/coercion lattice or universal Java
\texttt{T\_OPT}-\/-formal \texttt{T\_NORM} refinement. MK-T4 now
quantifies over every non-recursive payload of all 16 production
\texttt{OptimizedExpression} constructors; the Java matrix checks every
record field and possible target plan constructor, and executable
witnesses cover the 16/16 source universe. It still does not instantiate
the primitive operations with complete object/graph semantics or prove
every Java optimizer rewrite. The artifact also does not represent
nested collection values, the full typed \texttt{OCL\_val} induction,
Neo4j/Cypher semantics, or \texttt{AdapterAdequate} for arbitrary
deployments. Therefore kernel acceptance establishes that these selected
Lean propositions follow from their explicit premises; it does not
establish the end-to-end correctness of the prototype.

The historical machine-verification run at commit \texttt{7b7bdd25}
covers its contract-.3 checkpoint: clean baseline \texttt{f0938519},
311/311 selected Java tests, Lean 32/32, both six-mutant gates, and
successful proof/build artifact uploads. It is not a clean provenance
record for the current working copy. PO-19 remains discharged for that
historical CI artifact only; this does not expand the conditional
semantic scope. PO-20 remains open, out of correctness scope, and
local-only: paper sources and compiled outputs are ignored, not mirrored
in \texttt{verification/}, and not built or uploaded by GitHub Actions.

\begin{center}\rule{0.5\linewidth}{0.5pt}\end{center}

\section{20. Conclusion}\label{20-conclusion}

This report establishes a proof chain that is closed relative to its
explicit premises. Theorem 0 connects UML models with validation-visible
graph observations. Theorems 1-\/-3 preserve binding, Validation Algebra
abstraction, and normalization. Theorem 4 relates object-side and
graph-side VA interpretations. Theorem 5 realizes the graph
interpretation in the selected Cypher fragment. Theorem 6 composes these
results into equality of the source and returned-ID violation sets. Set
literals, \texttt{xor}, \texttt{union}, \texttt{intersection},
\texttt{asSet}, and \texttt{isUnique} are covered by dedicated typing,
typed-coercion, realization, and D8-\/-D11 arguments.

The result is conditional correctness of the formal specification, not
an unqualified correctness claim for every implementation. A concrete
deployment must supply an applicable \texttt{AdapterAdequate}
certificate, raw-AST/parser bridge evidence, scalar-profile closure,
bottom separation, and the probes for the selected Cypher runtime. PO-21
supplies the adapter certificate for the current canonical profile only;
it does not automatically cover another encoding or externally mutated
graph. This separation identifies precisely which obligations are
mathematical results and which require implementation conformance
evidence.

\end{document}
